% !TEX program = xelatex

%**start of header
\documentclass[11pt]{amsart} 
\usepackage{amscd,amssymb,amsxtra}
\usepackage{mathrsfs}  % \mathrsfs{A}
\DeclareSymbolFontAlphabet{\mathrsfs}{rsfs}
\usepackage[mathscr]{eucal} % \mathscr{A} 
\usepackage{mathabx}\DeclareMathSymbol{\emptyset}{\mathord}{symbols}{59}
                    \DeclareMathSymbol{\boxtimes}{\mathbin}{AMSa}{"02}
\usepackage{comment}
\usepackage{color}
\usepackage{enumitem} 
\usepackage{marginnote}

% ===== 中译本：CJK 支持（XeLaTeX / Tectonic）=====
\usepackage{xeCJK}
\setCJKmainfont{SimSun}[BoldFont=SimHei, ItalicFont=KaiTi, BoldItalicFont=SimHei,
                        SmallCapsFont=SimHei]
\setCJKsansfont{SimHei}
\setCJKmonofont{FangSong}
\xeCJKsetup{CJKmath=true}
\XeTeXlinebreaklocale "zh"
\XeTeXlinebreakskip = 0pt plus 1pt

\setlength{\textwidth}{6.5truein} \setlength{\hoffset}{-.5truein} 
\setlength{\textheight}{8.9truein} \setlength{\voffset}{-.4truein} 
\setlength{\abovedisplayskip}{18pt plus4.5pt minus9pt}
\setlength{\belowdisplayskip}{\abovedisplayskip}
\setlength{\abovedisplayshortskip}{0pt plus4.5pt}
\setlength{\belowdisplayshortskip}{10.5pt plus4.5pt minus6pt}

\renewcommand\thesubsection{\S\thesection.\arabic{subsection}}

\makeatletter
\let\@secnumfont\bfseries
\def\section{\@startsection{section}{1}%
  \z@{4\linespacing\@plus\linespacing}{\linespacing}%
  {\bfseries\centering}}
\def\introsection{\@startsection{section}{1}%
  \z@{3\linespacing\@plus\linespacing}{\linespacing}%
  {\bfseries\centering}}
\def\subsection{\@startsection{subsection}{2}%
   \z@{1.25\linespacing\@plus.7\linespacing}{.5\linespacing}%
   {\normalfont\bfseries}}
\def\subsectionsinline{\def\subsection{\@startsection{subsection}{2}%
  \z@{1\linespacing\@plus.7\linespacing}{-.5em}%
  {\normalfont\bfseries}}}

\def\noremarks{\def\rem##1{\relax}\def\refrem##1{\relax}}

\def\thmqed{\qed\vskip1.5em}

\newenvironment{qedequation}{%
   \pushQED{\qed}%
   \incr@eqnum
   \mathdisplay@push
   \st@rredfalse \global\@eqnswtrue
   \mathdisplay{equation}%
}{%
   \endmathdisplay{equation}%
   \mathdisplay@pop
   \ignorespacesafterend
   \popQED\@endpefalse
}

\makeatother

%%%%%%%%%%%%%%%%%%%%%%%%%%%%%%%%%%%%%%%%%%%
\theoremstyle{definition}
\newtheorem{definition}[equation]{定义}
\newtheorem{example}[equation]{例}
\newtheorem{problem}[equation]{\color{blue}Problem}
\newtheorem{exercise}[equation]{练习}
\newtheorem{question}[equation]{\color{blue}Question}
\newtheorem{project}[equation]{\color{blue}Project}
\newtheorem{construction}[equation]{构造}
\newtheorem{task}[equation]{\color{blue}Task}
\newtheorem{otask}[equation]{\color{green}Obsolete Task}

\newtheorem*{definition*}{定义}
\newtheorem*{example*}{例}
\newtheorem*{problem*}{\color{blue}Problem}
\newtheorem*{exercise*}{练习}
\newtheorem*{question*}{\color{blue}Question}
\newtheorem*{project*}{\color{blue}Project}
\newtheorem*{construction*}{构造}

\theoremstyle{remark}
\newtheorem{note}[equation]{注记}
\newtheorem{notation}[equation]{记号}
\newtheorem{remark}[equation]{注}
\newtheorem{data}[equation]{数据}

\newtheorem*{note*}{注记}
\newtheorem*{notation*}{记号}
\newtheorem*{remark*}{注}
\newtheorem*{data*}{数据}

\theoremstyle{plain}
\newtheorem{theorem}[equation]{定理}
\newtheorem{corollary}[equation]{推论}
\newtheorem{lemma}[equation]{引理}
\newtheorem{proposition}[equation]{命题}
\newtheorem{conjecture}[equation]{猜想}
\newtheorem{claim}[equation]{断言}
\newtheorem{proposal}[equation]{提议}
\newtheorem{conclusion}[equation]{结论}
\newtheorem{hypothesis}[equation]{假设}
\newtheorem{assumption}[equation]{假定}

 
\newtheorem*{theorem*}{定理}
\newtheorem*{corollary*}{推论}
\newtheorem*{lemma*}{引理}
\newtheorem*{proposition*}{命题}
\newtheorem*{conjecture*}{猜想}
\newtheorem*{claim*}{断言}
\newtheorem*{proposal*}{提议}
\newtheorem*{conclusion*}{结论}
\newtheorem*{hypothesis*}{假设}
\newtheorem*{assumption*}{假定}

\newenvironment{proof*}[1][\proofname]{
  \begin{proof}[#1]}{  \renewcommand\qedsymbol{\relax}
\end{proof}} 

%\renewcommand{\qedsymbol}{\relax}

\numberwithin{equation}{section}

\newcommand{\rem}[1]{\marginnote%
  [$\bf\color{blue}\Rightarrow$]{$\bf\color{blue}\Leftarrow$}\relax[{\tiny 
\color{blue}\bf #1}]} 



\definecolor{refkey}{rgb}{0,.6,.4}
\newcommand{\refrem}[1]{\marginnote%
  [$\bf\color{refkey}\Rightarrow$]{$\bf\color{refkey}\Leftarrow$}\relax[{\tiny 
\color{refkey} \bf #1}]} 

\renewcommand{\:}{\colon}
%%%%%%%%%%%%%%%%%%%%%%%%%%%%%%%%%%%%%
\renewcommand{\AA}{{\mathbb A}}
\newcommand{\Ahat}{{\hat A}}
\DeclareMathOperator{\Aut}{Aut}
\newcommand{\CC}{{\mathbb C}}
\newcommand{\CP}{{\mathbb C\mathbb P}}
\DeclareMathOperator{\Conf}{Conf}
\DeclareMathOperator{\Diff}{Diff}
\newcommand{\EE}{\mathbb E}
\newcommand{\FF}{\mathbb F}
\DeclareMathOperator{\End}{End}
\DeclareMathOperator{\Free}{Free}
\newcommand{\HH}{{\mathbb H}}
\DeclareMathOperator{\Hom}{Hom}
\DeclareMathOperator{\id}{id}
\DeclareMathOperator{\Ker}{Ker}
\DeclareMathOperator{\Map}{Map}
\DeclareMathOperator{\Met}{Met}
\newcommand{\NN}{{\mathbb N}}
\newcommand{\PP}{{\mathbb P}}
\DeclareMathOperator{\pt}{pt}
\newcommand{\QQ}{{\mathbb Q}}
\newcommand{\RP}{{\mathbb R\mathbb P}}
\newcommand{\RR}{{\mathbb R}}
\newcommand{\TT}{\mathbb T}
\DeclareMathOperator{\Spin}{Spin}
\DeclareMathOperator{\Tor}{Tor}
\DeclareMathOperator{\Trs}{Tr\mstrut _s}
\DeclareMathOperator{\Tr}{Tr}
\DeclareMathOperator{\tr}{tr}
\DeclareMathOperator{\trs}{tr\mstrut _s}
\newcommand{\ZZ}{{\mathbb Z}}
\DeclareMathOperator{\ch}{ch}
\newcommand{\chiup}{\raise.5ex\hbox{$\chi$}}
\newcommand{\cir}{S^1}
\DeclareMathOperator{\coker}{coker}
\newcommand{\dbar}{{\overline\partial}}
\DeclareMathOperator{\free}{free}
%\DeclareMathOperator{\index}{index}
\DeclareMathOperator{\ind}{ind}
\newcommand{\inv}{^{-1}}
%\newcommand{\mstrut}{\rule[-165000sp]{0pt}{0pt}}
\DeclareRobustCommand{\mstrut}{^{\vphantom{1*\prime y\vee M}}}
\newcommand{\mlstrut}{_{\vphantom{1*\prime y}}}
\DeclareMathOperator{\rank}{rank}
\newcommand{\res}[1]{\negmedspace\bigm|\mstrut_{#1}}
\newcommand{\Res}[1]{\negmedspace\Bigm|\mstrut_{#1}}
\newcommand{\temsquare}{\raise3.5pt\hbox{\boxed{ }}}
\DeclareMathOperator{\tor}{tor}
\newcommand{\zmod}[1]{\ZZ/#1\ZZ}
\newcommand{\zn}{\zmod{n}}
\newcommand{\zt}{\zmod2}
%\DeclareMathOperator{\Im}{Im}
%\DeclareMathOperator{\Re}{Re}
%\let\germ\frak
 
\newcommand{\theprotag}[2]{#2~#1}
\newcommand{\longhookrightarrow}{\lhook\joinrel\longrightarrow}
 
\DeclareMathOperator{\SO}{SO}
\DeclareMathOperator{\Sp}{Sp}
\let\O\relax
\DeclareMathOperator{\O}{O}
\DeclareMathOperator{\U}{U}
\DeclareMathOperator{\SU}{SU}
\DeclareMathOperator{\GL}{GL}
\DeclareMathOperator{\SL}{SL}
\DeclareMathOperator{\PSL}{PSL}
\DeclareMathOperator{\PGL}{PGL}
\DeclareMathOperator{\PSU}{PSU}
\DeclareMathOperator{\PU}{PU}


\renewcommand{\thesubsection}{\thesection.\arabic{subsection}}
\usepackage[all,2cell]{xy}
\usepackage{longtable}
\usepackage{booktabs}\renewcommand{\cir}{\ensuremath{S^1}}
\usepackage[colorlinks,backref=page,citecolor=refkey]{hyperref}
\let\O\relax\DeclareMathOperator{\O}{O}
\DeclareMathOperator{\Bord}{Bord}
\DeclareMathOperator{\Cliff}{Cliff}
\DeclareMathOperator{\Det}{Det}
\DeclareMathOperator{\Ell}{Ell}
\DeclareMathOperator{\Euler}{Euler}
\DeclareMathOperator{\Fix}{Fix}
\DeclareMathOperator{\Fred}{Fred}
\DeclareMathOperator{\Imm}{Im}
\DeclareMathOperator{\Line}{Line}
\DeclareMathOperator{\Sign}{Sign}
\DeclareMathOperator{\Sym}{Sym}
\DeclareMathOperator{\Todd}{Todd}
\DeclareMathOperator{\Vect}{Vect}
\DeclareMathOperator{\Vol}{Vol}
\DeclareMathOperator{\ad}{ad}
\DeclareMathOperator{\curv}{curv}
\DeclareMathOperator{\hol}{hol}
\DeclareMathOperator{\image}{image}
\DeclareMathOperator{\ord}{ord}
\DeclareMathOperator{\pfaff}{pfaff}
\DeclareMathOperator{\pr}{pr}
\DeclareMathOperator{\sign}{sign}
\DeclareMathOperator{\spec}{spec}
\newcommand{\AP}{\ssA_{\sP}}
\newcommand{\BSO}{B\!\SO}
\newcommand{\Bd}{\Bord_{\langle d,d+1  \rangle}}
\newcommand{\Clpm}{\Cliff_{\pm n}}
\newcommand{\Cnm}{\Cliff_{-n}}
\newcommand{\Cnp}{\Cliff_{+n}}
\newcommand{\Cn}{\Cliff_{n}}
\newcommand{\Fn}{\Fred_n}
\newcommand{\Mn}{\mathbf{Man}_n}
\newcommand{\PH}{\PP\sH}
\newcommand{\RZ}{\RR/\ZZ}
\newcommand{\Rn}{\RR^n}
\newcommand{\Set}{\mathbf{Set}}
\newcommand{\Vk}{\Vect_k}
\newcommand{\bX}{\partial X}
\newcommand{\cP}{\sC_{\sP}}
\newcommand{\et}[1]{\eta \mstrut _{#1}}
\newcommand{\evo}[2]{e^{-i(t_{#1}-t_{#2})H/\hbar}}
\newcommand{\hooklongrightarrow}{\lhook\joinrel\longrightarrow} 
\newcommand{\op}{^{\textnormal op}}
\newcommand{\q}{^{(q)}}
\newcommand{\sA}{\mathcal{A}}
\newcommand{\sB}{\mathscr{B}}
\newcommand{\sC}{\mathscr{C}}
\newcommand{\sE}{\mathscr{E}}
\newcommand{\sG}{\mathscr{G}}
\newcommand{\sH}{\mathscr{H}}
\newcommand{\sL}{\mathcal{L}}
\newcommand{\sM}{\mathscr{M}}
\newcommand{\sO}{\mathcal{O}}
\newcommand{\sP}{\mathscr{P}}
\newcommand{\sS}{\mathscr{S}}
\newcommand{\scrF}{\mathscr{F}}
\newcommand{\sect}[1]{\sL(#1)}
\newcommand{\ssA}{\mathfrak{A}}
\newcommand{\tD}{\widetilde{D}}
\newcommand{\tE}{\widetilde{E}}
\newcommand{\tG}{\widetilde{G}}
\newcommand{\tX}{\widetilde{X}}
\newcommand{\tc}{\tilde{c}}
\newcommand{\tsH}{\widetilde{\sH}}
\newcommand{\xia}[1]{\xi \mstrut _{#1}}
\newcommand{\zD}{\zeta \mstrut _\Delta }
\renewcommand{\Im}{\Imm}
\renewcommand{\Re}{\text{Re}}
\renewcommand{\SS}{\mathbb{S}}




% ===== 中译本：本节重点（Key Takeaways）环境 =====
\newenvironment{keytakeaways}{%
  \par\medskip
  \begingroup\small
  \noindent\rule{\linewidth}{0.5pt}\par\nobreak\vspace{3pt}
  \noindent{\bfseries 本节重点（Key Takeaways）}\par\nobreak\vspace{3pt}
  \begin{itemize}[leftmargin=1.5em,itemsep=1pt,topsep=1pt,parsep=0pt]
}{%
  \end{itemize}\endgroup
  \noindent\rule{\linewidth}{0.5pt}\par\medskip
}

  \begin{document}

\abovedisplayskip18pt plus4.5pt minus9pt
\belowdisplayskip \abovedisplayskip
\abovedisplayshortskip0pt plus4.5pt
\belowdisplayshortskip10.5pt plus4.5pt minus6pt
\baselineskip=15 truept
\marginparwidth=55pt

\makeatletter
\renewcommand{\tocsection}[3]{%
  \indentlabel{\@ifempty{#2}{\hskip1.5em}{\ignorespaces#1 #2.\;\;}}#3}
\renewcommand{\tocsubsection}[3]{%
  \indentlabel{\@ifempty{#2}{\hskip 2.5em}{\hskip 2.5em\ignorespaces#1%
    #2.\;\;}}#3} 
\renewcommand{\tocsubsubsection}[3]{%
  \indentlabel{\@ifempty{#2}{\hskip 5.5em}{\hskip 5.5em\ignorespaces#1%
    #2.\;\;}}#3} 
%  \indentlabel{\hskip 4em#3}}
\def\@makefnmark{%
  \leavevmode
  \raise.9ex\hbox{\fontsize\sf@size\z@\normalfont\tiny\@thefnmark}} 
\def\multfoot{\textsuperscript{\tiny\color{red},}}
\def\footref#1{$\textsuperscript{\tiny\ref{#1}}$}
%             \footnote{Note\label{SYM}} ... in footnote~\footref{SYM}
\makeatother

\setcounter{tocdepth}{2}

%**end of header

\renewcommand{\proofname}{证明}
\renewcommand{\contentsname}{目录}
\renewcommand{\refname}{参考文献}
% 批注（附录 A 重排）中可能用到的宏别名
\providecommand{\Coker}{\coker}
\providecommand{\Pfaff}{\pfaff}
\providecommand{\Sect}{\sect}



% ===== 中译本：标题页 =====
\title[阿蒂亚--辛格指数定理]{阿蒂亚--辛格指数定理\\[6pt]
  {\large\mdseries The Atiyah--Singer Index Theorem}}
\author[弗里德]{Daniel S.~Freed（丹尼尔·S·弗里德）}
\thanks{本工作受美国国家科学基金会（NSF）资助，资助号 DMS-2005286。}
\address{德克萨斯大学奥斯汀分校数学系\\ Austin, TX 78712, USA}
\email{[邮箱隐去]}
\dedicatory{谨以此文纪念 Michael Atiyah（迈克尔·阿蒂亚）}
\date{2021 年 7 月 7 日}
\renewcommand{\abstractname}{摘要}
\begin{abstract}
阿蒂亚--辛格指数定理是 20 世纪 60 年代初的一项里程碑式成就，它把分析、几何与拓扑中的思想融为一炉。
我们追述若干前驱与动机，陈述这一定理的各种形式，并讨论它的若干推论——这些推论一直延伸到今天。
\end{abstract}
\maketitle



{\small
\makeatletter
\def\reftext{参考文献}
\renewcommand{\tocsection}[3]{%
  \begingroup
   \def\tmp{#3}%
   \ifx\tmp\reftext
  \indentlabel{\phantom{1}\;\;} #3%
  \else
  \indentlabel{\@ifempty{#2}{\hskip1.5em}{\ignorespaces#1 #2.\;\;}}#3%
  \fi\endgroup}
\makeatother
\tableofcontents
}



\section*{译者说明}

\begin{description}[leftmargin=4.2em,style=nextline,itemsep=3pt,topsep=2pt]
  \item[原文] Daniel S. Freed, \emph{The Atiyah--Singer Index Theorem},
    arXiv:2107.03557v1 [math.HO]（2021 年 7 月 8 日），52 页；为纪念 Michael Atiyah 而作。
    原文以 \texttt{amsart} 排版，本文沿用其版式（6.5\,in 版心、原编号体系）。
  \item[翻译原则] 以“无损”为唯一标准：不增删论证、不改动作者的行文逻辑与语气、不添加原文没有的物理解释。
    全部数学公式、定理环境、编号、引用键与原 LaTeX 源码逐字符一致；仅自然语言译为中文。
    原文若有笔误，一律照录原样，另以批注指出（见附录 A）。
  \item[结构] 每一节开头附「本节重点（Key Takeaways）」，为译者所加阅读导引（原文无此内容）。
    正文中形如 \texttt{\% [Tutor Note: \dots]} 的行为译者审校批注，属 LaTeX 注释，编译后不显示；
    为便于通读，全部批注另汇编于附录 A。
  \item[关于图表] 经逐页核对，原文全文\textbf{没有任何插图或图表}：52 页中 Figure 数为 $0$，
    全部“可视元素”仅为 127 个 \texttt{equation} 环境（与定理/注记共用编号计数器）与若干定理/注记环境。
    因此本译本不含图、也不存在图注翻译；此系原文事实，并非缺漏。
  \item[术语] 译名以中文学界通行译法为准（如 gauge theory/规范场论、holonomy/和乐、anomaly/反常），
    首次出现时附英文；核心术语对照见附录 B。
  \item[编译] 单文件、可直接编译：\texttt{xelatex}（或 \texttt{lualatex}）编译两遍即可；
    需要支持中文的 XeTeX 环境（xeCJK 与 SimSun 等系统字体），或使用 Tectonic。
    原文件亦可由 \texttt{pdflatex} 编译英文版。
\end{description}


% ==================== §1 ====================

\section{引言}\label{sec:1}

\begin{keytakeaways}
  \item \textbf{出发点}：在黎曼球面（Riemann sphere）$\CP^1$ 上给定有限个点
    $z_i$ 与非零整数 $n_i$，求以 $z_i$ 为零点（$n_i>0$）或极点（$n_i<0$）的
    亚纯函数；此时全部解只相差一个常数因子，构成一维复向量空间。
  \item \textbf{阻碍的出现}：换成亏格（genus）为正的闭黎曼曲面（Riemann
    surface）后，同样的计数问题出现阻碍。椭圆曲线的情形归结为约束
    $\sum n_iz_i\in \ZZ+\ZZ\tau $（且须 $\sum n_i=0$）；对一般曲面，约束构成维数为亏格 $g$ 的
    向量空间，而 $g$ 是拓扑不变量。
  \item \textbf{分析视角}：亚纯函数是柯西–黎曼（Cauchy–Riemann）方程的解，
    后者是线性椭圆偏微分方程；这一方程的解与阻碍由拓扑“计数”，这正是
    经典的黎曼–罗赫（Riemann–Roch）定理。
  \item \textbf{主要定理}：把它推广到任意维紧致流形上的任意椭圆算子，即
    阿蒂亚–辛格（Atiyah–Singer）指数定理（1962–3）；这一定理用算子与底
    流形的拓扑不变量来计算 Fredholm 指数 $\dim\ker-\dim\coker$。
  \item \textbf{全文结构}：\S\ref{sec:2} 与 \S\ref{sec:3} 分别追溯代数几何/拓扑
    与分析两方面的渊源；\S\ref{sec:4} 给出基本定理及配边、$K$ 理论两种
    证明；\S\ref{sec:5} 讨论等变、参数族、带边等推广；
    \S\ref{sec:6}–\S\ref{sec:7} 转向热方程、局部指数定理与狄拉克（Dirac）
    算子的新几何不变量；\S\ref{sec:8} 转向物理中的反常。
  \item \textbf{写作立场}：本文是献给阿蒂亚的纪念之作，自然只聚焦于他的
    工作与影响的若干方面；即便如此，也只能触及这一丰富主题的表层。
\end{keytakeaways}

考虑黎曼球面 $\CP^1=\CC\cup\{\infty \}$。设 $\{z_i\}_{i=1}^N\subset \CP^1$
为有限集，并对每个 $i\in \{1,\dots ,N\}$ 给定一个非零整数 $n_i$。一个
经典问题要求给出一个亚纯函数 $f$，它在每个 $z_i$ 处有零点或极点。若
$n_i>0$，则 $z_i$ 是重数为 $n_i$ 的零点；若 $n_i<0$，则 $z_i$ 是阶为
$|n_i|$ 的极点。
% [Tutor Note: 易错——“解构成一维复向量空间”依赖于重数被精确指定：若要求 $z_i$ 恰为 $n_i$ 重零点（当 $n_i>0$）或恰为 $|n_i|$ 阶极点（当 $n_i<0$），则全部解只相差常数因子 $c$；但黎曼–罗赫定理的“至少”提法取 $D=\sum n_iz_i$、$\sect D=H^0(\CP^1,\sO(D))$，此时当 $\deg D=\sum n_i\ge0$ 有 $\dim\sect D=\deg D+1=1+\sum n_i$，一般并不等于 $1$（对照式 \eqref{eq:1}）。原文亦未声明 $\sum n_i=0$：若 $\sum n_i\neq0$，$f$ 还在 $\infty$ 处有零点或极点，解空间仍是一维（相差一个常数因子）。]
解答直截了当。即 $f(z)=c\prod(z-z_i)^{n_i}$ 是有理函数，在相差一个常数
因子 $c\in \CC\setminus \{0\}$ 的意义下唯一。换言之，若现在还允许
$f\equiv 0$，则解构成一个一维复向量空间。若把 $\CP^1$ 换成亏格为正的
闭黎曼曲面，则具有指定零点与极点的亚纯函数之存在性便会有阻碍。
% [Tutor Note: 易错——“唯一约束”一句漏掉了次数条件：按阿贝尔（Abel）定理，$E=\CC/(\ZZ+\ZZ\tau )$ 上的除子 $D=\sum n_iz_i$ 是主除子当且仅当 $\deg D=\sum n_i=0$ 且 $\sum n_iz_i\in \ZZ+\ZZ\tau $（后者即除子的阿贝尔–雅可比像在雅可比簇 $J(E)$ 中为零）。原文只写出后一条件；例如取 $n_1=1$、$z_1=0$（$0$ 表 $E$ 中的原点），此时 $\sum n_iz_i=0\in \ZZ+\ZZ\tau $ 但 $\deg D=1\neq0$，并不存在以 $0$ 为唯一零点、别无极点的亚纯函数。对亏格为 $g$ 的一般曲面，“更多约束”正是阿贝尔定理的 $g$ 个周期条件。]
例如，椭圆曲线可实现为复直线由 $1,\tau $ 生成的满格所作的商
$E=\CC/(\ZZ+\ZZ\tau )$，其中 $\tau \in \CC$ 满足 $\Im\tau >0$。$E$ 上的
亚纯函数提升为 $\CC$ 上的双周期函数，而亚纯函数的零点与极点所受的唯一
约束是 $\sum n_iz_i\in \ZZ+\ZZ\tau $。从 $E$ 推进到一般的闭黎曼曲面
$X$，我们会遇到更多约束。
% [Tutor Note: 含糊——“约束构成维数为亏格的向量空间”的精确含义是：对每个全纯微分 $\omega\in H^0(X,K_X)$ 都有周期条件 $\sum_i n_i\int_{z_0}^{z_i}\omega\equiv0$（模周期），共 $g=\dim H^0(X,K_X)=\dim\sect K$ 个独立复条件（由式 \eqref{eq:1} 取 $D=\sO$ 得 $\dim\sect K=g$，见 \S\ref{subsec:2.1}）。故“约束空间”指全纯微分空间的对偶 $H^0(X,K_X)^*$，它是 $g$ 维复向量空间；条件本身的取值落在 $g$ 维复环面（雅可比簇，Jacobian）$J(X)=H^0(X,K_X)^*/H_1(X;\ZZ)$ 中，固定度数的除子类构成 $J(X)$ 上的一个主齐性空间（torsor），并非线性空间。可查 [H1]。]
事实上，这些约束构成一个向量空间，其维数为 $X$ 的亏格，而亏格是一个
拓扑不变量。亚纯函数是柯西–黎曼方程的解，而后者是一个线性椭圆偏微分
方程。
% [Tutor Note: 含糊——“由拓扑来‘计数’”的实质是 Fredholm 指数：$\dbar$ 作用在线丛 $\sO(D)$ 的光滑截面上是椭圆算子，其指数 $\dim\ker\dbar-\dim\coker\dbar=\dim\sect D-\dim\sect{K-D}$ 恰为黎曼–罗赫公式 \eqref{eq:1} 的左端，只依赖 $\deg D$ 与 $g$ 这两个拓扑量；而单个空间 $\dim\sect D$ 的维数依赖于复结构。指数定理正是把“两个维数之差可以拓扑地计算”这一现象推广到一般椭圆算子；详见 \S\ref{subsec:2.1} 与 \S\ref{subsec:3.2}。]
这一椭圆偏微分方程的解与阻碍由拓扑来“计数”。在更一般的形式下，这就是
经典的黎曼–罗赫定理。

% [Tutor Note: 含糊——“于 1962–3 年提出并证明”是压缩的叙述：据本文 \S\ref{sec:2} 末尾的回忆，阿蒂亚与辛格的合作始于 1962 年 1 月（辛格访牛津），当年 3 月二人已得到狄拉克算子与指标公式，约九个月后完成第一证明；定理于 1963 年在 [AS1] 中宣布，第一证明（配边路线，仿照希策布鲁赫的号差与黎曼–罗赫证明）未由二人撰写，而发表于帕莱（Palais）主编的会议录 [Pa]；第二证明（$K$ 理论路线，仿照格罗滕迪克）于 1968 年发表于 [AS2]。此外还有以伪微分算子与热方程为工具的解析路线（[Se2,Se3]，本文 \S\ref{sec:6}，帕托迪（Patodi）、吉尔基（Gilkey））。按发表年份引用时，宜区分“1963 年宣布”与“1968 年发表证明”。]
阿蒂亚–辛格指数定理于 1962–3 年提出并证明，它是对任意维紧致流形上任意
椭圆算子的巨大推广。
% [Tutor Note: 术语——Fredholm 指数定义为 $\dim\ker-\dim\coker$，但“有限维”并不是定义的一部分：闭流形上的椭圆算子经椭圆正则性（Sobolev 空间上的 Fredholm 性质）其核与余核才有限维，且与光滑截面上的核、余核一致；标准出处见本文 \S\ref{subsec:3.2} 以及 [Se2,Se3]、[Ho1]。余核须在适当的 Sobolev（或光滑）范畴中取，才能由正则性保证其有限维性。]
这里涉及的 Fredholm 指数是线性椭圆算子的核的维数减去其余核的维数。
阿蒂亚–辛格定理用算子的拓扑不变量与底流形的拓扑不变量来计算该指数。
该定理把代数拓扑、代数几何、微分几何与线性分析中的概念和结果交织在
一起；其影响远远超出这些领域，延伸至数论、表示论、算子代数、非线性
分析与理论物理。此外，指数理论本身就是一个不断扩张的庞大事业。基本的
阿蒂亚–辛格定理催生了众多推广与新的路径。本文——一篇献给迈克尔·阿蒂亚
的纪念文章——自然聚焦于他的工作及其影响的若干方面。即便作了这样的限制，
我们也只能触及这个丰富故事的表层。

指数定理的渊源，一方面来自代数几何与拓扑，另一方面来自分析。我们将
依次在 \S\ref{sec:2} 与 \S\ref{sec:3} 中讨论这两方面。基本的阿蒂亚–辛格
定理是 \S\ref{sec:4} 的主题。
% [Tutor Note: 术语——cobordism 与 bordism 均译“配边”：本文 \S\ref{sec:2} 的脚注指出托姆（Thom）、希策布鲁赫等人用“cobordism”代替“bordism”，阿蒂亚 [A10] 澄清了二者的关系，故文中两种拼法并存，不必据拼法区分理论。又，希策布鲁赫的“号差定理”即 signature theorem（号差 = signature，切勿与“符号”symbol 混同），它与黎曼–罗赫定理都是指数定理的特例。]
第一个证明基于配边（cobordism），其大要遵循希策布鲁赫（Hirzebruch）对
其号差定理与黎曼–罗赫定理的证明。
% [Tutor Note: 含糊——“受格罗滕迪克的黎曼–罗赫定理启发”的实质：格罗滕迪克把黎曼–罗赫定理推广为相对版本（对复流形间的 proper 映射，见 \S\ref{subsec:2.2} 与 [H1]），其 $K$ 理论形式用 $f_!$ 与 Todd 类、陈特征的关系表达；阿蒂亚–辛格的第二证明把这一相对形式搬到微分拓扑中，在 $K^0(T^*X,T^*X\setminus0)$ 上使用 Thom 同构与差丛（difference bundle）构造，见 \S\ref{subsec:4.4} 与 [AS2]。]
第二个证明基于 $K$ 理论；它受到格罗滕迪克（Grothendieck）的黎曼–罗赫
定理的启发。在 \S\ref{sec:5} 中，我们讨论基本定理的一些推广与变体。
这些包括等变指数定理、算子参数族的指数定理、带边流形的指数定理，以及
另外几项。至此，我们的论述从一般线性椭圆算子的整体拓扑不变量转向狄拉克
算子的局部几何不变量。
% [Tutor Note: 术语——“局部指数定理”（local index theorem）指热核 $e^{-t\Delta }$ 当 $t\to0$ 时的对角渐近展开给出指数密度，其积分等于指标；它与“整体拓扑不变量”的形式相对，但“局部”仅指公式的逐点形式，不表示指标本身可以局部化。本文 \S\ref{sec:6} 指出该定理的特殊情形由帕托迪、一般情形由吉尔基在 1970 年代初证明。]
热方程方法是 \S\ref{sec:6} 的主题，其第一个应用是局部指数定理。狄拉克
算子的新几何不变量出现在 \S\ref{sec:7} 中。在 \S\ref{sec:8} 中，我们转向
物理学：1980 年代中期之后，物理学成为阿蒂亚的关注重心，而它也为围绕
指数定理的那一圈思想提供了一个始料未及的试验场。我们聚焦于量子理论中
的反常，阿蒂亚与辛格曾对这一主题作出早期贡献。

本文的每一节都有更多的导论性材料，我们建议即使只是随手翻阅的读者也不
妨一读。
% [Tutor Note: 术语——[F1] 是 Freed 本人 2021 年 4 月的同名讲座视频（文献表中给出链接），并非论文；正文“a lecture based on this paper may be viewed at [F1]”即指此条目，引用时不要把它当作印刷文献。]
此外，以本文为基础的一场讲座可在 \cite{F1} 处观看。

迈克尔具有非凡的数学魅力与人格魅力。他的文字捕捉到了他的活力，他的
讲座亦然，其中一些讲座可在网上观看。他写了许多关于指数定理的精彩综述
文章，尤其是早期阶段的文章；读者会乐于细读它们。

我热忱感谢唐纳森（Simon Donaldson）、里德（Charlie Reid）与西格尔
（Graeme Segal）细致审读本文的早期版本并提出意见。

% ==================== §2 ====================

\section{来自代数几何与拓扑的先驱与动机}\label{sec:2}

% [Tutor Note: 易错——源文件的小节标号与出现顺序不一致：按出现顺序依次是 \label{subsec:2.1}（黎曼–罗赫）、\label{subsec:2.3}（希策布鲁赫）、\label{subsec:2.2}（格罗滕迪克）、\label{subsec:2.4}（整性定理），而正文第 428、431 行的 \S\ref{subsec:2.3}、\S\ref{subsec:2.2} 正是按这套标号引用的。译文保留原标号、不按顺序重编，否则交叉引用会指向错误小节；排版时编号由出现顺序自动给出，仍是 2.1–2.4。]

\begin{keytakeaways}
  \item 本节是索引定理的“前史”：代数几何中黎曼–罗赫定理及其高维形式，
        拓扑中 Hirzebruch 号差定理、Grothendieck 的 $K$ 理论与相对黎曼–罗赫，
        再经 Atiyah–Hirzebruch 用 $K$ 理论给出 $\Ahat$ 整性的新证明与纯拓扑解释。
  \item 关键工具：层上同调与 Euler 示性数、Todd 类与陈特征、$L$-类与
        $\Ahat$-类，以及带 Thom 类的 $K$ 理论（“shriek”映射 $f_!$）。
  \item 主要公式：$\chi(X,V)=\Todd(X)\ch(V)[X]$（\eqref{eq:7}）、
        $\Sign(X)=L(X)[X]$（\eqref{eq:10}）、Grothendieck–Riemann–Roch
        \eqref{eq:12}、Atiyah–Hirzebruch 的 \eqref{eq:21}，以及推论
        $\Ahat(X)[X]\in\ZZ$。
  \item 逻辑链条：示性数先验只是有理数，代数几何先保证它是整数；Grothendieck
        把黎曼–罗赫变成相对情形下的自然变换；Atiyah–Hirzebruch 再把它搬到
        拓扑，用 $K$ 理论 Thom 类给出整性的纯拓扑解释。
  \item 遗留问题：\eqref{eq:24}“整数 $\Ahat(X)[X]$ 究竟是什么？”——这正是
        Dirac 算子与索引定理的出发点。
\end{keytakeaways}

代数几何中的计数问题常常给出具有拓扑解释的整数。一个经典的例子是黎曼–罗赫
（Riemann–Roch）公式，它是我们在 \S\ref{subsec:2.1} 中的出发点。其高维推广由
弗里茨·希策布鲁赫（Fritz Hirzebruch）在 1950 年代初给出，我们在
\S\ref{subsec:2.3} 中叙述。几年之后，亚历山大·格罗滕迪克（Alexander
Grothendieck）把希策布鲁赫的定理推广到相对情形，即推广到复流形的逆紧映射。
在此过程中，他引入了层的 $K$ 理论。他的思想在 \S\ref{subsec:2.2} 中简要介绍，
十年之后在阿蒂亚–辛格（Atiyah–Singer）指数定理的各种变体中起了根本作用。

% [Tutor Note: 术语——"stable homotopy groups of Lie groups" 指经典群的稳定同伦群 $\pi_{n+k}(U(n))$、$\pi_{n+k}(O(n))$（$n\to\infty$ 的稳定范围），即 Bott 周期性所计算的对象，不是"每个李群自身的同伦群"。第 753–755 行明确说这里算的是酉群，两处指同一件事；Bott 的 {\it periodicity theorem} 正是由此得名。]

更直接地说，正如格雷姆·西格尔（Graeme Segal）在本卷中所写~\cite{Seg4}，
阿蒂亚与希策布鲁赫
把格罗滕迪克的 $K$ 理论搬到了代数拓扑中。拉乌尔·博特（Raoul Bott）对李群的
稳定同伦群的计算，与希策布鲁赫和格罗滕迪克关于黎曼–罗赫定理的工作发生在同一
时期，它是他们理论的基石。对指数定理至关重要的是由此得到的整性定理，我们在
\S\ref{subsec:2.4} 中提到其中几个。这引出了一个问题——为什么 $\Ahat$ 亏格对
spin 流形是整数？——它在 1962 年初成为阿蒂亚与辛格合作的直接催化剂。

\subsection{黎曼–罗赫定理}\label{subsec:2.1}

% [Tutor Note: 术语——divisor 是整系数形式和 $D=\sum_x n_xx$（$n_x=\ord_x(D)\in\ZZ$，只有有限多个非零），并非集合意义上的"点集"；原句"每点附一个整数"是简化说法。 $\deg D=\sum_x\ord_x(D)$。另外 $\sect D$ 就是 $\sL(D)=H^0(X,\sO(D))$：源文件的宏定义为 \sect{#1}=\sL(#1)（第 265 行），故第 488 行写成的 $\sL(K)$ 与 $\sect K$ 完全同义。]

设 $X$ 是 $\CC$ 上的光滑射影曲线，即复射影空间中的一维闭复子流形。
除子（divisor）$D$ 是 $X$ 上的一个有限点集，并在每个点 $x\in D$ 上附有一个
整数 $\ord_x(D)$。一个除子决定 $X$ 上的一个全纯线丛；用 $\sect D$ 表示该丛的
全纯截面空间。我们可以把 $\sect D$ 描述为 $X$ 上这样的亚纯函数空间：它在每个
$x\in X$ 处有阶数 $\le \ord_x(D)$ 的极点。曲线论中的一个基本问题是：计算
$\sect D$ 的维数。虽然这个问题一般是相当困难的，但对于
$\dim\sect D - \dim\sect{K-D}$ 存在一个{\it 拓扑的}公式，其中 $K$ 是 $X$ 的
一个\emph{典范除子}。（全纯 $1$-形式的零点集，按零点的阶数加权，就是一个
典范除子。）

% [Tutor Note: 易错——上句"至多有 $\ord_x(D)$ 阶极点"必须按整系数读：$\ord_x(D)>0$ 时是极点阶数的上界；$\ord_x(D)<0$ 时应读作"$x$ 处至少有 $-\ord_x(D)$ 阶零点"；而 $x\notin D$ 时 $\ord_x(D)=0$，即要求全纯。漏掉负系数情形就理解不了后面 $K-D$ 这类除子为何可能没有截面。]

\begin{theorem}[Riemann-Roch]\label{thm:2}
设 $X$ 是光滑射影曲线，$D$ 是 $X$ 上的除子。则
  \begin{equation}\label{eq:1}
     \dim\sect D - \dim\sect{K-D} = \deg(D) -g + 1.
  \end{equation}
\end{theorem}

% [Tutor Note: 易错——定理左端其实是 $\sO(D)$ 的层上同调 Euler 示性数：由 Serre 对偶 $H^1(X,\sO(D))^\vee\cong H^0(X,K-D)$，即 $\dim\sect{K-D}=h^1(D)$，故左端 $=\chi(X,\sO(D))$。符号要点：对偶写的是 $K-D$；若误写成 $K+D$，$h^0(K+D)$ 与 $h^1(D)$ 不再对偶，(1) 即不成立。]

\noindent
这里 $g$ 是曲线 $X$ 的亏格，即其基本的拓扑不变量，定义为
$\frac{1}{2}\rank H^1(X;\ZZ)$。又 $\deg(D)=\sum \ord_x(D)$ 是定义除子 $D$ 的
那些整数之和。若 $\deg(D) > 2g-2$，可以证明 $\sect{K-D}=0$，于是此时
\eqref{eq:1} 给出了计算 $\dim\sect D$ 这一问题的完整解答。
定理~\ref{thm:2} 就是经典的黎曼–罗赫\footnote{黎曼~\cite{Ri} 证明了不等式
$\dim\sect D \ge \deg(D) -g + 1$，随后罗赫~\cite{Ro} 证明了更精确的
\eqref{eq:1}。令人惋惜的是，罗赫在 26 岁时死于肺结核，距 39 岁的黎曼死于
肺结核仅数月。}公式。阿蒂亚–辛格指数定理是 \eqref{eq:1} 的一个巨大推广，
我们将看到这一点。

% [Tutor Note: 易错——$\deg(K)=2g-2$，故 $\deg(D)>2g-2\iff\deg(K-D)<0$；负次数除子没有非零全纯截面（$\deg E<0\Rightarrow H^0(E)=0$），这才是 $\sect{K-D}=0$ 的理由。原文用 "it can be shown" 带过了这一步。另外 $g=\frac12\rank H^1(X;\ZZ)$ 用的是自由秩（$\dim_\RR H^1=2g$），不要把 $g$ 与 $\dim\sect K$ 混为一谈：后者等于 $g$ 是下面用 $D=\sO$ 代入 (1) 推出的。]

让我们立刻来注意黎曼–罗赫公式的一个推论。取 $D=\sO$ 为不含任何点的平凡
除子。则 $\sect{\sO}$ 是常数函数空间，而 $\sect K$ 是全纯微分空间。由
\eqref{eq:1} 我们推出后者的维数为 $g$。由此可知 $g$ 是整数，即
$\rank H^1(X;\ZZ)$ 是偶数。因此，欧拉数 $\Euler(X)$ 的一半是整数，这是我们
遇到的第一个{\it 整性定理}。这个证明值得注意：$1-\Euler(X)/2$ 是整数，因为
它是某个向量空间的维数，即 $\sL(K)$ 的维数。

% [Tutor Note: 易错——这里的 $\Euler(X)$ 是实二维闭曲面的欧拉示性数 $2-2g$，故 $1-\Euler(X)/2=g$。整性链条是：$\dim\sect K=g\in\ZZ$（层论/Serre 对偶）$\Rightarrow$ $\rank H^1=2g$ 为偶数 $\Rightarrow$ $\Euler/2\in\ZZ$。三步实为同一组等式，核对时注意不要把这个逻辑顺序读反。]

在 19 世纪的最后十年，诺特（Noether）、恩里克斯
（Enriques）与卡斯泰尔诺沃（Castelnuovo）把黎曼–罗赫不等式与等式推广到
代数曲面；见下文 \eqref{eq:4}。

% [Tutor Note: 含糊——曲面情形的推广（即 \eqref{eq:4} 的 Noether 公式）是关于 Euler 示性数 $\chi(\sO_X)$ 的等式，而不是逐项 $h^0$ 的公式；19 世纪末这三人的结果要到 Hirzebruch 的 \eqref{eq:7} 才被统一为"示性类积分"的形式。]

\subsection{希策布鲁赫的黎曼–罗赫定理与号差定理}\label{subsec:2.3}

我们远远跳到 1945–1954 年，跳到年轻的希策布鲁赫的工作；它基于几何中的两项
重要发展。第一项由勒雷（Leray）开创，是层论。第二项是托姆（Thom）学位论文中
的结果，特别是关于光滑流形配边群\footnote{托姆、希策布鲁赫以及许多其他人用
`cobordism' 代替 `bordism'；阿蒂亚~\cite{A10} 澄清了两者的关系。}的那些结果。
我们陈述希策布鲁赫的两个主要结果，它们在~\cite{H1} 中有叙述。

% [Tutor Note: 术语——这条脚注讲的是术语源流：老文献（Thom、Hirzebruch 等）把 "bordism" 说成 "cobordism"。不同作者用法不一：bordism 多指流形在配边关系下的等价类与配边群 $\Omega_*^{SO}$，cobordism 有时指具体的配边 $W$（$\partial W=M\sqcup N$），有时也指同一等价关系。 Atiyah~\cite{A10} 澄清了这一用法。本节其余各处按中文惯例一律译作"配边"。]

设 $X$ 是复维数为 $n$ 的非异射影簇，$V\to X$ 是全纯向量丛。（在曲线情形
我们用的是除子；回忆一下，一个除子决定一个全纯线丛，这就联系到我们这里的
表述。）于是上同调群 $H^q(X,V)$ 通过层论定义：$H^0(X,V)$ 是 $V\to X$ 的全纯
截面空间，而对 $q\ge1$，$H^q(X,V)$ 由 $V\to X$ 的全纯截面层的消解导出。
这些上同调群是有限维的，这可以用椭圆微分算子理论与多尔博（Dolbeault）定理
证明。（见 \S\S\ref{subsec:3.1}--\ref{subsec:3.2}。）{\it 欧拉示性数}定义为
交错和
  \begin{equation}\label{eq:3}
     \chi (X,V) = \sum\limits_{q=0}^n\, (-1)^q \dim H^q(X,V).
  \end{equation}
如同黎曼曲面的情形 $n=1$ 一样，人们常常想计算 $\dim H^0(X,V)$，但一般来说
$\dim H^0(X,V)$ 所依赖的不止是拓扑数据。另一方面，欧拉示性数 $\chi (X,V)$
确实有一个用陈类 $c_j(X)$ 与 $c_k(V)$ 表示的拓扑公式。特殊情形
$\dim X = \rank V = 1$ 就是经典的黎曼–罗赫公式 \eqref{eq:1}。当 $X$ 是光滑
射影代数曲面（$n=2$）而 $V\to X$ 是秩为 $1$ 的平凡丛时，所得结果通常称为
诺特公式：
  \begin{equation}\label{eq:4}
     \chi (X) = \frac{1}{12}\bigl( c_1^2(X) + c_2(X)\bigr) [X].
  \end{equation}
在 \eqref{eq:4} 中，陈类在由自然定向给出的 $X$ 的基本类上取值。分母中出现的
$12$ 给出了射影曲面的陈数的一个整性定理。

% [Tutor Note: 含糊——"可以用椭圆算子理论与 Dolbeault 定理证明"省略的实质是 Hodge 理论：$H^q(X,V)\cong H^{0,q}_{\bar\partial}(X,V)\cong\mathcal H^{0,q}(X,V)$ （调和 $(0,q)$-形式空间），而椭圆算子 $\dbar+\dbar^*$ 的核有限维（详见 \S\S\ref{subsec:3.1}--\ref{subsec:3.2}）。原文一句话带过了整个 Hodge–Dolbeault 同构，需要时可查 [H1] 与本文 \S\ref{subsec:3.1}。]

% [Tutor Note: 术语——$\chi(X,V)$ 是层上同调的 Euler 示性数（交错和），不是流形的 Euler 数 $\Euler(X)$：例如 $n=1$、$V=\sO$ 时 $\chi(X,\sO)=1-g$，而 $\Euler(X)=2-2g$。\eqref{eq:4} 中的 $\chi(X)$ 是 $\chi(X,\sO_X)$ 的简写，尽管字母相同，含义与 $\Euler$ 不同。]

% [Tutor Note: 易错——\eqref{eq:4} 是 Noether 公式 $\chi(\sO_X)=\frac1{12}(c_1^2+c_2)[X]$，其中 $c_1^2$ 表示 $c_1(X)\smile c_1(X)$，$c_2[X]=\Euler(X)$。由 $p_1=c_1^2-2c_2$、 $\Sign=p_1/3$ 可把它改写成 $\chi=(\Sign+\Euler)/4$。用 $K3$ 曲面核对：$\Sign=-16$、 $\Euler=24$，得 $\chi=(-16+24)/4=2$，与 $(0+24)/12=2$ 一致。若把分母 12 记成 24，就会与 $\Ahat$ 的 $-p_1/24$ 混淆。]

对所有 $X,V$ 解出黎曼–罗赫问题——即计算 \eqref{eq:3}——是希策布鲁赫最显著的
成就之一。希策布鲁赫的公式用 Todd 多项式与陈特征表出。设切丛
$TX=L_1\oplus \dots \oplus L_n$ 分裂为线丛之和，并令
$y_i = c_1(L_i) \in H^2(X;\ZZ)$。则{\it Todd 类}为
  \begin{equation}\label{eq:5}
     \Todd(X) = \prod_{i=1}^n \frac{y_i}{1-e^{-y_i}}.
  \end{equation}
这是一个（混合的）偶次上同调类。类似地，若 $V=K_1\oplus \dots \oplus K_r$
是线丛之和，$x_i = c_1(K_i)$，则{\it 陈特征}为
  \begin{equation}\label{eq:6}
     \ch(V) = \sum\limits_{i=1}^r \,e^{x_i}.
  \end{equation}
示性类理论中的分裂原理允许我们把这两个定义推广到并非线丛之和的 $TX\to X$
与 $V\to X$。

% [Tutor Note: 易错——Todd 类的逐因子展开为 $\frac{y}{1-e^{-y}}=1+\frac y2+\frac{y^2}{12}+\cdots$，故 $\Todd(X)=1+\frac{c_1}{2}+\frac{c_1^2+c_2}{12}+\cdots$。本节出现三个互不相同的归一化，务必分清：Todd 的 4 维项 $\frac1{12}(c_1^2+c_2)$（见 \eqref{eq:4}）、 $L$-类的 4 维项 $\frac{p_1}{3}$（见 \eqref{eq:9}）、$\Ahat$-类的 4 维项 $-\frac{p_1}{24}$（见 \eqref{eq:22}）。$\frac12$ 因子（$c_1/2$）来自 $\frac{y}{1-e^{-y}}$ 与 $\Ahat$ 的换算，不是笔误。]

% [Tutor Note: 含糊——"分裂原理"的实质：存在旗流形 $\pi:F(X)\to X$ 使 $\pi^*TX$ 分裂为线丛之和，且 $\pi^*:H^*(X)\hookrightarrow H^*(F(X))$ 为单射；于是先在分裂情形定义示性类/示性数，再证明结果与所取分裂无关。$\ch$ 的加性（$\ch(V\oplus W)=\ch V+\ch W$）与 $\Todd$ 对直和的乘性正是在此用到。]

\begin{theorem}[Hirzebruch-Riemann-Roch]\label{thm:1}
设 $X$ 是射影复流形，$V\to X$ 是全纯向量丛。则
  \begin{equation}\label{eq:7}
     \chi (X,V) = \Todd(X)\ch(V)[X].
  \end{equation}
\end{theorem}

希策布鲁赫的第二个主要定理是定理~\ref{thm:1} 证明中的一步，现在称为希策
布鲁赫号差定理。设 $X$ 是维数为 $4k$ 的闭定向实可微流形，其中 $k$ 为某个
正整数。则中间上同调 $H^{2k}(X;\RR)$ 上有一个非退化对称双线性配对，它由
杯积再在基本类上取值给出：
  \begin{equation}\label{eq:8}
     \begin{aligned} H^{2k}(X;\RR) \otimes H^{2k}(X;\RR) &\longrightarrow
     \quad\;\;\,\RR\\ \alpha _1 \quad \;\; \otimes \quad\;\; \alpha_2 \quad\quad
     &\longmapsto 
     (\alpha _1 
     \smallsmile \alpha_2 )[X]\end{aligned}
  \end{equation}
该配对的号差 $\Sign(X)$ 称为 $X$ 的号差。（在较老的文献中用 `index' 代替
`signature'。）希策布鲁赫把 $L$-类定义为由形式表达式
  \begin{equation}\label{eq:9}
     L(X) = \prod_{i=1}^{2k} \frac{y_i}{\tanh y_i},
  \end{equation}
确定的 $X$ 的 Pontryagin 类的多项式，其中 $y_i, -y_i$ 是复化切丛的\emph{陈根}
（Chern roots）。\footnote{总 Pontryagin 类
$p(X) = 1 + p_1(X) + p_2(X)+\cdots $ 由表达式 $\prod(1 + y_i^2)$ 定义。}这
类似于 \eqref{eq:5}：人们首先在 $TX\otimes \CC\to X$ 分裂为复线丛之和的情形
定义 $L(X)$。

% [Tutor Note: 易错——中间杯积配对要在 $H^{2k}(X;\RR)$（实系数）上谈对称性与非退化性；因 $2k$ 是偶数，$\alpha_1\smile\alpha_2=\alpha_2\smile\alpha_1$， $\Sign$ 才有定义。若 $\dim X\equiv2\pmod4$，中间配对是交错的，"号差"恒为 0，故号差定理只对 $4k$ 维陈述；原文的 $4k$ 不可省。]

% [Tutor Note: 术语——原文特意注明较老文献把 signature 叫作 "index"，这正是 Hirzebruch 号差定理旧称 "index theorem" 的由来；它与第 4 节起讨论的 Fredholm 指数 $\ind D$ 是两个概念，二者要等到指标定理才统一（号差等于号差算子的指标）。]

% [Tutor Note: 易错——\eqref{eq:9} 的乘积上限是 $2k$：实 $4k$ 维流形的复化切丛 $TX\otimes\CC$ 分裂为 $4k$ 个复线丛，其 Chern 根成对出现 $y_i,-y_i$ （$i=1,\dots,2k$），$p_1=\sum y_i^2$。展开得 $L(X)=1+\frac{p_1}{3} +\frac{7p_2-p_1^2}{45}+\cdots$，故 $\dim X=4$ 时 $\Sign=p_1/3$。若上限误取 $k$ 或漏掉根的另一半，度数与符号都会错。]

\begin{theorem}[Hirzebruch Signature Theorem]\label{thm:3}
闭定向光滑流形 $X$ 的号差为
  \begin{equation}\label{eq:10}
     \Sign(X) = L(X)[X].
  \end{equation}
\end{theorem}

\noindent
希策布鲁赫的证明本质地使用了托姆的配边理论~\cite{T1}。\eqref{eq:10} 的两边
都在定向配边下不变，并且都是乘性的；其中“号差在定向配边下不变”是托姆的一个定理
~\cite[\S IV]{T2}。因此只需在（有理）定向配边环的一组生成元上验证
\eqref{eq:10}，而这组生成元已由托姆算出。偶射影空间 $\CP^{2n}$ 提供了一组
方便的生成元，证明以如下观察结束：$L$-类恰好由它在这些生成元上取值 $1$ 所
刻画。Todd 类以类似的方式进入定理~\ref{thm:1} 的证明——它在所有射影空间
$\CP^n$ 上的值都是 $1$，并且由这一性质刻画。

% [Tutor Note: 含糊——"只需在生成元上验证"省略了两件事：(i) Thom 定理给出 $\Omega_*^{SO}\otimes\QQ=\QQ[\CP^2,\CP^4,\dots]$，故偶射影空间是有理生成元（\cite{T1}）；(ii) 号差在定向配边下不变是 Thom 的定理 [T2，\S IV]，而号差对乘积是乘性的（$\Sign(M\times N)=\Sign M\cdot\Sign N$，由 Künneth 公式与中间次数配对得到）。两条合起来，才能从生成元推广到所有闭定向流形。]

\subsection{格罗滕迪克的黎曼–罗赫定理}\label{subsec:2.2}

黎曼–罗赫–希策布鲁赫定理在 1957 年被格罗滕迪克~\cite{BS} 沿一个新的方向推广。
决定性的一步是格罗滕迪克在代数几何中引入 $K$ 理论。设 $X$ 是光滑代数簇。
把 $K(X)$ 定义为由 $X$ 上的凝聚代数层生成的自由阿贝尔群，模去如下等价关系：
若有短正合列 $0 \to \scrF' \to \scrF \to \scrF'' \to 0$，则
$\scrF \sim \scrF' + \scrF''$。在此定义中可以把“凝聚代数层”换成“全纯向量丛”，
而一个基本结果是群 $K(X)$ 不变。于是陈类与陈特征对 $K(X)$ 的元素有定义。
（格罗滕迪克把它们精化为取值于 $X$ 的{\it 周环}（Chow ring）中的类。）若
$f\:X\to Y$ 是簇的态射，$\scrF$ 是 $X$ 上的层，则 $R^qf_*(\scrF)$ 是与预层
$U\mapsto H^q(f\inv (U),\scrF)$ 相伴的 $Y$ 上的层。对应
  \begin{equation}\label{eq:11}
     f_! \: \scrF\longmapsto \sum (-1)^q R^qf_*(\scrF) \in K(Y)
  \end{equation}
延拓为阿贝尔群的同态 $f_!\: K(X)\longrightarrow K(Y)$，这从层上同调的长正合列
可以看出。

% [Tutor Note: 含糊——把"凝聚代数层"换成"全纯向量丛"而 $K(X)$ 不变，依赖 $X$ 光滑：此时每个凝聚层有有限长度的局部自由（向量丛）消解，故两者的 Grothendieck 群一致（Serre；[BS] 即 Borel–Serre 的黎曼–罗赫论文）。在一般概形上这一替换不成立。]

% [Tutor Note: 术语——此处的 $K(X)$ 是凝聚层生成的自由阿贝尔群模去由短正合列生成的子群（Grothendieck 群），不是拓扑 $K$ 理论；下一段的 $f_!$ 是层论直接像的交错和。 \eqref{eq:11} 中的 $R^qf_*$ 是"高次直接像层"，不是同调群。]

% [Tutor Note: 含糊——$f_!$ 良定义除需要短正合列给出的长正合列（保证 $\sum(-1)^qR^qf_*$ 只依赖 $K$ 类）之外，还需 $R^qf_*=0$ 对 $q\gg0$ 成立，这要求 $f$ 真（下一段因此明确假设 $f$ proper）；非逆紧态射的高次直接像未必凝聚，$f_!$ 未必落在 $K(Y)$ 中。]

现在设 $f\:X\to Y$ 是非异不可约拟射影簇之间的{\it 真}态射。上同调（或周环）
中存在前推 $f_*$。

\begin{theorem}[Grothendieck-Riemann-Roch]\label{thm:4}
对 $z\in K(X)$ 有
  \begin{equation}\label{eq:12}
     \Todd(Y)\ch\bigl(f_!(z)\bigr)  = f_*\bigl( \Todd(X)\ch(z)\bigr).
  \end{equation}
\end{theorem}

\noindent
取 $Y$ 为一点、$z$ 为全纯向量丛的 $K$ 理论类，它就约化为希策布鲁赫的
定理~\ref{thm:1}。

通向 Todd 类的一条途径是如下特殊情形：$f\:X\to Y$ 是除子的包含，$z$ 是结构层
$\sO_X$ 的类。此时对 $q\ge 1$ 有 $R^qf_*(\sO_X)=0$，而 $R^0f_*(\sO_X)$ 是把
$\sO_X$ 沿零延拓到 $Y$ 所得的层。设 $L\to Y$ 是该除子 $X$ 所定义的线丛。
注意 $f^*(L)\to X$ 是 $X$ 在 $Y$ 中的法丛。层正合列
  \begin{equation}\label{eq:13}
     0 \longrightarrow L\inv \longrightarrow \sO_Y \longrightarrow
     f_!\sO_X \longrightarrow 0
  \end{equation}
导出等式
  \begin{equation}\label{eq:14}
     f_!\sO_X = \sO_Y - L\inv
  \end{equation}
在 $K(Y)$ 中。令 $y=c_1(L)$。则由 \eqref{eq:14}，
  \begin{equation}\label{eq:15}
     \ch\bigl( f_!(\sO_X)\bigr) = 1 - e^{-y}.
  \end{equation}
另一方面
  \begin{equation}\label{eq:16}
     f_*\bigl(\ch(\sO_X)\bigr) = f_*(1) = y.
  \end{equation}
于是 $f_*\circ \ch = \ch \circ f_!$ 至多相差 $L$ 的 Todd 类。为在此情形验证
定理~\ref{thm:4}，利用 $X$ 上向量丛的正合列
  \begin{equation}\label{eq:127}
     0\longrightarrow TX\longrightarrow f^*TY\longrightarrow
     f^*L\longrightarrow 0
  \end{equation}
以及 Todd 类的乘性来改写 \eqref{eq:12}：
  \begin{equation}\label{eq:128}
     \begin{split} \ch\bigl(f_!(z) \bigr) &= f_*\left(
      \frac{\Todd(X)}{f^*\Todd(Y)}\,\ch(z) \right) \\ &= f_*\left(
      \frac{1}{f^*\Todd(L)}\,\ch(z) \right) \\ &=
      \frac{1}{\Todd(L)}\,f_*\bigl(\ch(z) \bigr).\end{split}
  \end{equation}
这正是我们对 $z=[\sO_X]$ 在 \eqref{eq:15} 与 \eqref{eq:16} 中所验证的。

% [Tutor Note: 易错——\eqref{eq:13} 中 $L=\sO_Y(X)$，故 $L\inv=\sO_Y(-X)$， \eqref{eq:14} 给出 $f_!\sO_X=\sO_Y-L\inv$；若把 $L\inv$ 写成 $L$，符号就反了。 \eqref{eq:16} 的 $f_*(1)=y$ 是 Gysin 前推：$f_*:H^0(X)\to H^2(Y)$ 把 $1$ 送到法丛的 Thom 类 $c_1(L)=y\in H^2(Y)$（Euler 类 $f^*c_1(L)\in H^2(X)$ 则是它的拉回）。把 \eqref{eq:15} 与 \eqref{eq:16} 对比可见 $f_*\circ\ch\ne\ch\circ f_!$，相差的因子恰是 $\Todd(L)=\frac{y}{1-e^{-y}}$，即 \eqref{eq:128} 末行。]

% [Tutor Note: 笔误——\eqref{eq:127}、\eqref{eq:128} 的标签命名是源文件修订时的旧编号残留（label 名与排版编号无关，不影响引用，故译文照录原标号）。另注意 \eqref{eq:128} 第一步是把 \eqref{eq:12} 除以 $Y$ 上的 $\Todd(Y)$：$\Todd(Y)$ 的常数项为 1、故可逆， $\frac{\Todd(X)}{f^*\Todd(Y)}$ 中分母是拉回到 $X$ 的类。]

在此阶段，考虑秩为 $k$ 的向量丛 $\pi \:E\to X$ 中零截面的包含 $f\:X\to E$
是有启发性的。此时层 $f_!\sO_X=R^0f_*\sO_X $ 位于 $E$ 上的层正合列
  \begin{equation}\label{eq:17}
     0 \longrightarrow \pi ^*{\textstyle\bigwedge} ^kE^* \longrightarrow \pi
     ^*{\textstyle\bigwedge} ^{k-1}E^* \longrightarrow \cdots \longrightarrow \pi
     ^*E^* \longrightarrow \sO_E \longrightarrow f_!\sO_X \longrightarrow 0
  \end{equation}
中。（比较 \eqref{eq:13}。）这里 $E^*$ 是 $\pi \:E\to X$ 的对偶丛（的截面层），
而 \eqref{eq:17} 中在 $e\in E$ 处的箭头是与 $e$ 作缩并。于是在 $K(E)$ 中有
  \begin{equation}\label{eq:18}
     f_!(\sO_X) = {\textstyle\bigwedge} ^{\bullet }(E^*),
  \end{equation}
其中
${\textstyle\bigwedge} ^{\bullet }(E^*)=\sum\limits_{}(-1)^k{\textstyle\bigwedge} ^kE^*$
是 $K$ 理论中的和。注意 $\pi \:E\to X$ 是 $X$ 在 $E$ 中的法丛。

% [Tutor Note: 易错——\eqref{eq:17} 是 Koszul 复形（比较 \eqref{eq:13} 的除子情形），映射为与 $e$ 的缩并 $\iota(e)$，$k=\rank E$，复形在 $e\ne0$ 处正合，故其 $K$ 类以零截面 $X$ 为支撑。\eqref{eq:18} 的 $\wedge^\bullet(E^*)$ 是交错和 $\sum_{k=0}^{\rank E}(-1)^k\wedge^kE^*$；原文 $\sum\limits_{}$ 的求和下标留空，是省略记法而非笔误。漏掉 $(-1)^k$ 得到的是和 $\sum_k[\wedge^kE^*]$（其陈特征为 $\prod_i(1+e^{-y_i})$），不再是 Thom 类。]

\subsection{拓扑中的整性定理}\label{subsec:2.4}

% [Tutor Note: 术语——genus 在本节有两种用法：(i) 曲线/曲面的亏格 $g$（拓扑不变量）；(ii) 示性类意义下的"亏格"，即上同调类在基本类上取值的数。第 710–711 行的 "Todd 亏格 $\Todd(X)[X]$" 属 (ii)：$\Todd(X)$ 是类（Todd 类），$\Todd(X)[X]$ 才是数（Todd 亏格）；同理有 $\Ahat$-类与 $\Ahat$ 亏格、$L$-类与 $L$ 亏格。注意 $\Ahat$ 亏格与 Todd 亏格是两个不同的不变量，仅在近复流形上由 \eqref{eq:23} 相联系。]

黎曼–罗赫–希策布鲁赫定理~\ref{thm:1} 的一个推论是：\eqref{eq:7} 右端的示性数
{\it 先验地}是有理数，实际上却是整数。这个整数被左端识别为向量空间维数之和与差。
另一方面，右端对任何近复流形都有定义。希策布鲁赫因此被引向这样一个问题
（早在 1954 年）：近复流形（更不用说非代数复流形）的{\it Todd 亏格}
$\Todd(X)[X]$ 是否为整数~\cite{H3}。他还对实流形提出了类似的问题。把实流形
$X^{4k}$ 的{\it $\Ahat$-类}\footnote{希策布鲁赫此前定义过一个 $A$-类，它与
$\Ahat$-类相差一个 $2$ 的幂，$\Ahat$ 的记号即由此而来。}定义为形式表达式
  \begin{equation}\label{eq:22}
     \Ahat(X) = \prod_{i=1}^{2k} \frac{y_i/2}{\sinh y_i/2},
  \end{equation}
其中 $y_i$ 是陈根。这是 Pontryagin 类的多项式。于是近复流形的 Todd 类可表为
  \begin{equation}\label{eq:23}
     \Todd(X) = e^{c_1(X)/2} \Ahat(X).
  \end{equation}
特别地，$\Todd(X)$ 只依赖于 Pontryagin 类与第一陈类。有理由推测，正是
\eqref{eq:23} 促使希策布鲁赫引入 $\Ahat$-类。此外，由于第二 Stiefel–Whitney
类 $w_2$ 是 $c_1$ 的模 $2$ 约化，希策布鲁赫问道：若实流形 $X^{4k}$ 满足
$w_2(X)=0$，即 $X$ 是 spin 流形，那么 $\Ahat(X)[X]$ 是否为整数？\footnote{在
~\cite{H3} 中希策布鲁赫只问了一个不那么精确的整除性问题（该文的问题~7）。更
精确的形式是后来才出现的，同时出现的还有更一般的问题：若闭实流形 $X$ 有一个
元素 $c\in H^2(X;\ZZ)$，其模 $2$ 约化为 $w_2(X)$，而 $\Todd(X)$ 由
\eqref{eq:23} 定义（以 $c$ 代替 $c_1(X)$），则 $\Todd(X)[X]$ 是否为整数？}这一
问题由博雷尔（Borel）与希策布鲁赫~\cite{BH3} 在 1950 年代末利用米尔诺（Milnor）
关于配边的结果~\cite{Mi1} 证明为真（最初在~\cite{BH2} 中只证到相差一个 $2$ 的幂）。

% [Tutor Note: 易错——归一化核对：\eqref{eq:22} 的 $\Ahat(X)=\prod_{i=1}^{2k}\frac{y_i/2}{\sinh y_i/2}$ 展开为 $1-\frac{p_1}{24}+\cdots$。$K3$ 曲面：$p_1=-48$，故 $\Ahat(X)[X]=-p_1/24=2$（即 Dirac 算子的指数），而 $L(X)[X]=\Sign=p_1/3=-16$。对 spin 四维流形 $p_1=3\Sign$，于是 $\Ahat[X]=-\Sign/8$；Rokhlin 定理 $\Sign\equiv0\pmod{16}$ 等价于 $p_1\equiv0\pmod{48}$，正是它保证 $\Ahat[X]\in\ZZ$（且为偶数）。文献中若见到 $-p_1/48$ 之类系数，那是把 $p_1/48$ 的整除性误当成 $\Ahat$ 的归一化，或 $A$ 与 $\Ahat$ 旧记号混用所致；引用前须按 \eqref{eq:22} 逐项换算。]

% [Tutor Note: 含糊——"$w_2$ 是 $c_1$ 的模 2 约化"只对近复流形成立（此时 $w_2(X)\equiv c_1(X)\pmod2$）；一般实流形没有 $c_1$，故脚注把问题放宽为"存在 $c\in H^2(X;\ZZ)$ 使 $c\equiv w_2(X)\pmod2$"。这个 $c$ 正是后文 \eqref{eq:19} 中"定向"条件的一般形式，也是 \eqref{eq:23} 中 $e^{c/2}$ 的来源。]

% [Tutor Note: 易错——\eqref{eq:23} 的 $c_1/2$ 因子不是任意的：逐因子恒等式 $\frac{y}{1-e^{-y}}=e^{y/2}\frac{y/2}{\sinh(y/2)}$（因 $1-e^{-y}=2e^{-y/2}\sinh\frac y2$）说明 $e^{c_1/2}$ 恰好补上 Todd 的 $y_i$ 与 $\Ahat$ 的 $y_i/2$ 之间的差别；写成 $e^{c_1}$ 会差一个 $e^{c_1/2}$。校核到 4 次项： $e^{c_1/2}\Ahat=1+\frac{c_1}{2}+\frac{c_1^2+c_2}{12}$，与 \eqref{eq:5} 一致（用 $p_1=c_1^2-2c_2$）。]

整性既已证明，显而易见的问题便提了出来：
  \begin{equation}\label{eq:24}
     \textnormal{整数~$\Ahat(X)[X]$ 是什么？}
  \end{equation}

对这个问题的第一个回答来自代数拓扑自身内部，尽管不是来自传统的
Eilenberg–MacLane 上同调理论。当阿蒂亚与希策布鲁赫了解到格罗滕迪克的工作时，
他们立即着手研究它在拓扑中可能的推广。第一步是为任意 CW 复形 $X$ 定义
$K$ 理论~\cite{AH1}。定义方式与代数簇情形相同，只是用“拓扑向量丛”代替
“凝聚代数层”。拓扑的基本构件是球面，而 $K(S^n)$ 的计算很快归约为酉群的稳定
同伦群的计算。幸运的巧合是，博特刚刚（1957 年）算出了这些同伦群~\cite{B1},
~\cite{B2}。他的{\it 周期性定理}成为新的拓扑 $K$ 理论的基石。由此得到的是一个
上同调理论，它满足 Eilenberg–MacLane 公理中除维数公理之外的所有公理。于是
“广义上同调”诞生了。$K$ 理论正是本卷中西格尔文章的主题~\cite{Seg4}。

% [Tutor Note: 易错——$K(S^n)$ 的约化部分同构于 $\pi_{n-1}(U)$ 而非 $\pi_n(U)$： $\tilde K(S^n)\cong[S^n,BU\times\ZZ]\cong\pi_n(BU)\cong\pi_{n-1}(U)$。Bott 给出 $\pi_{n-1}(U)=\ZZ$（$n$ 偶）、$0$（$n$ 奇），周期性由此而来。原文"归约为酉群的稳定同伦群"不可读成"球面的同伦群由酉群给出"。]

% [Tutor Note: 术语——"extraordinary cohomology" 即广义（上）同调：满足 Eilenberg–MacLane 公理中除维数公理外全部公理的上同调理论（$K$ 理论、配边理论等）；中文文献亦译"异常上同调"。维数公理失效正是它有别于奇异上同调之处。]

回到格罗滕迪克的纲领，阿蒂亚与希策布鲁赫对光滑流形表述了一版黎曼–罗赫定理
~\cite{AH2}, ~\cite{H2}。设 $f\:X\to Y$ 是可微流形之间的光滑映射，并设 $f$
是“定向的”，意即存在元素 $c\in H^2(X;\ZZ)$ 使
  \begin{equation}\label{eq:19}
     c\equiv w_2(X) - f^*w_2(Y) \pmod2.
  \end{equation}
回忆格罗滕迪克的定理 \eqref{eq:12} 是用映射 $f_!\:K(X)\to K(Y)$ 表述的。在
拓扑范畴中，我们不能像代数范畴中推前层那样推前向量丛，因此需要新的构造。
\footnote{$f_!$ 的定义在原始论文~\cite{AH2} 中没有给出；所缺的是 $K$ 理论中的
Thom 类，而它与 Dirac 算子的象征密切相关。Dirac 算子在阿蒂亚与辛格的合作中
进入故事（\S\ref{subsec:4.1}），随后 $K$ 理论 Thom 类与 Thom 同构出现在
~\cite[\S12]{ABS} 中。另见~\cite[\S1]{Seg4} 中的讨论。}这里我们只限讨论复流形
的嵌入，以简化表述。\footnote{一般情形：把闭流形 $X$ 嵌入球面，从而把任意映射
$f\:X \to Y$ 分解为嵌入后接投影：复合 $X \to S^N\times Y \to Y$。Bott 周期性
计算“叹号映射”（shriek map）$K(S^N\times Y) \to K(Y)$。对于实流形的嵌入
（法丛带有定向），旋量上的 Clifford 乘法取代 \eqref{eq:20}。}于是
\eqref{eq:14} 与 \eqref{eq:18} 启发了 $f_!$ 的定义。设 $\pi \:E \to X$ 是
$X$ 在 $Y$ 中的法丛。由管状邻域定理，我们可以把 $E$ 与 $Y$ 中 $X$ 的邻域
$U$ 等同。{\it Thom 复形} ${\textstyle\bigwedge} ^{\bullet }E^*\to E$ 在 $E$
的全空间上由缩并定义（比较 \eqref{eq:17}）：
  \begin{equation}\label{eq:20}
     0 \longrightarrow \pi ^*{\textstyle\bigwedge} ^kE^*
     \xrightarrow{\;\;\iota (e)\;\;} \pi 
     ^*{\textstyle\bigwedge} ^{k-1}E^* \longrightarrow \cdots \longrightarrow
     \pi ^*E^* \xrightarrow{\;\;\iota (e)\;\;} E\times \CC \longrightarrow 0.
  \end{equation}
注意 \eqref{eq:20} 对 $e\not= 0$ 正合，因此所得的 $K$ 理论元素以 $X$ 为支撑。
由管状邻域定理它也在 $U$ 上有定义，而零延拓给出所需的元素 $f_!(1) \in K(Y)$。
若 $V \to X$ 是向量丛，则把 \eqref{eq:20} 与 $\pi ^*V$ 张量积便定义了
$f_!(V)$。

% [Tutor Note: 术语——这里的"定向"不是流形的定向（只需 $w_1=0$），而是 $K$ 理论 Thom 类存在所需的 $w_2$-定向：$c\equiv w_2(X)-f^*w_2(Y)\pmod2$。近复流形取 $c=c_1(X)$ 即满足（$w_2$ 是 $c_1$ 的模 2 约化）；spin 流形取 $c=0$。$c$ 的作用是在 \eqref{eq:23} 中用 $e^{c/2}$ 代替 $e^{c_1/2}$（见第 734–736 行脚注）。]

% [Tutor Note: 含糊——\eqref{eq:20} 就是 $K$ 理论 Thom 类的构造：$f_!(1)\in K(Y)$ 是法丛的 Thom 类，"零延拓"一步用到管状邻域定理，与 \eqref{eq:17} 的 Koszul 复形是同一件事（把 $E$ 上的复形限制/延拓）。张量 $\pi^*V$ 即得 $f_!(V)$。下文"证明是一个练习：比较 $K$ 理论与上同调的 Thom 同构"实质是要验证 $\ch(U_K)=\Todd(E)\smile U_H$（$U_K,U_H$ 分别为两种 Thom 类），细节见 [ABS，\S12]。]

阿蒂亚–希策布鲁赫关于光滑流形的黎曼–罗赫定理断言
  \begin{equation}\label{eq:21}
     \ch\bigl(f_!(z)\bigr) \Todd(Y) = f_*\bigl( \ch(z)\Todd(X)\bigr), \qquad
     z\in K(X).
  \end{equation}
$f_!$ 一旦定义好，证明就是一个练习：比较 $K$ 理论与上同调中的 Thom 同构。
现在特殊化到 $Y=\pt$，并设 $w_2(X)=0$。取定向类 $c\in H^2(X;\ZZ)$ 为零。
于是在 \eqref{eq:21} 中取 $z=0$，借助 \eqref{eq:23} 我们推出
$\Ahat(X)[X] = f_!(1)\in K(\pt)\cong \ZZ$ 是整数。阿蒂亚–希策布鲁赫的这一
论证给出了 $\Ahat$ 整性定理的新证明，也给出了整数 $\Ahat(X)[X]$ 的一个拓扑
解释，因而是对 \eqref{eq:24} 的第一个回答。

% [Tutor Note: 笔误——第 812 行 "for $z=0$ in~\eqref{eq:21}" 疑应为 $z=1$ （$K(X)$ 的单位元）：$z=0$ 时 \eqref{eq:21} 两端都是 0，推不出任何整性；只有 $z=1$ 时左端才给出 $\ch(f_!(1))$，再由 $Y=\pt$、$c=0$ 及 \eqref{eq:23} 得 $\ch(f_!(1))=f_*(\Todd(X))=\Ahat(X)[X]$。译文照录原文 $z=0$，此处指出。]

% [Tutor Note: 易错——顺着上一条：$K(\pt)\cong\ZZ$，而 $\ch:K(\pt)\to H^*(\pt;\QQ)$ 是恒等映射，故 $f_!(1)$ 这个整数与 $\Ahat(X)[X]$ 的相等是"数"层面的等式；但 $f_!(1)$ 作为 $K$ 类存在，比"示性数恰为整数"强，这正是这一论证被称为"给出了 $\Ahat[X]$ 的拓扑解释"的原因。]

尽管如此，这一解释仍被认为不能令人满意。正如阿蒂亚~\cite{A1} 所述，希策
布鲁赫意识到号差是调和微分形式空间的维数之差，于是他为 $\Ahat$ 亏格
$\Ahat(X)[X]$ 寻求类似的分析解释。因此，当辛格于 1962 年 1 月抵达牛津作
学术休假时，阿蒂亚问他的第一个问题是：“为什么 A-roof 亏格对 spin 流形是
整数？”辛格~\cite{S1} 回答：“迈克尔，你为什么要问我这个问题？你知道答案的。”
但阿蒂亚在寻求更深的东西，而且他立刻让辛格上了钩。到三月，这对搭档已经掌握
了 Dirac 算子与指标公式。随后，在最初那次谈话九个月之后，阿蒂亚与辛格完成了
他们那以二人命名的指数定理的第一个证明。

% [Tutor Note: 含糊——"号差是调和形式空间的维数之差"省略了 Hodge 理论： $\Sign(X)$ 等于杯积形式限制在 $H^{2k}$ 的调和形式空间上所得对称型的正、负特征子空间维数之差（Hodge 分解 + 中间次数的 Hodge 星算子）。Hirzebruch 要的是 $\Ahat$ 对应的解析刻画，即某椭圆算子的核与余核维数之差，答案就是 Dirac 算子的指数 $\ind D$。]

% [Tutor Note: 术语——"A-roof genus" 是 $\Ahat$ 亏格的口语读法（$\hat A$ 形似屋顶，读作 "A-roof" 或 "A-hat"），并非另一个不变量；它与 \eqref{eq:24} 的问句是同一个问题，只是这里要求的是分析/指标解释，而非用 Atiyah–Hirzebruch $K$ 理论给出的拓扑解释。]

% ==================== §3 ====================

\section{分析中的先驱}\label{sec:3}
\begin{keytakeaways}
  \item de Rham 定理断言：由外微分定义的 de Rham 上同调与拓扑的实
  上同调自然同构；Hodge 定理则断言闭黎曼流形上每个上同调类有唯一的
  调和代表元。
  \item 由此调和形式空间维数的交错和
  $\sum_{q}(-1)^q\dim\sH^q(X)=\Euler(X)$ 可写成椭圆算子
  $d+d^*\:\Omega^{\textnormal{even}}(X)\to\Omega^{\textnormal{odd}}(X)$
  的 Fredholm 指数；同一机制还给出号差算子。
  \item Dolbeault 定理把层上同调 $H^q(X,V)$ 实现为 $\dbar$-复形的
  上同调；在凯勒流形上，黎曼–罗赫定理把全纯欧拉示性数表为拓扑量
  $\chi(X,V)=\Todd(X)\ch(V)[X]$。
  \item 椭圆算子由主象征在 $T^*X\setminus0$ 上逐点可逆来定义；椭圆
  理论保证 Sobolev 空间上的延拓为 Fredholm 算子，指数
  $\ind P=\dim\ker P-\dim\coker P$ 为整数。
  \item 指数在紧扰动与连续形变下不变，故只依赖主象征的同伦类——这
  正是阿蒂亚–辛格指数定理所要计算的对象。
  \item 1920 年诺特的 Toeplitz 算子指数公式与 1960 年盖尔范德提出
  的拓扑公式问题，构成指数定理在分析方面的前史。
\end{keytakeaways}

阿蒂亚–辛格（Atiyah–Singer）指数定理把代数几何与代数拓扑的世界同
微分几何与整体分析的世界联系起来。我们在 §\ref{subsec:3.1} 中对后
者的介绍，从关于调和微分形式及其与上同调之关系的基本定理开始。黎
曼流形上的几何椭圆微分算子处于核心地位。我们在
§\ref{subsec:3.2} 中讨论更一般的椭圆算子，并回顾 Fredholm 算子的
基本事实。Fredholm 指数是 Fredholm 算子的整值形变不变量，也正是
指数理论得名所系的主角。在 §\ref{subsec:3.3} 中，我们让读者略窥
1920--1963 年间围绕椭圆算子指数所展开的活动。

\subsection{de Rham、Hodge 与 Dolbeault}\label{subsec:3.1}

我们从 de Rham 定理与 Hodge 定理开始，它们体现了椭圆线性微分方程
与拓扑之间的关系。设 $X$ 为光滑 $n$ 维流形，考虑微分形式复形
  \begin{equation}\label{eq:25}
     \Omega ^0(X) \xrightarrow{\;\;d\;\;} \Omega ^1(X)
     \xrightarrow{\;\;d\;\;} \cdots \xrightarrow{\;\;d\;\;} \Omega ^n(X),
  \end{equation}
其中 $d$ 是埃利·嘉当（Elie Cartan）的外微分。
% [Tutor Note: 术语——对 $q\ge1$，(3.2) 的分子 $\ker[d:\Omega^q\to\Omega^{q+1}]$ 与分母 $\image[d:\Omega^{q-1}\to\Omega^q]$ 都是无限维向量空间（$q=0$ 时分子为常数、分母为 $0$），取商之后才有限维；有限性可由 Hodge 理论（(3.4) 的椭圆性）或 de Rham 定理证明，无法从 (3.2) 直接看出。]
de Rham 上同调向量空间定义为商
  \begin{equation}\label{eq:26}
     H^q_{\text{DR}}(X) = \frac{\quad\,\ker \,[d\: \Omega ^q(X) \to \Omega
     ^{q+1}(X)]}{\image \,[d\: \Omega ^{q-1}(X) \to \Omega ^q(X)]},\qquad
     0\le q\le n.
  \end{equation}
德·拉姆（de Rham）在其 1931 年学位论文 \cite{deR} 中证明的定理断
言：对每个 $q$，de Rham 上同调与由奇异上链定义的实上同调之间存在
自然同构
$H^q_{\text{DR}}(X)\xrightarrow{\;\;\cong \;\;} H^q(X;\RR)$。
（这是现代语言；德·拉姆实际证明的是：给定周期，存在闭形式，且模
恰当形式后唯一。）注意，$H^q_{\text{DR}}(X)$ 是利用微分算子定义的，
而 $H^q(X;\RR)$ 来自拓扑。
% [Tutor Note: 含糊——“存在唯一的‘最佳’形式”所略去的是 Hodge 分解 $\Omega^q(X)=\sH^q(X)\oplus\image d\oplus\image d^*$：存在性由在闭流形上极小化 $L^2$ 范数得到（椭圆算子 $\Delta$ 的 Fredholm 性及紧性），唯一性由 $\langle\Delta\omega,\omega\rangle=\|d\omega\|^2+\|d^*\omega\|^2$ 得到。原始文献为 [Hod]；椭圆理论的证明见 [Ho3]。]
霍奇（Hodge）受代数几何中问题的驱动，证明了：在{\it 闭黎曼\/}流
形上，每个上同调类中存在唯一的“最佳”形式。
% [Tutor Note: 易错——Hodge 星号 $*$ 的归一化决定 (3.3) 的写法：若取 $\alpha\wedge *\beta=\langle\alpha,\beta\rangle\,\mathrm{vol}$，则伴随为 $d^*=(-1)^{np+n+1}*d*$（$p$ 为形式次数），故 $d(*\omega)=0$ 与 $d^*\omega=0$ 相差一个可逆的 Hodge 星号算子（可能差符号），两种写法等价；但若 $*$ 的约定整体差一个常数或符号，则 $d^*$ 与 $\Delta$ 的整体符号须重新核对（(3.4) 取 $\Delta=dd^*+d^*d$ 使 $\Delta$ 非负；几何分析文献有时把非负拉普拉斯算子记作 $-\Delta$，与 (3.4) 相差一个整体符号）。]
具体说来，在定向黎曼流形 $X$ 上，霍奇定义了对偶运算
$*\: \Omega ^q(X)\to \Omega ^{n-q}(X)$，并对闭流形论证 \cite{Hod}：
每个 de Rham 上同调类中存在唯一的形式 $\omega$ 满足
  \begin{equation}\label{eq:27}
     d\omega =0,\qquad d(*\omega )=0.
  \end{equation}
% [Tutor Note: 易错——“在闭流形上 (3.4) 等价于 (3.3) 中的两个方程”依赖分部积分：由 $\int_X\langle\Delta\omega,\omega\rangle\,\mathrm{vol}=\|d\omega\|^2+\|d^*\omega\|^2$ 知 $\Delta\omega=0$ 当且仅当 $d\omega=d^*\omega=0$，而 $d^*\omega=0$ 与 $d(*\omega)=0$ 相差一个可逆的 Hodge 星号算子（可能差符号）。闭流形（无边）保证分部积分无边界项；带边流形上须补加边界条件。]
这些{\it 调和\/}微分形式是椭圆 Hodge–Laplace 方程
  \begin{equation}\label{eq:31}
     \Delta \omega = (dd^*+d^*d)\omega =0,
  \end{equation}
的解，而该方程在闭流形上等价于 \eqref{eq:27} 中的两个方程。
% [Tutor Note: 术语——$\sH^q(X)$ 是调和 $q$-形式构成的空间（由 $\Delta\omega=0$ 定义），与上同调群 $H^q(X;\RR)$（由闭形式模恰当形式定义）是不同对象；“更强”的同构 $\sH^q(X)\xrightarrow{\;\cong\;}H^q(X;\RR)$ 才是 Hodge 定理的完整内容，$\dim\sH^q=b_q$ 只是其推论。]
线性无关解的个数——即解空间 $\sH^q(X)$ 的维数——等于一个拓扑不
变量，即贝蒂数 $b_q=\dim H^q(X;\RR)$。还有更强的论断，即存在同构
$\sH^q(X)\xrightarrow{\;\cong \;}H^q(X;\RR)$。
% [Tutor Note: 易错——“这两个论断都不能推广”指对一般的椭圆算子：$\dim\ker P$ 不再是某个贝蒂数，$\ker P$ 也不再与某个上同调群同构；保留下来的只有交错和 $\ind P=\dim\ker P-\dim\coker P$（核与余核一般也不等维，见 (3.6)），这就是 §3.2 的 Fredholm 指数。]
这两个论断都不能推广到一般的椭圆微分算子；不过，在这种情形下，
指数定理所计算的正是调和形式空间维数的交错和，而这是一个人们熟
悉的拓扑不变量：
  \begin{equation}\label{eq:32}
     \sum\limits_{q=0}^n (-1)^q\dim\sH^q(X) = \Euler(X),
  \end{equation}
% [Tutor Note: 易错——(3.5) 的 $\Euler(X)$ 是常系数实上同调的交错和，对应平凡系数 $\RR$ 的 de Rham 复形；而 (2.3) 的全纯欧拉示性数即使取平凡全纯线丛也给出 $\Todd(X)[X]$，一般不等于 $\Euler(X)$。作者“(Compare (2.3))”只是提醒同一“交错和”格式，不可读成数值相等。]
其中 $\Euler(X)$ 是 $X$ 的欧拉数。（比较 \eqref{eq:3}。）

我们可以把 \eqref{eq:32} 的左端写成某个椭圆算子的指数，即
  \begin{equation}\label{eq:77}
     d+d^*\:\Omega ^{\textnormal{even}}(X)\longrightarrow \Omega
     ^{\textnormal{odd}}(X).
  \end{equation}
% [Tutor Note: 易错——把 (3.6) 的余核与伴随的核等同，需要 $d+d^*$ 为 Fredholm 算子（闭值域、有限维核与余核），此时才有 $\coker T\cong\ker T^*$；对一般算子这仅在闭值域时成立，否则 $\ker T^*=(\image T)^\perp$ 不再同构于 $\coker T$。此处由椭圆理论保证。]
它的形式伴随是
$d+d^*\:\Omega ^{\textnormal{odd}}(X)\to \Omega ^{\textnormal{even}}(X)$，
而我们把 \eqref{eq:77} 的余核与伴随的核等同起来。
% [Tutor Note: 易错——脚注中的 $\tau$ 是对合：由 $*^2=(-1)^{p(n-p)}$（$n=2\ell$）可得 $\tau^2\omega=(-1)^{p(n-p)}\,i^{p(p-1)+q(q-1)+2\ell}\omega=\omega$（$q=n-p$，且 $i$ 的幂为 $(-1)^p$），故 $\Omega^\pm$ 是 $\tau$ 的 $\pm1$ 特征空间。若把 $\tau$ 换成 $-\tau$（或把 $i$ 的幂次整体乘 $-1$），$\Omega^\pm$ 互换，(3.7) 的指数变号；此 $\zt$-分级不同于 (3.6) 的偶/奇次分级。核对 [AS3, §6] 的约定。]
若 $\dim X$ 为偶数，则在复微分形式
$\Omega ^{\bullet }(X,\CC)$ 上取另一种
$\zt$-分级\footnote{分级
$\Omega ^{\bullet }(X,\CC)= \Omega ^+(X,\CC)\oplus \Omega ^-(X,\CC)$
是对合 $\tau $ 在 $\Omega ^{\bullet }(X,\CC)$ 上的特征空间分解，其中
$\tau (\omega )= i^{p(p-1)+\ell }*\omega $，$\dim X=2\ell $。}会得到
另一个椭圆算子 \cite[\S6]{AS3}
  \begin{equation}\label{eq:78}
     d+d^*\:\Omega ^+(X,\CC)\longrightarrow \Omega ^-(X,\CC),
  \end{equation}
% [Tutor Note: 含糊——(3.7) 的指数等于号差，其表述要求 $\dim X$ 为 $4$ 的倍数（中间维数 $2k$ 为偶数，(2.9) 的配对才有非零号差）；若 $\dim X\equiv2\pmod4$，(2.9) 的配对在中间维数是交替的、号差为 $0$，(3.7) 的指数也应为 $0$。此外定向与 $\Omega^\pm$ 的选取须与 (2.9) 一致，否则差一个负号。]
即\emph{号差算子}，它在维数被 $4$ 整除的闭流形上的指数等于配对
\eqref{eq:8} 的号差。

% [Tutor Note: 术语——复流形上 $\Omega^{p,q}$ 的第一个指标是全纯次数、第二个是反全纯次数，$\dbar$ 只增大第二个指标；故 (3.8) 是 $(0,q)$ 行，Dolbeault 上同调记作 $H^{0,q}_{\dbar}(X,V)$，与层上同调 $H^q(X,V)$ 的 $q$ 对应。若与 de Rham 的 $H^q_{\text{DR}}$ 混用次数，会差一位。]
% [Tutor Note: 易错——此处 $n$ 改指复维数（$X$ 是 $n$ 维复流形），与本节开头 $n=\dim_{\RR}X$ 的用法不同；(3.8) 的末项 $\Omega^{0,n}$ 与 (3.9) 的求和上限 $n$ 都应读作复维数，请结合 $\dim X=1$ 的退化情形核对。]
设 $X$ 为闭 $n$ 维复流形，$V\to X$ 为全纯向量丛。那么
\S\ref{subsec:2.3} 中所用的层上同调群 $H^q(X,V)$ 同构于
$\dbar$-复形
  \begin{equation}\label{eq:33}
     \Omega ^{0,0}(X,V)\xrightarrow{\;\;\dbar\;\;}\Omega
     ^{0,1}(X,V)\xrightarrow{\;\;\dbar\;\;}\cdots \xrightarrow{\;\;\dbar\;\;}\Omega
     ^{0,n}(X,V),
  \end{equation}
的上同调群，这一点由多尔博（Dolbeault）证明。若 $\dim X=1$，则
\eqref{eq:33} 退化为单个椭圆算子；当 $V\to X$ 是 $X$ 上除子 $D$ 相
伴的线丛时，\S\ref{subsec:2.1} 中的向量空间 $\sL(D)$ 自然同构于
$\dbar$ 的核。
% [Tutor Note: 含糊——“Hodge 理论表明……”分两步：(i) 对任意紧 Hermitian 复流形，$\dbar$-Laplacian $\Delta_{\dbar}=\dbar\dbar^*+\dbar^*\dbar$ 的 Hodge 理论（$\dbar$ 复形是椭圆复形）已给出 $H^{0,q}(X,V)\cong\sH^{0,q}_{\dbar}(X,V)$；(ii) 凯勒条件用于 Kähler 恒等式 $\Delta_d=2\Delta_{\dbar}$，从而 $\dbar$-调和的 $(0,q)$-形式恰为 de Rham 意义下的复调和形式，并有 $H^r(X;\CC)=\bigoplus_{p+q=r}H^{p,q}(X)$。见 [H1]。]
若 $X$ 是凯勒（Kähler）流形，则 Hodge 理论表明，Dolbeault 上同调
向量空间同构于复调和微分形式构成的向量空间。把这些定理合在一起，
我们推出：在凯勒流形上，全纯欧拉示性数 \eqref{eq:3} 是各调和形式
空间维数的交错和。希策布鲁赫（Hirzebruch）的黎曼–罗赫定理
\ref{thm:1} 是关于这个解析量的拓扑公式：
% [Tutor Note: 笔误——(3.9) 左端原文作 $\sum_{q=0}^n\dim\sH^{0,q}(X)$，但按上文“交错和”的表述与 (2.3) 的定义，疑应为 $\sum_{q=0}^n(-1)^q\dim\sH^{0,q}(X,V)$：既缺 $(-1)^q$，又缺系数丛 $V$，否则左端不可能等于 $\chi(X,V)=\Todd(X)\ch(V)[X]$。译文照录原文，未改 (3.9)。]
  \begin{equation}\label{eq:35}
     \sum\limits_{q=0}^n \dim\sH^{0,q}(X) = \chi (X,V)=\Todd(X)\ch(V)[X].
  \end{equation}
正如 \S\ref{sec:2} 末尾所提到的，号差定理也可以用调和微分形式来
诠释。

\subsection{椭圆微分算子与 Fredholm 指数}\label{subsec:3.2}

我们在闭 $n$ 维流形 $X$ 上建立指数问题。设 $E^0,E^1\to X$ 是 $X$
上的向量丛，并设
$P\:C^{\infty}(X,E^0) \to C^{\infty}(X,E^1)$ 是 $m$ 阶线性微分算
子。在 $X$ 的局部坐标 $x^1,\dots ,x^n$ 下，对 $E^0\to X$ 的光滑截
面 $u\in C^{\infty}(X,E^0)$，有
  \begin{equation}\label{eq:28}
     Pu = a^{i_1i_2\dots i_m} \frac{\partial^mu}{\partial x^{i_1}\partial
     x^{i_2}\cdots \partial x^{i_m}} + \text{低阶项},
  \end{equation}
其中 $a^{i_1i_2\cdots i_m}$ 是对指标 $i_j$ 对称的丛映射
$E^0\to E^1$，我们并对指标 $1\le i_j\le n$ 求和。这个最高阶项在
坐标变换下按张量方式变换，因而定义了一个整体的丛映射
  \begin{equation}\label{eq:29}
     \sigma (P)\: \Sym^m(T^*X) \otimes E^0\longrightarrow E^1,
  \end{equation}
% [Tutor Note: 术语——原文所谓 $P$ 的 symbol $\sigma(P)$ 实为主象征（principal symbol），只记录最高阶项；(3.10) 的全象征还含各低阶项。下文“symbols 的连续路径”也指主象征的路径。]
称为 $P$ 的{\it 象征}。把 $\sigma (P)$ 视为 $T^*X$ 上取值于
$\Hom(E^0,E^1)$ 的 $m$ 次齐次多项式。
% [Tutor Note: 易错——“象征可逆”是逐点条件：对每个 $x$ 与每个非零 $\theta\in T^*_xX$，$\sigma(P)(\theta,\dots,\theta)$ 可逆；它在 $\theta=0$ 处不可逆，故准确说法是 $\sigma(P)$ 在 $T^*X\setminus0$ 上给出丛同构。这要求 $\rank E^0=\rank E^1$；椭圆性是开条件，对主象征的 $C^0$ 扰动稳定。]
若微分算子 $P$ 的象征可逆，则称 $P$ 为{\it 椭圆的}；也就是说，若
对每个 $x\in X$ 与非零的 $\theta \in T^*_xX$，线性映射
$\sigma (P)(\theta ,\dots ,\theta )\:E^0_x\to E^1_x$ 可逆。
% [Tutor Note: 含糊——“由椭圆理论可知核与余核有限维（用到紧致性）”所略去的是：椭圆正则性（$Pu\in C^\infty$ 推出 $u\in C^\infty$，见 [Se2] 与 [Ho3]）加上 Rellich–Kondrachov 定理（$X$ 紧时 $H^s\hookrightarrow H^{s-m}$ 为紧嵌入），由此 $P:H^s\to H^{s-m}$ 为 Fredholm 算子。$X$ 非紧时紧嵌入失效，核或余核可以是无限维的。]
由椭圆理论可知，$P$ 的核与余核都是有限维的。（这里用到 $X$ 的
\emph{紧致性}。）$P$ 的（\emph{Fredholm}）\emph{指数}为
  \begin{equation}\label{eq:30}
     \ind P = \dim\ker P - \dim\coker P.
  \end{equation}

% [Tutor Note: 含糊——“适当 Sobolev 空间”指对每个实数 $s$，$P$ 延拓为 $P_s:H^s(X,E^0)\to H^{s-m}(X,E^1)$ 且为 Fredholm 算子；其核在 $s$ 充分大时与 $C^\infty$ 核一致，指数与 $s$ 无关。细节见 [AS2, §1] 与 [Ho3]。]
椭圆理论证明：$P$ 到适当 Sobolev 空间上的延拓是\emph{Fredholm 算
子}。
% [Tutor Note: 术语——Fredholm 算子的定义关键是“闭值域”（不是“闭算子”）：$T$ 为 Fredholm 算子当且仅当它在 Calkin 代数（有界算子模紧算子构成的商代数）中可逆（Atkinson 定理），此时 $\coker T\cong\ker T^*$。只有有限维核与余核而不闭值域者不算 Fredholm 算子。]
回顾：希尔伯特空间之间的有界线性算子 $H^0\to H^1$，若具有闭值
域、有限维核与有限维余核，则称为 Fredholm 算子。（该定义可推广到
巴拿赫空间乃至更广的范围。）
% [Tutor Note: 含糊——(3.12) 的符号脚注实质是行列式的约定问题：若在次数 $0,1$ 上写两项复形 $0\to V^0\to V^1\to0$，则 $\sum(-1)^q\dim V^q=\dim V^0-\dim V^1=\dim\ker-\dim\coker$ 恰为 (3.12)，故数值上并无歧义；作者的顾虑（“符号本应相反”）来自 $T\in V^1\otimes(V^0)^*$ 的对偶性以及 $\det T$ 与 $\det T^*$ 在行列式线意义下互为倒数，这在后文行列式线丛的讨论中会再次出现。]
Fredholm 算子的指数由 \eqref{eq:30} 定义\footnote{有限维向量空间
之间的线性映射 $T\:V^0\to V^1$ 是 $(V^1)\otimes (V^0)^*$ 的元素，
因此按理说 \eqref{eq:30} 中的符号本应相反：对偶——或者说负号——
落在定义域一侧，而不是余定义域一侧。惯用的符号约定日后会带来麻
烦，例如在行列式理论中。另一方面，人们也可以为惯用的符号约定辩
护：把单个算子改写为位于 $0$ 次与 $1$ 次的两项复形，于是指数的惯
用符号（如 \eqref{eq:3} 中那样）就归结为 \eqref{eq:30}。}。
% [Tutor Note: 笔误——脚注末句原文作“the usual sign for the index (as in (3)) reduce to (30)”，其中 reduce 疑应为 reduces（主语是 the usual sign）；按印刷号即 (2.3) 与 (3.12)，译文按单数理解。]
连续线性映射构成的空间 $\Hom(H^0,H^1)$ 以算子范数赋予巴拿赫空间
结构，而 Fredholm 算子构成的开子空间
$\Fred(H^0,H^1)\subset \Hom(H^0,H^1)$ 具有任意高阶的非平凡同伦
群。
% [Tutor Note: 含糊——Fredholm 算子空间“有无界阶的非平凡同伦群”指 Atiyah–Jänich 定理：$\Fred(H^0,H^1)$ 同伦等价于 $\ZZ\times BU$（复 K 理论的分类空间），故由 Bott 周期性（周期 $2$）各阶同伦群在无穷多阶非平凡，且 $\pi_0\cong\ZZ$ 正由 (3.13) 的指数给出。见 [A4]；这正是 §4–§5 用 K 理论证明指数定理的出发点。]
特别地，指数
  \begin{equation}\label{eq:34}
     \ind \:\pi _0\Fred(H^0,H^1)\longrightarrow \ZZ
  \end{equation}
是一个同构。换言之，数值指数是单个 Fredholm 算子的完全形变不变
量。
% [Tutor Note: 含糊——“紧算子平移保持 Fredholm 子空间且指数不变”即 Atkinson 定理的推论：$T$ Fredholm、$K$ 紧则 $T+K$ Fredholm 且 $\ind(T+K)=\ind T$；这也是下段“低阶项相对最高阶项为紧”能推出“指数只依赖象征”的原因。]
此外，紧算子通过平移作用保持 Fredholm 算子子空间，因此指数在此
平移下不变。

% [Tutor Note: 含糊——“低阶项相对最高阶项为紧”即 Rellich 定理：$m$ 阶算子与 $(m-1)$ 阶算子之差 $H^s\to H^{s-m}$ 为紧算子（经由紧嵌入 $H^{s-m+1}\hookrightarrow H^{s-m}$）；等价地，$P$ 与只保留主象征的算子相差一个紧算子，故指数相同。参见 [Se2] 与 [Ho3]。]
椭圆理论表明：光滑流形上椭圆微分算子 \eqref{eq:28} 的低阶项相对
于最高阶项是紧的，而最高阶项本质上就是象征 \eqref{eq:29}。由此可
知，椭圆微分算子的指数是其象征的不变量。此外，椭圆微分算子的连
续路径诱导出 Fredholm 算子与象征的连续路径（在适当的拓扑下）。指
数在这样的形变下不变；事实上，它只依赖于象征的同伦类。
% [Tutor Note: 术语——“index”（指数）在此指解析的 Fredholm 指数 (3.12)；§4 将证明它等于拓扑指数（由主象征的同伦类经 K 理论推前给出）。注意 §2 提到旧文献把 signature 也称作 index，那与 Fredholm 指数不是同一概念，勿混。]
阿蒂亚–辛格指数定理用象征的同伦类给出了指数的公式。

\subsection{椭圆算子的指数问题}\label{subsec:3.3}

在 20 世纪上半叶，线性椭圆问题的指数定理有着
悠久而丰富的历史。其中许多被阿蒂亚–辛格指数定理及其到带边流形的
推广（\S\ref{subsec:5.2}）所涵盖。阿格拉诺维奇（Agranovich）
\cite{Agr} 与西利（Seeley）\cite{Se1} 的文章是了解这段历史的极好
向导。
% [Tutor Note: 术语——Fritz Noether（弗里茨·诺特）是 Emmy Noether（埃米·诺特）之弟，两人中文均译“诺特”，此处须保留名“弗里茨”以免混淆；[E,AM] 都把 Fredholm 指数的首创归于其 1920 年 Math. Ann. 论文 [N]。]
第一个指数定理包含在弗里茨·诺特（Fritz Noether）1920 年的一篇
论文 \cite{N} 中。（现代文献如 \cite{E,AM} 中已确认这一归属。）此
外，这篇论文似乎正是 Fredholm 指数的源头。
% [Tutor Note: 术语——旧文献的“Noether property”（诺特性质）即今之 Fredholm 性（闭值域、有限维核与余核），而旧称“Fredholm”专指指数为 $0$ 者；“Fredholm alternative”（Fredholm 二择一）指 $Tu=f$ 要么对每个 $f$ 唯一可解、要么齐次方程有非平凡解，恰是指数为 $0$ 的 Fredholm 算子的性质。]
% [Tutor Note: 笔误——原文此句“in the older literature the following terminology that is sometimes used”语法不完整，疑脱漏 is（应为“the following terminology is sometimes used”）；译文按其本意译出。]
事实上，较早的文献中有时使用如下术语：若线性算子是 Fredholm 的，
就说它满足“诺特性质”；此时“Fredholm”一词专指指数为零的算子
（这样的算子满足“Fredholm 二择一”）。

诺特的工作中有一个情形，即 Toeplitz 算子的指数公式。设
$\cir\subset \CC$ 为单位圆，$f\:\cir\to \CC^\times $ 为光滑的非零
复值函数。
% [Tutor Note: 易错——$\sH$ 若理解为圆盘上的 Bergman 空间（$L^2$ 全纯函数），它并不能按傅里叶级数嵌入 $L^2(\cir)$；下一句“By Fourier series $\sH$ sits as a subspace in $\tsH$”表明这里实际是 Hardy 空间 $H^2$（边界值实现，Fourier 系数只含非负频率）。$i$ 为包含映射、$\pi$ 为正交投影（Szegő 投影）。]
Toeplitz 算子 $T_f\:\sH\to \sH$ 定义在希尔伯特空间 $\sH$ 上，
$\sH$ 由单位圆盘 $\Omega \subset \CC$ 的闭包 $\overline{\Omega }$
上的 $L^2$ 全纯函数组成。借助傅里叶级数，$\sH$ 作为子空间落在
$\tsH=L^2(\cir,\CC)$ 中。设 $i\:\sH\hookrightarrow \tsH$ 为包含映
射，$\pi \:\tsH\twoheadrightarrow\sH$ 为正交投影。于是
$T_f=\pi \circ M_f\circ i$ 是乘法算子 $M_f\:\tsH\to \tsH$ 到 $\sH$
的压缩。

% [Tutor Note: 易错——定理的符号约定为 $\ind T_f=-\mathrm{wind}(f)$（绕数逆时针为正）；用 $f_n(z)=z^n$ 核对：$n>0$ 时 $T_{z^n}$ 是单侧移位的 $n$ 次幂，核为 $0$、余核为 $n$ 维，指数为 $-n=-\mathrm{wind}(z^n)$。若把 $\sH$ 取在圆盘外部，或把 $T_f$ 定义为 $i\circ M_f\circ\pi$，符号会改变。]
  \begin{theorem}[Noether, 1920]\label{thm:32}
 $T_f$ 是 Fredholm 算子，其指数等于 $f$ 的绕数的相反数。
  \end{theorem}

\noindent
读者或许愿意对 $f_n(z)=z^n$、$n\in \ZZ$ 显式计算该指数。定理
\ref{thm:32} 后来被戈赫贝格–克莱因（Gohberg–Krein）\cite{GK} 重
新发现；任意维数情形也有 Toeplitz 算子的指数定理，例如布泰·德·
蒙韦尔（Boutet de Monvel）\cite{BdM} 的工作。

1960 年，盖尔范德（Gelfand）\cite{G} 注意到指数是同伦不变量，并
提出了一般问题：为指数计算一个拓扑公式。看来，阿蒂亚与辛格在踏上
通往指数定理的旅程时，并不了解俄国发生的这些事件；不过，在斯梅尔
（Smale）访问牛津 \cite{A1} 期间，他们得知了这些情况。盖尔范德的
论文，以及其中一些解决指数问题特殊情形的先行工作，在阿蒂亚–辛格
宣布其一般指数定理的短文 \cite{AS1} 的开头被引用。

% ==================== §4 ====================

\section{指数定理及其证明}\label{sec:4}

\begin{keytakeaways}
  \item 本节讨论单个椭圆算子的阿蒂亚–辛格（Atiyah–Singer）指数定理：闭流形上
    椭圆算子的解析指数 $\ind P$ 等于由象征的上同调类算出的拓扑指数
    （定理 \ref{thm:5}）。
  \item 中心工具是 Clifford 代数与旋量丛。手征 Dirac 算子的指数
    $\dim\sH^+(X)-\dim\sH^-(X)=\Ahat(X)[X]$ 把 $\Ahat$ 亏格实现为
    解析不变量（式 \eqref{eq:42}）。
  \item 第一个证明走配边路线：把问题化到被向量丛扭化的号差算子，
    对扭化号差算子证明 \eqref{eq:44} 形式的两端均为配边不变量，
    再在配边类所成有理向量空间的一组基上验证公式。
  \item 伪微分算子是技术支柱：主象征映射在椭圆伪微分算子上满射且
    连续，使每个椭圆象征类可由算子实现，象征的同伦也可由算子的
    同伦实现。
  \item 第二个证明采用 $K$ 理论：在 $\Ell(X)=K^0(T^*X,T^*X\setminus 0)$
    上定义解析指数 a-ind 与拓扑指数 t-ind，用一组公理唯一刻画二者
    并证明 a-ind = t-ind（定理 \ref{thm:8}），其基础是 $K$ 理论的
    Thom 同构与 Bott 周期性。
  \item 早期应用包括希策布鲁赫号差定理、希策布鲁赫–黎曼–罗赫定理
    以及标量椭圆算子指数为零等结论。
\end{keytakeaways}

% [Tutor Note: 笔误——源文件中 \label{subsec:4.5} 与 \label{subsec:4.4} 互换：文档顺序第 4 小节（应用）带 \label{subsec:4.5}，第 5 小节（K 理论）带 \label{subsec:4.4}。因此本节开头“第二个证明（\S\ref{subsec:4.4}）”编译后印作 4.5，而“\S\ref{subsec:4.5} 列出应用”印作 4.4。译文按源码文档顺序保留原 label，不重排、不改名；印刷号与交叉引用本身均正确，仅为 label 命名与所在小节不符。]
% [Tutor Note: 术语——"announced in~\cite{AS1} in~1963" 指 Bull. Amer. Math. Soc. 69 (1963) 的公告；第一个证明由 Palais 编成 Annals of Math. Studies 57 (1965) 的 Seminar 文集 [Pa]（撰稿人含 Atiyah、Borel、Floyd、Seeley、Shih、Solovay），并非 Atiyah–Singer 本人执笔；第二个证明见 Ann. of Math. 87 (1968) 系列第一篇 [AS2]。核对年代时注意“宣布年”与“正式发表年”的差别。]

我们现在来到单个椭圆算子的阿蒂亚–辛格指数定理。该定理于 1963 年
在~\cite{AS1} 中公布。他们的第一个证明
（§\ref{subsec:4.2}）以希策布鲁赫（Hirzebruch）对其号差定理与
黎曼–罗赫定理的配边证明为范本，他们本人并未写出，而是发表在
帕莱斯（Palais）所编的~\cite{Pa} 中，由多位作者分篇写成；这卷文集
至今仍是宝贵的参考文献。第二个证明（§\ref{subsec:4.4}）更多地以
格罗滕迪克（Grothendieck）为范本，于 1968 年发表在阿蒂亚与辛格
系列论文的第一篇~\cite{AS2} 中。其后的论文处理各种变形与推广。
我们从 §\ref{subsec:4.1} 黎曼几何中的 Dirac 算子开始。它是洛伦兹
（Lorentz）几何中 Dirac 算子的类比物，在几何与物理的许多场合以及
一般指数理论中都居于中心地位。另一方面，伪微分算子在指数定理的两个
证明中都起重要作用；我们在 §\ref{subsec:4.3} 中作一简要介绍。在
§\ref{subsec:4.5} 中我们列出指数定理的几个早期应用。

% [Tutor Note: 术语——Dirac 原文 [D] 在 Minkowski 时空（号差 $(1,3)$）上工作，其方程是一阶的，其平方为波算子；本节把时间方向解析延拓 $t\mapsto i x^n$（Wick 旋转），得到欧氏空间上的椭圆算子。原文说“the analogue”，要点正在于这一转动把双曲方程变为椭圆方程，从而进入 Fredholm 指数理论的适用范围。]

  \subsection{Dirac 算子}\label{subsec:4.1}

% [Tutor Note: 易错——(4.1)–(4.4) 的符号约定是配套的：取几何学家常用的 $\Delta=-\sum_i\partial_i^2$ 时，Clifford 关系必须是 (4.4) 的 $\gamma^i\gamma^j+\gamma^j\gamma^i=-2\delta^{ij}$（即 $\Cliff_{-n}$）；若改用 $\Delta=+\sum_i\partial_i^2$，右端要换成 $+2\delta^{ij}$（即脚注中的 $\Cnp$）。两处符号必须同时翻转，否则 $D^2=\Delta$ 不成立。]

1928 年，狄拉克（Dirac）~\cite{D} 在他的电子相对论理论中提出了他的
方程。狄拉克是在闵可夫斯基（Minkowski）时空上工作的。把狄拉克的探究
路线类比到带标准坐标 $x^1,\dots ,x^n$ 的欧氏空间~$\EE^n$ 上，就是要找
一个一阶微分算子
  \begin{equation}\label{eq:36}
     D = \gamma ^1\frac{\partial }{\partial x^1} +\cdots+\gamma
     ^n\frac{\partial }{\partial x^n} 
  \end{equation}
使其平方为拉普拉斯算子
  \begin{equation}\label{eq:37}
     \Delta =-\left\{ \left(\frac{\partial }{\partial x^1}\right)^2 + \cdots
     + \left(\frac{\partial }{\partial x^n}\right)^2 \right\} . 
  \end{equation}
（在闵可夫斯基时空中，椭圆型拉普拉斯算子~\eqref{eq:37} 被双曲型波
算子取代。）设 $\gamma ^1,\dots ,\gamma ^n$~是 $\EE^n$ 上的常函数。
那么微分方程
  \begin{equation}\label{eq:38}
     D^2=\Delta 
  \end{equation}
等价于代数方程组
  \begin{equation}\label{eq:39}
     \gamma ^i\gamma ^j + \gamma ^j\gamma ^i = -2\delta ^{ij} = \begin{cases}
     -2,&i=j;\\\phantom{-}0,&i\neq j,\end{cases}\qquad 1\le i,j\le n. 
  \end{equation}
当~$n\ge2$ 时没有标量解，但有矩阵解。

% [Tutor Note: 术语——$\Cnm=\Cliff_{-n}$、$\Cnp=\Cliff_{+n}$ 是原文的记号（见脚注与 §5.4，源码 label 为 subsec:5.6）；实 Clifford 代数 $\Cliff_{\pm n}$ 的类型以 $n\bmod 8$ 为周期（[ABS] 的周期表），这决定了旋量表示 $\SS^0,\SS^1$ 在 $n\bmod 8$ 不同时的实/复/四元数性质。不要把 $\Cnm$ 的下标读成“负维数”。]
% [Tutor Note: 含糊——“the spin group is a double cover of $\SO_n$, and a subgroup of the units in $\Cnm$”一句跳过两点：(i) 对 $n\ge3$，$\Spin_n\to\SO_n$ 是万有覆盖（$\pi_1(\SO_n)=\ZZ/2$），$n=2$ 时 $\Spin_2\cong U(1)$；(ii) 更准确地说 $\Spin_n$ 位于 $\Cnm$ 偶部分 $\Cliff^0_{-n}$ 的可逆元群中，这正是旋量丛可由主 $\Spin_n$-丛相伴构造的原因。详见 [ABS] 与 [Ch]。]

设\footnote{要定义密切相关的 Clifford 代数~$\Cnp$，只需把关系式
\eqref{eq:39} 中的 $-$~号改成 $+$~号。两种 Clifford 代数都在
§\ref{subsec:5.6} 中出现。} $\Cnm$ 为\emph{Clifford
代数}~\cite{Cl,Ca,BW,Ch,ABS}，即由 $\gamma ^1,\dots ,\gamma ^n$ 生成、
服从关系式~\eqref{eq:39} 的 $\RR$ 上带单位元的代数。\eqref{eq:39} 的
一个矩阵解定义了一个 $\Cnm$-模。\emph{旋量群}~$\Spin_n$ 是特殊正交群
$\SO_n$ 的二重覆盖，并且 $\Spin_n$~是 $\Cnm$ 中可逆元群的一个子群，
正如 $\SO_n$~是 $n\times n$~实矩阵代数中可逆元群的子群一样。
Clifford 模限制到~$\Spin_n$ 上得到特殊的表示，称为\emph{自旋}表示或
\emph{旋量}表示。

% [Tutor Note: 笔误——(4.5)（源码 \eqref{eq:40}）原文作 \frac{\sinh y_i/2}{1/2}；按 TeX 的结合规则，$\sinh y_i/2$ 读作 $(\sinh y_i)/2$，再除以分母 $1/2$ 得 $\sinh y_i$，并不是旋量特征标之差。作者本意应为 $\sinh(y_i/2)/(1/2)=2\sinh(y_i/2)$，因为 $\SS^0-\SS^1$ 的特征标等于 $\prod_i(e^{y_i/2}-e^{-y_i/2})=\prod_i 2\sinh(y_i/2)$。译文照录原文，未改公式。]
% [Tutor Note: 笔误——紧接 (4.5) 的正文说 $y_1,\dots ,y_n$ 是 $\Spin_n$ 极大环面特征标的一组基，但 $\Spin_n$（$n$ 偶）的极大环面秩为 $n/2$，(4.5) 的乘积也只到 $n/2$；疑应为 $y_1,\dots ,y_{n/2}$。可对照 (2.24) 的 $\Ahat$ 亏格：那里 $X^{4k}$ 有 $2k=n/2$ 个 Chern 根。]
% [Tutor Note: 易错——(4.5) 并不是 $1/\Ahat$：由 (2.24)，$\Ahat=\prod_i (y_i/2)/\sinh(y_i/2)$，其倒数的分母应为 $y_i/2$ 而非 $1/2$。准确地说 (4.5) 等于 $\prod_i 2\sinh(y_i/2)$，即旋量特征标之差；不要把“upside down variant”理解为严格取倒数。]
% [Tutor Note: 含糊——“the spin condition was known to be related to the integrality of the $\Ahat$-genus, as explained in~\S\ref{subsec:2.4}”：所指是 Borel–Hirzebruch [BH3]（早期 [BH2] 差一个 2 的幂）利用 Milnor [Mi1] 的配边结果证明：若 $w_2(X)=0$，则 $\Ahat(X)[X]\in\ZZ$；更一般地，若 $w_2$ 是某个整上同调类的模 2 约化，则相应的 Todd 数为整数。见 §2.4 的脚注与 [H3]。]

旋量群为何会进入确定整数 $\Ahat(X)[X]$ 的探索？一个关键在于
$\Ahat$ 亏格的公式~\eqref{eq:22}。设 $n$~为偶数，群~$\Spin_n$ 有两个
特殊的、互不等价的不可约复表示~$\SS^0,\SS^1$，而 $\SS^0$ 与 $\SS^1$
的特征标之差是 $\Ahat$ 亏格~\eqref{eq:22} 的一个上下颠倒的变形，即
  \begin{equation}\label{eq:40}
     \prod\limits_{i=1}^{n/2}\frac{\sinh y_i/2}{1/2}, 
  \end{equation}
其中现在 $y_1,\dots ,y_n$ 是 $\Spin_n$ 的极大环面的一组特征标基。其次，
spin 条件已知与 $\Ahat$ 亏格的整性有关，这一点在~\S\ref{subsec:2.4}
中已有说明。或许正是这些考虑促使阿蒂亚与辛格在黎曼 spin 流形上
构造 Dirac 算子。还要注意，在黎曼流形上，一阶微分算子 $d+d^*$ 的平方
就是 Hodge–Laplace 算子~\eqref{eq:31}，因此在这个意义上 $d+d^*$ 已经
是一个 Dirac 算子。在凯勒流形上，$\dbar + \dbar^*$ 也是如此。

% [Tutor Note: 术语——“spin structure is a double cover $\Spin(X)\to\SO(X)$ together with a compatible principal $\Spin_n$-bundle structure”是 spin 结构的几何定义：它等价于把主 $\SO_n$-丛提升为主 $\Spin_n$-丛；存在性障碍是 $w_2(X)=0$，而两个 spin 结构之差由 $H^1(X;\ZZ/2)$ 分类。不要把主丛 $\Spin(X)$ 与旋量丛 $\SS_X$ 混为一谈：前者是主 $\Spin_n$-丛，后者是相伴的复向量丛。]
% [Tutor Note: 易错——(4.6) 的 $D_X=c\circ\nabla$ 只是手征 Dirac 算子 $\SS^0_X\to\SS^1_X$，它不（形式）自伴；（形式）自伴的是作用在 $\SS^0_X\oplus\SS^1_X$ 上、以 $D_X$ 与 $D_X^*$ 为两个非对角块的分块算子。其自伴性来自 $c$ 关于度量是斜伴随、$\nabla$ 保度量，以及 $X$ 紧闭（可用 Stokes 定理分部积分验证）。指数按 (3.12) 取 $\dim\ker D_X-\dim\coker D_X$，与 (4.7) 的左端对应。]
% [Tutor Note: 含糊——“The Levi-Civita connection on $X$ induces a covariant derivative $\nabla$ on each spinor bundle”省略的实质是 spin 联络：把 Levi-Civita 联络的 $\SO_n$-联络形式经 $\Spin_n\to\SO_n$ 的微分（同构 $\mathfrak{spin}_n\cong\mathfrak{so}_n$）提升到旋量丛，得到协变导数 $\nabla$；Dirac 算子的局部公式是 $D=\sum_i c(e_i)\nabla_{e_i}$（$e_i$ 为局部正交标架）。标准处理见 [ABS] 与 [Ch]。]
% [Tutor Note: 术语——“only determined uniquely up to tensoring by a line”：$\Spin_n$ 的两个不可约复旋量表示只在相差一个一维表示的意义下唯一。张量一个线表示不改变手征分解的结构，却改变旋量丛与 Dirac 算子（相当于用该线丛扭化）；§5.4（源码 label 为 subsec:5.6）的 Clifford 线性 Dirac 算子在 $\Cliff_n$-模层面消除了这一非典范性。]

\emph{那个} Dirac 算子在偶数维黎曼 spin 流形~$X^n$ 上是指数定理的
主角，因此我们在此给出定义。设 $\SO(X)\to X$ 是以 $\SO_n$ 为主群的
主丛，其在~$x\in X$ 处的纤维由定向等距 $\RR^n\to T_xX$ 组成。spin
结构是一个二重覆盖 $\Spin(X)\to\SO(X)$，连同复合映射
$\Spin(X)\to\SO(X)\to X$ 上的一个相容的主 $\Spin_n$-丛结构。利用
复旋量表示~$\SS^0,\SS^1$ 构造相伴的复向量丛
$\SS^{0}_X,\SS^{1}_X\to X$；这些向量丛的截面称为\emph{旋量}或
\emph{旋量场}。$X$ 上的列维–奇维塔（Levi-Civita）联络在每个旋量丛上
诱导一个协变导数~$\nabla $。Dirac 算子是
  \begin{equation}\label{eq:41}
     D_X = c\circ \nabla \:C^{\infty}(X,\SS^0_X)\longrightarrow
     C^{\infty}(X,\SS^1_X),  
  \end{equation}
其中 $c$~是 Clifford 乘法，由一个 $\Spin_n$-等变线性映射
$\RR^n\otimes \SS^{0}\to\SS^{1}$ 诱导而来。此外还有 Clifford 乘法
$\RR^n\otimes \SS^1\to \SS^0$，两者合起来使 $\SS^0\oplus \SS^1$ 具有
$\Cliff_n$-模的结构。这一构造并不完全典范，因为不可约表示
$\SS^0,\SS^1$ 只在相差张量一个线表示的意义下唯一确定。在
§\ref{subsec:5.6} 中我们将介绍 Clifford 线性 Dirac 算子，它是典范的。

% [Tutor Note: 含糊——“solutions $\psi$ to $D\psi=0$”中的 $D$ 是 $\SS^0_X\oplus\SS^1_X$ 上的自伴 Dirac 算子，而非 (4.6) 的手征算子 $D_X$；因此“调和旋量”按手征性自然分成 $\sH^+$ 与 $\sH^-$，它们分别对应 $D_X$ 的核与余核（经由 $D_X^*$ 的核）。这是 (4.7) 左端取 $\dim\sH^+-\dim\sH^-$ 的原因。]
% [Tutor Note: 易错——(4.7) 的正负号依赖于 $\SS^0,\SS^1$ 的编号约定：交换两者会使左端整体变号。与 (4.7) 配套的约定是 $\SS^0$ 取“正”手征、$\Ahat$ 取 (2.24) 的 $\prod(y_i/2)/\sinh(y_i/2)$；若采用 $\Delta=+\sum\partial_i^2$ 或相反的 $\gamma$ 符号约定，需重新核对符号。参见 [ABS] 对 $\SS^\pm$ 的约定。]
% [Tutor Note: 笔误——原文说“the DPhil thesis [Hi] of Atiyah's student Nigel Hitchin”，但文献表 [Hi] 是 N. J. Hitchin, Harmonic spinors, Adv. Math. 14 (1974) 1–55 的论文，并非学位论文本身（该文源自其 1972 年牛津 DPhil 论文）。引用时可注明“论文正式发表版”。]

de Rham 算子~\eqref{eq:25} 与 Dolbeault 算子~\eqref{eq:33} 有类似的
描述，于是 \eqref{eq:32}~与 \eqref{eq:35} ~启发人们猜测 Dirac 算子的
指数公式。设 $\sH^{\pm}(X)$ 为\emph{调和旋量}构成的复向量空间，即
$D\psi =0$ 的解~$\psi $。猜测的指数公式为
  \begin{equation}\label{eq:42}
     \dim\sH^+(X) - \dim\sH^-(X) = \Ahat(X)[X]. 
  \end{equation}
关于调和旋量诸方面的研究，可参看阿蒂亚的学生奈杰尔·希钦
（Nigel Hitchin）的博士（DPhil）论文~\cite{Hi}。

  \subsection{第一个证明：配边}\label{subsec:4.2}

% [Tutor Note: 术语——analytic index（解析指数）指 (3.12) 的 Fredholm 指数 $\dim\ker P-\dim\coker P$；topological index（拓扑指数）在这里先以上同调形式给出，即 $\Todd(TX\otimes\CC)\smallsmile\phi^{-1}\ch\sigma(P)$ 的示性数；§4.5 的 $K$ 理论形式用 $K$ 理论 Thom 同构定义 t-ind。阅读时注意作者的措辞层次：$\ind P$、上同调拓扑指数、t-ind 在定理中相等。]

设 $X$~为闭定向流形；$E^0,E^1\to X$ 为复向量丛；
$P\:C^{\infty}(X,E^0)\to C^{\infty}(X,E^1)$ 为椭圆微分算子。$P$ 的
\emph{解析指数}就是 \eqref{eq:30} 中定义的 Fredholm 指数。
\emph{拓扑指数}——以上同调形式给出时——如下所述。椭圆性蕴涵
\eqref{eq:29} 中的象征~$\sigma (P)$ 限制到非零向量 $T^*X\setminus 0$
上是同构。因此，$\sigma (P)$~定义了相对 $K$ 理论中的一个类，属于
$K^0(T^*X,\,T^*X\setminus 0)$。陈特征把 $K$ 理论映到有理上同调，
然后 Thom 同构的逆
  \begin{equation}\label{eq:43}
     \phi \:H^{\bullet }(X;\QQ)\xrightarrow{\;\;\cong \;\;} H^{\bullet +\dim
     X}(T^*X,\,T^*X\setminus 0;\,\QQ) 
  \end{equation}
把我们带回底流形~$X$ 的上同调。\emph{拓扑指数}就是复化切丛的 Todd 类
与 $\sigma (P)$ 在陈特征和 Thom 同构下的像的杯积。指数定理断言解析
指数与拓扑指数相等。

% [Tutor Note: 易错——(4.10) 中必须是复化切丛的 Todd 类 $\Todd(TX\otimes\CC)$，它与 $\Ahat$ 的关系是 $\Todd(TX\otimes\CC)=\Ahat(X)^2$（复化后 Chern 根成对为 $\pm y_i$，故 $\Todd=\prod_i[(y_i/2)/\sinh(y_i/2)]^2$）。注意不要与复流形的 Todd 类混淆：后者按 (2.25) 为 $\Todd(X)=e^{c_1(X)/2}\Ahat(X)$，二者相差一个因子。另注意 $\phi$ 是 Thom 同构、$\phi\inv$ 才是回到 $X$ 上同调的映射。]

  \begin{theorem}[Atiyah-Singer Index Theorem~\cite{AS1}]\label{thm:5}
 $P$ 的指数为
  \begin{equation}\label{eq:44}
     \ind P = \bigl(\Todd(TX\otimes \CC)\smallsmile \phi \inv \ch \sigma
     (P)\bigr)[X].  
  \end{equation} 
  \end{theorem}

\noindent
 在本节余下的部分，我们略述定理~\ref{thm:5} 的证明所涉及的主要思想。

% [Tutor Note: 术语——注 4.11 说“$K$ 理论公式（\S\ref{subsec:4.4}）更自然”：此处的 $K$ 理论公式指 §4.5（源码 label 为 subsec:4.4）末尾的 K 理论定理（定理 4.18，a-ind = t-ind）；由于 label 互换，\ref{subsec:4.4} 印刷为 4.5，但指向的正是 $K$ 理论小节。]

  \begin{remark}[]\label{thm:6}
 指数的 $K$ 理论公式（§\ref{subsec:4.4}）更为自然，并且便于作种种
推广。这与阿蒂亚的哲学相符：以线性代数为基础的 $K$ 理论比上同调
更为初等。诚然，在代数拓扑的各种形式中，$K$ 理论与线性微分算子
最为贴近。
  \end{remark}

% [Tutor Note: 含糊——“the cyclic module $K^0(X)\sigma_0$ ... is a subgroup of finite index”省略了论证的关键：解析指数与拓扑指数都是 $\Ell(X)\to\ZZ$ 的群同态；若子群 $K^0(X)\sigma_0$ 在 $\Ell(X)$ 中指数有限，则只需在该子群上验证等式（对任意 $\sigma$，有 $m\neq0$ 使 $m\sigma\in K^0(X)\sigma_0$，同态性给出 $m(\ind\sigma-\text{top}\ \sigma)=0$）。而 $K^0(X)\sigma_0$ 中的元素正是扭化号差算子的象征，故问题化归到 $D_W$。]

现在进入证明。在~\S\ref{subsec:3.2} 末尾我们指出，解析指数只依赖于
象征~$\sigma (P)$ 的同伦类。阿蒂亚–辛格引入椭圆象征类构成的阿贝尔群
  \begin{equation}\label{eq:45}
     \Ell(X) = K^0(T^*X,\,T^*X\setminus 0) 
  \end{equation}
它通过张量积成为环~$K^0(X)$ 上的模，而由号差算子的象征类生成的
\emph{循环}模 $K^0(X)\sigma _0$ 是一个指数有限的子群。这就把固定流形
$X$~上的问题化归为被向量丛 $W\to X$ 扭化的号差算子。然而，要完成这一
论证，需要：$\Ell(X)$ 的每个元素都是某个椭圆算子的象征；并且，若椭圆
算子~$P_0,P_1$ 的象征定义 $\Ell(X)$ 中的同一个元素，则存在同伦
$P_0\sim P_1$。\emph{微分}算子对于这些性质来说过于刚性，证明中的一个
关键步骤是引入\emph{伪微分}算子，我们在下一节中简要讨论。

% [Tutor Note: 含糊——“Bott periodicity, as realized in $K$-theory”：指复 $K$ 理论的 Bott 周期性 $K^0(X)\cong K^2(X)$（[B1]、[B2]），它给出 $K$ 理论的 Thom 同构 $K^0(X)\cong K^0(T^*X,T^*X\setminus 0)$（在约化 $K$ 群中即 suspension 同构），从而在 spin 流形上 $\Ell(X)$ 由 Dirac 象征循环生成。实情形的周期 8 见 [ABS] 的 Clifford 模表。]

  \begin{remark}[]\label{thm:7}
 若 $X$~是 spin 流形，则 $\Ell(X)$ 是由 Dirac 算子的象征生成的循环
$K^0(X)$-模。这一事实是 Bott 周期性在 $K$ 理论中的一种体现。
  \end{remark}

% [Tutor Note: 含糊——“This is straightforward for the cohomological formula”省略了配边不变性的来源：右端是示性数，其配边不变性来自 $\Todd(TX\otimes\CC)$、$\Ahat$ 等示性类在配边下的自然性（在配边 $Y$ 上由 Stokes 定理化为边界上的积分）。左端解析指数的配边不变性则靠下述“在 $Y$ 上补一个算子”的论证；细节见 [AS3] 与 [AB3]。]
% [Tutor Note: 易错——“local elliptic boundary condition”指 Atiyah–Bott 意义下的局部（Lopatinski–Shapiro）边界条件，而非 APS 的整体谱投影条件；局部条件使 $D_U$ 在 $Y$ 上形式自伴且边界条件自伴，故 $\ind_YD_U=0$；而 (i) 式 $\ind_YD_U=\ind_XD_W$ 来自内部算子与其边界算子的指数关系（粘合/分离论证）。参见 [AB3]、[AS3]。]

配边论证在此阶段进入证明。对被向量丛 $W\to X$ 扭化的号差算子，
\eqref{eq:44} 的两端都被视为对 $(X,W)$ 的函数，其中我们只用到等价类
$[W]\in K^0(X)$。关键一步是证明两端都是 $(X,W)$ 的配边不变量。对右端
的上同调公式来说这是直接的。对解析指数，设 $(X,W)$ 是某一对 $(Y,U)$
的边界，其中 $Y$~是紧致定向流形，$U\to Y$ 是复向量丛。
阿蒂亚–辛格在 $Y$~上引入一个椭圆微分算子~$D_U$，其在 $X$~上的边界
算子就是扭化号差算子。他们指定一个\emph{局部}椭圆边界条件，并证明
(i)~$\ind_YD_U = \ind_X D_W$ 与 (ii)~$\ind_YD_U=0$。有了配边不变性，
剩下的就是计算 $(X,W)$ 的配边类所成有理向量空间的一组基，并对这些
基元素验证~\eqref{eq:44}。

从这一粗略梗概走到完整证明的旅程中，充满了分析、几何与拓扑方面的
有趣迂回。

  \subsection{伪微分算子}\label{subsec:4.3}

% [Tutor Note: 术语——(4.14) 中的 $p(x,\xi)$ 是伪微分算子的全象征（total symbol），而 (3.11) 与 (4.15) 中的 $\sigma(P)$ 是主象征（principal/top order symbol）；原文在“symbol”一词上两义并用，阅读时须按上下文区分。全象征依赖于坐标与截断，主象征才是流形上整体定义的。]
% [Tutor Note: 术语——“Schwartz kernel smooth away from the diagonal”：伪微分算子的核 $K(x,y)$ 在对角线外光滑，在对角线附近有 $|x-y|^{-n-m}$ 型奇性；微分算子的核是 $\delta(x-y)$ 及其导数，支集在对角线上，故“off the diagonal 为零”。这一性质保证伪微分算子把 $C^\infty$ 映到 $C^\infty$，并在 Sobolev 空间上有界。]

设 $P$~是 $\RR^n$ 上 $m$ 阶微分算子，如~\eqref{eq:28} 所示。它在紧
支集光滑函数~$u$ 上的作用，用傅里叶（Fourier）变换~$\hat{u}$ 来写
较为方便：
  \begin{equation}\label{eq:46}
     Pu(x) = (2\pi )^{-n}\int_{{\RR^n}^*} p(x,\xi )\,\hat{u}(\xi )\,e^{i\langle
     x,\xi \rangle} \,d\xi ,\qquad x\in \RR^n, 
  \end{equation}
其中对每个~$x\in \RR^n$，函数 $p(x,\xi )$ 是~$\xi \in {\RR^n}^*$ 的
$m$ 次多项式。一个可追溯到米赫林（Mikhlin）~\cite{Mik} 与
卡尔德隆–齐格蒙德（Calderón–Zygmund）~\cite{CZ} 的推广，允许更一般的
（全）象征~$p$，只要当 $|\xi |\to\infty $ 时 $p$ 及其导数具有一致界。
一个动机来自椭圆算子：拟基本解——一种近似逆——就是这种形式的线性
算子。此外，这类算子的施瓦茨（Schwartz）核在离开对角线的地方是
光滑的，这带来了良好的性质。（微分算子的施瓦茨核在对角线之外为零。）
伪微分算子理论及其向光滑流形的整体化，见于西利（Seeley）~\cite{Se2,Se3}、
赫尔曼德（Hörmander）~\cite{Ho2,Ho3}、科恩–尼伦伯格
（Kohn–Nirenberg）~\cite{KN} 以及帕莱斯–西利（Palais–Seeley）~\cite{PS}
等人的论文。这只是浩繁文献中极小的一部分。

% [Tutor Note: 易错——(4.15) 用极限定义主象征，要求 $p$ 满足 $p(x,\lambda\xi)/\lambda^m$ 当 $\lambda\to\infty$ 有极限；对一般的全象征这并非自动成立（例如对数项会破坏齐次性）。因此“限制到使极限存在的伪微分算子”不是技术性小条件，而是主象征有定义的前提；主象征仅保留最高阶的齐次部分，低阶项与有界项都被丢弃。]

若 $P$~是 $m$ 阶微分算子，从而对每个~$x$，$p(x,\xi )$~是~$\xi $ 的
$m$ 次多项式，则 $P$~的\emph{主}象征或\emph{最高阶}象征是由
  \begin{equation}\label{eq:47}
     \sigma (P)(x,\xi ) = \lim\limits_{\lambda \to\infty }\frac{p(x,\lambda
     \xi )}{\lambda ^m}. 
  \end{equation}
给出的 $m$ 次齐次多项式。我们只考虑象征~$p$ 使极限~\eqref{eq:47}
存在的伪微分算子。此时主象征有定义，其在光滑流形上的整体形式为
\eqref{eq:29}，而椭圆性一如既往是主象征的可逆性。椭圆伪微分算子上
主象征映射的满射性与连续性，是指数定理证明中的关键要素。

% [Tutor Note: 含糊——“the surjectivity and continuity of the principal symbol map ... are crucial”的实质是两条定理：(i) 每个椭圆象征类都可由某个伪微分算子实现（满射性，需要 $X$ 紧致）；(ii) 算子同伦与象征同伦相容（主象征映射的连续性/稳定性），从而解析指数只依赖象征的同伦类。这正是 §4.2 末尾所提两条性质的来源，见 [AS2] §1 与 [AB1]。]

我们指出，在阿蒂亚–博特（Atiyah–Bott）~\cite{AB1} 与阿蒂亚–辛格~\cite{AS2,AS4}
关于指数理论的论文中，伪微分算子的整体理论得到了进一步展开。

  \subsection{若干应用}\label{subsec:4.5}

% [Tutor Note: 术语——(4.7) 的 $\Ahat(X)[X]$ 是 $\Ahat$ 亏格的示性数：$\Ahat(X)$ 按 (2.24) 取 $\prod_i(y_i/2)/\sinh(y_i/2)$，其中 $p_1$ 的系数是 $-1/24$（即 $\Ahat=1-p_1/24+\cdots$）。若采用把 $\Ahat$ 换成与 $L$ 亏格混用的约定（$L=p_1/3$，$\Ahat=1-L/8+\cdots$），公式会差一个因子；核对时以 (2.24) 的 $\prod_i(y_i/2)/\sinh(y_i/2)$ 为准，参见 [H1]。]

阿蒂亚–辛格指数定理~\ref{thm:5} 的一个直接推论是 Dirac 算子的
指数公式~\eqref{eq:42}。这给出了 spin 流形的 $\Ahat$ 亏格的一种
\emph{解析}解释，从而回答了~\eqref{eq:24}。

% [Tutor Note: 易错——“the more flexible techniques of Atiyah-Singer prove it for arbitrary closed complex manifolds, a powerful generalization from projective algebraic manifolds”：希策布鲁赫–黎曼–罗赫定理 2.7（thm:1）原本对射影代数流形成立，而 AS 的椭圆算子方法把它推广到任意闭复流形（无需射影/代数）。这是“解析指数=拓扑指数”的直接应用：$\dbar+\dbar^*$ 的指数就是 $\chi(X,V)$。]
% [Tutor Note: 含糊——“given appropriate inequalities between $\dim X$ and $\rank E^0=\rank E^1$”原文未写出具体不等式；其意为：当 $E^0,E^1$ 为平凡丛时，象征是球面丛到 $\GL_k(\CC)$（$k=\rank E^0$）的映射，在 $\dim X$ 与 $k$ 满足适当不等式时该象征类平凡，指数因而为零。最清楚的特例是单个函数（$k=1$）。精确陈述请回到 [AS1] 核对维数与秩的条件。]

若干其他推论在~\cite{AS1} 中有描述。如前所述，希策布鲁赫号差定理
~\ref{thm:3} 是一个特例。希策布鲁赫–黎曼–罗赫定理~\ref{thm:1} 也是
特例，但阿蒂亚–辛格更为灵活的技术对任意闭复流形证明了它，这是
从射影代数流形出发的一个有力推广。最后，对 $X$~上的椭圆算子组——
即 $E^0,E^1\to X$ 为平凡向量丛的情形——只要 $\dim X$ 与
$\rank E^0=\rank E^1$ 之间满足适当的不等式，指数就为零；这一结果把
定理~\ref{thm:5} 与当时的偏微分方程文献联系起来。特别地，作用在
单个函数上的椭圆算子的指数为零。

随着新的指数定理不断涌现，应用也日益增多，我们将在~\S\ref{sec:5}
中看到。

  \subsection{第二个证明：$K$ 理论}\label{subsec:4.4}

% [Tutor Note: 术语——pushforward $f_!$ 在格罗滕迪克的黎曼–罗赫定理（(2.13)）中是层的前推 $f_!:K(X)\to K(Y)$；在拓扑范畴中 Atiyah–Hirzebruch 用 $K$ 理论 Thom 同构定义拓扑推前，使其在 $X$ 为复流形、$f$ 为复映射时与代数 $f_!$ 一致（见 (2.27) 附近的讨论）。]

回顾格罗滕迪克版本的黎曼–罗赫定理中出现的推前~\eqref{eq:11}。设
$X$~为紧射影簇。对唯一映射 $f\:X\to\pt$ 这一特殊情形，推前
$f_!\:K(X)\to K(\pt)\cong \ZZ$ 计算全纯向量丛 $V\to X$ 的 Euler
示性数~\eqref{eq:3}。根据 Dolbeault 定理，层上同调群同构于
$\dbar$-复形~\eqref{eq:33} 的上同调群。对光滑流形与光滑映射，
阿蒂亚–希策布鲁赫在他们关于 $K$ 理论的最早工作之一中找到了一个
\emph{拓扑推前}，其设计正是为了在这一情形下与格罗滕迪克的~$f_!$ 相
匹配。（参看 \eqref{eq:19} 附近的正文。）这里还有一种解析解释。即，
在凯勒流形~$X$ 上，Dolbeault 上同调群~$H^{0,q}_{\dbar}(X,V)$
同构于 $(0,q)$-形式上的椭圆算子 $\dbar + \dbar^*$ 的核，因此 Euler
示性数~\eqref{eq:3} 就是
  \begin{equation}\label{eq:48}
     \dbar+\dbar^*\:\Omega ^{0,\textnormal{even}}(X,V) \longrightarrow \Omega
     ^{0,\textnormal{odd}}(X,V). 
  \end{equation}
的指数。这个指数是 $V\to X$ 在映射~$f$ 下的一个\emph{解析推前}。
Riemann–Roch–de Rham–Hodge–Dolbeault–Hirzebruch 诸定理蕴涵解析推前与
拓扑推前相等。

% [Tutor Note: 含糊——(4.17) 的箭头只标了 $\ind$，而紧接着的正文定义了两个同态 a-ind 与 t-ind；可把 (4.17) 理解为二者共同的“定义域与值域”式，两个同态在此后分别命名。此处不是笔误，只是表述从简。下文定理 4.18 写 a-ind = t-ind，可证此读法。]

阿蒂亚–辛格指数定理的 $K$ 理论形式是对任意椭圆（伪）微分算子的
推广。回顾~\S\ref{subsec:4.2}：光滑流形~$X$ 上椭圆微分算子~$P$ 的
象征~$\sigma (P)$ 的“同伦类”是元素
$[\sigma (P)]\in K^0(T^*X, T^*X\setminus 0)$。阿蒂亚–辛格定义两个
同态
  \begin{equation}\label{eq:49}
     \ind\:K^0(T^*X, T^*X\setminus 0)\longrightarrow \ZZ. 
  \end{equation}
\emph{解析指数} a-ind 把象征类~$\sigma $ 送到任一满足
$[\sigma (P)]=\sigma $ 的椭圆伪微分算子~$P$ 的指数~$\ind P$。
\emph{拓扑指数} t-ind 与拓扑推前类似。它基于 $K$ 理论中的 Thom
同构，而后者又依赖于 Bott 周期性。

% [Tutor Note: 术语——$K$ 理论的 Thom 同构给出 $\Ell(X)=K^0(T^*X,T^*X\setminus 0)\cong K^0(X)$（“符号映射”把相对 $K$ 群与底流形的 $K$ 群等同起来）；于是 t-ind 可写成 $\Ell(X)\xrightarrow{\cong}K^0(X)\to K^0(\pt)=\ZZ$ 的复合，其中第二步是映到一点的推前。这正是 §4.2 上同调公式的 $K$ 理论对应物，其存在性依赖 Bott 周期性（即 $K$ 群中的 suspension/Thom 同构）。]
% [Tutor Note: 含糊——定理 4.18（a-ind = t-ind）标为“Atiyah-Singer~1967”，而本节上文说第二个证明于 1968 年发表于 [AS2]（Ann. of Math. 87）；1967 或指宣布/投稿/预印年份，不宜判为笔误。核对文献年份时以 [AS2] 为准。]

  \begin{theorem}[Atiyah-Singer~1967]\label{thm:8} a-ind = t-ind. 
  \end{theorem}

\noindent
 这与上同调形式的定理~\ref{thm:5} 等价，但正如我们将在~\S\ref{sec:5}
中看到的，$K$ 理论表述与证明的自然性使它们便于作许多推广。

% [Tutor Note: 含糊——“Atiyah-Singer uniquely characterize 'index homomorphisms' ... by a short list of axioms”：所谓短公理表包括正规化（在一点或平凡象征上的取值）、切除（excision）与乘性（multiplicativity）等；证明的核心是用 $K$ 理论的 Thom 同构与 Bott 周期性把任意象征类化归到这些基本情形。细节见 [AS2] §1 与 §2。]

定理~\ref{thm:8} 是~\cite{AS2} 的主题。该文 §1 中有关于证明思路的
简明概述，不过实际证明走的是略有不同的路线。本质上，阿蒂亚–辛格
用一份简短的公理表唯一地刻画了“指数同态”~\eqref{eq:49}，并验证
a-ind 与 t-ind 都满足这些公理。除正规化公理之外，指数的两条主要性质
在证明中起作用。一条\emph{切除公理}把指数延拓到任意（可能非紧的）
流形~$U$ 上具紧支集的象征，至少当 $U$~可嵌入紧流形时如此；它断言：
把 $U$~上具紧支集的象征零延拓为周围紧流形上的象征，指数不依赖于
延拓的选取。一条\emph{乘性公理}给出（扭化）乘积象征的乘积公式。
证明进行到此，需要一套稳健的伪微分算子整体理论。

% ==================== §5 ====================

\section{主题的各种变奏}\label{sec:5}

\begin{keytakeaways}
  \item 等变指数定理把经典的莱夫谢茨不动点公式推广到椭圆复形：指数成为紧群 $G$ 的
        虚拟表示，其特征标由不动点处的局部贡献求和给出。
  \item 带边流形上的指数定理需要\emph{局部}椭圆边界条件；这类边界条件的存在性障碍
        落在边界象征的 $K$ 群中；障碍消失时拓扑指数可延拓到相对情形，
        并等于解析指数。
  \item 实指标定理涵盖模 $2$ 指标与整指标（如 Rokhlin 型）两类情形；前者由
        $KO^{-n}(\pt)$ 的 $\ZZ/2$ 分量刻画，后者经四维 spin 流形的 Rokhlin 定理与
        号差整性相联系。
  \item Clifford 线性 Dirac 算子以 $\Cliff_{\pm n}$ 模取代不可约旋量表示，
        给出典范的 $KO^{-n}(\pt)$ 类；其关于 $n\bmod 8$ 的周期表等价于
        Bott 周期性。
  \item 族指数定理把指数从点上的 $K$ 群提升为底空间上的 $K^{\bullet }(S)$ 类，
        其最低阶分量是到 $H^1(S;\ZZ)$ 的行列式映射。
  \item 覆盖空间上的 von Neumann 指标把 $L^2$ 维数解释为实值指数，开启了非紧空间与
        奇异空间上的指标理论。
\end{keytakeaways}

在阿蒂亚–辛格（Atiyah–Singer）的初始工作之后，指数理论朝多个方向分枝发展。
阿蒂亚处于其中许多发展的中心，而本节只能略加触及。我们从 \S\ref{subsec:5.1} 开始，
介绍阿蒂亚–博特（Atiyah–Bott）把经典的莱夫谢茨（Lefschetz）不动点定理推广到
椭圆复形的工作。它的应用十分广泛，其中包括紧李群理论中外尔（Weyl）特征标公式的
一个几何证明。把阿蒂亚–辛格的 $K$ 理论指数框架与其先前的椭圆算子指数理论联系
起来的，是阿蒂亚–辛格定理到紧带边流形的推广；这项工作与博特合作完成，我们在
\S\ref{subsec:5.2} 中叙述。基本的阿蒂亚–辛格定理适用于复椭圆算子。对它还有一个
重要的精化（\S\ref{subsec:5.3}），即推广到实椭圆算子；这自然地联系到
阿蒂亚此时引入的 $K$ 理论的实版本。黎曼 spin 流形上存在一个\emph{典范的}实
Dirac 算子——它或许并未像它应得的那样被经常使用——因此我们借 \S\ref{subsec:5.6}
之机把它展示出来。格罗滕迪克（Grothendieck）的箴言是：在底空间上、而不只是在
一点上做几何；由此引出椭圆算子族的指数定理。正如格罗滕迪克的哲学所许诺的，这一
在 \S\ref{subsec:5.4} 中描述的定理在几何与物理中屡屡出现。最后，在
\S\ref{subsec:5.5} 中我们介绍阿蒂亚的一个指数定理，它开辟了这一课题的一个全新
分支，尽管阿蒂亚本人并未沿此分支继续。


  \subsection{等变指数定理}\label{subsec:5.1}

% [Tutor Note: 笔误——式 \eqref{eq:51} 左端求和指标为 $x\in \Fix(f)$，被加项却写作 $\epsilon _X$（下标为大写 $X$），按式 \eqref{eq:50} 应为 $\epsilon _x$。此式照录原文，引用时请按 $\epsilon _x$ 理解。]
% [Tutor Note: 术语——"带符号计数（count with signs）"指按 $\sign\det(1-df_x)$ 加权；这里的 $f$ 必须是横截的（脚注假设 $1$ 不是 $df_x$ 的本征值，即 $\det(1-df_x)\neq0$），否则须借助等变局部化或热核方法处理非横截不动点。]
回顾经典的莱夫谢茨不动点定理。设 $X$~为光滑紧致流形，$f\:X\to X$ 为具有孤立
不动点的微分同胚。此时子集 $\Fix(f)\subset X$ 有限，我们的目标或许是计算它的基数。
然而基数并非形变不变量——在形变之下，成对的不动点会诞生与消亡——因此不存在拓扑
公式。但是，若带符号计数，即对每个 $x\in \Fix(f)$ 以下式加权\footnote{在假定
对一切 $x\in \Fix(f)$，$df_x$~都没有非零不动向量的前提下}
  \begin{equation}\label{eq:50}
     \epsilon _x=\sign\det(1-df_x)\,\in \{-1,+1\}, 
  \end{equation}
则莱夫谢茨不动点定理断言
  \begin{equation}\label{eq:51}
     \sum\limits_{x\in \Fix(f)}\epsilon _X = \sum\limits_{q=0}^{\dim X}
     (-1)^q\Tr\Bigl( f^*\res{H^q(X;\RR)} \Bigr) , 
  \end{equation}
其中右端是 $f$ 作用在第 $q$ 个实上同调上的迹的交错和。

% [Tutor Note: 易错——式 \eqref{eq:52} 的分子是\emph{分级}迹 $\Tr(f^0)-\Tr(f^1)$，分母取绝对值 $|\det(1-df_x)|$；符号信息全在分子的分级迹里。核验：对 de Rham 复形 $E^q=\Omega^q$，有 $\sum_q(-1)^q\tr\Lambda^q(df_x^*)=\det(1-df_x^*)=\det(1-df_x)$，故 $\nu_x=\sign\det(1-df_x)=\epsilon_x$，与式 \eqref{eq:50}–\eqref{eq:51} 一致。若误把分母写成不带绝对值的 $\det(1-df_x)$，则比值为 $1$ 而非 $\epsilon_x$，即逐点乘上 $\epsilon_x$（它可在同一连通分支内变号），并非仅在分支上整体差一个符号。]
在 1965 年，阿蒂亚与博特对椭圆算子与椭圆复形表述并证明了一个推广。设 $X$~为
闭流形；$E^0,E^1\to X$ 为向量丛；$f\:X\to X$ 为具有孤立不动点的微分同胚；
$f^i\:E^i\to E^i$，$i=0,1$，为 $f$ 的提升；而 $P\:C^{\infty}(X,E^0)\to
C^{\infty}(X,E^1)$ 为与 $f^i$ 交换的椭圆微分算子。定义
  \begin{equation}\label{eq:52}
     \nu _x = \frac{\Tr\Bigl( f^0\res{E^0_x} \Bigr) - \Tr\Bigl(
     f^1\res{E^1_x} \Bigr) }{\bigl| \det(1-df_x) \bigr| },\qquad x\in
     \Fix(f). 
  \end{equation}

  \begin{theorem}[Atiyah-Bott~\cite{AB1}]\label{thm:9}
 我们有
  \begin{equation}\label{eq:53}
     \sum\limits_{x\in \Fix(f)}\nu _x = \Tr\Bigl(f^0\res{\ker P}\Bigr) -
     \Tr\Bigl( f^1\res{\coker P}\Bigr). 
  \end{equation}
  \end{theorem}

% [Tutor Note: 含糊——作者以"不评论证明"带过，其实质是 Atiyah–Bott 在 [AB1]（Ann. of Math. 86 (1967)）中用伪微分算子把 $P$ 局部化到不动点附近，并对分级迹作迹公式；所需技术是伪微分算子代数上的迹与象征演算。完整的证明见 [AB1]；椭圆\emph{复形}（而非单个算子）情形还需对每个 $f^i$ 逐项跟踪分级。]
\noindent 
我们在此不讨论其证明——证明也涉及伪微分算子——而只叙述若干应用。

% [Tutor Note: 术语——"等变指数"指指数作为 $R(G)$ 中的虚拟表示；$g\in G$ 给出虚拟特征标值 $\ind_g(P)\in\CC$，即"特征标值指标"，且 $g=e$ 时 $\ind_e(P)=\ind P\in\ZZ$。注意它与"指标丛"（families 情形，取值于 $K^{\bullet }(S)$）是两回事：前者底空间是一点、对称群为 $G$，后者底空间是族参数空间 $S$。]
% [Tutor Note: 含糊——"等变 $K$ 理论中的局部化定理"即 Atiyah–Segal 的局部化定理（[Seg1]，亦见 [ASeg]）：等变 $K$ 类由它在不动点集上的限制（或按法向权展开的贡献）决定。作者未写出公式；具体地，$\chi(g)=\ind_g(P)$ 可写成 $g$ 的不动点集 $\Fix(g)$ 上的求和，各项含 $\det(1-dg)$ 型的因子，这正是式 \eqref{eq:53} 的等变来源。]
% [Tutor Note: 易错——把定理~\ref{thm:9} 用于外尔特征标公式时必须取 $T$ 的\emph{一般}元素（generic element）：一般性保证不动点孤立且 $1$ 不是 $dt_x$ 的本征值，即定理~\ref{thm:9} 的横截条件成立。若 $t$ 落在环面的真子群中，不动点集非离散，须改用等变局部化公式。]
在仿照格罗滕迪克（\S\ref{subsec:4.5}）的指数定理证明中，阿蒂亚–辛格~\cite{AS2}
把\emph{紧}李群 $G$ 对紧致流形 $X$ 上椭圆算子 $P$ 的作用纳入其中。当阿蒂亚–博特
框架中的微分同胚 $f$ 生成一个（紧）微分同胚环面时，两种框架有重叠。在阿蒂亚–辛格
的情形，由于 $G$ 与 $P$ 交换，它作用在 $\ker P$ 与 $\coker P$ 上，因此形式差——
即指数——是 $G$ 的虚拟表示。于是它定义了表示环 $R(G)$ 中的一个类。另一方面，
象征 $\sigma (P)$ 也与 $G$ 交换，因而它定义了\emph{等变 $K$ 理论}中的一个类；
等变 $K$ 理论是拓扑 $K$ 理论的一个重要推广，由阿蒂亚的学生格雷姆·塞加尔
（Graeme Segal）在其博士论文~\cite{Seg1} 中研究。在等变 $K$ 理论中施行的拓扑
指数 t-ind 也给出 $R(G)$ 的一个元素。定理~\ref{thm:8} 的等变推广把这些 $R(G)$
元素等同起来。阿蒂亚–塞加尔~\cite{ASeg} 利用等变 $K$ 理论中的局部化定理
~\cite{Seg1} 导出莱夫谢茨型公式，并特别在 $f$ 生成一个紧李群（环面）微分同胚群
时重新得到定理~\ref{thm:9}。

% [Tutor Note: 术语——"$h$-配边（$h$-cobordism）"指嵌入 $W$ 满足两侧包含映射都是同伦等价；"透镜空间（lens space）"是 $S^3$ 对有限循环群自由作用的商。此处结论的强度在于把拓扑的 $h$-配边关系升级为黎曼几何的等距，所用工具是号差算子的 Lefschetz 公式，见 [AB2]。]
% [Tutor Note: 含糊——从 Lefschetz 公式到外尔特征标公式的推导被一句"leads to"带过。实质：旗流形 $G/T$ 上全纯线丛 $\mathscr{L}$ 的上同调实现 $G$ 的不可约表示（Borel–Weil），而 $t\in T$ 的 Lefschetz 数 $\sum_q(-1)^q\tr(t^*|H^q(G/T;\mathscr{L}))$ 正是该表示的特征标在 $t$ 处的值；不动点即 Weyl 群元素对应的点。详见 [AB2] 与 [AS3]。]
% [Tutor Note: 易错——"$H^q(X;\sO_X)=0$ 对 $q>0$" 的全纯映射必有不动点：这是用 Lefschetz 数等于 $1$（只有 $H^0$ 贡献）而非用拓扑不动点性质；请核对 $\sO_X$ 是结构层、$H^0(X;\sO_X)=\CC$（$X$ 连通紧致）。]
这些不同的莱夫谢茨公式有许多应用；与定理同时期实现的那些见~\cite{AB2}
与~\cite{AS3}。这里给出一个小样本。设 $X$~为连通闭复流形，且对 $q>0$ 有
$H^q(X;\sO_X)=0$；则任意全纯映射 $f\:X\to X$ 都有不动点。换一个方向，设 $G$~为
紧连通李群，$T\subset G$ 为极大环面。把定理~\ref{thm:9} 应用于 $T$ 的一个一般
元素在旗流形上的全纯线丛 $\mathscr{L}\to G/T$ 上的作用，便导致外尔特征标公式。
在紧致黎曼流形 $X$ 上，借助号差算子可以推出广义莱夫谢茨定理的若干强结论。例如，
阿蒂亚–博特–米尔诺（Atiyah–Bott–Milnor）证明两个 $h$-配边的透镜空间是等距的。
关于号差算子的莱夫谢茨公式的某些数论侧面，在阿蒂亚–博特与阿蒂亚–辛格的文中有所
描述。希策布鲁赫–扎吉尔（Hirzebruch–Zagier）~\cite{HZ} 把这些思想大大推进，他们
不仅对透镜空间、而且对射影空间、Brieskorn 簇与代数曲面探讨了这一莱夫谢茨号差
定理。他们发现了它与余切和、Dedekind $\eta$ 函数、模形式以及实二次域等数论对象
之间的深刻联系。莱夫谢茨型不变量与数论之间的这一联系，是 1970 年代诸发展的核心，
我们在 \S\ref{sec:7} 中论及。


  \subsection{带边流形上的指数定理}\label{subsec:5.2}

% [Tutor Note: 术语——"局部边界条件"指条件只依赖未知函数在边界点处的有限阶 jet：Dirichlet 条件只涉及 $0$ 阶 jet（$u(x)$），Neumann 条件涉及法向导数（$1$ 阶 jet）。与之相对的是"整体边界条件"（如 §7.2 的 APS 条件），后者涉及边界上算子的谱投影，因而不是局部的。]
% [Tutor Note: 含糊——"Lopatinski 判别法"即今称的 Lopatinski–Shapiro 条件（亦称覆盖条件）：要求边界象征在法向的零空间维数恒定，且边界象征的各因子在法向"覆盖"靶向量丛。作者只给出 [Lo]（1953）与 [Ho1]（其 §20.1 疑有误，见下条）的指引，未写条件本身；核对时请注意不同文献对"椭圆边界条件"的定义在光滑性要求上略有差异。]
% [Tutor Note: 笔误——原文 \cite[\S20.1]{Ho1} 疑误：[Ho1] 是 Hörmander 1965 年 CPAM 短文（pp. 501–517），不可能有 §20.1；疑应为 Hörmander 专著 [Ho3]（*The analysis of linear partial differential operators III*），请核对文献与节号。]
正如 \S\ref{subsec:3.3} 中所讨论的，线性椭圆方程的研究传统上在具有光滑边界的
区域 $\Omega \subset \AA^n$ 上进行，此时必须施加\emph{椭圆边界条件}。（$\AA^n$~是
标准的实仿射 $n$ 维空间。）对 $\overline{\Omega }$ 上的二阶拉普拉斯算子
~\eqref{eq:37}，\emph{Dirichlet 问题}是方程组
  \begin{equation}\label{eq:54}
     \begin{aligned} \Delta u&=0 \\ u\res{\partial \Omega }&=f\end{aligned} 
  \end{equation}
其中 $f$ 是 $\partial \Omega $ 上给定的函数。\eqref{eq:54} 中的第二个方程是
$\Delta $ 的\emph{局部}边界条件：在每一点 $x\in \partial \Omega $ 处它只依赖于
$u(x)$。更一般地，局部边界条件可以依赖于 $u$ 在 $x$ 处的有限个导数。
洛帕京斯基（Lopatinski）~\cite{Lo} 给出了局部边界条件为椭圆的一般判别法；例如
见~\cite[\S20.1]{Ho1}。

% [Tutor Note: 易错——圆盘上 $\dbar$ 的核无限维，根源是 Cauchy–Riemann 方程作为边值问题是\emph{不适定}的（Lopatinski 条件失败），而不是 $\dbar$ 不椭圆。$\dbar$ 的主象征在余切丛的非零部分是可逆的（复线性），所以它是椭圆算子；毛病完全出在边界。]
% [Tutor Note: 含糊——障碍的精确表述：作者说"$\sigma$ 到 $K^0(T^*X|_{\partial X})$ 的限制"，但该限制所在的群、以及"提升到适当的相对 $K$ 理论群"是哪个群，原文未写出。按 [AB3]（Bombay 1964 报告）与 [A2]（Seminar 附录 I）的叙述，障碍 $=\sigma$ 的边界限制类 $\in K^0(T^*X|_{\partial X})$；该障碍消失时可提升到相对群 $K^0(T^*X,T^*X|_{\partial X})$，其消失当且仅当存在局部椭圆边界条件。审校时请核对 [AB3] 的定义。]
% [Tutor Note: 含糊——"Atiyah–Bott–Singer" 并无独立的三人论文：上文引的是 [AB3]（Atiyah–Bott，Bombay 报告）与 [A2]（Atiyah，Seminar 附录 I），可理解为二者的合并称呼；经典的 K 理论证明见 [AS2]。请勿据三人合名另找文献。]
% [Tutor Note: 术语——此处 "Bott 周期性定理" 指复 $K$ 理论的 $2$ 周期性。Atiyah–Bott 用带边流形的相对 $K$ 群给出的新证明把周期性化归为边界椭圆算子的指数定理，见 [AB4]（Acta Math. 112 (1964)）。]
% [Tutor Note: 术语——本小节只处理\emph{局部}椭圆边界条件：此时拓扑指数可无修正项地计算解析指数。若改用\emph{整体}边界条件（如 §7.2 的 APS 条件），指标公式须加上边界修正项（$\eta$ 不变量，其几何原型即号差亏损 signature defect），见 §7.3 与 [APS1]–[APS3]；两种边界问题不可混为一谈。]
并非所有椭圆算子都容许局部椭圆边界条件。考虑单位圆盘 $\Omega \subset \CC$ 上
函数 $u\:\overline{\Omega }\to\CC$ 的一阶 $\dbar$ 算子。$\dbar$ 的核由
$\overline\Omega $ 上的全纯函数组成，这是一个无限维向量空间，正如在缺少椭圆边界
条件时所预期的那样。对该算子不存在\emph{局部}椭圆边界条件：圆盘上的
Cauchy–Riemann 方程是不适定的。（正如我们将在 \S\ref{subsec:7.2} 中看到的，存在
\emph{整体}边界条件。此外，若考虑向量丛截面上的 $\dbar$ 算子——例如在研究辛流形
中边界落在一个拉格朗日（Lagrangian）子流形上的全纯圆盘的形变时——则确实存在
局部边界条件。）局部椭圆边界条件存在性问题的拓扑本质，在阿蒂亚–博特
~\cite{AB3,A2} 的工作中被揭示出来。$K$ 理论再次成为障碍的自然栖身之所。事实上，
他们的工作导致了 Bott 周期性定理~\cite{AB4} 的一个新证明。设 $X$~为紧带边流形，
并设 $\sigma \in K^0(T^*X,\,T^*X\setminus 0)$ 为椭圆象征类。局部椭圆边界条件的
障碍是 $\sigma $ 到~$K^0(T^*X\res{\partial X})$ 的限制；若该障碍消失，则局部椭圆
边界条件把 $\sigma $ 提升到某个适当的相对 $K$ 理论群。阿蒂亚–博特–辛格
（Atiyah–Bott–Singer）证明了带局部椭圆边界条件的椭圆算子的指数定理。拓扑指数
容易延拓到提升后的相对椭圆象征，而该提升象征的拓扑指数计算了椭圆边值问题的
解析指数。


  \subsection{实椭圆算子}\label{subsec:5.3}

% [Tutor Note: 易错——反自伴 Fredholm 算子的整数指数必为零，因为 $\ker P^*=\ker P$ 且 $\coker P\cong\ker P^*$，正负指标相消；因此这里只能取模 $2$ 指标。这不是笔误。另注意：$D$（周期）与 $D'$（反周期）的核分别是 $1$ 维与 $0$ 维，故其模 $2$ 指标分别为 $1$ 与 $0$，恰好区分 $S^1$ 的两个 spin 结构。]
% [Tutor Note: 术语——$\cir$ 上恰有两个 spin 结构：周期（Lie 群）结构与反周期（扭化、bounding）结构；相应的旋量丛截面可等同于周期或反周期函数。见 \S\ref{subsec:4.1} 与 [LM] I.3。学者常说的"两个 spin 结构"与 $\zt$ 指标的对应，正是 $KO^{-1}(\pt)\cong\zt$ 的几何体现。]
考虑作用在 $\RR/2\pi \ZZ$ 上实值函数、即 $2\pi $ 周期函数 $u\:\RR\to\RR$ 上的
算子 $D=d/dx$。分部积分表明 $D$~是（形式上的）反自伴算子。它是一个\emph{实}
Dirac 算子的例子。一个密切相关的 Dirac 算子是 $D'=d/dx$，它作用在圆上的扭化
函数上；等价地说，它作用在满足对一切 $x\in \RR$ 有 $u(x+2\pi )=-u(x)$ 的函数
$u\:\RR\to\RR$ 上。注意 $\ker D\cong \RR$ 由常数函数组成，而 $\ker D'=0$。
反自伴算子在~\eqref{eq:30} 的意义下的整数指数为零。但对反自伴 Fredholm 算子~$P$，
模~$2$ 维数 $\dim \ker P\pmod2$ 是形变不变量：反自伴性意味着非零本征值在 $P$~
形变时成对移动，因此特别地 $\dim\ker P$ 只能按偶数跳跃。\footnote{Fredholm 算子
的谱未必离散；但关于核的断言仍然成立。} 对圆上的 Dirac 算子~$D,D'$，这一
\emph{模~$2$ 指数}区分了两个 spin 结构，我们在此已用具体的形式把它们表示出来。

% [Tutor Note: 术语——大写的 "Real $K$ 理论"（[A3]，*K-theory and reality*, 1966）指 Atiyah 的 KR 理论：带对合 $\tau$ 的空间上的 $K$ 理论，其元素是"实"丛在对合下的复共轭配对；不要与实 $K$ 理论 $KO$ 混为一谈。本节的实指标定理最终用 $KO$ 表述（见式 \eqref{eq:81}），KR 只是导出它的框架；二者的关系见 [A3]。]
% [Tutor Note: 易错——"四维实 Dirac 算子的自然整数不变量是整数指数的一半"：复指标为 $\ind_\CC D=\Ahat=-\sigma/8$；实 Dirac 算子的 $KO^{-4}\cong\ZZ$ 指标与此相差因子（源自 $\Cliff_{-4}\cong\HH(2)$ 的模维数），故其自然整数不变量为 $\Ahat/2$。于是"该不变量为整数"等价于 $\Ahat$ 为偶数，即 $\sigma/8$ 为偶数，即 $16\mid\sigma$，这正是 Rokhlin 定理。核对 [AS5] 与 [LM] II.4 的归一化，别把因子 $2$ 弄丢。]
\emph{实}椭圆算子 $P$ 在 $\RR^n$ 上的象征 $p(x,\xi )$，作为 $\xi \in {\RR^n}^*$
的函数，本质上是某个实值函数的 Fourier 变换；见~\eqref{eq:46}。因此 $p$~是满足
如下"实性条件"的复值函数
  \begin{equation}\label{eq:55}
     p(x,-\xi )=\overline{p(x,\xi )}.
  \end{equation}
流形 $X$ 上的整体象征 $\sigma (P)$ 是复向量丛间的映射~\eqref{eq:29}，而实性条件
~\eqref{eq:55} 整体化为这些向量丛之间的一个复共轭同构，它覆盖 $T^*X$ 上的
$(-1)$ 对合并与 $\sigma (P)$ 交换。这一思想促使阿蒂亚发展了 Real $K$ 理论
~\cite{A3}；见~\cite[\S8]{Seg4}。在此基础上，阿蒂亚–辛格~\cite{AS5} 证明了实椭圆
算子的指数定理。该指数定理涵盖模~$2$ 指标（如上文提到的 $\cir$ 上的不变量）以及
其他情形。例如，四维实 Dirac 算子的自然整数不变量是整数指数的一半，由此得到
Rokhlin 定理~\cite{Roh} 的解析证明：闭 spin 四维流形的号差可被~$16$ 整除。这些
实指数定理由 Clifford 代数得到阐明，如 \S\ref{subsec:5.6} 所述。

% [Tutor Note: 含糊——作者自陈"对模 $2$ 指标不知道其他证明"，这凸显了 $KO$ 理论方法在 $\zt$ 指标上的不可替代性；热方程/$\eta$ 不变量方法（§6、§7）只直接给出整指标（或 $\eta$ 的实值修正），模 $2$ 指标需额外的拓扑输入。引用时请勿把这句话误读为"模 $2$ 指标没有证明"。]
  \begin{remark}[]\label{thm:10}
 按设计，$K$ 理论证明（\S\ref{subsec:4.4}）适合证明定理~\ref{thm:8} 的诸推广，
如实指数定理。对模~$2$ 指标，我不知道其他证明；而对整数指标，还有其他途径，
例如我们在 \S\ref{sec:6} 中综述的热方程方法。
  \end{remark}


  \subsection{Clifford 线性算子的指数定理}\label{subsec:5.6}

% [Tutor Note: 笔误——本小节在文档中排第 4 个（按编号应为 5.4），却带 label subsec:5.6；随后"族指数定理"排第 5 个（应为 5.5）而带 label subsec:5.4，"覆盖与 von Neumann 代数"排第 6 个（应为 5.6）而带 label subsec:5.5：三个 label 名与位置错位。LaTeX 按位置编号，故交叉引用仍指向正确小节、印出的编号也正确，问题仅在 label 名与位置不符；为不破坏全篇引用，此处照录源码 label 未改。]
% [Tutor Note: 易错——式 \eqref{eq:79} 的 $\pm$ 约定决定了 $\Cliff_{\pm n}$ 的符号：$\Cliff_{+n}$ 中 $\gamma_i^2=+1$，$\Cliff_{-n}$ 中 $\gamma_i^2=-1$。Bott 周期表（$n\bmod 8$）与 $KO^{-n}$ 的取值都依赖这一约定；若把两个符号对调，$\Cliff_{\pm n}$ 的矩阵代数类型会移位，周期表随之整体平移。核对 [ABS]（*Clifford modules*, Topology 3 (1964)）与 [LM] I.4 的约定。]
% [Tutor Note: 术语——"混合构造（mixing construction）"即纤维丛的相伴丛构造 $P\times_G V$：用主 $\Spin_n$ 丛 $\Spin(X)$ 把表示 $\Cnp$ 扭化为向量丛。它不是新的几何操作，只是把表示论数据升级为丛的标准手段。]
% [Tutor Note: 易错——"$\Cnm$ 是 $\Cnp$ 的反代数"这一步用的是分级反代数同构 $\Cliff_{-n}\cong(\Cliff_{+n})^{\text{graded op}}$（Koszul 符号反乘法 $a\cdot b=(-1)^{|a||b|}ba$；复化后也可用 $\gamma ^i\mapsto i\gamma ^i$ 实现，参见 [LM] I.1），因此右 $\Cnp$ 模自动成为左 $\Cnm$ 模。这正是"Clifford 线性"的含义：算子与左 $\Cnm$ 作用分级交换，而单个 $\Spin_n$ 表示被整套 Clifford 生成元的作用取代。]
设 $X$~为 $n$ 维黎曼 spin 流形，$\Spin(X)\to X$ 为刻画该黎曼 spin 结构的
主 $\Spin_n$ 丛；见 \S\ref{subsec:4.1}。回顾 Clifford 代数~$\Clpm$，其生成元为
$\gamma ^1,\dots ,\gamma ^n$，关系为
  \begin{equation}\label{eq:79}
     \gamma ^i\gamma ^j + \gamma ^j\gamma ^i = \pm 2\delta ^{ij} = \begin{cases}
     \pm 2,&i=j;\\\phantom{\pm}0,&i\neq j,\end{cases}\qquad 1\le i,j\le n. 
  \end{equation}
用"混合构造"（mixing construction）构成向量丛
  \begin{equation}\label{eq:80}
     S_X=\Spin(X)\times \mstrut _{\Spin_n}\Cnp\longrightarrow X, 
  \end{equation}
其中 $\Spin_n\subset \Cnp$ 通过左乘作用在~$\Cnp$ 上。$S_X\to X$ 的纤维是
$\zt$ 分级的\emph{右} $\Cnp$ 模，因为右乘与左乘交换。于是，由于 $\Cnm$~是
$\Cnp$ 的（$\zt$ 分级）反代数，$S_X\to X$ 的纤维是 $\zt$ 分级的\emph{左}
$\Cnm$ 模。此外，$S_X\to X$ 从 Levi-Civita 联络继承了一个协变导数。Dirac 算子
~\eqref{eq:41} 作用在 $S_X\to X$ 的截面上，其中我们用嵌入
$\RR^n\hookrightarrow \Cnp$（把 $\RR^n$ 实现为 $\gamma ^1,\dots ,\gamma ^n$
张成的子空间）来定义 Clifford 乘法。于是 $D_X$~是（在适当 Sobolev 空间上的）反
自伴 Fredholm 算子，它关于该分级为奇的，并与左 $\Cnm$ 作用分级交换。

% [Tutor Note: 术语——第 2 条"旋量表示只确定到与线丛张量积"（原文只在 §4.1 提过）：因为 $\Spin_n$ 的中心（$n$ 偶时为 $\zt$）与复化后矩阵代数的自同构允许把 $\SS^{\pm}$ 与一维表示作张量；故不能单靠表示论定义典范旋量丛，须额外选 spin 结构。$\Cliff$ 模构造的优点是：$\Cnp$ 作为 $\Spin_n$ 模虽是可约的，但带有典范的 Clifford 生成元作用，因此"典范"到无需选择。]
% [Tutor Note: 术语——"Clifford 线性"意为算子与左 $\Cnm$ 模结构相容（分级交换），故核与余核都是 $\Cnm$ 模；因此指标落在 $\Cliff$ 模的 Grothendieck 群中，这由 [ABS] 等同于 $KO^{-n}(\pt)$。这一"以模代替表示"的技巧是 Atiyah–Singer III 之后的标准语言，见 [LM] II.7。]
  \begin{remark}[]\label{thm:14}
\ 

 \begin{enumerate}%[label=\textnormal{(\roman*)}]
 \item 这一\emph{Clifford 线性 Dirac 算子}出现在~\cite[\S
II.7]{LM} 中，归功于阿蒂亚–辛格。

 \item 如同任何不可约表示，\S\ref{subsec:4.1} 中的旋量表示~$\SS^{0},\SS^1$ 只
确定到与一个线丛作张量积，那里已经指出过这一点。相比之下，$S_X\to X$ 与~$D_X$
的构造对黎曼 spin 流形是典范的。我们用典范表示~$\Cnp$ 连同其上一个典范的交换
算子代数，换掉了不可约表示~$\SS^0,\SS^1$。

 \item 可以把 $\Cnp$ 复化，从而把 $S_X\to X$ 复化，以构造一个典范的复 Dirac
算子；但实 Dirac 算子含有更精细的信息。

 \item 若 $V\to X$ 是带协变导数的实向量丛，我们可以构成 $S_X\otimes V\to X$
截面上的扭化 Dirac 算子~$D_X(V)$；它也与左 $\Cnm$ 作用交换。

 \end{enumerate}
  \end{remark}

% [Tutor Note: 术语——"解析前推"与"拓扑前推"：$f_!$ 指沿 $f$ 的 umkehr（Gysin）映射；实 $K$ 理论中 $f\:X\to\pt$ 的定向由一个 spin 结构给出（$M\Spin$ 定向），见式 \eqref{eq:19} 的复情形类比。不要把这里的 $f_!$ 与前推 $f_*$（同调/上同调）混淆。]
% [Tutor Note: 易错——式 \eqref{eq:81} 的周期表请按 $n\bmod 8$ 核对：$KO^{-0}=\ZZ$，$KO^{-1}=\zt$，$KO^{-2}=\zt$，$KO^{-3}=0$，$KO^{-4}=\ZZ$，$KO^{-5}=0$，$KO^{-6}=0$，$KO^{-7}=0$。表中 $n\equiv 1,2$ 给 $\zt$ 而 $n\equiv 0,4$ 给 $\ZZ$，这对应 $\pi _{n-1}O_\infty$ 的取值（$\pi_0O=\zt$，$\pi _1O=\zt$，$\pi _3O=\ZZ$，$\pi _7O=\ZZ$）。注意表只对 $n\ge1$ 陈述，$n=0$ 需另行补上 $KO^0(\pt)\cong\ZZ$。]
% [Tutor Note: 易错——"$n=1$ 的模 $2$ 指标是特例"：$KO^{-1}(\pt)\cong\zt$，由 $\dim\ker D\bmod2$ 给出；对 $S^1$ 的两个 spin 结构分别给 $1$ 与 $0$（见 \S\ref{subsec:5.3}），正与表中 $n\equiv1$ 的 $\zt$ 相符。]
$D_X(V)$ 的核与余核是有限维 $\Cnm$ 模，因此由~\cite{ABS}，它们的形式差——即
指数——表示实 $K$ 理论群~$KO^{-n}(\pt)$ 中的一个元素。这就是 $V$ 的\emph{解析
前推}。该 spin 结构为实 $K$ 理论定向了映射 $f\:X\to \pt$——复 $K$ 理论中的类似
情形见~\eqref{eq:19}——而\emph{拓扑前推}是 $f_![V]\in KO^{-n}(\pt)$。Dirac 算子
的指数定理断言这些前推相等。点的 $KO$ 群由 Bott 周期性算出。设 $O_\infty $~为
稳定正交群。则对~$n\ge1$，有
  \begin{equation}\label{eq:81}
     KO^{-n}(\pt)\cong \pi _{n-1}O_\infty \cong \begin{cases} \ZZ,&n\equiv
     0\pmod8;\\ \zt,&n\equiv 1\pmod8;\\ \zt,&n\equiv 2\pmod8;\\ 0,&n\equiv
     3\pmod8;\\ \ZZ,&n\equiv 4\pmod8;\\ 0,&n\equiv 5\pmod8;\\ 0,&n\equiv
     6\pmod8;\\ 0,&n\equiv7\pmod8.\end{cases} 
  \end{equation}
（$KO^0(\pt)\cong \ZZ$ 也是对的。）\S\ref{subsec:5.3} 中讨论的 $n=1$ 的模~$2$
指标是一个特例。

% [Tutor Note: 术语——此处的"谱（spectrum）"是同伦论意义下的谱（$\Omega$-谱），与算子的谱毫无关系。同一节里两种"谱"并存是易混点：$\Fn(\sH_n)$ 构成的谱代表 $KO$ 或 $K$ 的 $\Omega$-谱，而"算子的谱"指本征值集。]
% [Tutor Note: 易错——$\Fn(\sH_n)$ 的定义要求算子\emph{奇}且反自伴；去掉"奇"或去掉"反自伴"都会得到不同的分类空间（例如自伴情形对应 $KR$ 的其他分量）。这里 $\Cn$ 作用与"odd"条件共同保证 $\Fred$ 恰是 $KO^{-n}$（或 $K^{-n}$）的分类空间；详见 [AS6] 定理 A。]
在 1969 年写成的~\cite{AS6} 中，阿蒂亚–辛格研究了空间~$\Fn(\sH_n)$，$n\in \ZZ$，
其元素是作用在带有 $\Cn$~作用的实或复 $\zt$ 分级希尔伯特空间~$\sH_n$ 上的奇
反自伴 Fredholm 算子；这些算子与 $\Cn$ 作用分级交换。Clifford 线性 Dirac 算子
是一个例子。主要结果是：空间~$\Fn(\sH_n)$ 在同伦论的意义下构成一个谱，而
阿蒂亚–辛格在实与复两种情形下分别把它等同于 $KO$-~与 $K$ 理论的谱。这推广了
$K$ 理论与 Fredholm 算子之间的关系，后者出现在 Jänich~\cite{J} 的早期工作中。
它也给出了 Clifford 线性 Dirac 算子解析指数的另一种构造。最后，Clifford 代数的
周期性~\cite{ABS} 导致了 Bott 周期性的又一个证明；另见~\cite{Kar,Wo}。


  \subsection{族指数定理}\label{subsec:5.4}

% [Tutor Note: 术语——"逆紧（proper）纤维丛"指纤维紧致、$f$ 逆紧；族指数定理要求纤维紧致以保证逐纤维的 Fredholm 性，而底空间 $S$ 可以非紧（此时 $K^{\bullet }(S)$ 用紧支撑或 $K$ 理论的一般定义）。不要把 proper 与"正常"或"真"混淆。]
% [Tutor Note: 术语——"族指数（family index）"与"指标丛（index bundle）"：族指标是 $K^{\bullet }(S)$ 中的一个类；当 $\ker$、$\coker$ 的维数沿 $S$ 跳跃时它一般不是真实向量丛，而是以 $\ker-\coker$ 为纤维的\emph{虚拟}丛（K 理论类）。只有在维数局部常定的开集上才能谈论具体的指标丛；这正是必须用 $K$ 理论而非普通丛来陈述族指数定理的原因。]
% [Tutor Note: 易错——"Fredholm 算子构成 $K$ 理论的分类空间"对复 $K$ 理论给 $\Fred\simeq \ZZ\times BU$（$K^0$），带分级的奇算子给 $K^{-1}$；因此族指标取值于 $K^{\bullet }(S)$ 而非单个 $K^0(S)$。实情形对应 $KO$ 与 8 周期，见 [AS6] 与 [J]。]
回顾格罗滕迪克在其黎曼–罗赫定理~\ref{thm:4} 中考虑的是逆紧态射 $f\:X\to S$，
而不只是单个簇 $X\to\pt$。类似地，阿蒂亚–辛格~\cite{AS4} 把定理~\ref{thm:8} 推广
到逆紧纤维丛 $f\:X\to S$，其上带有沿纤维的椭圆（伪）微分算子族~$P$。在适当的
函数空间上，$P$~是由 $S$ 参数化的连续 Fredholm 算子族，而解析指数是这个族的
同伦类。由于 Fredholm 算子构成 $K$ 理论的一个分类空间~\cite{J,A4,AS6}，解析指数
落在~$K^{\bullet }(S)$ 中。在纤维丛~$f$ 中施行的拓扑指数构造也给出
$K^{\bullet }(S)$ 的一个元素。\emph{族指数定理}就是解析指数与拓扑指数的相等。

\S\ref{subsec:5.6} 的 Clifford 线性故事也有族版本，我们加以说明。

% [Tutor Note: 易错——纤维取\emph{奇}维是关键的：复 Bott 周期性给出 $K^{-n}(S)\cong K^{-1}(S)$（$n$ 奇），而 $K^{-1}(S)=[S,U(\infty)]$ 经 $\det:U(\infty)\to U(1)$ 映到 $H^1(S;\ZZ)\cong[S,S^1]$，这才有"行列式映射"。若 $n$ 为偶数，最低阶分量将是 $K^0$ 的秩（映到 $H^0$），而不是行列式。]
% [Tutor Note: 含糊——"相对黎曼 spin 结构"的定义放在脚注里，但族 Dirac 算子的自伴/反自伴性与 Fredholm 性需要"相对黎曼结构"（内积 + 水平分布）才能逐纤维定义；作者未给出细节。可查 [AS4]（*The index of elliptic operators IV*）与 [APS3] §3（自伴表述）。]
% [Tutor Note: 易错——"$H^1(S;\ZZ)$ 无挠"才使"由在 $\cir$ 上的周期决定类"成立：一般地 $H^1(S;\ZZ)\cong\Hom(H_1(S),\ZZ)$ 确实无挠，故周期（对 $S^1$ 的配对）足以确定类；但对 $H^2$ 一般有挠部分，不能用 de Rham/周期方法决定。基变换（base change）即沿 $\varphi $ 拉回纤维丛与向量丛。]
% [Tutor Note: 术语——"最低阶片"指 $K$ 理论类在 Atiyah–Hirzebruch 谱序列中的最低非零分量；行列式映射 $K^{-1}(S)\to H^1(S;\ZZ)$ 由 $\det:U(\infty)\to U(1)$ 诱导，把 $K$ 类送到 $U(1)$-类（即 $S\to S^1$ 的绕数/相位；等价于 $S\times S^1$ 上行列式线丛的第一陈类），勿与 $K^0\to H^2$ 的第一陈类混同。注意它\emph{不是}单射：核包含挠信息与高阶信息。]
  \begin{example}[]\label{thm:24}
 设 $f\:X\to S$ 为逆紧纤维丛，纤维的维数为奇数~$n$。假设还给定一个相对的
黎曼 spin 结构\footnote{\emph{相对 spin 结构}是竖直切丛 $T(X/S)\to X$ 上的一个
spin 结构。\emph{相对黎曼结构}是该向量丛上的一个内积，连同纤维丛~ $f$ 上的一个
水平分布。}。在\emph{复} $K$ 理论中有前推
  \begin{equation}\label{eq:82}
     f_!\:K^0(X)\longrightarrow K^{-n}(S).
  \end{equation}
若 $V\to X$ 是带协变导数的复向量丛，我们构成由 $S$ 参数化的复 Clifford 线性
Dirac 算子族~$D_{X/S}(V)$。解析指数是由这些 Dirac 算子给出的映射
$S\to\Fred_{-n}$ 的同伦类，其中 $\Fred_{-n}$ 是 \S\ref{subsec:5.6} 末尾引入的
Fredholm 算子空间，它是 $K^{-n}$ 的分类空间。指数定理断言解析指数等于
$f_![V]$。指数的"最低阶片"（lowest piece）由与一个自然映射
  \begin{equation}\label{eq:83}
     K^{-n}(S)\longrightarrow H^1(S;\ZZ), 
  \end{equation}
的复合给出，该映射可视为行列式映射。回顾 $H^1(S;\ZZ)$ 中的类由它在映射
$\varphi \:\cir\to S$ 上的周期决定，换言之由它在 $\varphi _*[\cir]\in H_1(S)$
上的值决定（因为 $H^1(S;\ZZ)$ 无挠）；于是我们通过 $\varphi $ 作基变换，得到带
向量丛 $V_\varphi \to X_\varphi $ 的纤维丛 $X_\varphi \to\cir$，以此计算
$\ind D_{X/S}(V)$ 在~\eqref{eq:83} 下的像。此时 $X_\varphi $~是偶数维的闭黎曼
spin 流形，而我们所要的整数周期就是数值指数 $\ind D\mstrut _{X_\varphi
}(V_\varphi )\in \ZZ$。
  \end{example}

% [Tutor Note: 术语——$\RZ$ 在此是 Eilenberg–MacLane 空间 $K(\ZZ,1)\simeq S^1$，故 $H^1(S;\ZZ)\cong[S,\RZ]$；把同伦类"提升为具体映射"指 $\eta$ 不变量给出规范代表元（而不仅是同伦类），这正是"几何不变量"的含义。]
  \begin{remark}[]\label{thm:15}
 \ 
 \begin{enumerate}%[label=\textnormal{(\roman*)}]

 \item Clifford 线性 Dirac 算子族 $D_{X/S}(V)$ 给出无分级希尔伯特空间上的一族
自伴 Dirac 算子，这种表述出现得更为频繁；例如见 \cite[\S3]{APS3}。

 \item $H^1(S;\ZZ)$ 中的一个上同调类是映射 $S\to
\RZ$ 的同伦类。族~$D_{X/S}(V)$ 的一个\emph{几何}不变量——Atiyah–Patodi–Singer
$\eta $-不变量——把这个同伦类提升为一个具体的映射，正如我们将在
\S\ref{subsec:7.3} 中看到的。

 \item 尽管我们是对 Dirac 算子陈述这些拓扑构造的，它们以多种方式推广到椭圆伪
微分算子族。

 \end{enumerate}
  \end{remark}

% [Tutor Note: 术语——"模空间，或栈"：模空间 $\sM=F\inv (c)/\sG$ 一般是奇异的（当 $\sG$ 作用有稳定化子时），"栈"指保留稳定化子信息的精细版本；此时"虚拟切丛"由族的指数给出（形变复形的上同调），在基点处等于 $\ker dF_b-\coker dF_b$。]
% [Tutor Note: 易错——原文 "its linearizations $dF_b$ at $b\in F\inv (c)$ fit together to a family ... parametrized by $\sB$"：严格地说族应由 $F\inv (c)/\sG$ 参数化，且只有当 $F$ 在 $\sG$-轨道横截方向上 Fredholm 时指标才有定义；作者用 $\sB$ 参数化是为简化。审校时可保留原文措辞，但须知这里隐含了"模 $\sG$ 的 Fredholm"假设。]
遇到\emph{线性}椭圆方程族的一种情形，是对\emph{非线性}椭圆方程的解作线性化。
抽象的设定是无限维希尔伯特流形之间的非线性 Fredholm 映射 $F\:\sB\to\sC$。
有时 $F$~关于某个无限维李群~$\sG$ 的作用是等变的，且只是模~$\sG$ 的 Fredholm
映射。此时人们感兴趣的是对某个 $c\in \sC$ 的模空间（或栈）~$\sM=F\inv (c)/\sG$。
若 $F$~是闭流形上的非线性椭圆算子、在函数空间之间映射，则它在
$b\in F\inv (c)$ 处的线性化~$dF_b$ 拼成由~$\sB$ 参数化的线性椭圆算子族。该族的
指数计算 $\sM$ 的（虚拟）切丛，并给出关于~$\sM$ 的有用信息。这一总体方案被
阿蒂亚–希钦–辛格（Atiyah–Hitchin–Singer）~\cite{AHS} 用来研究四维流形上的瞬子
方程（见本卷中 Donaldson 的文章~\cite[\S2]{Do}），此后又被用于几何分析中的许多
其他问题。

% [Tutor Note: 含糊——"\eqref{eq:12} 的拓扑类比"指族情形的 Riemann–Roch 型公式 $\ch(f_!V)=f_*(\ch V\cdot \Todd(T(X/S)))$ 之类；作者未写出公式。热方程方法能算出 $\ch$ 像是因为热核的局部展开给出示性形式的积分，而要恢复 $K$ 理论中的挠信息（如指标丛本身）则超出热核方法的能力，见 [Bi2]。]
  \begin{remark}[]\label{thm:11}
 族指数在陈特征 $\ch\:K^{\bullet
}(S)\to H^{\bullet }(S;\QQ)$ 下的像，是比 $K$ 理论指数更粗糙的不变量，但往往
含有令人感兴趣的信息。它可由~\eqref{eq:12} 的一个拓扑类比来计算，并且可以借助
热方程方法~\cite{Bi2} 得到；而据我们所知，更强的 $K$ 理论指数则不能。
  \end{remark}

阿蒂亚–希钦–辛格的工作是族指数定理在量子场论与弦论、以及由这些物理产生的
数学问题中众多应用的第一例。我们在 \S\ref{sec:8} 中再举一个例子。


  \subsection{覆盖与 von Neumann 代数}\label{subsec:5.5}

% [Tutor Note: 术语——"II 型 von Neumann 代数"（type II von Neumann algebra）的关键特征是存在实值的维数函数（Murray–von Neumann 维数），可以取连续值；这正是把指数推广到实数的代数基础。与 $\ssA$ 上的正规迹 $\tau$ 配合，投影的 $\Gamma$-维数 $\dim_\Gamma=\tau(p)$ 就是 $L^2$ 意义下的"维数"。]
至此应当已经清楚，阿蒂亚把指数定理当作向多个方向展开数学探索的发射台。其中
特别富有成果的一个方向涉及冯·诺依曼（von Neumann）代数~\cite{A5}。
阿蒂亚~\cite{A1} 评论道：

  \begin{quote}
 我尤其从辛格那里（他在泛函分析方面有很强的背景）得知了 II 型 von Neumann 代数
及其奇特的实值维数。我们意识到 $K$ 理论与指数理论可以朝这个方向推广，但起初
并不清楚这样的推广是否真的会有意义。然而在一个特别简单的场合，即基本群为无限
群的流形，von Neumann 代数的想法变得相当自然，并导致了具体的非平凡结果。这就是我
在纪念昂利·嘉当（Henri Cartan）的会议上的报告~\cite{A5} 的内容。由于我并非
von Neumann 代数方面的专家，我在这次报告中试图对这些结果给出一个简单、初等且
基本自足的论述。后来，在这一领域的权威阿兰·孔涅（Alain Connes）手中，这些简单
的想法被极大地推广，发展成关于叶状结构的整个线性分析理论。
  \end{quote}

\noindent 
而在孔涅、卡什帕罗夫（Kasparov）以及许多其他人的手中，这些想法则发展成
$C^*$ 代数上的指数理论与 $K$ 理论。这确实是一次有影响力的会议报告！

% [Tutor Note: 易错——von Neumann 迹的归一化是本节最容易出错处：$\Gamma$ 无限时 $\ssA$ 为 II$_\infty$ 型，$\tau(1)=\infty$，不存在 $\tau(1)=1$ 的迹；应取 $\tau$ 使 $\dim_\Gamma(L^2(\Gamma))=1$（即把基本域维数归一化为 $1$），此时 $\dim_\Gamma(\ker\tD)=\tau(p)$ 是"每个基本域的维数"。若不归一化，定理会差一个体积因子，$\ind_\Gamma\tD=\ind D$ 不再成立。核对 [A5] 的约定。]
% [Tutor Note: 术语——"$L^2$ 指数"与"$\Gamma$ 维数"：对投影 $p$（到 $\ker\tD$ 的正交投影），$\dim_\Gamma(\ker\tD)=\tau(p)\in[0,\infty)$ 可为非整数；这就是 $L^2$-index。它与普通指数 $\ind D\in\ZZ$ 的区别只在归一化与适用范围（非紧覆盖上），定理断言两者数值相等。]
% [Tutor Note: 术语——"无分叉 Galois 覆盖"：拓扑覆盖自动无分叉，"Galois"即正规覆盖，其 deck 变换群 $\Gamma=\pi_1(X)/\pi_1(\tX)$ 自由作用，"unramified"是为区别于代数几何中带分歧的覆盖。定理所需的是 $\Gamma$ 自由作用且 $X=\tX/\Gamma$ 紧致。]
% [Tutor Note: 易错——定理~\ref{thm:12} 断言 $\ind_\Gamma\tD=\ind D$：左端按基本域计数、右端是 $X$ 上的整数指数。等号成立依赖 $\Gamma$ 自由作用与 $X$ 紧致；deck 群总是自由作用，挠元同样可以自由作用（如 $\ZZ_2$ 对径作用于 $S^n$），故"$\Gamma$ 有挠元"并不等于"作用不自由"。若作用确有不动点，须改用轨道叶状空间上的指标理论（Connes 的叶状结构理论），见 [A5] 及后续文献。]
\cite{A5} 中的情形是一个无分叉 Galois 覆盖 $\pi \:\tX\to X$，其 Galois 群
$\Gamma $ 自由地作用在~$\tX$ 上，且 $\tX$ 上有一个 $\Gamma $ 不变的椭圆算子
$\tD$。$X$ 上有诱导的椭圆算子~$D$，我们假定 $X$~紧致。例如，$X$~可以是闭黎曼
曲面，$\tX$~是它的万有覆盖；若 $X$ 的亏格~$\ge2$，则 $\tX$~同构于单位圆盘
$\Omega $。若 $\tD$~是 $\dbar$ 算子，则 $\ker\tD$~是 $\Omega $ 上全纯函数的
\emph{无限维}空间。一般地，若 $\Gamma $~无限，则 $\ker\tD$ 与 $\coker\tD$~是
无限维的，而 $\ker D$~与 $\coker D$~是有限维的——后者因为假定 $X$~紧致。因此
$\ind D$~是良定义的。阿蒂亚在 $\tX$ 上引入一个 $\Gamma $ 不变的测度，以及
$L^2(\tX,\tE^0)$ 上有界线性算子的 von Neumann 代数~$\ssA$，其中 $\tE^0\to\tX$ 是
使 $\tD$~定义在其截面上的向量丛。到 $\ker\tD$ 的正交投影属于该 von Neumann 代数，
其 von Neumann 迹是一个实数，被定义为 $\ker\tD$ 的 $\Gamma $~维数。对 $\coker\tD$
重复此过程，阿蒂亚定义了一个实值指数 $\ind_\Gamma \tD$。

  \begin{theorem}[Atiyah~\cite{A5}]\label{thm:12}
 $\ind_\Gamma \tD = \ind D$. 
  \end{theorem}

\noindent
如前所述，阿蒂亚对这一结果的阐述，成为非紧空间、奇异空间以及更广范围内指标
理论的催化剂。

% ==================== §6 ====================

\section{热方程证明}\label{sec:6}

\begin{keytakeaways}
  \item 1960 年代末起，由阿蒂亚–辛格指数定理生发出来的思想圈不断
    扩大，研究重心从椭圆算子的核转向更高的特征值；这先引出指数定理的
    \emph{局部}版本，再引出局部的\emph{几何}不变量。
  \item 把全部特征值汇总为单一函数的两个基本工具，是 $\zeta$ 函数
    $\zD(s)=\Tr\Delta^{-s}$ 与热算子 $H_t=e^{-t\Delta}$ 的迹；二者由
    梅林变换相联系。
  \item 对闭流形上的非负自伴二阶椭圆算子，热核当 $t\to0$ 时在对角线上
    有渐近展开 $h_t(x,x)\sim t^{-n/2}\sum_{k}A_k(x)t^k$，其中 $A_k$ 为
    Seeley–DeWitt 系数，只依赖象征在一点的有限阶射流。
  \item 由 $P^*P$ 与 $PP^*$ 的非零谱一致，得 McKean–Singer 公式
    $\ind P=\Tr e^{-tP^*P}-\Tr e^{-tPP^*}$（对一切 $t>0$）：取
    $t\to\infty$ 给出 $\dim\ker P^*P-\dim\ker PP^*$，取 $t\to0$ 给出
    局部指标密度。
  \item 局部指数定理把 $t\to0$ 的极限识别为一个示性形式；关键输入是
    Gilkey 的相消定理：自然、正则、非负权齐次的微分形式必为
    Pontryagin 类的 Chern–Weil 形式的多项式。
  \item 1980 年代的新证明（Getzler 的标度论证、Witten 与 Atiyah 的
    环路空间平稳相、Bismut 的维纳测度、Berline 与 Vergne 的标架丛
    计算）从不同来源导出 $\Ahat$ 亏格 $\prod_i\frac{y_i/2}{\sinh(y_i/2)}$。
\end{keytakeaways}

% [Tutor Note: 术语——本节的“局部”指 $t\to0$ 时在每一点 $x$ 处成立的恒等式（局部指标密度），与“整体拓扑”（示性类在基本类 $[X]$ 上的配对）相对；下一节讨论的“几何不变量”指 $\eta$ 不变量、号差亏损等次级不变量，勿与本节的热系数混淆。]
自 1960 年代末起，从基本的阿蒂亚–辛格指数定理生发出来的思想圈不断
扩大，并转向更为解析的方面。关注点从椭圆算子的核扩展到更高的特征值。
这首先引出了指数定理的\emph{局部}版本，随后又引出了局部的\emph{几何}
不变量（与\emph{整体拓扑}不变量相对）。前者在本节讨论，后者在下一节
讨论。

我们从 §\ref{subsec:6.1} 开始，先给出两个基本构造，用以把更高的特征值
汇总成一个函数：$\zeta$ 函数与热算子之迹。（它们是解析数论中若干基本
对象的类比物。）局部指数定理先由维贾伊·帕托迪（Vijay Patodi）在特殊
情形下证明，后由彼得·吉尔基（Peter Gilkey）在一般情形下证明，这是
§\ref{subsec:6.2} 的主题。这项工作完成于 1970 年代初。1980 年代，局部
指数定理的若干新证明使人们对 Dirac 算子指标公式中 $\Ahat$ 亏格的来源
有了更深入的理解。在 §\ref{subsec:6.3} 中，我们简要概述下列诸人的
贡献：埃兹拉·盖茨勒（Ezra Getzler）、爱德华·威滕（Edward Witten）、
让-米歇尔·比斯穆特（Jean-Michel Bismut）、妮科尔·贝尔利纳（Nicole
Berline）、米歇尔·韦尔涅（Michele Vergne），以及当然还有迈克尔·阿蒂亚
（Michael Atiyah）。

  \subsection{热算子、$\zeta$ 函数与指数}\label{subsec:6.1}

% [Tutor Note: 术语——\emph{热算子}是算子 $H_t=e^{-t\Delta}$（\eqref{eq:56}），\emph{热核}是它的积分核 $h_t(x,y)$（\eqref{eq:59}）。二者不可互换：只有热核才有逐点的 $t\to0$ 渐近展开，而热算子本身在 $t\to0$ 时只是趋于恒等算子。]
设 $\Delta $~是希尔伯特空间~$\sH$ 上的非负自伴算子。我们试图定义
\emph{热算子}
  \begin{equation}\label{eq:56}
     H_t=e^{-t\Delta },\qquad t\in \RR^{>0}, 
  \end{equation}
与\emph{$\zeta $ 函数}
  \begin{equation}\label{eq:57}
     \zD(s)=\Tr\Delta ^{-s},\qquad s\in \CC. 
  \end{equation}
% [Tutor Note: 易错——\eqref{eq:57} 的 $\Tr\Delta^{-s}$ 在无穷维时是形式写法：先取谱截断、再解析延拓（$\Delta$ 须有离散谱且特征值按 Weyl 律增长，闭流形上的椭圆算子满足）。零特征值的约定尤须注意：若把 $\ker\Delta$ 上的 $0^{-s}$ 排除（只对 $\lambda>0$ 求和），$\zeta$ 就少一项 $\dim\ker\Delta$；作者在正文（$s=0$ 处的取值 $=\dim\sH^0-\dim\sH^1$）用的是计入零特征值、取 $0^0=1$ 的约定。两种归一化相差 $\dim\ker\Delta$，引用时须核对。]
当 $\sH$ 为有限维时，两者都有良好定义，且由梅林变换相联系
  \begin{equation}\label{eq:58}
     \Tr\Delta ^{-s} = \frac{1}{\Gamma
     (s)}\int_{0}^{\infty}\frac{dt}{t}\,t^s\,\Tr\bigl(e^{-t\Delta } \bigr). 
  \end{equation}
% [Tutor Note: 易错——\eqref{eq:58} 的被积式是 $\frac{dt}{t}\,t^s=t^{s-1}dt$（不是 $t^{-s}$），收敛域也不是“仅 $\Re s>0$”：小 $t$ 端由 \eqref{eq:60} 的 $t^{-n/2}$ 控制，要求 $\Re s>n/2$（否则 $t\to0$ 处发散）；大 $t$ 端若 $\ker\Delta\neq0$，则 $\Tr e^{-t\Delta}\to\dim\ker\Delta$ 而不趋于 $0$，须先减去零特征值投影 $P_0$（改用 $\Tr(e^{-t\Delta}-P_0)$）才收敛。前置的 $1/\Gamma(s)$ 消去 $s=0,-1,\dots$ 处的极点；$\zD$ 的解析延拓即由此得到。]
% [Tutor Note: 术语——$\zeta$ 正规化（zeta regularization）指用 $\zD(s)$ 的亚纯延拓给发散量赋值；例如 $\det\Delta:=e^{-\zeta_\Delta'(0)}$（$\zeta$ 正规化行列式）。这是取有限部分的\emph{定义}，不是普通行列式；本节只用到 $\zD$ 的亚纯延拓，尚未用到行列式。]

% [Tutor Note: 笔误——原文作“the heat operator~$H_t$ in~\eqref{eq:56} is exists by basic elliptic theory”，多一“is”；译文按 exists 理解。]
% [Tutor Note: 含糊——“basic elliptic theory”的实质是：闭流形上非负自伴椭圆算子的谱定理（本质自伴、离散谱）与半群理论给出 $e^{-t\Delta}$ 作为强连续压缩半群，且其核是热方程初值问题的基本解。完备叙述见 \cite{Se4} 与 Berline–Getzler–Vergne《Heat Kernels and Dirac Operators》第 2 章（本书无该文献键）。]
% [Tutor Note: 术语——\emph{磨光算子}（smoothing operator）指核 $h_t$ 光滑、从而把分布映为光滑函数；它必是迹类紧算子，但反之不然，“紧”不蕴含“磨光”。]
若 $\Delta $~是闭流形~$X$ 上非负自伴的二阶椭圆算子，则由基本椭圆理论，
\eqref{eq:56} 中的热算子~$H_t$ 存在。它是\emph{磨光算子}：$H_t$~把分布
映为光滑函数。
% [Tutor Note: 易错——\eqref{eq:59} 约定 $h_t(x,y)=(e^{-t\Delta}\delta_y)(x)$，配合热方程 $u_t=-\Delta u$（传播子 $e^{-t\Delta}$，$\Delta$ 非负）；若改用 $e^{+t\Delta}$，则 $t$ 换成 $-t$，渐近展开的方向随之反转。又 $h_t(x,y)$ 关于 $x,y$ 对称，但须以同一黎曼测度 $|dx|$、$|dy|$ 为前提。]
例如，若 $\delta _y$~是~$y\in X$ 处的 Dirac
$\delta $~分布，则\footnote{我们的记号假定 $\Delta $ 是作用在函数上的
算子；稍作修改便可纳入向量丛的情形。}
  \begin{equation}\label{eq:59}
     h_t(x,y) = \left( e^{-t\Delta }\delta _y \right)\!(x),\qquad x\in X, 
  \end{equation}
是~$t,x,y$ 的光滑函数，称为\emph{热核}。若 $\Delta $~是~$X$ 上某个黎曼
度量所附的标量拉普拉斯算子——即拉普拉斯–贝尔特拉米算子——则直观上
$h_t(x,y)$~是初始热量分布为~$\delta _y$ 时、经过时间~$t$ 后~$x$ 处的
热量。热的传播速度为无穷大，且瞬时扩散：对一切~$t,x,y$ 有
$h_t(x,y)>0$。物理热流的这些性质，为一般的非负自伴二阶微分算子
~$\Delta $（作用在向量丛 $E\to X$ 的截面上）的热算子在大时间与小时间下
的行为提供了直观。
% [Tutor Note: 易错——两端极限的拓扑不同：$t\to0$ 时 $e^{-t\Delta}\to I$ 仅在\emph{强}算子拓扑下成立（当 $\spec\Delta$ 无界时 $\|I-e^{-t\Delta}\|=1$，范数收敛失败），而 $t\to\infty$ 时在\emph{一致}（算子范数）拓扑下收敛到 $\ker\Delta$ 上的投影，速率由第一个正特征值控制。原文特意区分这两种拓扑，不可混用。]
当 $t\to\infty $ 时，热算子~$e^{-t\Delta }$ 收敛（在一致拓扑下）到向
$\ker \Delta $ 的投影。当 $t\to0$ 时，热算子收敛（在强算子拓扑下）到
恒等算子。
% [Tutor Note: 易错——原文说 $x\neq y$ 时 $h_t(x,y)$ “指数式地收敛到零”；精确速率由 Varadhan 公式给出（本书无该文献键）：$\log h_t(x,y)\sim -d(x,y)^2/4t$，$d$ 为黎曼距离。故“指数式”应读作 $e^{-d^2/4t}$ 这一尺度，而非任意的 $e^{-c/t}$。]
小~$t$ 行为的一个更精确版本，是米纳克希孙达拉姆（Minakshisundaram）
与普莱耶尔（Pleijel）一篇有影响的 1948 年论文~\cite{MP} 的主题（当
$\Delta $~为拉普拉斯–贝尔特拉米算子时）；西利（Seeley）~\cite{Se4} 把
他们的结果推广到更一般的椭圆伪微分算子。
% [Tutor Note: 笔误——原文作“Minakshisundarum”，而 \cite{MP} 条目作“Minakshisundaram”（S. Minakshisundaram），少一字母“a”；又原文称该文为 1948 年的论文，而 \cite{MP} 为 Canad. J. Math. \textbf{1} (1949), 242–256，年份不符。]
当 $t\to0$ 时，若~$x\neq y$，热核~$h_t(x,y)$ 指数式地收敛到零；而在
对角线上有渐近展开
  \begin{equation}\label{eq:60}
     h_t(x,x)\sim t^{-n/2}\sum\limits_{k=0}^{\infty}A_k(x)t^i\qquad
     \textnormal{当 $t\to0$ 时}, 
  \end{equation}
% [Tutor Note: 笔误——\eqref{eq:60} 的求和项写作 $A_k(x)t^i$，指数 $t^i$ 中的 $i$ 未定义（且与 \eqref{eq:72} 的下标 $i$ 冲突），按标准热核展开应为 $t^k$：$h_t(x,x)\sim t^{-n/2}\sum_k A_k(x)t^k$。照录原文，特此指出。]
% [Tutor Note: 易错——\eqref{eq:60} 的归一化：$t^{-n/2}$ 与 $A_0=(4\pi)^{-n/2}$（见 §\ref{subsec:6.2} 首段）配套，等价于 $h_t(x,x)\sim(4\pi t)^{-n/2}\sum_k a_k(x)t^k$、$a_0=1$；若采用 $(2\pi)^{-n/2}$ 等别的约定，全部 $A_k$ 会差一个常数因子。核对 \cite{MP}、\cite{MS} 的归一化。]
其中 $A_k$~是~$X$ 上的光滑函数。对~$x\in X$，$A_k(x)\in \End E_x$ 的值
只依赖于微分算子~$\Delta $ 的全象征在~$x$ 处的有限阶射流。若 $\Delta $~由
黎曼度量典范地给出，则 $A_k(x)$~只依赖于度量在~$x$ 处的有限阶射流；射流
的阶随~$k$ 增长。

% [Tutor Note: 含糊——“渐近展开与亚纯延拓等价”正是梅林变换的实质：把 $\int_0^\infty t^{s-1}\bigl(\Tr e^{-t\Delta}-\text{有限项}\bigr)dt$ 逐项积分，$t\in(0,1)$ 上的 $t^{s-1}t^{m}$ 项给出 $\frac{1}{s+m}$ 型极点，$t\in(1,\infty)$ 上给出整函数；于是 \eqref{eq:60} 等价于 $\zD$ 在全 $s$ 平面的亚纯延拓。极点的位置与留数（即 Seeley–DeWitt 系数）的严格证明见 \cite{MP,Se4}。]
对于上一段所讨论的那类紧流形上的椭圆算子~$\Delta $，$\zeta$ 函数
\eqref{eq:57} 存在，且当 $\Re (s)>\!>0$ 时 $\zD$~是~$s$ 的全纯函数。热核
~$h_t$ 的渐近展开~\eqref{eq:60} 经梅林变换~\eqref{eq:58} 等价于：$\zD$ 可
亚纯延拓到整个复 $s$ 直线——事实上 \cite{MP,Se4} 所证明的正是这一点。

现在设 $P\:C^{\infty}(X,E^0)\to C^{\infty}(X,E^1)$ 是闭 $n$ 维流形~$X$ 上的
一阶椭圆算子，$X$ 上配有复向量丛 $E^0,E^1\to X$。处处假定已给定度量，使
形式伴随 $P^*\:C^{\infty}(X,E^1)\to C^{\infty}(X,E^0)$ 有定义；它同样是一阶
椭圆算子。于是
  \begin{equation}\label{eq:61}
     \ind P = \dim\ker P^*P - \dim\ker PP^*, 
  \end{equation}
% [Tutor Note: 易错——\eqref{eq:61} 用到 $\ker P^*P=\ker P$ 与 $\ker PP^*\cong\coker P$；这需要 $P$ 为 Fredholm 算子（椭圆性保证闭值域与有限维核、余核）。若不假定 $P^*$ 为形式伴随，右端第二项应直接写成 $\dim\coker P$。]
且 $P^*P,\,PP^*$ 都是非负自伴的二阶椭圆微分算子。对~$\lambda \ge0$，设
$\sE^0_\lambda \subset C^{\infty}(X,E^0)$ 是~$P^*P$ 的 $\lambda $~特征子空间，
$\sE^1_\lambda \subset C^{\infty}(X,E^1)$ 是~$PP^*$ 的 $\lambda $~特征子空间。
于是对~$\lambda >0$，
  \begin{equation}\label{eq:62}
     P\res{\sE^0_\lambda }\:\sE^0_\lambda \longrightarrow \sE^1_\lambda 
  \end{equation}
是同构。
% [Tutor Note: 易错——\eqref{eq:62} 只在 $\lambda>0$ 时成立；$\lambda=0$ 处 $P$ 把 $\ker P^*P=\ker P$ 映入 $0$，指数的全部信息恰在 $\lambda=0$ 这一项。这正是 McKean–Singer 把“核的维数差”改写为“谱上的和之差”的出发点。]
因此，对任何函数 $\chi \:\RR^{\ge0}\to\RR^{\ge0}$，只要 $\chi (0)=1$ 且当
$\lambda \to\infty $ 时 $\chi (\lambda )\searrow0$ 足够快，就有
  \begin{equation}\label{eq:63}
     \ind P = \sum\limits_{\lambda \in \spec P^*P}\chi (\lambda ) -
     \sum\limits_{\lambda \in \spec PP^*}\chi (\lambda ). 
  \end{equation}
% [Tutor Note: 易错——\eqref{eq:63} 要求 $\chi(0)=1$ 以保住 $\lambda=0$ 的项，并要求 $\chi$ 衰减足够快以使两个和收敛；取 $\chi(\lambda)=\lambda^{-s}$ 时需 $\Re s\gg0$，才能先对谱截断求和再取极限（即交换求和与取迹）。]
对 $\chi (\lambda )=\lambda ^{-s}$（$\Re(s)>\!>0$），我们得到出现在
阿蒂亚–博特~\cite[\S8]{AB1} 中的一个指数公式：
  \begin{equation}\label{eq:64}
     \ind P = \Tr \zeta \mstrut _{P^*P}(s) - \Tr \zeta \mstrut _{PP^*}(s). 
  \end{equation}
% [Tutor Note: 术语——$\zD$ 在 \eqref{eq:57} 中已含迹（$\zD=\Tr\Delta^{-s}$），而 \eqref{eq:64} 又对 $\zeta_{P^*P}$ 取迹，可见此处的 $\zeta_{P^*P}(s)$ 指算子 $(P^*P)^{-s}$（复幂），记号前后并不完全一致；勿据此认为原文重复取迹。]
事实上，由于 $\zeta$ 函数的亚纯延拓，\eqref{eq:64}~对一切~$s\in \CC$ 成立。
阿蒂亚–博特指出，鉴于 \cite{Se4} 给出了用~$P$ 的象征表示的显式积分公式，
$s=0$ 是特别好的一个取值。（取~$s=0$ 的另一动机是：对有限维希尔伯特空间
之间的算子 $P\:\sH^0\to\sH^1$，$s=0$ 处的取值就是
$\dim \sH^0 - \dim \sH^1$。）
% [Tutor Note: 含糊——“explicit integral formulas~\cite{Se4} in terms of the symbol”指 Seeley 的复幂构造：$(P^*P)^{-s}$ 的核由象征 $\sigma(P^*P)^{-s}$ 的振荡积分给出，故 $\zeta(s)$ 是象征的积分。但正因如此，这些公式含象征的高阶导数，与 \eqref{eq:44}（Chern–Weil 形式只用曲率的低阶导数）形成作者所说的“mismatch”。]
但这些显式公式虽是局部的，却涉及象征的高阶导数；而指数的示性类公式
\eqref{eq:44} 若用曲率的 Chern–Weil 多项式写出，只涉及少数几阶导数。
正是这一不匹配，在此后数年间一直是个谜。

代替 $\zeta$ 函数，人们常在 \eqref{eq:63} 中使用热核的迹。这对应
$\chi (\lambda )= e^{-t\lambda }$、$t>0$。于是对一切~$t>0$，有
  \begin{equation}\label{eq:65}
     \ind P = \Tr e^{-tP^*P} - \Tr e^{-tPP^*}. 
  \end{equation}
% [Tutor Note: 易错——\eqref{eq:65} 两端与 $t\in(0,\infty)$ 无关：$t>0$ 时 $\Tr e^{-tP^*P}$ 与 $\Tr e^{-tPP^*}$ 均可微（热核光滑、迹类），求导后由迹的循环性得 $\Tr\bigl(P^*P e^{-tP^*P}\bigr)=\Tr\bigl(PP^*e^{-tPP^*}\bigr)$，故导数为零。不能以直接取 $t\to0$ 或 $t\to\infty$ 的极限代替这一论证：两端的收敛性都须单独证明。]
事实上，可以求导证明右端与~$t$ 无关，而取~$t\to\infty $ 的极限就重新得到
\eqref{eq:65}。
% [Tutor Note: 笔误——原文作“evaluation as $t\to\infty$ reproduces~\eqref{eq:65}”，但 $t\to\infty$ 时 $e^{-tP^*P}$ 与 $e^{-tPP^*}$ 分别收敛到 $\ker P^*P$、$\ker PP^*$ 上的投影，得到的是 \eqref{eq:61}（即 $\ind P=\dim\ker P^*P-\dim\ker PP^*$），而不是 \eqref{eq:65} 本身。疑应为 \eqref{eq:61}。]
另一方面，令~$t\to0$ 并利用渐近展开~\eqref{eq:60}，得到
  \begin{equation}\label{eq:66}
     \ind P = \int_{X}\tr\left[ A^0_{n/2}(x) - A^1_{n/2}(x) \right] \,|dx|, 
  \end{equation}
% [Tutor Note: 易错——\eqref{eq:66} 中的 $A^i_{n/2}(x)$ 要求 $n$ 为偶数（热系数 $A_k$ 只对非负整数 $k$ 有定义）；对奇数维 $n$ 此式需另行解释，原文未加声明。]
其中 $A^0_k(x),A^1_k(x)$~分别是~$P^*P,PP^*$ 作用在向量空间~$E^0_x,E^1_x$ 上
的热系数。正如对 $\zeta$ 函数那样，$A^i_{n/2}(x)$~的公式涉及~$P$ 的象征在
~$x$ 处的许多导数，因此似乎难以用来证明指数公式。
% [Tutor Note: 术语——此处“热系数”即 Seeley–DeWitt 系数；\eqref{eq:66} 的被积式 $\tr\bigl[A^0_{n/2}-A^1_{n/2}\bigr]$ 是“局部指标密度”（local index density），积分后给出指数。注意区分：局部指数定理断言该密度存在并等于某个示性形式；“指标密度”只是被积式本身。]

  \subsection{局部指数定理}\label{subsec:6.2}

对 $n$ 维黎曼流形~$X$ 上的拉普拉斯–贝尔特拉米算子~$\Delta $，\eqref{eq:60}
中第一个系数~$A_0$ 是常函数 $1/(4\pi )^{n/2}$。这反映了欧氏空间中经典热
方程的解，并且蕴含关于~$\Delta $ 特征值渐近增长的外尔（Weyl）定律，后者只
依赖~$n$ 与 $\Vol(X)\sim\int_{X}A_0(x)\,|dx|$。外尔定律不依赖热核展开，它
是马克·卡茨（Mark Kac）~\cite{K} 于 1966 年提出下述问题的动机之一：
$\Delta $ 的特征值在多大程度上决定黎曼流形~$X$？卡茨关注的是欧氏平面
~$\EE^2$ 中的区域，不过更一般的问题已隐含其中。

麦基恩–辛格（McKean–Singer）~\cite{MS} 立即着手研究这一问题，他们确定
了热核展开中接下来的几个系数~$A_1,A_2$，从而证明了卡茨–普莱耶尔
（Kac–Pleijel）的一个猜想。（同期还有德布鲁因（de Bruijn）、阿诺德
（Arnold）与贝热（Berger）的独立结果。）特别地，$A_1$ 是标量曲率的倍数。
对~$n=2$，McKean–Singer 观察到在~$X$ 的每一点都成立的相消，并猜测类似
结果在一切维数下成立。即，设 $\Delta \q$~表示微分 $q$ 形式上的拉普拉斯
算子，$h_t\q$~为相应的热核，$A_k\q$ 为热系数。
% [Tutor Note: 易错——\eqref{eq:67} 的交错和 $\sum_q(-1)^q\tr h_t\q$ 就是超迹 $\Trs=\Tr(\gamma\,\cdot)$：在 $\Omega^{\bullet}$ 上取 $\gamma=(-1)^{\deg}$（形式次数分级），即 de Rham 算子的手征性算子；若采用 $\Omega^{\pm}$（偶/奇形式）分级的 $\gamma=\pm1$ 约定，分项写法不同而交错和相同。若与 §2 的手征约定符号相反，则 $\Trs$ 整体变号，务必核对。]
McKean–Singer 猜想是：对一切~$x\in X$，极限
  \begin{equation}\label{eq:67}
     \lim\limits_{t\to0}\;\sum\limits_{q=0}^n\,(-1)^q\tr h_t\q(x,x) 
  \end{equation}
存在；此外，对偶数~$n$ 它等于高斯–博内–陈（Gauss–Bonnet–Chern）
被积式，后者积分后给出~$X$ 的欧拉数。
% [Tutor Note: 术语——“Gauss–Bonnet–Chern integrand”即欧拉形式 $\Euler(TX)$ 的被积式：$\chi(X)=\int_X \Euler(TX)$；原文以连字符写作“Gauss-Bonnet-Chern”，规范写法为 Gauss–Bonnet–Chern。对奇数维 $n$，McKean–Singer 证明该极限为零。]
由~\eqref{eq:60}，极限的存在性等价于相消
  \begin{equation}\label{eq:68}
     \sum\limits_{q=0}^n(-1)^q\tr A_k\q(x)=0,\qquad k<\frac n2,\quad x\in
     X. 
  \end{equation}
对~$k=n/2$ 的交错和等于极限~\eqref{eq:67}。McKean–Singer 证明了极限
存在、且在 $n$~为奇数时为零，并计算了~$n=2$ 时的极限。

% [Tutor Note: 含糊——Patodi 的“virtuoso direct computation”指逐项算出 Seeley–DeWitt 系数 $A_k\q$ 并直接验证 $k<n/2$ 时交错和为零，本质上是曲率张量多项式的组合恒等式硬算（极为繁重）；见 \cite{P1} 以及 \cite{ABP} 的相应部分。]
1970 年，帕托迪（Patodi）~\cite{P1} 通过一次大师级的直接计算证明了
McKean–Singer 猜想。他随即~\cite{P2} 把其方法用于证明凯勒流形上的
黎曼–罗赫定理~\ref{thm:1}。几年后，吉尔基（Gilkey）在其博士论文
~\cite{Gi1} 中用不同的方法——标度在其中起关键作用——证明了扭化号差算子
的相应定理。
% [Tutor Note: 含糊——“由标准的拓扑论证，这蕴含 Atiyah–Singer 指数定理”的实质是配边论证（即本文 §\ref{subsec:4.2} 的“第一证明”）：扭化号差算子情形的指数定理经配边不变性与乘法性，可化归为一般椭圆算子的情形。细节见 \cite{ABP}；配边（第一）证明的原始文本见 \cite{Pa}（据本文 §\ref{sec:4} 开头的说明，Atiyah–Singer 的第一个证明未由他们自己写出，而是收入 Palais 编的 Seminar 卷；1968 年的 \cite{AS2} 是第二个、K 理论证明）。]
由标准的拓扑论证，这蕴含阿蒂亚–辛格指数定理~\ref{thm:5}；见
§\ref{subsec:4.2}。随后，阿蒂亚–博特–帕托迪（Atiyah–Bott–Patodi）
~\cite{ABP} 给出了吉尔基定理的一个证明，以及由此得到的指数定理的证明。
下面叙述吉尔基的主要结果。

% [Tutor Note: 术语——原文把 functor 与 sheaf 并列（“the functors (sheaves)”）：$\Met$ 与 $\Omega^q$ 都是 $\Mn\op\to\Set$ 的函子，且是层；范畴 $\Mn$ 的态射取\emph{局部微分同胚}，正是为了使度量的拉回只在局部有定义，从而得到层而非仅为预层。]
吉尔基研究由黎曼度量典范地构造出的微分形式。为节省笔墨而使用现代术语：
设 $\Mn$~是光滑 $n$ 维流形与局部微分同胚构成的范畴。考虑函子（层）
  \begin{equation}\label{eq:69}
     \begin{aligned} \Met\:\Mn\op&\longrightarrow \Set \\ \Omega
      ^q\:\Mn\op&\longrightarrow \Set\end{aligned} 
  \end{equation}
其中若 $M$~是光滑 $n$ 维流形，则 $\Met(M)$~是 $M$ 上黎曼度量的集合，
$\Omega ^q(M)$~是微分 $q$ 形式的集合。我们寻求自然变换
  \begin{equation}\label{eq:70}
     \omega \:\Met\longrightarrow \Omega ^q. 
  \end{equation}
粗略地说，这些是把微分形式指派给黎曼度量、且在坐标变换下协变的法则。
更诗意地说，它们是 $\Met$ 上的微分形式。即便对~$q=0$，分类问题也是棘手
的：标量曲率的任何光滑函数都给出黎曼度量的一个自然函数。现在引入标度。
% [Tutor Note: 易错——\eqref{eq:71} 的自变量是 $\lambda^2g$（而非 $\lambda g$）；常值缩放 $g\mapsto\lambda^2g$ 不改变 Levi-Civita 联络（Christoffel 符号不变），故曲率 2-形式 $\Omega$ 与由它构造的 Pontryagin/Chern–Weil 形式都不变，只有标量曲率按 $\lambda^{-2}$ 缩放。于是 Pontryagin 形式的任何多项式权均为 $0$；再结合定理~\ref{thm:13}（非负权、正则、自然的微分形式必为 Pontryagin 形式的多项式），即得“非零形式的权为零”，同时可见 Hirzebruch 的 $L$ 类 \eqref{eq:9} 与 $\Ahat$ 类 \eqref{eq:22} 的权为 $0$。]
称 $\omega $~是\emph{权为 $k$ 的齐次形式}，若
  \begin{equation}\label{eq:71}
     \omega (\lambda ^2g)=\lambda ^k\omega (g)\qquad \textnormal{对一切
     $\lambda \in \RR^{>0}$}. 
  \end{equation}
称 $\omega $~是\emph{正则的}，若在任意局部坐标系 $x^1,\dots ,x^n$ 中它取
如下形状
  \begin{equation}\label{eq:72}
     \omega (g)(x) = \sum\limits_{I}\sum\limits_{\alpha
     }^{\textnormal{有限}}\sum\limits_{i,j=1}^n \omega ^{i,j}_{I,\alpha
     }(x)\, \frac{\partial ^{|\alpha |}g_{ij}}{\partial x^{\alpha
     _1}\cdots\partial x^{\alpha _n}} \,dx^{i_1}\wedge \cdots\wedge
     dx^{i_q}, 
  \end{equation}
其中 $I=(i_1,\dots ,i_q)$ 满足 $1\le i_1<\cdots< i_q\le n$，而
$\alpha =(\alpha _1,\dots ,\alpha _n)$ 满足 $\alpha _k\in \ZZ^{\ge0}$。
$\omega ^{ij}_{I,\alpha }$~是光滑函数。关键是：只有有限个 $\alpha $ 出现。
% [Tutor Note: 笔误——\eqref{eq:72} 中被求和的系数写作 $\omega ^{i,j}_{I,\alpha }$（带逗号），紧接着的正文却写作 $\omega ^{ij}_{I,\alpha }$（无逗号），二者指同一对象；照录原文，仅指出写法不一致。]

  \begin{theorem}[Gilkey~\cite{Gi1}]\label{thm:13}
 满足 \eqref{eq:70} 的自然微分形式，若正则且为非负权齐次，则是 Pontryagin
 类的 Chern–Weil 形式的多项式。
  \end{theorem}

\noindent
 非零形式的权为零。

% [Tutor Note: 含糊——“Weyl's theorem~\cite{W} on invariants of the orthogonal group”即正交群不变量的\emph{第一基本定理}：$O(n)$ 不变多项式环由内积 $g_{ij}$ 与体积形式 $\epsilon$ 生成。Gilkey 用它把 \eqref{eq:72} 中出现的全部 $g$ 的导数项化归为这些不变量的多项式，再借“权”的齐次性消去除 Pontryagin 形式之外的项；完整论证见 \cite{ABP}。]
\cite{ABP} 中吉尔基定理~\ref{thm:13} 的证明使用了外尔关于正交群
不变量的定理~\cite{W}。阿蒂亚–博特–帕托迪把定理~\ref{thm:13} 应用于
黎曼流形~$X$ 上的号差算子~$P$。沿用 §\ref{subsec:6.1} 的记号，正权形式
的消失蕴含
  \begin{equation}\label{eq:73}
     \tr\bigl[A^0_k(x) - A^1_k(x) \bigr]=0,\qquad k<\frac n2,\quad x\in X. 
  \end{equation}
如~\eqref{eq:68} 那样，这一相消结果蕴含
  \begin{equation}\label{eq:74}
     \lim\limits_{t\to0}\;\bigl[\tr h^0_t(x,x) - \tr h^1_t(x,x) \bigr], 
  \end{equation}
的存在性；而定理~\ref{thm:13} 指出该极限是 Pontryagin 形式的多项式。
精确的多项式——即 $L$ 亏格~\eqref{eq:9}——像在希策布鲁赫原来的证明中
那样，通过计算足够多的例子来确定。从这里通往\emph{整体}指数定理的道路
沿既定的线索展开。真正引向日后发展的，是\emph{局部}指数定理——即极限
~\eqref{eq:74} 的存在性与识别。
% [Tutor Note: 术语——“局部指数定理”含两层内容：(i) 极限 \eqref{eq:74} 存在；(ii) 它等于某个局部示性形式（此处为 $L$ 类，对 Dirac 算子则为 $\Ahat$ 形式）。被取极限的量 $\tr h^0_t-\tr h^1_t$ 是\emph{局部指标密度}；“局部指数定理”断言密度与示性形式恒等，而“整体指数定理”是其积分结果。]

  \subsection{后记：$\Ahat$ 亏格从何而来？}\label{subsec:6.3}

对相消~\eqref{eq:73} 与极限值~\eqref{eq:74} 的不同概念性理解，是在 1980
年代前半期获得的。背景是（广义的）Dirac 算子，其中 Clifford 代数的基本
性质给出相消。极限是 Chern–Weil $\Ahat$ 形式，即 Pontryagin 形式的多项
式，它的出现可从各种来源导出。在这些工作中，$\Ahat$ 亏格是靠直接论证
出现的。我们给出简要概述。

% [Tutor Note: 含糊——Mehler 公式（\cite[p.~19]{GJ}）给出谐振子热核的显式高斯核 $K_t(x,y)$；Getzler 的要点在于：把位似后的热核当 $t\to0$ 时直接与这一显式核比较，从而把渐近展开化为算子代数的极限。细节见 \cite{Ge2}。]
通往 $\Ahat$ 亏格的一条路径经过谐振子热核的梅勒（Mehler）公式
~\cite[p.~19]{GJ}，盖茨勒（Getzler）~\cite{Ge2} 在其局部指数定理的证明中
使用了它。他用一个位似变换，它不仅缩放时间与空间，还缩放 Dirac 算子
~\eqref{eq:41} 中的 Clifford 代数变量。
% [Tutor Note: 易错——Getzler 的位似须同时缩放三类对象：$t\mapsto\varepsilon^2 t$、空间坐标按 $\varepsilon$ 缩放、Clifford 变量（奇变量）按 $\varepsilon$ 缩放（各版本指数略有出入，须核对 \cite{Ge2} 的约定）。若只缩放时空而不缩放 Clifford 变量，重标度后的热核不收敛，也得不到 $\Ahat$ 形式。]
他的技巧部分受到当时用超对称量子力学给出的指数定理物理证明
~\cite{Wi1,Ag,FW,Ge1} 的启发。

% [Tutor Note: 术语——“Duistermaat–Heckman exactness of stationary phase”指平稳相近似在等变（环面作用）情形下\emph{精确}：积分恰等于驻点局部贡献之和，无高阶修正项，即 DH 局部化公式。Atiyah 在 \cite{A6} 中把这一精确性形式地移置到自由环路空间上。]
在纪念洛朗·施瓦茨（Laurent Schwartz）的一次会议上，阿蒂亚~\cite{A6}
阐述了威滕的想法：将杜斯特马特–赫克曼（Duistermaat–Heckman）平稳相的
精确性定理~\cite{DH} 应用于紧黎曼流形的自由环路空间，从而导出指数定理。
（同一时期，阿蒂亚–博特~\cite{AB5} 把 Duistermaat–Heckman 的结果置于
等变上同调局部化的框架中。）在这一证明中，$\Ahat$ 亏格是通过正规化某个
无穷乘积而进入的，正如它在超对称量子力学的证明~\cite[\S1.2.4]{Wi2} 中
那样。

受阿蒂亚这一叙述的启发，比斯穆特（Bismut）~\cite{Bi1} 用环路空间上的
维纳（Wiener）测度与马拉万（Malliavin）微积分实现了指数定理的一个证明。
这样他就严格地处理了环路空间上的积分。借助利希内罗维奇（Lichnerowicz）
公式（它用协变拉普拉斯算子表示 Dirac 拉普拉斯算子），热核被用维纳测度
表示出来。在~$t\to0$ 时向点环路的局部化，在这一概率进路中是自然的。
变量~$t$ 表示布朗轨道存在的总时间；当时间趋于零时，只有常值环路有显著
的出现概率。对点环路上积分的计算，借助保罗·莱维（Paul Lévy）~\cite{L}
的一个公式完成：他考虑平面中一条在时间~$2\pi $ 后闭合的布朗曲线。于是
随机曲线所围面积~$S$ 的特征函数（$e^{izS}$ 的期望值，$z\in \CC$）为
$\pi z/\sinh \pi z$。
% [Tutor Note: 术语——此处的“特征函数”（characteristic function）是概率论意义（$\mathbb{E}[e^{izS}]$），与示性类（characteristic class）无关；中文两者都作“特征”，极易混淆。]
% [Tutor Note: 易错——Lévy 公式 $\pi z/\sinh\pi z$ 与 $\Ahat$ 的对应靠 $2\pi$ 归一化：取 $z_i=y_i/2\pi$（$y_i$ 为曲率特征值），则 $\prod_i\frac{\pi z_i}{\sinh \pi z_i}=\prod_i\frac{y_i/2}{\sinh(y_i/2)}$，正是 \eqref{eq:22} 的 $\Ahat$ 类。注意此处 $\sinh$ 取实变量，$z$ 与曲率的对应含因子 $2\pi$，漏掉该因子会得到错误的归一化。]
同一计算出现在 Bismut 的工作中，只不过那里~$X$ 的曲率代替了~$z$，于是
再一次得到 $\Ahat$ 亏格。

% [Tutor Note: 含糊——“a standard formula in the theory of Lie groups”指：对 $\exp(a+tX)$ 在 $t=0$ 展开可得 $d\exp_a=\frac{1-e^{-\ad a}}{\ad a}$（Maurer–Cartan 形式展开）；此处 $J(x)=(1-e^{-x})/x$，其级数从 $1$ 开始，故可逆。推导与约定见 \cite{BV}。]
在贝尔利纳与韦尔涅（Berline–Vergne）~\cite{BV} 给出的指数定理证明中，
$\Ahat$ 亏格以相当不同的方式出现。设 $G$~是李群，李代数为~$\frak{g}$。于是
李群理论中的一个标准公式断言：指数映射~$\exp\:\frak{g}\to G$ 在
~$a\in \frak{g}$ 处的微分是
  \begin{equation}\label{eq:76}
     d\exp_a = \frac{1-e^{-\ad a}}{\ad a} = J (\ad a),
  \end{equation}
% [Tutor Note: 笔误——公式标签顺序与出现顺序不符：\eqref{eq:76} 在正文中先于 \eqref{eq:75} 出现，按排版编号二者将互换。正文引用无误，仅标签编号顺序错位，照录原文。]
其中定义~$J$ 的幂级数是定义 Todd 亏格~\eqref{eq:5} 的幂级数的乘法逆。
Todd 亏格出现在~\eqref{eq:76} 中是否与指数定理有关，曾是一个谜。
% [Tutor Note: 易错——$\Ahat$ 与 $J$ 的关系正是此谜的答案：$J(x)=\frac{1-e^{-x}}{x}$，故 $\Ahat(x)=\bigl(J(x)J(-x)\bigr)^{-1/2}=\frac{x/2}{\sinh(x/2)}$（与 \eqref{eq:22} 核对）。于是 Berline–Vergne 公式中由 Todd 级数之逆 $J$ 产生的因子，经曲率代入后恰好给出 $\Ahat$ 亏格；注意其中取了几何平均 $\bigl(J(x)J(-x)\bigr)^{-1/2}$，符号约定须与定向一致。]
Berline–Vergne 注意到：若 $X$~是黎曼流形，$\O(X)$~是正交标架主丛，则
$\O(X)$ 上黎曼指数映射的微分由一个类似的公式给出。确切地说，经由
列维–奇维塔联络，存在自然同构 $T_p\O(X)\cong \Rn\oplus\frak{o}(n) $，而
指数映射 $\exp\:\Rn\oplus\frak{o}(n)\to \O(X)$ 在~$p\in \O(X)$ 处的微分，
在~$a\in \frak{o}(n)$ 上的取值是
  \begin{equation}\label{eq:75}
     \begin{aligned} &d\exp_a\res{\Rn} = \exp(-a)\,J \bigl( \langle
     \Omega _p/2,a\rangle\bigr).\\ &d\exp_a\res{\frak{o}(n)}=J
     (\ad a); 
     \end{aligned} 
  \end{equation}
% [Tutor Note: 易错——\eqref{eq:75} 有三处约定须核对 \cite{BV}：(i) 曲率项写作 $\Omega_p/2$，这个 $1/2$ 正是 $\Ahat$ 级数中 $x/2$ 的来源；(ii) $\exp(-a)$ 的符号取决于指标准架与 Levi-Civita 联络的约定，取 $\exp(+a)$ 会使 $J$ 的自变量变号；(iii) 与 $a$ 的缩并用 Killing 型，结果是 $\mathfrak{o}(n)$ 值的 2 形式，还须与 $\Omega$ 配对。]
在此公式中，取值于~$\frak{o}(n)$ 的黎曼曲率~$\Omega $ 用基灵（Killing）型
与~$a$ 缩并。结果是一个 2 形式，可等同于~$\frak{o}(n)$ 的一个元素。为证明
指数定理，Berline–Vergne 在标架丛~$\O(X)$ 而非底流形~$X$ 上工作。为补偿
纤维方向上多出的自由度，他们必须研究热核沿纤维的行为。
% [Tutor Note: 含糊——“研究热核沿纤维的行为”指在 $\O(X)$ 上作等变（沿纤维的 $\O(n)$ 作用）热核的渐近展开并对纤维方向作局部化；正是在 $t\to0$ 的极限中，纤维方向的 Jacobian 取 \eqref{eq:75} 的形式，从而产生 $\Ahat$ 因子。细节见 \cite{BV}。]
正是在这一阶段，在小时间极限下，\eqref{eq:75} 出现了。归根结底，
$\Ahat$ 亏格就是以这种方式进入他们的证明的。

% ==================== §7 ====================

\section{Dirac 算子的几何不变量}\label{sec:7}

\begin{keytakeaways}
  \item \textbf{主题转向}：在此之前，指数理论给出的是椭圆算子的
\emph{整体拓扑}不变量，其自然归宿是拓扑 $K$ 理论；自 1973 年
阿蒂亚–帕托迪–辛格（Atiyah–Patodi–Singer）的工作起，焦点转向由热方程
构造的\emph{局部微分几何}不变量，且对象限于 Dirac 算子。
  \item \textbf{号差亏损}：带边 4-流形上的
$\alpha(Y)=\Sign(X)-\int_X\omega$ 只依赖边界 $Y$，它由 $\eta$ 不变量给出
（$\alpha=-\eta$），其典型实例之一是 Hilbert 模曲面上 Shimizu $L$ 函数在 $s=1$
的值。
  \item \textbf{整体边界条件}：局部椭圆边界条件受 Atiyah–Bott 拓扑障碍
限制；APS 改用边界自伴算子 $A_Y=\sigma\inv D_Y$ 的谱投影来定义整体边界
条件，$\dbar$ 在单位圆盘上的 Fourier 截断是其原型。
  \item \textbf{APS 指标定理}：$\ind D_X=\int_X\omega-(\et{\bX}+h)/2$，
其中 $h=\dim\ker D_{\bX}$；对号差算子它化为
$\Sign(X)=\int_X\omega-\et{\bX}$。
  \item \textbf{次级不变量与行列式}：陈–西蒙斯不变量与整上同调相配，
$\eta$ 不变量与 $K$ 理论相配；再进一步是 $\zeta$ 正规化行列式、
Quillen 度量与行列式线丛，其和乐由绝热极限下的 $\eta$ 给出，曲率为
$2\pi i\int_{X/S}\omega$。
\end{keytakeaways}

% [Tutor Note: 术语——标题中的「几何不变量」指依赖度量（对 Dirac 算子还依赖 spin 结构与联络）的不变量，与 §2–§6 中由示性数和 $K$ 理论给出的拓扑不变量相对；本节按三个层次展开：号差亏损（普通实数或 $\RZ$ 值）、$\eta$ 不变量、以及行列式线丛（带度量与联络的几何对象）。]
至此，指数理论所产生的都是椭圆算子的\emph{整体拓扑}不变量；
它们的自然归宿是拓扑 $K$ 理论。自阿蒂亚–帕托迪–辛格
（Atiyah–Patodi–Singer）于 1973 年公布的工作开始，指数理论转向了
\emph{局部微分几何}不变量。\footnote{几何不变量的自然归宿是\emph{微分}
$K$ 理论，但那是更晚近的发展，超出本文范围。}此外，关注点也从一般的
椭圆伪微分算子转向 Dirac 算子。热方程方法提供了构造不变量的基本工具。

% [Tutor Note: 含糊——脚注说几何不变量的归宿是微分 $K$ 理论「a more recent development」：其实际内容是把初级（示性数）与次级（陈–西蒙斯、$\eta$）不变量统一为同一个上同调理论中的类的 Cheeger–Simons 微分特征理论（本文 [ChS]、[F2]、[HS]），其算子代数版本见 [FL]；作者在此只提及而不展开。]
% [Tutor Note: 含糊——"the index theorem for flat bundles" 指 APS 系列论文中的平坦丛指数定理（见 [APS2]、[APS3] 的相关讨论），它把带边指标定理推广到平坦 $\Gamma$-丛的情形；作者未注出处，阅读时不要把它与 \eqref{eq:92} 的 Dirac 算子指标定理混为一谈。]
阿蒂亚–帕托迪–辛格的三篇论文~\cite{APS1,APS2,APS3} 引入了
Dirac 算子的 $\eta$ 不变量。它的定义与分析数论中 $L$ 函数的定义相平行。
他们的第一个主要定理（我们在 \S\ref{subsec:7.3} 中叙述）是紧致带边
黎曼 spin 流形上 Dirac 算子的指数定理。其中的一个关键要素是新的
\emph{整体}椭圆边界条件（\S\ref{subsec:7.2}）；局部椭圆边界条件在多数
情形下受到阻碍，这一点阿蒂亚–博特（Atiyah-Bott）在十年前就已发现。
$\eta$ 不变量解决了一个曾是他们工作动机之一的问题，即号差亏损的计算；
这正是我们在 \S\ref{subsec:7.1} 的出发点。阿蒂亚–帕托迪–辛格的论文
还包含许多其他重要定理，例如平坦丛的指数定理，我们在此不予涉及。

$\eta$ 不变量的另一种视角是 \S\ref{subsec:7.4} 的主题。闭定向流形上
向量丛的示性数——即把向量丛的陈类与 Pontryagin 类的乘积与底流形的
基本类配对所得的整数——是初级的整值拓扑不变量。与之相伴的 $\RZ$ 值
次级微分几何不变量，已在几年前由陈省身与西蒙斯（Chern–Simons）引入。
类似地，整值的 $K$ 理论示性数（按指数定理它们就是 Dirac 算子的指数）
是初级的拓扑不变量。与之相伴的次级微分几何量就是阿蒂亚–帕托迪–辛格的
$\eta$ 不变量。Dirac 算子（或 Dirac 算子族）的下一个几何不变量是行列式。
相关的理论在 1980 年代发展起来，我们在 \S\ref{subsec:7.5} 中叙述。
它是我们在 \S\ref{sec:8} 中所处理的物理学应用的一个重要成分。



  \subsection{号差亏损}\label{subsec:7.1}

% [Tutor Note: 易错——\eqref{eq:84} 的归一化：全 Gauss–Bonnet 定理写作 $\int_X K\,d\mu_X=2\pi\chi(X)$，故被积式取 $\frac{K}{2\pi}$ 才恰给出 Euler 数；把 $\Euler$ 理解为 Euler 形式的积分时，这个 $2\pi$ 不可漏掉。]
回顾经典的 Gauss–Bonnet 定理。设 $X$ 是闭黎曼 2-流形，$K\:X\to\RR$
是其高斯曲率。则 $X$ 的 Euler 数是曲率积分
  \begin{equation}\label{eq:84}
     \Euler(X) = \int_{X}\frac{K}{2\pi }\,d\mu _X, 
  \end{equation}
其中 $d\mu _X$ 是黎曼测度。若 $X$ 紧致带边，则还有来自测地曲率
$\kappa \:\bX\to\RR$ 的边界贡献，即
  \begin{equation}\label{eq:85}
     \Euler(X) = \int_{X}\frac{K}{2\pi }\,d\mu _X + \int_{\bX}\frac{\kappa
     }{2\pi }\,d\mu _{\bX}. 
  \end{equation}
% [Tutor Note: 易错——\eqref{eq:85} 中 $\kappa$ 的符号取决于边界的外法向约定：对平面圆盘（$K=0$、$\kappa=1/r$，取外法向）有 $\Euler=1$；若改取内法向，边界项整体变号，须与 $X$ 的定向一并固定后再核对。]
% [Tutor Note: 术语——"a neighborhood ... isometric to the cylinder $[0,\epsilon )\times \bX$ with its product metric" 即度量在边界附近为乘积度量（领状邻域）。这是本节所有带边公式的标准假设：只有此时边界算子 $A_Y$ 才与 $t$ 无关，\eqref{eq:90} 的分解才成立。]
若 $X$ 中 $\bX$ 的某个邻域与柱面 $[0,\epsilon )\times \bX$ 带乘积度量
等距（对某个 $\epsilon >0$），则边界项消失。

% [Tutor Note: 易错——\eqref{eq:86} 中的 $\omega$ 是 $p_1/3$ 的 Chern–Weil 4-形式，而 4 维时 $L$ 亏格的最低次项正是 $L_1=p_1/3$，故 \eqref{eq:86} 与 \eqref{eq:94} 在 4 维一致。切勿把它与 $\Ahat$ 的整性归一化混用：$\hat A_1=-p_1/24$（\eqref{eq:24} 与 \eqref{eq:99}），两者相差因子 $-8$。]
现在设 $X$ 是闭定向黎曼 4-流形。号差定理~\ref{thm:3} 给出
  \begin{equation}\label{eq:86}
     \Sign(X) = \int_{X}\omega , 
  \end{equation}
其中 $\omega$ 是有理示性类 $p_1/3$ 的 Chern–Weil 4-形式。
% [Tutor Note: 易错——这里是本小节的关键对比：乘积度量使 \eqref{eq:85} 的边界项消失，却不能让 \eqref{eq:86} 成立。原因是 4 维的「边界项」并非某个 3-形式的局部积分（下文说 $\alpha$ 不是 $\int_Y\eta$ 的形式），而是本质上的谱量（非局部）；切勿照搬 Gauss–Bonnet 的直观。]
若 $X$ 紧致带边，即便假定黎曼度量在边界附近为乘积度量（下文我们确实
如此假定），式~\eqref{eq:86} 也未必成立。令 $Y=\bX$。于是\emph{号差亏损}
（signature defect）~\cite[\S10.3]{A7}
  \begin{equation}\label{eq:87}
     \alpha (Y)=\Sign(X) - \int_{X}\omega 
  \end{equation}
只依赖于闭定向黎曼 3-流形 $Y$，这由~\eqref{eq:86} 容易推出。
% [Tutor Note: 含糊——"as follows easily from \eqref{eq:86}" 的实质：若 $Y=\bX=\bX'$，把 $X$ 与 $-X'$ 沿 $Y$ 粘合得闭 4-流形 $Z$，对 $Z$ 用 \eqref{eq:86} 即得 $\Sign(X)-\Sign(X')=\int_X\omega-\int_{X'}\omega$，故 $\alpha$ 只依赖 $Y$；此处用到号差在粘合下的可加性与积分的可加性。]
阿蒂亚–帕托迪–辛格~\cite{APS1} 论证了 $\alpha$ 是黎曼度量的光滑函数，
在定向反向下变号，并且\emph{不}具有 $\int_{Y}\eta $ 的形式（其中 $\eta $
是某个由度量自然给出的 3-形式），因为号差亏损在有限覆盖下不具乘性。
% [Tutor Note: 易错——号差亏损 $\alpha$ 在定向反向下变号：$\Sign(\bar X)=-\Sign(X)$ 且 $\int_{\bar X}\omega=-\int_X\omega$，故 $\alpha(\bar Y)=-\alpha(Y)$。$\eta$ 不变量同样在流形反向下变号（其谱关于原点翻转，见 \eqref{eq:93}），故 $-\eta=\alpha$ 与这一奇偶性相容；对照 \eqref{eq:94} 时符号不要弄反。]
% [Tutor Note: 含糊——"not of the form $\int_Y\eta$" 的论据是「号差亏损在有限覆盖下不具乘性」：若 $\alpha(Y)=\int_Y\eta$ 且 $\eta$ 由度量局部给出，则 $n$ 重覆盖 $Y'\to Y$ 将给出 $\alpha(Y')=n\,\alpha(Y)$；而 [APS1]（亦见 [A7] §10.3）说明号差亏损没有这种乘性。严格核对需计算 $\eta$ 在覆盖下的行为。]

% [Tutor Note: 术语——Hilbert 模曲面 $X^0(K)=(\HH\times\HH)/\PSL_2(\sO)$ 是 Hilbert 模形式所对应的模曲面；它有 orbifold 点（由 $\sO$ 中的有限阶单位引起）和与理想类一一对应的端。截断端、挖去 orbifold 点邻域后得到的 $X(K)$ 是带边光滑流形，其边界正是号差亏损的计算对象。]
希策布鲁赫（Hirzebruch）~\cite{H4} 在 1970 年代早期研究的一个具体
号差亏损实例，曾是阿蒂亚–帕托迪–辛格工作的主要动机之一。设
$K=\QQ(\sqrt d)$ 为实二次数域——其中 $d\in \ZZ^{>1}$ 假定无平方因子——
并设 $\sO\subset K$ 为整数环。$d$ 的两个平方根给出两个嵌入
$K\hookrightarrow \RR$，从而给出一个嵌入
  \begin{equation}\label{eq:88}
     G=\PSL_2(\sO)\hooklongrightarrow \PSL_2(\RR)\times \PSL_2(\RR). 
  \end{equation}
设 $\HH$ 为上半平面。商
  \begin{equation}\label{eq:89}
     X^0(K) = \bigl(\HH\times \HH\bigr)\bigm/ G 
  \end{equation}
是\emph{Hilbert 模曲面}。群 $G$ 的作用可能有有限稳定化子，即 $X^0(K)$
可能是轨形（orbifold）而非光滑流形。它是非紧的；$X^0(K)$ 有有限多个端，
与 $K$ 的理想类群一一对应。截断每个端，可得一个紧致带边轨形；每个边界
分量是以 $\cir$ 为底、以 $\cir\times \cir$ 为纤维的纤维丛。此外，挖去
每个轨形点的一个邻域，可得光滑紧致带边流形 $X(K)$；新增的边界分量是
透镜空间。透镜空间边界处的号差亏损~\eqref{eq:87} 已由等变号差
定理~\S\ref{subsec:5.1} 知道~\cite[\S10.3]{A7}。
% [Tutor Note: 含糊——透镜空间边界上的号差亏损「was known from the equivariant signature theorem §5.1」：实质是对轨形点的循环群稳定化子用 [A7] §10.3 的等变号差定理计算 $\Sign$（所得边界为透镜空间 $L(p,q)$，公式含 Dedekind 和）；作者未写出公式，查 [A7] §10.3 与 §5.1 的等变号差定理。]
希策布鲁赫~\cite[p.~222]{H4} 通过边界锥上尖点奇点的解消
（desingularization）计算了其余边界处的号差亏损。该公式是奇点所对应的
连分数的一个简单表达式。
% [Tutor Note: 含糊——「desingularization of the cusp singularity in the cone on the boundary」指 Hirzebruch 用尖点奇点的解消（例外除子链的相交数由 Hirzebruch–Jung 连分数给出）把号差亏损写成连分数表达式（[H4] p.~222）；"continued fraction" 就是该解消的连分数。]
另一方面，清水（Shimizu）~\cite{Sh} 引入了与对应该端的理想类相联系的
$L$ 函数。希策布鲁赫~\cite[p.~231]{H4} 证明了尖点处的号差亏损等于
该 $L$ 函数在 $s=1$ 处的值，至多相差一个数值因子。
% [Tutor Note: 含糊——「up to a numerical factor」中的因子需按 [H4] p.~231 的 $L$ 函数归一化核对（是否含 $\pi$、类数以及 $\sqrt d$ 的幂）；这一步是 $\eta$ 不变量与数论 $L$ 函数对应的核心，不要把它当作纯记号问题。]
这一图景可推广到任意次的全实数域。在一般情形下，希策布鲁赫猜测尖点
奇点处的号差亏损仍是清水 $L$ 函数的一个值。该猜想由阿蒂亚–唐纳利–辛格
（Atiyah–Donnelly–Singer）~\cite{ADS} 与米勒（Müller）~\cite{Mu} 于
1982–3 年独立证明。
% [Tutor Note: 术语——号差亏损 $\alpha$ 与 $\eta$ 不变量的关系是 $\alpha(Y)=-\eta(Y)$，其中 $\eta(Y):=\eta(0)$ 是解析延拓值（本节末 \eqref{eq:94} 之后）。作者先用 $\alpha$ 的几何/数论性质（[H4]、[Sh]、[ADS]、[Mu]）引出问题，再用 $\eta$ 不变量给出其解析表达，这条线索不要读反。]


  \subsection{整体边界条件}\label{subsec:7.2}

% [Tutor Note: 含糊——局部椭圆边界条件的「拓扑障碍」见 [AB3]（Atiyah–Bott, The index problem for manifolds with boundary）与本文 §5.2：符号类 $\sigma\in K^0(T^*X,T^*X\setminus0)$ 在边界 $K^0(T^*X|_{\bX})$ 上的限制一般非零；只有当它消失（即 $\sigma$ 提升为相对 K 群中的类）时才存在局部边界条件。对号差算子这一限制非零，因此必须改用整体条件。]
为一般地计算号差亏损~\eqref{eq:87}，自然的想法是在紧致带边流形上
考虑号差算子~\eqref{eq:78}。但我们必须施加一个椭圆边界条件；而对号差
算子而言，\emph{局部}椭圆边界条件（\S\ref{subsec:5.2}）的拓扑障碍非零。
阿蒂亚–帕托迪–辛格用一种新颖的手法克服了这一障碍：他们引入了
\emph{整体}椭圆边界条件，它对任何广义 Dirac 算子（包括号差算子）都存在。
这些整体边界条件如今在带边流形上 Dirac 算子的理论与应用中无处不在。

% [Tutor Note: 术语——$\dbar$ 是 Cauchy–Riemann 算子（$\partial/\partial\bar z$ 的一半），其核是全纯函数而非调和函数，故为无穷维；这正是需要边界条件把核截成有限维的原因。与调和形式的 Hodge 理论对照阅读更清楚。]
作为第一个例子，考虑单位圆盘闭包 $\Omega \subset \CC_z$ 上的 $\dbar$
算子，如 \S\ref{subsec:5.2} 那样。其核由 $\overline\Omega $ 上的全纯函数
组成；一个稠密子空间是多项式空间，即 $\{z^n:n\in \ZZ^{\ge0}\}$ 的线性
张成。这是一个无穷维向量空间。椭圆边界条件必须把它截成有限维子空间。
固定 $a\in \RR\setminus \ZZ^{\ge0}$。设 $\sH_a$ 表示光滑函数
$u\:\overline\Omega \to\CC$ 的子空间，使得 $u\res{\partial \Omega }$ 的
Fourier 展开在 $m>a$ 时 $e^{im\theta }$ 的 Fourier 系数为零，其中我们在
$\partial \Omega $ 上写 $z=e^{i\theta }$。
% [Tutor Note: 易错——$a\in \RR\setminus \ZZ^{\ge0}$ 且条件取 $m>a$：若 $a$ 取非负整数，则 $m=a$ 的 Fourier 系数的归属需另行规定（谱在 $a$ 处有本征值，谱投影 $\lambda<a$ 与 $\lambda\le a$ 给出不同的边界条件）；这正是作者要求 $a$ 避开 $\ZZ^{\ge0}$ 的原因，核对 §5.2 与 [APS1] §1 的相应约定。]
则 $\dbar$ 到 $\sH_a$ 的限制具有有限维的核与余核，并且它在适当的
Sobolev 完备化上延拓为 Fredholm 算子。换言之，限制到 $\sH_a$ 是一个
椭圆边界条件。但由于其定义中用到 Fourier 变换，它不是局部的；
比较~\eqref{eq:54}。
% [Tutor Note: 含糊——"it extends to a Fredholm operator on suitable Sobolev completions" 的实质：在 $H^s(\Omega)$ 上取 $\sH_a$ 的闭包（即边界 Fourier 系数在 $L^2$ 中的闭包），并验证该闭扩张的核、余核有限维且有闭值域；这正是 [APS1] §1 对「椭圆边界条件」所要求的两条性质。]

% [Tutor Note: 易错——\eqref{eq:90} 中 $\sigma$ 是 $dt$ 的 Clifford 乘法，若 $dt$ 取单位 1-形式并采用黎曼约定则 $\sigma^2=-1$，从而 $\sigma\inv=-\sigma$；$A_Y=\sigma\inv D_Y$ 的自伴性、以及 \eqref{eq:91} 中谱 $\lambda$ 的符号都依赖这一约定，核对 §4 的 Clifford 约定。]
这个例子可推广到紧致带边黎曼流形 $X$ 上的 Dirac 算子 $D_X$。
假定度量在 $Y=\bX$ 附近为乘积度量，于是在边界附近有分解
  \begin{equation}\label{eq:90}
     D_X = \sigma \frac{\partial }{\partial t} + D_Y 
  \end{equation}
其中 $t$ 是垂直于 $Y=\bX$ 的测地线的长度坐标，代数算子 $\sigma$ 是
$dt$ 的 Clifford 乘法，$D_Y$ 是 $Y$ 上的 Dirac 算子。于是算子
$A_Y=\sigma \inv D_Y$ 是自伴的。设
  \begin{equation}\label{eq:91}
     \bigoplus\limits_{\lambda \in \spec (A_Y)} \!\!\sE_\lambda 
  \end{equation}
是 $Y$ 上旋量的谱分解。对每个 $a\in \RR\setminus \spec(A_Y)$，
阿蒂亚–帕托迪–辛格整体边界条件把算子限制到 $X$ 上这样的旋量子空间：
其在 $Y$ 上的限制落在 $\oplus _{\lambda <a}\sE_\lambda $ 的完备化中。
事实上，阿蒂亚–帕托迪–辛格选取 $a=0$，我们在下文也如此。若
$0\in \spec(A_Y)$，则须像下文~\eqref{eq:92} 那样单独考虑 $\ker D_Y$。
% [Tutor Note: 易错——取 $a=0$ 时边界条件是「负谱空间」（$\lambda<0$）的闭包；若改取正谱空间，相当于取伴随算子或反转定向，指标公式 \eqref{eq:92} 中 $(\et{\bX}+h)/2$ 的符号要相应改变。用 \eqref{eq:92} 前先固定这个约定，并与 [APS1] 定理 4.14 对照。]
% [Tutor Note: 含糊——"one must take into account $\ker D_Y$ separately" 的实质：零特征空间的归属（取谱投影 $P_{<0}$ 还是 $P_{\le0}$）会改变指标，其效应正是 \eqref{eq:92} 分子中的 $h=\dim\ker D_{\bX}$；所以 $h$ 项不是修正项，而是边界条件定义的一部分。]


  \subsection{Atiyah–Patodi–Singer $\eta$ 不变量}\label{subsec:7.3}

% [Tutor Note: 含糊——"The problem splits into two pieces: a cylinder near $\bX$ and $X\setminus\bX$" 与 "Gluing ... via a partition of unity" 的实质是 [APS1] §4 的两步：柱面上把 $D_X$ 化为以谱参数为变量的常微分算子族并显式计算其指标；内部用热核方法（如 §6.2 的局部指数定理）计算，最后用单位分解与热核的边界层渐近展开把两部分拼接。真正困难的是边界附近热核的渐近展开。]
有了椭圆边界条件，阿蒂亚–帕托迪–辛格进而计算紧致带边流形上一般
Dirac 算子的 $\ind D_X$。问题分为两部分：$\bX$ 附近的柱面区域与
$X\setminus \bX$。在柱面上，他们利用谱分解~\eqref{eq:91} 与乘积度量
分解~\eqref{eq:90}，把 $D_X$ 转化为以 $\spec(A_Y)$ 为参数的一族常微分
算子。在 $\bX$ 的补集上，他们使用热核方法，如局部指数定理
（\S\ref{subsec:6.2}）中那样。用单位分解把两个区域粘合，他们证明了
如下结果。

  \begin{theorem}[Atiyah-Patodi-Singer~\cite{APS1}]\label{thm:16}
 设 $X$ 为紧致带边黎曼流形，并假定黎曼度量在 $\bX$ 的一个邻域内为
乘积度量。设 $D_X$ 为广义 Dirac 算子。则关于整体边界条件，有
  \begin{equation}\label{eq:92}
     \ind D_X = \int_{X}\omega \;-\; \frac{\et{\bX} + h\mstrut _{\bX}}{2}, 
  \end{equation}
其中 $\omega$ 是 $\Ahat$ 亏格的 Chern–Weil 形式，$\et{\bX}$ 是
$\eta$ 不变量，而 $h\mstrut _{\bX}=\dim\ker D_{\bX}$。
  \end{theorem}

% [Tutor Note: 易错——\eqref{eq:92} 的归一化：分母的 $2$ 来自把 $\eta$ 的解析延拓与谱投影（正/负谱空间）拼成 Fredholm 指标，$(h+\eta)/2$ 中的 $h=\dim\ker D_{\bX}$ 是零模的贡献；若把 $h$ 取成别的量（如 $\dim\ker D_X^-$）或把边界条件取成另一侧投影，系数会变成 $(\eta-h)/2$ 或整体反号。核对 [APS1] 定理 4.14 的约定。]
\noindent
 为在 $Y=\bX$ 上定义 $\eta$ 不变量，令 $A_Y=\sigma\inv  D_Y$，并构造
  \begin{equation}\label{eq:93}
     \eta (s) = \sum\limits_{\lambda \in \spec(A_Y)\setminus \{0\}}
     (\sign\lambda )\,|\lambda |^{-s},\qquad \Re(s)>\!>0. 
  \end{equation}
% [Tutor Note: 易错——\eqref{eq:93} 只对非零谱求和：零特征值若出现在求和中，$\sign 0$ 无定义且 $|\lambda|^{-s}$ 项发散，这正是 \eqref{eq:92} 中另设 $h$ 的原因。另注意 $\eta(0)$ 是解析延拓值，不是级数在 $s=0$ 处的和（级数只在 $\Re(s)$ 充分大时收敛）。]
这是 $L$ 函数的黎曼版本，是 Hilbert 模曲面上号差亏损中数论 $L$ 函数的
回响，也是 $\zeta$ 函数的黎曼版本~\eqref{eq:57} 的一个变体。\eqref{eq:93}
中的无穷和在 $\Re(s)>\!>0$ 时收敛，并可亚纯延拓\footnote{值得注意的是，
$\eta (s)$ 在 $\Re(s)>-1/2$ 上全纯。}到复 $s$ 平面，且 $s=0$ 是正则点。
定义 $\et Y=\eta (0)$。我们指出，~\cite{APS1} 中的主要定理适用于更一般的
一阶椭圆微分算子。
% [Tutor Note: 含糊——解析延拓的存在性、有限性与脚注中「在 $\Re(s)>-1/2$ 上全纯」的断言依赖 Seeley 的复幂理论（见 [Se4]，以及 [Se1]、[Se2]）与 [APS1] §2 的定理；一般的亚纯延拓可以有极点，$s=0$ 处正则并非自动成立，而是 APS 定理的实质内容之一。作者以一句「there is a meromorphic continuation」带过。]
% [Tutor Note: 易错——$\eta$ 在流形反向时变号：反向使 $A_Y$ 变为 $-A_Y$（或复共轭），故 $\eta(-A_Y)=-\eta(A_Y)$，于是 $\eta(\bar Y)=-\eta(Y)$。这与 \eqref{eq:94} 中 $\Sign(X)=\int_X\omega-\et{\bX}$ 的奇偶性一致：$\Sign$ 与 $\int\omega$ 都变号时，$-\eta$ 保持同一变换规律。]

% [Tutor Note: 易错——\eqref{eq:94} 中边界算子的公式 $(-1)^{k+q+1}(*d-d*)$ 须按 [APS2] 的约定理解：$*$ 为 Hodge 星、$d$ 为外微分，$k$ 来自 $X$ 的维数 $4k$（故 $\dim\bX=4k-1$），$q$ 为形式的次数。不同文献把它写成 $\pm(*d-d*)$ 或等价的 $\pm(d*+*d)$ 等形式，符号与 $k,q$ 的平移约定不同；用 \eqref{eq:94} 前请核对 [APS2] §4 与 [APS1] 的 $\eta$ 定义。]
指数定理~\ref{thm:16} 对号差算子~\eqref{eq:78} 有更简单的形式。设 $X$ 为
紧致定向黎曼 $4k$-流形，且度量在 $\bX$ 附近为乘积度量。对称双线性
型~\eqref{eq:8} 限制在 $H^{2k}(X,\bX;\RR)$ 在 $H^{2k}(X;\RR)$ 中的像上
是非退化的，而 $\Sign(X)$ 就是它的号差。

  \begin{corollary}[Atiyah-Patodi-Singer~\cite{APS1}]\label{thm:17}
 在此情形下，
  \begin{equation}\label{eq:94}
     \Sign(X) = \int_{X}\omega \;-\;\et{\bX}, 
  \end{equation}
其中 $\omega$ 是 $L$ 亏格的 Chern–Weil 形式，而 $\et{\bX}$ 是
$\Omega ^{\textnormal{偶数}}(\bX)$ 上由 $(-1)^{k+q+1}(*d-d*)$ 在
$\Omega ^q(\bX)$ 上给出的自伴算子的 $\eta$ 不变量。
  \end{corollary}

% [Tutor Note: 含糊——\eqref{eq:94} 与 \eqref{eq:92} 相比不见了 $h$ 项（\eqref{eq:92} 为 $(\et{\bX}+h)/2$）。这一「简化」依赖号差算子的特殊约定（零特征空间与谱投影的处理方式由号差算子的奇偶性决定），作者未加说明；严格核对须看 [APS1] 定理 4.14 与 [APS2] §4 中 $h$ 的定义，以及 $-\eta$ 即号差亏损 \eqref{eq:87} 这一事实。]
\noindent
 特别地，$-\eta$ 就是号差亏损~\eqref{eq:87}。注意号差亏损是 $\bX$ 上
一个自然微分算子的\emph{谱不变量}，这一性质从其定义并不能直接看出。
% [Tutor Note: 术语——「谱不变量」指只依赖 $\bX$ 上号差算子谱（特征值）的量，而不像示性数那样是局部曲率多项式的积分。号差亏损的定义 \eqref{eq:87} 用的是 4-流形的号差与 $p_1$ 积分，二者都不是谱量，但差值却是谱量——这正是 APS 定理的惊人之处。]

  \begin{remark}[]\label{thm:18}
 号差亏损在阿蒂亚后来的两篇论文~\cite{A8,A9} 中也扮演了重要角色。
  \end{remark}

% [Tutor Note: 易错——\eqref{eq:95} 定义的是 $\RZ$ 值映射 $\xia{Y/S}$（模 $1$ 的等价类），不是实值函数；只有选定实值提升后才能谈 \eqref{eq:96} 的微分。$(\et{Y/S}+h)/2$ 的模 $1$ 不确定性恰好是族指标定理所给的整性，故 \eqref{eq:95} 把 $S\to\RZ$ 的\emph{同伦类}精化为具体的几何映射。]
$\eta$ 不变量是 Dirac 算子的\emph{几何}不变量的第一个例子。为说明
这一点，回顾例~\ref{thm:24}。设 $Y\to S$ 是相对维数为奇数的逆紧
（proper）纤维丛，并带有相对黎曼 spin 结构。由这些几何数据我们得到
以 $S$ 为参数的一族自伴 Dirac 算子。拓扑指数的最低阶部分是一个映射
$S\to\RZ$ 的\emph{同伦类}。表达式
  \begin{equation}\label{eq:95}
     \xia{Y/S} = \frac{\et{Y/S} + h\mstrut _{Y/S}}{2}\!\!\!\!\pmod1\:S\longrightarrow \RZ 
  \end{equation}
来自~\eqref{eq:92}，它把该同伦类精化为一个具体映射，即所论的几何
不变量。此外，定理~\ref{thm:16} 蕴含~\eqref{eq:95} 的微分是
  \begin{equation}\label{eq:96}
     d\xia{Y/S} = \int_{Y/S}\omega , 
  \end{equation}
其中 $\omega$ 是~\eqref{eq:92} 中的微分形式。这是几何不变量
$\xia{Y/S}$ 的一种「曲率」公式；Dirac 算子的其他几何不变量也有类似的
对应物。
% [Tutor Note: 易错——\eqref{eq:96} 中 $d\xia{Y/S}$ 是对实值提升取的微分，而 $\xia{Y/S}$ 只定义到模 $1$；两者相容是因为 $\int_{Y/S}\omega$ 的积分在换提升时不变（整数部分被模 $1$ 吸收）。与 \eqref{eq:111} 比较时注意这里的 $d$ 是 $S$ 上的外微分，而非流形 $X$ 上的。]
% [Tutor Note: 含糊——"Atiyah-Patodi-Singer [§7 APS3], in collaboration with Lusztig" 给出的另一种解析计算：谱流 = 绕数。其严格证明需要族的谱在 $\cir$ 上构成 $\cir\times\RR$ 的闭子集（由自伴性与逆紧性保证），从而相交数有定义；见 [APS3] §7（该文与 Lusztig 合作的结果）以及 [APS1] 中 $\eta$ 的解析性质。]
阿蒂亚–帕托迪–辛格~\cite[\S7]{APS3} 与卢斯蒂格（Lusztig）合作，给出了
映射 $S\to\RZ$ 的同伦类的另一种解析计算。对任意环路 $\cir\to S$，
他们证明复合映射 $\cir\to S\to \RZ$ 的绕数等于由 $\cir$ 参数化的拉回
Dirac 算子族的\emph{谱流}。谱流带符号地计数 $\xi$ 不变量~\eqref{eq:95}
在我们绕 $\cir$ 一周时发生的整数跳变。等价地说，诸 Dirac 算子的谱之并
是闭子集 $C \subset \cir\times \RR$；谱流就是 $C$ 与 $\cir\times \{0\}$
的相交数。
% [Tutor Note: 术语——谱流与指标：谱流是 $\cir$ 参数族 Dirac 算子的整数不变量，可理解为 $\cir\times Y$ 上的族指标（或映射环面上的指标），也是 \eqref{eq:96} 的「曲率」在 $1$ 维底空间上的积分；不要把它与单个算子的 Fredholm 指标 \eqref{eq:34} 混为一谈。]
  \begin{example}[]\label{thm:25}
 最简单的非平凡谱流出现在复自伴 Dirac 算子族
  \begin{equation}\label{eq:100}
     D_s = \sqrt{-1} \frac{d}{dx} + s,\qquad 0\le s\le 1, 
  \end{equation}
的情形，它们作用在带坐标 $x$ 的 $\RR/2\pi \ZZ$ 上。算子 $D_1$ 与 $D_0$
同构：用乘以 $e^{\sqrt{-1}x}$ 的乘法算子共轭即可。谱之并
$C\subset \cir\times \RR$ 是一条螺旋线。
  \end{example}
% [Tutor Note: 易错——本例谱为 $\{s+n:n\in \ZZ\}$，$s$ 从 $0$ 走到 $1$ 时恰有一个分支穿过 $0$（$n=0$ 分支由 $0$ 移到 $1$，同时 $n=-1$ 分支由 $-1$ 移到 $0$），故谱流为 $1$（「从负到正」取正号）。不要因为 $s=0$ 与 $s=1$ 两端各有一个零特征值就误判为 $0$ 或 $2$。]


  \subsection{次级几何不变量}\label{subsec:7.4}

% [Tutor Note: 术语——初级不变量（示性数、指标）是整值拓扑量；次级不变量是其「积分」的依赖联络的修正，取 $\RZ=\RR/\ZZ$ 值。总测地曲率 $\int\kappa/2\pi$ 就是最简单的次级不变量：它与和乐角只差 $2\pi$ 的整数倍，故模 $\ZZ$ 后等于 Levi-Civita 联络的和乐，这解释了为什么 CS 型次级不变量天然是 $\RZ$ 值的。]
最初等的次级不变量是黎曼 2-流形中曲线的总测地曲率；它出现在
Gauss–Bonnet 公式~\eqref{eq:85} 中。与之相伴的初级拓扑不变量是 Euler 数。
下面的推广更接近于总测地曲率模 $\ZZ$ 的约化，后者（至多相差一个符号）
就是 Levi-Civita 联络的和乐。

% [Tutor Note: 易错——Chern–Weil 的曲率约定：$\Omega = d\Theta + \frac{1}{2}[\Theta,\Theta]$ 与 $\Omega = d\Theta + [\Theta,\Theta]$ 两种写法取决于 $\Theta$ 的取值是否含因子；$p(\Omega)$ 的符号与系数随之改变。对照本文 §3 与 [CS] 时务必先固定 $\Theta = g^{-1}dg$（左不变）还是 $\Theta = dg\,g^{-1}$（右不变）的约定。]
设 $G$ 为只有有限多个连通分支的 Lie 群，$\pi \:P\to X$ 为主 $G$-丛，
$\Theta \in \Omega ^1(P;\mathfrak{g})$ 为联络形式。设
$p\in \bigl(\Sym^k\mathfrak{g}^* \bigr)^G$ 为 Lie 代数 $\mathfrak{g}$ 上的
Ad-不变多项式。正如我们已经用过的，Chern–Weil 理论对这批数据给出一个
闭微分形式 $p(\Omega )\in \Omega ^{2k}(X)$，它只依赖于 $\Theta $ 的曲率
$\Omega $。此外，它对联络 $\Theta $ 是自然的。其 de Rham 上同调类
$\bigl[p(\Omega )\bigr]\in H^{2k}(X;\RR)$ 与 $\Theta $ 无关，因而是主丛
$\pi $ 的不变量。若 $G$ 紧致（我们此后如此假定），则该不变量来自 $G$
的分类空间上同调中的一个示性类 $c_p\in H^{2k}(BG;\RR)$。1972 年，
陈省身–西蒙斯~\cite{CS} 引入了附着于 $c_p$ 到整上同调类
$c\in H^{2k}(BG;\ZZ)$ 的精细化的一个次级几何不变量。（这种精细化只有当
$c_p$ 的周期为整数时才存在。）
% [Tutor Note: 含糊——"Refinements exist only if the periods of $c_p$ are integers"：这是必要条件，其充分性等价于 $c_p$ 落在 $H^{2k}(BG;\ZZ)\to H^{2k}(BG;\RR)$ 的像中（可查 [ChS] 的 Cheeger–Simons 理论）。作者未说明 "periods" 指在闭 $2k$ 维流形上的配对值，也未给出构造：$d\,\mathrm{CS}(\Theta) = p(\Omega)$，其中 $\mathrm{CS}$ 由沿 $\Theta_t=(1-t)\Theta_0+t\Theta_1$ 的积分给出。]
% [Tutor Note: 笔误——原文作「1972 年」陈–西蒙斯引入该不变量，但 [CS] 于 1974 年才刊出（1971 年提交）；此年份应指提交/预印本，引用时以 [CS] 的发表年 1974 为准。]
在此情形下，初级的 $\ZZ$ 值不变量是 $2k$ 维闭定向流形上主丛
$\pi \:P\to X$ 的示性数。次级的 $\RZ$ 值\emph{陈–西蒙斯不变量}则对带
联络 $\Theta $ 的 $\pi \:Q\to Y$ 定义，其中 $Y$ 是闭定向 $(2k-1)$ 维
流形。次级 $\RZ$ 值不变量依赖于联络，而初级 $\ZZ$ 值不变量是拓扑的。

% [Tutor Note: 笔误——\eqref{eq:97} 写作 $p_1(W)[W]$，但本段刚把 4-流形记作 $X$，$W$ 未定义；\eqref{eq:99} 又写成 $p_1(TW)[W]$。疑原文应以 $X$（或 $TX$）替代 $W$；照录原文，请核对 [H4] 或 §2 中示性数记号。]
  \begin{example}[]\label{thm:19}
 对 $N\in \ZZ^{\ge3}$，设 $p_1\in H^4(\BSO_N;\ZZ)$ 为主 $\SO_N$-丛的
万有第一 Pontryagin 类。内蕴地（即用流形的切几何）来说，闭定向
4-流形 $X$ 的 $\ZZ$ 值初级不变量是
  \begin{equation}\label{eq:97}
     p_1(W)[W]. 
  \end{equation}
闭定向黎曼 3-流形 $Y$ 上的次级不变量 $\Gamma (Y)$ 是其 Levi-Civita 联络
的陈–西蒙斯不变量。在此情形下~\cite[\S6]{CS}，$\Gamma (Y)$ 是共形不变量，
并且是存在共形浸入 $Y\to\EE^4$ 的障碍，其中 $\EE^4$ 是欧氏 4-空间。
  \end{example}
% [Tutor Note: 含糊——$\Gamma (Y)$ 的共形不变性与「共形浸入 $Y^3\to\EE^4$ 的障碍」出自 [CS] §6；这里 $\Gamma$ 是 Levi-Civita 联络的 CS 不变量（3 维、$\mathrm{SO}_3$ 丛），与 \eqref{eq:98} 的 $\xia Y$ 相差一个归一化因子（见例~\ref{thm:22} 中 $-24\xia Y$ 与 $\Gamma(Y)$ 的比较）。]

  \begin{remark}[]\label{thm:20}
 陈–西蒙斯不变量在\emph{微分上同调}~\cite{ChS,F2,HS} 中有自然的表达，
 后者把初级与次级不变量统一在同一个框架内。
  \end{remark}
% [Tutor Note: 术语——微分上同调（differential cohomology）的类同时携带上同调类（初级）、联络/曲率（次级）与挠信息；[ChS] 是 Cheeger–Simons 微分特征，[F2] 是 Freed 的经典 Chern–Simons 理论，[HS] 是 Hopkins–Singer 关于二次函数与微分上同调的表述。与 \eqref{eq:98} 的 $\RZ$ 值 $\eta$ 不变量对照，可见它应由微分 $K$ 理论（而非微分上同调）承载，见注~\ref{thm:23}(2)。]

% [Tutor Note: 易错——\eqref{eq:98} 中 $\pmod1$ 的归一化：$\eta$ 与 $h$ 都是实数，除以 $2$ 后模 $1$；若改用 $\eta/2$（不除 $h$）或把 $\eta$ 的定义 \eqref{eq:93} 换成不取 $\sign\lambda$ 的版本，结果会整体相差 $h/2$ 或符号。核对 [APS1]/[APS2] 的 $\eta$ 归一化。]
阿蒂亚–帕托迪–辛格的 $\eta$ 不变量是类似于陈–西蒙斯不变量的次级
不变量，但它属于指数理论而非示性类理论。设 $Y$ 为 spin 黎曼流形，
$\pi \:Q\to Y$ 为带联络的主 $G$-丛。陈–西蒙斯的故事从整上同调类
$c\in H^{\bullet }(BG;\ZZ)$ 出发，而定义 $\eta$ 不变量则从复\footnote{对实
表示与实 $KO$ 理论有相应的精细化。}线性表示 $\rho \:G\to\Aut(V)$ 出发；
$\rho $ 的同构类是等变 $K$ 理论群 $K^0_G(\pt)$ 中的一个元素。从一开始
就能看出：陈–西蒙斯不变量属于整上同调，而 $\eta$ 不变量属于 $K$ 理论。
在 $Y$ 上构造与带继承协变导数的相伴向量丛 $V_Q\to Y$ 耦合的 Dirac 算子。
于是 $\RZ$ 值不变量
  \begin{equation}\label{eq:98}
     \xia Y(V) = \frac{\et Y(V) + h\mstrut _Y(V)}{2}\pmod1 
  \end{equation}
光滑依赖于黎曼度量与 $\pi $ 上的联络。\footnote{在号差算子的特殊情形，
其核——调和形式——具有上同调意义，$\RR$ 值不变量是光滑的。（它出现在
~\eqref{eq:94} 中。）}相应的初级不变量是 $\dim Y+1$ 维流形上带主 $G$-丛的
扭化 Dirac 算子的 $\ZZ$ 值指标。
% [Tutor Note: 易错——"depends smoothly" 的断言依赖族指标定理（[APS3]）：$\eta$ 不变量本身只是 $\RR/\ZZ$ 值，其光滑性指模 $1$ 后的光滑映射；脚注中号差算子情形的 $\RR$ 值光滑性来自调和形式的上同调解释，二者不同，不要混淆。]

  \begin{remark}[]\label{thm:21}
 为得到非拓扑的不变量，必须要求 $\dim Y$ 为奇数。$\eta$ 不变量在偶数维
 也导致新的、有趣的拓扑不变量，例如在带 pin 结构的非定向流形
 上~\cite{Gi2}。
  \end{remark}

  \begin{example}[]\label{thm:22}
 考虑如同例~\ref{thm:19} 的 $\dim Y=3$，但现在假定 $Y$ 是闭 spin 黎曼
流形。对标准 Dirac 算子（没有主 $G$-丛），表达式~\eqref{eq:98} 是闭
spin 4-流形 $X$ 上 Dirac 算子指数的 $\RZ$ 值次级不变量，而按指数
定理~\eqref{eq:42}，该指数为
  \begin{equation}\label{eq:99}
     \ind D_X = -\frac{1}{24}\,p_1(TW)[W]. 
  \end{equation}
它与~\eqref{eq:97} 的差别在于有理因子。$\Ahat$ 亏格的整性
（回忆~\eqref{eq:24}）蕴含~\eqref{eq:99} 是整数。转到次级不变量
$\Gamma (Y)$ 与 $\xia Y$，由于它们是 $\RZ$ 值的，我们不能把
$\Gamma (Y)$ 乘以像 $-1/24$ 这样的非整有理数；只能先清分母，再比较
$-24\xia Y$ 与 $\Gamma (Y)$。在所论情形下二者一致；在更一般情形下的
类似比较给出一个 spin 配边不变量。无论如何，我们看到 $\eta$ 不变量是比
陈–西蒙斯不变量更精细的不变量。换个说法，基于 $K$ 理论的次级不变量
包含了比基于整上同调的次级不变量更精细的信息。这与从初级不变量的整性
所得出的更强拓扑信息相呼应；见 \S\ref{subsec:2.4}。
  \end{example}
% [Tutor Note: 易错——\eqref{eq:99} 的系数 $-1/24$ 来自 $\hat A_1=-p_1/24$（与 \eqref{eq:24} 的整性定理配套）；不要把它与 \eqref{eq:86} 的 $p_1/3$ 相混。$\hat A$ 亏格的整性（在 spin 流形上）是这里「$-24\xia Y$ 与 $\Gamma(Y)$ 可比较」的根据。]
% [Tutor Note: 含糊——"a similar comparison in more general circumstances leads to a spin bordism invariant"：作者未指出该配边不变量的确切群与计算；其含义是次级不变量的比较在（一般维数或带 $G$-丛的）配边群上给出 $\RZ$ 值不变量（见 [APS2] 及 [Gi2]；注意 $\Omega^{Spin}_3=0$，3 维情形本身不产生非平凡同态）。清分母这一步正是 $\RZ$ 值与有理系数示性数不能直接相乘的技术后果。]

  \begin{remark}[]\label{thm:23}
 \ 


 \begin{enumerate}

 \item 作为这种额外威力的具体例证，阿蒂亚–帕托迪–辛格~\cite[\S4]{APS2} 用
号差算子的 $\eta$ 不变量精化了陈–西蒙斯关于共形嵌入 $Y^3\to\EE^4$ 的
障碍。他们还说明了如何用 $\eta$ 不变量构造 Adams $e$ 不变量——标架
配边（framed bordism）的一个不变量。

 \item 初级的指标与次级的 $\eta$ 不变量在微分 $K$ 理论的框架中得到统一；
比较注~\ref{thm:20}。见~\cite{FL} 及其中的参考文献。

 \item 阿蒂亚–帕托迪–辛格的 $\eta$ 不变量出现在几何乃至更广的许多语境中。
它还与拓扑与几何指数理论中的其他特征标一起，在理论物理中多次出现：
量子场论、弦论与凝聚态理论。

 \end{enumerate}
  \end{remark}
% [Tutor Note: 术语——Adams $e$ 不变量是稳定同伦论中标架配边 $\Omega^{fr}_*$ 的不变量，[APS2] §4 说明它可由 $\eta$ 不变量构造；而 "framed bordism" 与本文 §2 的配边记号一致，勿与 spin 配边混淆。]


  \subsection{Dirac 算子的行列式}\label{subsec:7.5}

% [Tutor Note: 易错——\eqref{eq:101} 只是形式记号（原文用引号标出 $=$）：无穷乘积 $\prod\lambda_m$ 一般发散，且若 $\lambda_1=0$ 则平凡地为零；故必须（i）剔除零特征值，(ii) 用 $\zeta$ 正规化（\eqref{eq:104}）赋值。$X=\cir$ 时 $\prod m^2$ 发散这一点正说明需要正规化。]
在转向 Dirac 算子之前，先考虑闭流形 $X$ 上的二阶 Laplace 算子
$\Delta $，如 \S\ref{subsec:6.1} 中那样。则 $\Delta $ 有离散谱，由特征值
$0\le\lambda _1\le\lambda _2\le\cdots$ 组成（按重数计）。形式上，
$\Delta $ 的行列式是
  \begin{equation}\label{eq:101}
     \det\Delta \;\textnormal{``$=$''} \;\prod\limits_{m=1}^{\infty}\lambda _m. 
  \end{equation}
（假定 $\lambda _1>0$，或剔除零特征值以避免 $\det\Delta =0$。）当然，
这个无穷乘积发散。例如当 $X=\cir$、$\Delta $ 是通常的标量 Laplace 算子时，
在剔除零特征值后，该无穷乘积在相差一个常数的意义下是
$\prod_{m=1}^{\infty}m^2$。给这个无穷乘积赋值的一种办法，由雷–辛格
（Ray–Singer）~\cite{RS} 于 1971 年首创，并沿用了复分析中熟悉的
技巧~\cite{JL}，就是利用 $\zeta$ 函数~\eqref{eq:57} 的解析延拓，其定义为
  \begin{equation}\label{eq:102}
     \zD(s) = \sum\limits_{m=1}^{\infty}\lambda ^{-s},\qquad s\in \CC, 
  \end{equation}
其中 $\Re (s)>\!>0$。
% [Tutor Note: 笔误——\eqref{eq:102}、\eqref{eq:103} 中 $\lambda^{-s}$ 与 $\log\lambda$ 均未标出下标 $m$，而求和指标为 $m$；疑应为 $\lambda_m^{-s}$、$\log\lambda_m$。照录原文，请核对 [RS] 或 §6.1 的 $\zeta$ 函数定义。]
于是在无穷和绝对收敛的区域内，我们求导
  \begin{equation}\label{eq:103}
     -\zeta_{\Delta }'(s) = \sum_{m=1}^{\infty}\lambda ^{-s}\log\lambda , 
  \end{equation}
然后利用 $\zD$ 的解析延拓在 $s=0$ 处的正则性来定义
  \begin{equation}\label{eq:104}
     \det\Delta := e^{-\zeta_{\Delta }'(0)} .
  \end{equation}
% [Tutor Note: 含糊——$\zeta_\Delta(s)$ 的解析延拓存在、且 $s=0$ 为正则点，这依赖 Seeley 的复幂/热核理论（[Se4]）与 [RS] §1；作者以「use the regularity」一句带过。另注意 $\det\Delta$ 无尺度不变性：替换 $\Delta\mapsto c\Delta$ 会引入 $e^{+\zeta_\Delta(0)\log c}$ 型因子。]
对 $\cir$ 上的标量 Laplace 算子，\eqref{eq:102} 定义的椭圆 $\zeta$ 函数
在相差一个常数的意义下归结为黎曼 $\zeta$ 函数。

% [Tutor Note: 易错——行列式线 \eqref{eq:106} 的定义要求 $\dim V^0=\dim V^1$：此时 $\Det T$ 是 $\Det V^0\to\Det V^1$ 的映射（即一条线上的元素），而不是数；若 $V^0=V^1$，$\det T$ 才是通常的数值行列式。对一般的 Fredholm 算子，$V^0,V^1$ 无穷维，\eqref{eq:123} 不再有意义，必须改用 $\ker$ 与 $\coker$ 的有限维性。]
一阶 Dirac 算子~\eqref{eq:41} 不是自伴的；其定义域与值域不同。一个
有限维模型是不同向量空间 $V^0,V^1$ 之间的线性算子
  \begin{equation}\label{eq:105}
     T\:V^0\longrightarrow V^1 
  \end{equation}
在每阶外幂上都有诱导映射
${\textstyle\bigwedge} ^qT\:{\textstyle\bigwedge} ^qV^0\to {\textstyle\bigwedge} ^qV^1$。
若 $\dim V^0=\dim V^1$，则对 $q=\dim V^i$ 的诱导映射就是行列式
  \begin{equation}\label{eq:106}
     \det T\:\Det V^0\longrightarrow \Det V^1, 
  \end{equation}
其中 $\Det V^i = {\textstyle\bigwedge} ^{\dim V^i}(V^i)$ 是行列式线。若
$V^0=V^1$，则算子 $\det T$ 就是乘以数值行列式。但一般地\footnote{若
$\dim V^0\neq \dim V^1$，定义 $\det T=0$。}$\det T$ 是某条直线中的元素，
即一维向量空间
  \begin{equation}\label{eq:123}
     \Hom(\Det V^0,\Det V^1)
  \end{equation}
中的元素，这个一维空间称为\emph{行列式线}。行列式构造可推广到 Fredholm
算子~\eqref{eq:105}，此时 $V^0,V^1$ 通常是无穷维的。若 $V^0,V^1$ 无穷维，
式~\eqref{eq:123} 不再有意义；取而代之，行列式线要用核与余核的有限维性
来定义。奎伦（Quillen）~\cite{Q} 在 Fredholm 算子空间上构造了行列式线丛
$\pi \:\Det\to \Fred$，并带有连续截面 $\det\:\Fred\to \Det$。它是继数值
指标~\eqref{eq:34} 之后 Fredholm 算子的下一个拓扑不变量。
% [Tutor Note: 含糊——Quillen [Q] 在 $\Fred$ 上构造行列式线丛 $\pi:\Det\to\Fred$：对 $S$ 上的 Fredholm 算子族 $T:S\to\Fred$，拉回 $T^*\Det\to S$ 的纤维为 $\det(\ker T_s)\otimes\det(\coker T_s)^*$（在指标局部常值的分层上粘合），且 $\det$ 给出连续截面；该族的同构类正是 \eqref{eq:108} 中的 $\Det\in H^2(S;\ZZ)$（拓扑部分）。这里「next topological invariant」指该 $H^2$ 类，而非数值指标 \eqref{eq:34}。]

设 $f\:X\to S$ 为纤维维数为偶数 $n$ 的逆紧纤维丛，并假定 $f$ 带有
相对黎曼 spin 结构，如例~\ref{thm:24}。对带协变导数的复向量丛
$V\to X$，所得复 Clifford 线性 Dirac 算子族 $D_{X/S}(V)$ 的指标由前推
  \begin{equation}\label{eq:107}
     f_!\:K^0(X)\longrightarrow K^{-n}(S). 
  \end{equation}
计算。数值指标是 $f_![V]$ 在 $K^{-n}(S)\to H^0(S;\ZZ)$ 下的像。
与~\eqref{eq:83} 平行，拓扑指数的次低阶部分由自然映射
  \begin{equation}\label{eq:108}
     \Det\:K^{-n}(S)\longrightarrow H^2(S;\ZZ). 
  \end{equation}
计算。族指标定理蕴含 Fredholm 行列式线丛
  \begin{equation}\label{eq:109}
     \Det D_{X/S}(V)\longrightarrow S 
  \end{equation}
的\emph{拓扑}等价类是 $\Det f_![V]\in H^2(S;\ZZ)$。

% [Tutor Note: 易错——\eqref{eq:110} 的相位约定：$e^{-2\pi i\xia{X_\varphi(\epsilon)}(V)}$ 中 $\xi$ 取 \eqref{eq:98} 的 $\RZ$ 值，$2\pi i$ 因子使其与和乐的 $U(1)$ 值匹配；若把 $\xi$ 换成 $-\xi$ 或 $\eta$ 而非 $(\eta+h)/2$，会得到 $e^{+2\pi i\cdots}$ 或差一个符号。核对 [BiF2] 的符号约定。]
这一拓扑部分的指标有一个\emph{几何}精细化。其解析表达是行列式线
丛~\eqref{eq:109} 上的埃尔米特联络及相容的协变导数。\emph{Quillen
度量}~\cite{Q} 用 $\zeta$ 函数行列式~\eqref{eq:104} 构造，而协变导数则
用类似的——但更微妙的——$\zeta$ 函数~\cite{BiF1}。带协变导数的线丛的
同构类由绕环路的和乐决定。沿环路 $\varphi \:\cir\to S$ 的基变换给出
纤维丛 $X_\varphi \to \cir$ 与向量丛 $V_\varphi \to X_{\varphi }$。
$\cir$ 上的可延拓 spin 结构与 $X_\varphi /\cir$ 上的相对
spin 结构相结合，产生 $X_\varphi $ 上的 spin 结构。任选 $\cir$ 上的
度量 $g\mstrut _{\cir}$，并令 $X_\varphi (\epsilon )$ 为 $X_\varphi $ 配以
黎曼度量 $g\mstrut _{\cir}\,/\,\epsilon ^2\,\oplus \,g\mstrut _{X_\varphi /\cir}$
所得的流形。（直和是相对于相对黎曼结构中的水平分布取的。）则绕环路
$\varphi $ 的和乐为~\cite{BiF2}
  \begin{equation}\label{eq:110}
     \hol_\varphi \Det D_{X/S}(V) = \lim\limits_{\epsilon \to0}e^{-2\pi i\xia
     {X_\varphi (\epsilon )}(V)}, 
  \end{equation}
其中指数中的量是阿蒂亚–帕托迪–辛格的 $\eta$ 不变量~\eqref{eq:98}。这个和乐
公式受到 Witten 整体反常~\cite{Wi3} 的启发。行列式线丛的曲率是
  \begin{equation}\label{eq:111}
     \curv\Det D_{X/S}(V) = 2\pi i\int_{X/S}\omega , 
  \end{equation}
其中 $\omega$ 是代表 $\Ahat(X/S)\smallsmile\ch(V)$ 的 Chern–Weil 形式；
比较~\eqref{eq:96}。
% [Tutor Note: 含糊——"the \emph{adiabatic limit} in \eqref{eq:110} was introduced in [Wi3]"：绝热极限 $\epsilon\to0$ 使纤维方向的度量相对底空间变得巨大（$g_\cir/\epsilon^2$），谱的间隙行为决定极限存在；其严格分析见 [BiF2]（Bismut–Freed 的和乐定理），其他几何解释见 [Che]、[S2]。作者未说明极限为何收敛。]
% [Tutor Note: 易错——\eqref{eq:111} 的 $2\pi i$ 与 $\omega$ 的 Chern–Weil 归一化必须匹配：$\omega$ 代表 $\Ahat\ch$ 时，曲率形式按 $\curv = 2\pi i\,\omega$ 的约定给出第一 Chern 类；对照 \eqref{eq:96} 的 $d\xi=\int_{Y/S}\omega$（无 $2\pi i$）可见两处因子的来源不同。]

  \begin{remark}[]\label{thm:26}
 \ 
 \begin{enumerate}

 \item \eqref{eq:110} 中的\emph{绝热极限}由~\cite{Wi3} 引入。Witten 整体
反常公式的其他几何解释见~\cite{Che,S2}。

 \item 一旦知道线丛 $L\to \cir$ 的和乐就是 $D_{\cir}(L)$ 的指数化
$\eta$ 不变量，和乐公式~\eqref{eq:110} 就可看作联系几何不变量 $\eta$
与 $\Det$ 的 Fubini 定理。

 \item 如注~\ref{thm:23}(2) 那样，带度量与协变导数的行列式线丛的同构类
可由微分 $K$ 理论中的前推计算。这精化了（拓扑）指数定理——后者把
\eqref{eq:109} 的拓扑同构类计算为 $\Det f_![V]$。

 \end{enumerate}
  \end{remark}
% [Tutor Note: 术语——Quillen 度量与 [BiF1] 的协变导数用的是不同的 $\zeta$ 函数：前者用 $-\zeta'(0)$（\eqref{eq:104}），后者用沿族参数方向的 $\zeta$ 函数（即对族 Dirac 算子作超迹/热核展开），二者合起来才给出 Hermite 线丛的结构；不要把它们混为同一个 $\zeta$ 正规化。]

% ==================== §8 ====================

\section{反常与指数理论}\label{sec:8}

\begin{keytakeaways}
  \item 反常是量子理论射影本性的表现：经典对称群在量子系统中只能射影地作用，
        而反常正是把这一射影作用线性化（提升为中心扩张）的阻碍。
  \item 对旋量场，费米子路径积分的结果是狄拉克（Dirac）算子族的行列式线丛的
        截面；该线丛的第一陈类是族指标的拓扑部分
        （$\Det f_![V]\in H^2(S;\ZZ)$），其曲率是“局部反常”，和乐则是威滕
        （Witten）的“整体反常”。
  \item 阿蒂亚（Atiyah）的拓扑场论公理把拓扑场论定义为配边范畴
        $\Bd\to \Vk$ 的对称幺半函子，从而把配边理论“范畴化”。
  \item 综合：反常本身是一个\emph{可逆}的 $(n+1)$ 维场论；在稳定同伦论中可逆
        场论就是谱的映射，而旋量场反常的一般公式（\cite[Conjecture~9.70]{FH}）
        中用到阿蒂亚–博特–夏皮罗映射 $M\!\Spin\to KO$。
  \item 阅读提示：§\ref{subsec:8.1} 给出“射影性 $\Rightarrow$ 群扩张”的软性
        框架，§\ref{subsec:8.3} 给出其几何实现（族指标定理），
        §\ref{subsec:8.2} 提供公理化语言，§\ref{subsec:8.4} 把二者融合。
\end{keytakeaways}

% [Tutor Note: 术语——本节“anomaly（反常）”至少有四义，读时须分清：(i) 手征/微扰反常，由 $\frac{1}{2\pi}\int F$ 型积分表达的对称性破缺；(ii) 整体反常，即把射影对象提升为线性对象时的拓扑阻碍，由和乐测得（[Wi3]）；(iii) 引力反常，见 [AgW]；(iv) 反常流入，用高一维的经典作用量抵消低维边界上的反常。作者在本节把“反常”统一提升为量子理论射影性的度量，见 \eqref{eq:114} 与 §\ref{subsec:8.4}。]

部分受伊萨多·辛格（Is Singer）之倡导的启发——他主张理论物理是几何学的沃土——
迈克尔·阿蒂亚（Michael Atiyah）从 1970 年代末起把注意力转向量子场论、其后是
弦论中的几何问题。本卷中西蒙·唐纳森（Simon Donaldson）的论文 \cite{Do} 覆盖了
1970 年代末与 1980 年代初由杨–米尔斯方程引发的研究热潮。阿蒂亚–辛格指标定理
是这一故事的一部分，但我在本文中只限于叙述阿蒂亚 1980 年代后期关于反常以及
关于拓扑场论公理化的工作。

我们从 §\ref{subsec:8.1} 开始，论证反常是量子理论之射影本性（而非线性本性）
的一种表现。反常与指标理论之间的几何与拓扑联系，始于阿蒂亚–辛格的一篇合作
论文，我们在 §\ref{subsec:8.3} 中概述。阿蒂亚关于拓扑场论的诸公理、以及它们
与代数拓扑中配边理论的关系，是 §\ref{subsec:8.2} 的主题。我们在
§\ref{subsec:8.4} 中作结，提及关于反常的现代观点，它把上述两项阿蒂亚的工作
结合起来。对旋量场的特殊化还会引入拓扑与几何指标定理，它们都被综合在一个
一般的反常公式之中，该公式援引了我们讨论过的许多阿蒂亚论文。

   \subsection{量子力学中的射影性与对称性}\label{subsec:8.1}

人们常说，反常是经典对称性在相应量子系统中失效。更精确地说，量子对称性是
射影的，而反常就是线性化的阻碍。量子理论内在地是射影的，而反常——
在比“与对称性相联系”这一意义更广的意义上——编码了量子系统的射影性。

% [Tutor Note: 含糊——作者用三元组 $(\sP,p,H)$ 而不用希尔伯特空间来定义量子系统，正是为了把相位自由度从结构里剔除：$\sH$ 只是坐标选择，$\sP$ 才是物理态空间；\eqref{eq:112} 与 \eqref{eq:113} 两处的讨论都只依赖射线的相对相位。脚注提到的 $C^*$ 代数框架（代数量子力学）被作者以“较简单的语境已经足够”一句带过，实质是本节不需要算子代数的一般理论。]

% [Tutor Note: 含糊——“Define $\sP$ as the limit of $\PH$ over all $(\sH,\theta )\in \cP$”中的 limit 用词不严：$\cP$ 是连通的群胚，且其自同构在 $\PH$（以及 $\End\sH$）上作用平凡，故此图表上的 limit 与 colimit 同构，$\sP$ 按哪一种读都得到同一空间；作者沿用 limit 一词只是行文方便，不必按范畴论的严格意义理解。]

% [Tutor Note: 易错——\eqref{eq:112} 的跃迁概率只依赖两条射线的相对相位：把 $\psi _i$ 换成 $e^{i\alpha _i}\psi _i$ 不改变 $|\langle \psi _1,\psi _2\rangle|^2$，却改变振幅 \eqref{eq:116} 的相位，这正是射影性的数学内容；此外 $p$ 的定义要求 $\psi _i$ 取单位范数，否则须除以范数，归一化因子极易漏掉。]

一个量子力学系统由数据三元组 $(\sP,p,H)$ 确定\footnote{关于量子理论还有用
$C^*$ 代数表述的更一般框架，但就本文的叙述而言，较简单的语境已经足够。}。
空间 $\sP$ 是一个射影空间，即复可分希尔伯特空间 $\sH$ 的射影化 $\PH$，但
$\sH$ 本身并未被指定。定义 $\sP$ 的一种方式是：固定 $\sH_0$，令 $\cP$ 为如下
群胚，其对象是希尔伯特空间与射影线性同构 $\theta \:\PH_0\to \PH$ 构成的偶
$(\sH,\theta )$。态射 $(\sH,\theta )\to (\sH',\theta ')$ 是线性等距
$\sH\to \sH'$，其射影化与 $\theta ,\theta '$ 交换；它们构成酉群
$\TT\subset \CC^{\times }$（单位模标量）上的一个挠子。基点 $\sH_0$
\emph{不}是结构的一部分。定义 $\sP$ 为 $\PH$ 在一切 $(\sH,\theta )\in \cP$ 上
的极限。此外，定义 $*$ 代数 $\AP$（可观测量）为 $\End\sH$ 在
$(\sH,\theta )\in \cP$ 上的极限。（$\End\sH$ 是 $\sH$ 上有界线性算子构成的
代数；为简单起见，本文略去无界算子。）射影空间 $\sP$ 是量子系统的\emph{纯态}
空间。把 $\sP\hookrightarrow \AP$ 嵌入为秩一正交投影。于是混合态是纯态的凸
组合。函数
  \begin{equation}\label{eq:112}
     \begin{aligned} p\:\sP\times \sP&\longrightarrow \;\;[0,1] \\
      L_1,L_2&\longmapsto |\langle \psi _1,\psi _2
     \rangle|^2,\qquad \psi _i\in L_i\end{aligned}  
  \end{equation}
用来计算跃迁概率。为定义 $p$，取 $(\sH,\theta )\in \cP$ 并把 $\sP$ 与 $\PH$
等同；于是 $L_i\subset \sH$ 是一条线，而 $\psi _i\in L_i$ 是单位范数向量。
\emph{哈密顿量} $H$ 是 $\AP$ 中的一个自伴元素，通常假定其谱有下界。

% [Tutor Note: 含糊——“A basic theorem of Wigner asserts”省略的是：$\PH$ 上保持跃迁概率 $p$ 的任一（双射）变换必由 $\sH$ 的酉算子或反酉算子诱导，故 \eqref{eq:113} 中 $q$ 满射、\eqref{eq:114} 的扩张存在。本文参考文献未收 Wigner 原文；标准证明见 Wigner 1931 年专著 \emph{Group Theory}（中译本《群论与量子力学》），或 Weinberg \emph{The Quantum Theory of Fields} 卷一 §2.2。]

% [Tutor Note: 术语——“反常”在此定义为 \eqref{eq:114} 这个 $\zt$ 分次的群扩张，而不只是一个中心扩张：反酉（时间反演型）对称性的存在使扩张带有 $\zt$ 分次 $\epsilon $。若 $\epsilon $ 非平凡，$G_\sH$ 不再是 $G$ 上的中心扩张，而是“$\zt$ 分次中心扩张”，此时同构类不由 $H^2(G;\TT)$ 单独决定，还需记录 $\epsilon $。]

固定 $(\sH,\theta )\in \cP$。设 $G(\sH)$ 为无穷维李群，其单位连通分支是 $\sH$
的酉自同构群 $U(\sH)$，而非单位连通分支是反酉自同构构成的挠子。维格纳
（Wigner）的一个基本定理断言：李群同态的序列
  \begin{equation}\label{eq:113}
     1\longrightarrow \TT\longrightarrow
     G(\sH)\xrightarrow{\;\;q\;\;}\Aut(\PH,p)\longrightarrow 1 
  \end{equation}
是群扩张：$q$ 是满射。设 $G$ 是 $(\sP,p)$ 的一个对称李群，即一个同态
$G\to \Aut(\sP,p)$。通过拉回，我们得到群扩张
  \begin{equation}\label{eq:114}
     1\longrightarrow \TT\longrightarrow G_{\sH}\longrightarrow G\longrightarrow
     1 
  \end{equation}
以及一个 $\zt$ 分次 $\epsilon \:G\to \zt$。该 $\zt$ 分次与扩张 \eqref{eq:114} 的同构
类都与 $(\sH,\theta )$ 的选取无关。$\zt$ 分次的群扩张 \eqref{eq:114} 可以称为
\emph{反常}；它度量对称性的射影性。

   \begin{remark}[]\label{thm:27}
 \ 

 \begin{enumerate}%[label=\textnormal{(\roman*)}]

% [Tutor Note: 含糊——“the symmetry persists in the limit”的实质是：经典极限下对称群 $G$ 线性地作用，而量子系统中它只射影地作用，$G_\sH$ 就是这一射影作用的线性化；作者没有给出“经典极限存在且对称性保持”的判别法，这在一般情形下是困难的（需半经典分析）。]

 \item 若 $(\sP,p,H)$ 有经典极限，且对称性在该极限下保持，则作用在经典系统
 上的是李群 $G$。从 $G$ 过渡到 $G_\sH$，正是本节第一段所提到的。

% [Tutor Note: 易错——第 2 条里“把 $S$ 积掉”要求先选定 $S$ 上的测度，$L^2$ 截面空间才良定义；测度会改变所得量子系统，但反常（提升的阻碍）与测度无关，它只依赖射影丛结构。另注意“反常是阻碍”与“反常是特征”两种说法的区别：不积分 $S$ 时它恰是族指标定理中的那个类，这正是 §\ref{subsec:8.4} 中“反常本身是场论”的雏形。]

 \item 设给定一族以流形 $S$ 为参数的量子系统，特别地给出射影希尔伯特空间的
 一个纤维丛 $\sP\to S$。在某些物理情形中，人们希望对 $S$“积分”以构造一个新
 的量子力学系统。若把 $\sP\to S$ 写成希尔伯特空间向量丛 $\sE\to S$ 的射影化，
 则积分所得系统的态空间是 $\sE\to S$ 的 $L^2$ 截面空间（假定 $S$ 上有测度）。
 反常是把 $\sP\to S$ 提升为向量丛的阻碍，而这一提升是把 $S$ 积掉的必要的
 第一步。另一方面，若不对 $S$ 积分，则反常不是阻碍，而是一族量子系统的一个
 特征，在许多场合下有用。

% [Tutor Note: 笔误——注 \ref{thm:27}(3) 中原文作 “in quantum field ~(\S\ref{subsec:8.4})”，疑脱 theory 一词，应为 “in quantum field theory”（量子场论）；译文按补足后的意思翻译，原文照录。]

% [Tutor Note: 含糊——“The isomorphism class of the central extension lives in $H^2(G;\TT)$”是中心扩张的分类定理，作者未给证明：$G$ 离散时用抽象群上同调，$G$ 为李群时须用连续（光滑）上同调，否则会出现非连续的“怪”扩张——这正是“The type of cohomology depends on the type of group”一句的含义。作者以“假定 $G$ 有限”回避了这一点；§\ref{subsec:8.4} 把 $G$ 换成配边范畴时，群上同调就被替换为配边谱的上同调。]

 \item 中心扩张 \eqref{eq:114} 的同构类位于上同调群 $H^2(G;\TT)$ 中。（上同调的类型
 取决于群的类型。此处假定 $G$ 是有限群并不损失任何东西。）若把上同调的次数
 降低一次，则 $H^1(G;\TT)$ 是一维酉表示构成的群。因此，\eqref{eq:114} 所度量的射影性
 可由一个从 $G$ 在 $\sH$ 上的线性作用出发的两步手续得到：(1) 把 $\sH$ 换成
 一维向量空间；(2) 用上同调解释 $G$ 的一维表示，并把上同调次数提高 $1$。
 我们将在量子场论中看到类似的手续（§\ref{subsec:8.4}）。

 \end{enumerate}
   \end{remark}

% [Tutor Note: 易错——\eqref{eq:115} 的 $\hbar$ 只出现在演化算子 $\evo{k+1}k$ 中（它们通常含因子 $e^{-iH(t_{k+1}-t_k)/\hbar}$）；概率对 $\hbar$ 的依赖是间接的。别把 $\hbar$ 误写进 \eqref{eq:116} 的振幅相位，也别把 $\hbar$ 与 $2\pi$（例如瞬子相位 $e^{2\pi i\theta}$ 中的 $2\pi$）混为一谈。]

设 $t_0<t_1<\cdots< t_{n+1}$ 为实数，视作仿射时间轴上的点，并设
$A_1,\dots ,A_n\in \AP$。固定初始与末态的纯态 $L_0,L_{n+1}\in \sP$。量子力学
中一个基本的有趣量是概率\footnote{若线性算子
$\evo{n+1}n A_n\cdots\evo21 A_1\evo10$ 在 $L_0$ 上为零，则概率为零。}
  \begin{equation}\label{eq:115}
     p\bigl(L_{n+1}\,,\, \evo{n+1}n A_n\cdots\evo21 A_1\evo10 L_0 \bigr),
  \end{equation}
其中 $\hbar$ 是普朗克常数。物理问题可以用这些概率来表述。固定
$(\sH,\theta )\in \cP$，把 $\sP\approx \PH$ 等同，并选取单位范数向量
$\psi _0\in L_0$ 与 $\psi _{n+1}\in L_{n+1}$。这组数据的\emph{振幅}，或称
\emph{关联函数}，是复数
  \begin{equation}\label{eq:116}
     \bigl\langle \psi _{n+1}\,,\, \evo{n+1}n A_n\cdots\evo21 A_1\evo10 \psi _0
     \bigr\rangle_{\sH} .
  \end{equation}
更内蕴地说，数据 $L_0,L_{n+1},t_0,\dots ,t_{n+1},A_1,\dots ,A_n$ 决定一条
埃尔米特线 $\sL$，而振幅是 $\sL$ 中的一个元素，其范数是概率 \eqref{eq:115}。

% [Tutor Note: 术语——注 \ref{thm:28} 中“诸线 $\sL$ 构成数据参数空间上的线丛”把本节的射影性讨论与 §\ref{subsec:8.3} 接上：在后者中这些线丛具体化为狄拉克算子族的行列式线丛 \eqref{eq:122}。注意线丛作为拓扑对象只记 $c_1$，作为带度量与联络的线丛才记录反常的几何信息（[F3]、[BiF1]）。]

   \begin{remark}[]\label{thm:28}
  $(\sP,p,H)$ 的振幅落在一条没有区分基元的线 $\sL$ 中，这是量子力学射影性的
 又一侧面。诸线 $\sL$ 构成数据参数空间上的线丛。这些线丛是量子系统反常的
 一部分，是注 \ref{thm:27}(2) 中射影丛在关联函数方面的对应物。
   \end{remark}

% [Tutor Note: 含糊——“soft”指本节只用到射影空间、群扩张等不依赖具体分析的拓扑/范畴事实，类比于对一般 Fredholm 算子的指标理论（只需算子的拓扑性质即可定义指数）。作者预告“这不是类比”：一旦换成威克转动后的旋量场，反常就由族指标定理精确计算，见 \eqref{eq:118} 与 \eqref{eq:119}。]

本节勾勒的反常框架是软性的，正如对一般 Fredholm 算子的指标理论那样。量子
场论引入威克转动后时空的几何，而在这一语境中反常也展现出更多的几何，正如
对狄拉克算子的指标理论那样。对费米子场而言，这不止是类比；下面我们就转向
费米子场。

   \subsection{量子场论中的旋量场与反常}\label{subsec:8.3}

% [Tutor Note: 含糊——[AS7] 是《美国科学院院刊》上的四页短文，只陈述结果；“we first summarize their work”所省略的细节（族指标、超渡、行列式线丛）散见阿蒂亚–辛格的系列长文与 [S3]、[F3]；本节的 \eqref{eq:117}–\eqref{eq:119} 就是这些细节的提要。另注意 [Lt2] 是 Lott 的“轴反常的一个证明”，不要与 [Lt1] 的指标 gerbe 混淆。]

旋量场反常与指标理论之间的几何联系，是在阿蒂亚–辛格 1984 年的一篇论文
\cite{AS7} 中锻造出来的。（同时期的其他论文，如 \cite{AgW,AgG,Lt2}，也把
指标多项式引入反常理论。）其背景是物理学中的规范场论。（到那时，阿蒂亚已在
数学规范场论中取得许多结果；见本卷中唐纳森的论文 \cite{Do}。）我们先概述他们
的工作，再把它与反常的几何图像联系起来。

% [Tutor Note: 易错——为何要除以 $\sG_0$ 而非整个 $\sG$：若 $g\in\sG$ 固定联络 $A$ 且 $g(p_0)=1$，则由 $A$ 的平行性 $g$ 必恒等于 $1$，故 $\sG_0$ 在 $\sA$ 上自由作用；对整个 $\sG$，使 $A$ 不动的常值（中心）变换使作用不自由，$\sA/\sG$ 也不再是好的丛。由 $g\mapsto g(p_0)$ 得 $\sG/\sG_0\cong G$，于是 \eqref{eq:117} 中 $\pi $ 的纤维恰是 $G$。$\sG_0$ 依赖 $p_0$ 的选取，但 $\pi $ 的同构类与 $\ind$ 不依赖。]

% [Tutor Note: 术语——\eqref{eq:118} 中 $K^{-2n}\cong K^0$ 用的是复 $K$ 理论的 Bott 周期性（周期 $2$）：$-2n$ 为偶数故同构于 $K^0$。注意指数里出现 $2n$，与本节开头所设 $\dim X=n$ 不一致，详见下一条批注。]

设 $n$ 为偶正整数，$X$ 为 $n$ 维闭黎曼 $\Spin$ 流形。设 $P\to X$ 是以紧李群
$G$ 为结构群的主丛。设 $\sA$ 是 $P\to X$ 上联络构成的无穷维仿射空间，$\sG$
是规范变换群，即作为 $P\to X$ 的自同构而在 $X$ 上恒等作用者。则 $\sG$ 作用在
$\sA$ 上。假定 $X$ 连通，固定基点 $p_0\in P$，令 $\sG_0\subset \sG$ 为固定
$p_0$ 的规范变换构成的子群。则 $\sG_0$ 自由地作用在 $\sA$ 上，而在序列
  \begin{equation}\label{eq:117}
     \frac{\sA\times P}{\sG_0}\xrightarrow{\;\;\pi
     \;\;}\frac{\sA}{\sG_0}\times X
     \xrightarrow{\;\;\pr_1\;\;}\frac{\sA}{\sG_0} ,
  \end{equation}
中，$\pi $ 是主 $G$-丛，$\pr_1$ 是纤维为 $X$ 的乘积纤维丛。对 $G$ 的一个酉
表示，我们配以 $\pi $ 的一个向量丛。这些数据决定一个拓扑指数
  \begin{equation}\label{eq:118}
     \ind\in K^{-2n}(\sA/\sG_0)\cong K^0(\sA/\sG_0). 
  \end{equation}

% [Tutor Note: 易错——归一化：本文取陈特征 $\ch_1=\frac{i}{2\pi}\Tr F$，故第一陈类的实表示由 $\frac{i}{2\pi}\Tr F$ 给出、周期为整数；若改用物理学家的实值场强 $F$ 书写，则对应 $\frac{1}{2\pi}\int F\in \ZZ$。$2\pi$ 因子写错时反常只差整数倍，却会翻转“整体反常是否为零”的判定；此处的 $2\pi$ 与 $\hbar$ 无关，切勿混同。]

% [Tutor Note: 笔误——本段三处 $2n$ 与开头“$\dim X=n$”不相容：\eqref{eq:118} 的指数是 $K^{-2n}$，\eqref{eq:119} 的 $\omega $ 是 $(2n+2)$ 形式、超渡后是 $(2n+1)$ 形式。若改按 $\dim X=2n$ 读则全部自洽：$(2n+2)$ 形式沿 $2n$ 维纤维积分恰得基上的 $2$ 形式 $=\ch_1$，族指数落在 $K^{-2n}\cong K^0$。疑原文的 $n$ 与 $2n$ 混用（请核对 [AS7] 中流形维数所采用的记号），照录原文。]

% [Tutor Note: 含糊——“They also relate their result to determinants”一句跳过的实质是：$c_1$ 等于行列式线丛的同构类（拓扑层面），而要在几何层面得到该线丛的度量与联络，须用 [BiF1] 的 Bismut–Freed 联络与 [F3] 的分析；本文 §\ref{subsec:7.5} 已讨论这一构造。仅凭 \eqref{eq:119} 只能得到 $c_1$，得不到“平坦截面”的阻碍（见下段）。]

阿蒂亚–辛格如下构造出代表 $\ch(\ind)\in H^{\bullet }(\sA/\sG_0;\RR)$ 的闭
微分形式。固定 $G$ 上一个双不变黎曼度量。于是对每个联络 $A\in \sA$，流形 $P$
有一个黎曼度量，使得投影 $P\to X$ 成为黎曼淹没，且该联络的水平子空间与纤维
正交。用 $\sA$ 上的 $L^2$ 度量，在 $\sA\times P$ 上构造翘曲乘积度量。于是
$\sG_0\times G$ 以等距方式作用；在模 $\sG_0$ 的商上取 $G$ 轨道的正交补，便
得到 $\pi $ 上的一个联络。相伴向量丛继承一个协变导数，因此有以 $\sA/\sG_0$
为参数、作用在 $X$ 上的狄拉克算子族；族指标定理（§\ref{subsec:5.4}）蕴含
\eqref{eq:118} 中的 ind 等于其解析指数。Chern–Weil 手续从这些微分几何数据
产生所需的微分形式。特别地，$\ch_1(\ind)\in H^2(\sA/\sG_0;\RR)$ 由
  \begin{equation}\label{eq:119}
     \int_{X}\omega 
  \end{equation}
代表，其中 $\omega $ 是 $(\sA/\sG_0)\times X$ 上的 $(2n+2)$ 形式。阿蒂亚–
辛格把 $\omega $ 超渡为 $\sG_0\times X$ 上的 $(2n+1)$ 形式，而正是这个微分
形式在同时期的物理学文献中作为量子场论中旋量场的反常出现。他们还把结果与
行列式联系起来。特别地，他们把 \eqref{eq:119} 的上同调类
$c_1(\ind)\in H^2(\sA/\sG_0;\ZZ)$ 解释为以 $\sA/\sG_0$ 为参数的 $X$ 上狄拉克
算子族的行列式线丛的同构类。

% [Tutor Note: 易错——\eqref{eq:121} 的指数是 $-(\psi ,D_s\psi )$，既无 $i$ 也无 $\hbar$：这是已威克转动到黎曼号差的欧几里得作用量；回到洛伦兹号差时指数为 $e^{iS/\hbar}$，而“相位取 $2\pi i$ 的整数倍”这类量子化条件在欧几里得表述中表现为 $\det D$ 的相位/符号。不要把 $e^{iS/\hbar}$ 的 $i$ 与 $2\pi i$ 中的 $2\pi$ 的归一化混在一起，否则行列式的相位会差一个与 $\hbar$ 无关的整数倍。]

% [Tutor Note: 术语——\eqref{eq:121} 中的 $(\psi ,D_s\psi )$ 不是希尔伯特空间上的埃尔米特内积，而是旋量丛的实（Majorana 型、电荷共轭型）双线性型 $\bar\psi D_s\psi $；其反对称性等价于 $D_s$ 关于相应埃尔米特内积自伴、而该双线性型由电荷共轭 $C$ 给出（于是 $CD_s$ 反对称）。换用不同的实结构（或改变旋量丛定向）会使该型乘上符号，从而 $\pfaff$ 差一个符号——这是行列式线丛相位模糊的根源。]

狄拉克算子的行列式（更一般地，普法夫式）直接出现在含旋量场的威克转动量子
场论中；下面我们在一般框架中略述之。设 $n$ 为正整数，$G$ 为紧李群，$\rho $
为 $G$ 的一个表示。设
  \begin{equation}\label{eq:120}
     P\xrightarrow{\;\;\pi \;\;}X\xrightarrow{\;\;f\;\;}S 
  \end{equation}
是一个带联络的主 $G$-丛 $\pi $ 与一个相对维数为 $n$ 的逆紧纤维丛 $f$，后者
配有相对的黎曼 $\Spin$ 结构。如 §\ref{subsec:5.4} 那样，这些数据产生以 $S$
为参数的狄拉克算子族 $D_{X/S}$。对旋量场 $\psi $ 的费曼积分是标准高斯积分的
无穷维变形：
  \begin{equation}\label{eq:121}
     \int_{\sS_s} e^{-(\psi ,D_s\psi )}. 
  \end{equation}
对每个 $s\in S$，我们应当在旋量场构成的无穷维向量空间 $\sS_s$ 上积分。狄拉克
算子 $D_s$ 出现在指数中，该指数是 $\sS_s$ 上的反对称双线性型。\footnote{物理
学家从洛伦兹号差下闵可夫斯基时空上的相对论性理论出发，然后威克转动到黎曼流形。
自然出现的是指数中的双线性型。}

% [Tutor Note: 含糊——“the pfaffian reduces to a determinant”省略的是：当旋量丛按手征性分裂为 $\sS=\sS^+\oplus \sS^-$、狄拉克算子呈非对角形式 $D=\begin{pmatrix}0&D^-\\D^+&0\end{pmatrix}$ 时，$\pfaff D=\det D^+$（至多相差一个与实结构、定向有关的符号）；“阿蒂亚–辛格框架”之所以能作此简化，是因为那里的算子正是这种手征形式。若不指定实结构，$\pfaff$ 本身只在相差一个相位的意义下有定义。]

% [Tutor Note: 术语——行列式线丛的纤维是 $\det\Ker D_s\otimes(\det\Coker D_s)^{-1}$；本文的 \eqref{eq:121} 直接给出正规化行列式作为截面，而在一般的族上定义 $\det D_{X/S}$ 须先作 $\zeta$ 正规化（[F3]）并赋予 Quillen 度量与 [BiF1] 的联络。作者把 \eqref{eq:121} 当作定义，是为了避开这些分析细节。]

与有限维积分类比，\eqref{eq:121} 被定义为 $\pfaff D_s$，即狄拉克算子的普法夫式。在
特殊情形下，例如上述阿蒂亚–辛格框架，普法夫式化为行列式；我们在此作这一
简化。于是积分 \eqref{eq:121} 的结果是行列式线丛
  \begin{equation}\label{eq:122}
     \Det D_{X/S}\longrightarrow S, 
  \end{equation}
的截面 $\det D_{X/S}$，如 §\ref{subsec:7.5} 所述。

   \begin{remark}[]\label{thm:29}
 \ 

 \begin{enumerate}%[label=\textnormal{(\roman*)}]

% [Tutor Note: 易错——第 1 条中“就是反常”应理解为：反常的完整数据是线丛 + 度量 + 联络，而不只是它的拓扑类 $c_1$。只记 $c_1$ 只能看到“能否取非零截面”，看不到“能否取平坦截面”；后者才是（下一段 [F3] 意义下的）反常。]

 \item 积分 \eqref{eq:121} 是费曼（Feynman）进路的量子场论中关联函数的一个例子。它是
 复线中的一个元素而非一个复数，这一事实正是 §\ref{subsec:8.1} 中所讨论的反常
 振幅的情形。这里，行列式线丛 \eqref{eq:122}（连同其度量与协变导数）就是反常。

% [Tutor Note: 含糊——“Such integrals are problematic analytically”在注 \ref{thm:27}(2) 处已被点过一次：对度量/联络的积分没有勒贝格测度，路径积分在分析上未定义；但“反常”与“平凡化”可以严格定义为线丛及其截面，故作者能在数学上谈论它。这正是 §\ref{subsec:8.4} 把反常提升为场论的出发点。]

 \item 对固定的流形 $X$，万有的参数空间 $S$ 是 $X$ 上黎曼度量与 $G$-联络
 构成的空间。在量子场论中，人们常常要对度量或联络（或二者）积分。反常就是
 这样做的一个阻碍。

 \end{enumerate}
   \end{remark}

% [Tutor Note: 易错——局部与整体的区分在此是：局部反常用 \eqref{eq:119} 的 $2$ 形式（自然联络的曲率，闭而不精确）度量，整体反常用其和乐（沿参数空间中非平凡闭路的相位）度量。若讨论反常流入（anomaly inflow），方向是：$(n+1)$ 维体上的陈–西蒙斯型作用量的变分精确抵消 $n$ 维边界上的规范变分；把体与边的定向弄反，会使抵消变成加倍。本文不展开流入机制，可参 [Wi3]、[AgW]。]

反常是把 \eqref{eq:122} 平凡化的阻碍。阿蒂亚–辛格所讨论的指数的拓扑第一陈类（参看
\cite{S3}）阻碍非零截面的存在，但人们需要更精细的东西。在 \cite{F3} 中，
反常被识别为（相对于自然联络 \cite{BiF1} 的）\emph{平坦}截面的阻碍。这导致
把阿蒂亚–辛格算出的 2-形式 \eqref{eq:119} 解释为该自然联络的曲率。物理学家把这一对
平坦截面的局部阻碍称为\emph{局部反常}；其和乐就是威滕的\emph{整体反常}
\cite{Wi3}。事实上，§\ref{subsec:7.5} 中所述的几何指标理论的发展，正是直接
受这一反常几何的启发。

% [Tutor Note: 术语——gerbe（中文或作“格贝”，常保留原文）是线丛的下一级对象，可理解为“局部由线丛粘合而成的层”，其“曲率”是 $3$ 形式；把射影丛提升为向量丛的阻碍比第一陈类“高一级”，故由狄拉克格贝而非线丛记录。[Lt1]、[Bu] 讨论指标 gerbe 及其超渡；[FS] 是 Faddeev–Shatashvili 的论文（切勿误读为本文作者 Freed 与他人合著），[Seg2]、[NAg] 讨论其在物理中的实现。]

   \begin{remark}[]\label{thm:33}
  反常作为把射影空间丛提升为向量丛的阻碍（注 \ref{thm:27}(2)），由几何指标
 理论中行列式线丛之后的下一个不变量——狄拉克格贝（Dirac gerbe）
  \cite{Lt1,Bu}——来度量。我们请读者参看 \cite{FS,Seg2,NAg}。
   \end{remark}

   \subsection{拓扑场论与配边}\label{subsec:8.2}

阿蒂亚对量子场论的介入远不止于指标定理。一篇特别有影响的论文 \cite{A11}
给出了\emph{拓扑}量子场论的公理，与西格尔（Segal）先前为二维共形场论引入的
公理 \cite{Seg3} 平行。一个关键的推动是威滕的量子陈–西蒙斯理论 \cite{Wi4}，
它把纽结的琼斯（Jones）不变量置于一个明显三维的框架之中。但被阿蒂亚抽象进
公理的还有许多别的例子。他的论文 \cite{A11} 献给托姆（Thom），而配边理论
确实正处在最前沿。阿蒂亚 \cite{A12} 关于这篇论文写下如下的话：

% [Tutor Note: 含糊——引文出自 [A12]（阿蒂亚自述）。“the combinatorial approach is so far the only one that has been made to work (for the Jones polynomials)”指 Reshetikhin–Turaev/量子群给出的组合构造：在阿蒂亚写这段话的年代，陈–西蒙斯路径积分尚无严格定义，而组合构造可以严格化。对照 [Wi4] 中威滕用路径积分给出不变量，二者的关系即“公理 vs 构造”。]

 \begin{quote} 
 数学家对费曼积分心怀畏惧，又不熟悉物理学家的全套行话，因此在我看来
 有必要以对用户友好的方式向数学家解释拓扑量子场论究竟是什么。我给出了一个
 简单的公理化处理（这是数学家所钟爱的），并列举了由物理学产生的例子。数学家
 的任务于是成为：用任何可能的方法去构造一个符合公理的理论。我愿意把这比作
 上同调的艾伦伯格–斯廷罗德（Eilenberg–Steenrod）公理：人们可以用单纯方法、
 切赫（Čech）方法或德·拉姆方法来构造该理论。最后一种最接近物理学，但另外
 几种也有一些优点。在量子场论的语境中，事情要困难得多，组合进路迄今是唯一
 被做成功的进路（就琼斯多项式而言）。
 \end{quote}

\noindent
下面我们给出这些公理的一个简洁版本。

% [Tutor Note: 易错——号差为何是配边不变量：$\Sign(X)=\langle L(X),[X]\rangle$，$L$ 亏格由 Pontryagin 类给出，配边关系下该配对不变；更本质地，\eqref{eq:124} 是整配边群 $\Omega _{4k}(\SO)$ 上的 $\ZZ$ 值同态，而任何到 $\ZZ$ 的同态在挠类上自动为零，故号差只看到自由部分（$\Omega _{4k}(\SO)\otimes\QQ$ 由复射影空间的乘积生成）。定理 \ref{thm:3} 的证明正依赖于此。]

% [Tutor Note: 易错——这里的 $d$ 是空间维数，TQFT 的时空维数是 $d+1$；§\ref{subsec:8.3} 的 $n$ 是时空维数，故 $n=d+1$。后文 §\ref{subsec:8.4} 说反常是 $(n+1)$ 维场论，因为它比被反常化的 $n$ 维理论高一维。这三处维数极易混，读时务必对表：$d$（对象维数）、$d+1$（态射/时空维数）、$n=d+1$、$n+1$（反常理论维数）。]

作为动机，回忆维数为 $4k$（$k\in \ZZ^{\ge0}$）的闭定向流形的号差。正如在
定理 \ref{thm:3} 的证明中所关键用到的，号差是一个\emph{配边不变量}，即一个
阿贝尔群同态
  \begin{equation}\label{eq:124}
     \Sign\:\Omega _{4k}(\SO)\longrightarrow \ZZ, 
  \end{equation}
其中 $\Omega _{4k}(\SO)$ 是托姆的闭定向 $4k$ 流形配边群。拓扑场论是一个
“范畴化的配边不变量”。固定一个非负整数 $d$。（它与 §\ref{subsec:8.3} 中
$n$ 的关系是 $n=d+1$。）设 $\Bd$ 为下述范畴，它由米尔诺（Milnor）\cite{Mi2}
首先引入。其对象是闭 $d$ 流形。若 $Y_0,Y_1$ 是两个这样的流形，则一个态射
$Y_0\to Y_1$ 由紧 $(d+1)$ 流形 $X$ 代表，其边界分划为
$\partial X=Y_0\amalg Y_1$。换言之，$X$ 是从 $Y_0$ 到 $Y_1$ 的一个配边。在
边界固定意义下微分同胚的配边代表同一个态射。复合通过粘合配边实现，而流形的
不交并给出对称幺半结构。

   \begin{remark}[]\label{thm:30}
 \ 

 \begin{enumerate}%[label=\textnormal{(\roman*)}]

% [Tutor Note: 含糊——“categorifies”的确切含义是：$\Bd$ 的 $\pi_0$（对象模“存在态射”）给出 $\Omega _d$，而 $\Bd$ 还保留了对象之间的态射与复合，即保留了配边的全部信息。严格说 $\Omega _d$ 视切结构而定，见第 2 条。]

 \item 若我们宣布 $\Bd$ 的对象 $Y_0,Y_1$ 等价，当且仅当存在态射
 $Y_0\to Y_1$，则等价类的集合就是配边群 $\Omega _d$。在这个意义上，$\Bd$
 “范畴化”了 $\Omega _d$。

% [Tutor Note: 术语——tangential structure（切结构）包括定向、$\Spin$ 结构、标架等，它决定配边群的类型；正文未指定切结构处写的 $\Omega _d$ 严格应为非定向配边群 $\Omega _d(\O)$，而 \eqref{eq:124} 中有定向，写作 $\Omega _{4k}(\SO)$。后文 $M\!\Spin$ 对应的正是 $\Spin$ 切结构下的配边谱。]

 \item 稍作变化便得到带切结构的配边范畴，例如带定向者。

 \end{enumerate}
   \end{remark}

现在我们可以陈述公理了。设 $k$ 是域，$\Vk$ 是 $k$-向量空间与线性映射构成的
范畴。（在量子理论中 $k=\CC$。）张量积运算在 $\Vk$ 上定义对称幺半结构。

% [Tutor Note: 含糊——定义 \ref{thm:31} 只写了函子性，未写出全部相容性：还须要求 $F(\emptyset)=\mathbf{1}$（空流形对应基域）、$F(Y_0\amalg Y_1)\cong F(Y_0)\otimes F(Y_1)$（不交并对张量积）、以及粘合配边对应态射的复合；这些正是“symmetric monoidal functor”一词的内容。现代严格化版本要求 $\Bd$ 为 $(\infty,1)$-范畴、$F$ 为幺半 $(\infty,1)$-函子（配边假设），作者以“with suitable modifications and extensions it is believed to apply widely”一句带过。]

   \begin{definition}[Atiyah~\cite{A11}]\label{thm:31}
  \emph{拓扑场论}是一个对称幺半函子
  \begin{equation}\label{eq:125}
     F\:\Bd\longrightarrow \Vk. 
  \end{equation}
   \end{definition}

% [Tutor Note: 术语——Atiyah–Segal 公理系统中，Segal 的原始公理针对二维共形场论（[Seg3]，保留共形/度量结构），Atiyah 把它搬到拓扑情形（[A11]，只需配边结构，并要求 $F(Y)$ 有限维）。二者的差别在于是否保留度量/共形数据；拓扑情形相当于共形权为零的极端，故配分函数只剩拓扑信息。]

\noindent
这个定义有时被称为\emph{阿蒂亚–西格尔公理系统}；经过适当的修改与扩充，人们
相信它广泛适用于数学与物理中的场论。这些公理所采取的视角与物理学教科书中所
见的大不相同。阿蒂亚–西格尔公理系统为整整一代数学家提供了进入这一物理学的
入口，它们是许多数学发展所依托的结构，同时也阐明了物理学中量子场论的几何
侧面。

% [Tutor Note: 易错——法标架的选取会改变同痕类 $\partial X\approx\cir\times \cir$，从而改变 $F(X)$ 的取值：这正是“标架纽结不变量”与“无标架纽结不变量”的区别。琼斯多项式要成为无标架不变量，须再乘/除以一个依赖标架的因子（陈–西蒙斯理论中表现为一个相位）；作者说“of this type”是就标架依赖的意义而言。]

为说明定义 \ref{thm:31}，我们展示如何从 $d=2$ 的场论 $F$ 提取闭 3-流形 $M$
中带法标架的纽结 $K\subset M$ 的数值不变量。设 $X$ 是从 $M$ 中去掉 $K$ 的
一个开管状邻域所得的 3-流形。结果是配边
$X\:\partial X\to \emptyset ^2$：它从 $\partial X$ 走向空 2-流形。$K$ 的一个法
标架给出微分同胚 $\partial X\approx\cir\times \cir$ 的一个同痕类。因此场论
$F$ 在 $X$ 上的取值是一个线性映射
  \begin{equation}\label{eq:126}
     F(X)\:V\longrightarrow k, 
  \end{equation}
其中 $V=F(\cir\times \cir)$ 是附着于标准 2-环面的向量空间。对每个向量
$\xi \in V$——它可以看作附着于 $K$ 的一个“标签”——我们得到数值不变量
$F(X)(\xi )$。纽结的琼斯不变量就属于这一类。这是 \cite{Wi4} 中的关键观察
之一。

   \subsection{综合}\label{subsec:8.4}

最后，我们把阿蒂亚–辛格关于反常的工作（§\ref{subsec:8.3}）、量子场论的
阿蒂亚–西格尔公理系统（§\ref{subsec:8.2}）以及 §\S\ref{sec:4},
\ref{sec:5}, \ref{sec:7} 中的各项指标定理汇总起来作结。

% [Tutor Note: 含糊——“the trivializations of the anomaly must be coherent in the background data”是本节最关键也最含糊的一句：局部性要求反常的平凡化在配边的粘合下相容（即平凡化构成背景数据范畴上的一个函子性对象），这正是阿蒂亚–西格尔公理的内容；作者未证明“局部性 $\Leftrightarrow$ 函子性”，该等价是现代公理化场论（[FHT]、[FH]）的工作。]

% [Tutor Note: 术语——“可逆场论”的确切含义是：$F$ 经由线与可逆线性映射构成的 Picard 群胚 $\Line_k\subset \Vk$ 分解，于是所有态空间都是一维、所有态射都可逆，$F$ 只记录相位而无动力学自由度；物理上这对应可逆相/SPT 相。反常正是这种场论，但一般不是拓扑的：它依赖背景度量与联络（见 \eqref{eq:120} 的数据），故只能定义在带相应结构的流形上——这与 §\ref{subsec:8.1} 中“振幅落在未被平凡化的线里”是同一件事的两种说法。]

% [Tutor Note: 易错——脚注强调：$n$ 维理论的反常只定义在维数 $\le n$ 的流形上，故它不是完整的 $(n+1)$ 维理论。把它延拓为完整的 $(n+1)$ 维可逆理论需要额外数据（例如把 $n$ 维理论视为 $\sB$ 的表示）。作者在下一条注中称之为“once-categorified invertible $n$-dimensional field theory”，切勿与“$(n+1)$ 维理论”混为一谈。]

一个出发点是注 \ref{thm:29}(2)：它告诉我们，要通过对某些场（例如度量或联络）
积分来构造量子场论，必须先把反常平凡化。（这类积分在分析上是有问题的，但
反常与平凡化在数学上是良定义的。）现在，若要所得的量子场论是\emph{局部的}
——而局部性是量子场论的一个标志性特征——那么反常的平凡化必须在背景数据
\eqref{eq:120} 中相容。这种相容性正是阿蒂亚–西格尔公理系统及其各种扩充所
表达的内容。这是导致如下认识的一条推理路线：反常本身就是一个量子场论，尽管
是非常特殊的一类。\footnote{这并非普遍成立：$n$ 维场论的反常未必是一个完整的
$(n+1)$ 维理论，而可能只定义在维数 $\le n$ 的流形上。}对
§\ref{subsec:8.3} 中的旋量场，\eqref{eq:122} 的纤维中的诸行列式线是某个 $(n+1)$ 维
场论中的一维态空间。这个\emph{反常理论}是\emph{可逆的}，但不一定是拓扑的。
（可逆场论 \eqref{eq:125} 经由线与可逆线性映射构成的子群胚
$\Line_k\subset \Vk$ 分解。）

% [Tutor Note: 含糊——注 \ref{thm:34} 的“cohomological interpretation”指 [FHT] 的如下对应：可逆场论 = 谱的映射（或稳定同伦意义下的一个类），于是把“一维表示”提升为上同调类、把次数 $+1$，正好对应“取 $(n+1)$ 维配边谱的一个类”。作者说这是“roadmap”而非证明，实质是 [FHT] 的定理；这正是下一段用稳定同伦论重述的原因。]

   \begin{remark}[]\label{thm:34}
  沿着 §\ref{subsec:8.1} 的想法，我们会得到同样的图像。即，形如
 \eqref{eq:125} 的场论是配边的\emph{线性}表示，但量子理论是\emph{射影的}，
 而反常度量这种射影性。此外，注 \ref{thm:27}(3) 给出了定位这一度量的路线图。
 这里配边范畴扮演群 $G$ 的角色。第一，把 $\Vk$ 换成一维向量空间与可逆线性
 映射构成的范畴。所得的是一个\emph{可逆}场论。第二，我们用上同调解释可逆
 场论，并把上同调次数提高 $1$。上同调解释由 \cite{FHT} 引入；参看
 \cite[\S5]{FH}。我们得到同样的结论：一个场论的反常（即其射影性的度量）是
 一个可逆的 $(n+1)$ 维场论。\footnote{更精确地说，$n$ 维场论是配边范畴
 $\sB$ 的表示。按本段的逻辑，反常定义在\emph{同一}配边范畴 $\sB$ 上，所以
 它是上一脚注中所说的“一次范畴化的可逆 $n$ 维场论”。通常，反常会延拓为
 完整的可逆 $(n+1)$ 维场论，但这并非必需。}

   \end{remark}

% [Tutor Note: 易错——反常相位常写成 $e^{2\pi i\eta(D)}$，准确说是 $e^{2\pi i\xi}$，其中 $\xi(Y)=\frac{1}{2}\bigl(\eta(D)+\dim\ker D\bigr)$：对闭 $\Spin$ 流形 $Y$，$\xi(Y)$ 只在模 $1$ 意义下良定义（换度量或算子会差整数），故指数化后相位有意义；仅当 $Y$ 是某个 $\Spin$ 流形 $W$ 的边界、且维数使 $\int_W\Ahat\in\ZZ$ 时，APS 指标定理才给出 $\xi(Y)\equiv\int_W\Ahat\pmod 1$ 为整数（相位为 $1$）。若忘了 $\frac{1}{2}\dim\ker$ 或谱流修正而直接用 $e^{2\pi i\eta}$，相位不闭环——这正是整体反常的和乐。]

% [Tutor Note: 含糊——“conjectural as a mathematical assertion until more foundations are laid”指 [FH, Conjecture 9.70] 依赖于配边谱与路径积分的对应尚未完全严格化（需把 $\eta$ 不变量、指标密度、$M\!\Spin\to KO$ 的 Atiyah–Bott–Shapiro 映射（[ABS]）组装成一个可逆场论）。对旋量场，$M\!\Spin\to KO$ 把 $n$ 维 $\Spin$ 配边类送到 $KO^{-n}$ 中的类，其与 $\eta$ 的关系是 APS 定理的“谱侧”表述；作者说“Implicit in it are ... index theorems”，正是本节的落脚点。]

这已经把前面提到的阿蒂亚 1980 年代的工作综合起来了，但我们还可以走得更远。
一般场论是对称幺半范畴之间的函子，而可逆场论可以在稳定同伦论中表述为谱的
映射。其定义域不再是配边范畴，而是托姆引入的那一类配边谱
\footnote{此陈述是针对幺正理论的。}。（我们指出，阿蒂亚 \cite{A10} 把配边与
上配边置于广义同调理论的框架之中。）对旋量场，定义域谱是 $M\!\Spin$ 或与之
相近的变体。旋量场反常的一般公式 \cite[Conjecture~9.70]{FH}——在更多基础被
奠定之前，它作为数学断言仍是猜想性的——引入了阿蒂亚–博特–夏皮罗映射
$M\!\Spin\to KO$，以及上述一切要素。其中隐含了本文所综述的各种拓扑与几何
指标不变量与指标定理。

\clearpage
\renewcommand{\refname}{参考文献}
 \bigskip\bigskip
 
\providecommand{\bysame}{\leavevmode\hbox to3em{\hrulefill}\thinspace}
\providecommand{\MR}{\relax\ifhmode\unskip\space\fi MR }
% \MRhref is called by the amsart/book/proc definition of \MR.
\providecommand{\MRhref}[2]{%
  \href{http://www.ams.org/mathscinet-getitem?mr=#1}{#2}
}
\providecommand{\href}[2]{#2}
\begin{thebibliography}{AGW83}

\bibitem[A1]{A1}
M.~F. Atiyah, \emph{Commentary}, Michael Atiyah Collected Works, Index
  Theory:1, vol.~3, Oxford University Press, Walton Street, Oxford, 1987,
  pp.~1--9.

\bibitem[A2]{A2}
\bysame, \emph{The index theorem for manifolds with boundary}, Seminar on the
  {A}tiyah-{S}inger index theorem, With contributions by M. F. Atiyah, A.
  Borel, E. E. Floyd, R. T. Seeley, W. Shih and R. Solovay. Annals of
  Mathematics Studies, No. 57, Princeton University Press, Princeton, N.J.,
  1965, pp.~x+366. Appendix~I.

\bibitem[A3]{A3}
\bysame, \emph{{$K$}-theory and reality}, Quart. J. Math. Oxford Ser. (2)
  \textbf{17} (1966), 367--386.

\bibitem[A4]{A4}
\bysame, \emph{{$K$}-Theory}, second ed., Advanced Book Classics,
  Addison-Wesley, Redwood City, CA, 1989.

\bibitem[A5]{A5}
\bysame, \emph{Elliptic operators, discrete groups and von {N}eumann algebras},
  Colloque ``{A}nalyse et {T}opologie'' en l'{H}onneur de {H}enri {C}artan
  ({O}rsay, 1974), 1976, pp.~43--72. Ast\'{e}risque, No. 32--33.

\bibitem[A6]{A6}
\bysame, \emph{Circular symmetry and stationary-phase approximation},
  Ast{\'e}risque (1985), no.~131, 43--59. Colloquium in honor of Laurent
  Schwartz, Vol. 1 (Palaiseau, 1983).

\bibitem[A7]{A7}
\bysame, \emph{Elliptic operators and compact groups}, Lecture Notes in
  Mathematics, Vol. 401, Springer-Verlag, Berlin-New York, 1974.

\bibitem[A8]{A8}
\bysame, \emph{The logarithm of the {D}edekind {$\eta$}-function},
  \href{http://dx.doi.org/10.1007/BF01458075}{Math. Ann. \textbf{278} (1987)},
  no.~1-4, 335--380.

\bibitem[A9]{A9}
\bysame, \emph{On framings of {$3$}-manifolds}, Topology \textbf{29} (1990),
  no.~1, 1--7.

\bibitem[A10]{A10}
\bysame, \emph{Bordism and cobordism}, Proc. Cambridge Philos. Soc. \textbf{57}
  (1961), 200--208.

\bibitem[A11]{A11}
\bysame, \emph{Topological quantum field theories}, Inst. Hautes \'Etudes Sci.
  Publ. Math. (1988), no.~68, 175--186 (1989).

\bibitem[A12]{A12}
\bysame, \emph{A personal history}, The Founders of Index Theory, Reminiscences
  of Atiyah, Bott, Hirzebruch, and Singer (S.~T. Yau, ed.), International
  Press, 2003, pp.~5--13.

\bibitem[AB1]{AB1}
M.~F. Atiyah and R.~Bott, \emph{A {L}efschetz fixed point formula for elliptic
  complexes. {I}}, \href{http://dx.doi.org/10.2307/1970694}{Ann. of Math. (2)
  \textbf{86} (1967)}, 374--407.

\bibitem[AB2]{AB2}
\bysame, \emph{A {L}efschetz fixed point formula for elliptic complexes. {II}.
  {A}pplications}, \href{http://dx.doi.org/10.2307/1970721}{Ann. of Math. (2)
  \textbf{88} (1968)}, 451--491.

\bibitem[AB3]{AB3}
\bysame, \emph{The index problem for manifolds with boundary}, Differential
  {A}nalysis, {B}ombay {C}olloq., 1964, Oxford Univ. Press, London, 1964,
  pp.~175--186.

\bibitem[AB4]{AB4}
\bysame, \emph{On the periodicity theorem for complex vector bundles},
  \href{http://dx.doi.org/10.1007/BF02391772}{Acta Math. \textbf{112} (1964)},
  229--247.

\bibitem[AB5]{AB5}
\bysame, \emph{The moment map and equivariant cohomology},
  \href{http://dx.doi.org/10.1016/0040-9383(84)90021-1}{Topology \textbf{23}
  (1984)}, no.~1, 1--28.

\bibitem[ABP]{ABP}
M.~Atiyah, R.~Bott, and V.~K. Patodi, \emph{On the heat equation and the index
  theorem}, Invent. Math. \textbf{19} (1973), 279--330.

\bibitem[ABS]{ABS}
M.~F. Atiyah, R.~Bott, and A.~Shapiro, \emph{Clifford modules}, Topology
  \textbf{3} (1964), no.~suppl. 1, 3--38.

\bibitem[ADS]{ADS}
M.~F. Atiyah, H.~Donnelly, and I.~M. Singer, \emph{Eta invariants, signature
  defects of cusps, and values of {$L$}-functions},
  \href{http://dx.doi.org/10.2307/2006957}{Ann. of Math. (2) \textbf{118}
  (1983)}, no.~1, 131--177.

\bibitem[Ag]{Ag}
L.~Alvarez-Gaume, \emph{Supersymmetry and the {A}tiyah-{S}inger index theorem},
  \href{http://dx.doi.org/10.1007/BF01205500}{Commun. Math. Phys. \textbf{90}
  (1983)},
161.
%%CITATION = CMPHA,90,161;%%.

\bibitem[AgG]{AgG}
L.~Alvarez-Gaum{\'e} and P.~Ginsparg, \emph{The structure of gauge and
  gravitational anomalies}, Annals of Physics \textbf{161} (1985), no.~2,
  423--490.

\bibitem[AgW]{AgW}
L.~Alvarez-Gaum\'e and E.~Witten, \emph{Gravitational anomalies}, Nucl. Phys.
  \textbf{B234} (1983), 269.

\bibitem[Agr]{Agr}
M.~S. Agranovi\v{c}, \emph{Elliptic singular integro-differential operators},
  Russian Math. Surveys \textbf{20} (1965), no.~5 (125), 3--120.

\bibitem[AH1]{AH1}
M.~F. Atiyah and F.~Hirzebruch, \emph{Vector bundles and homogeneous spaces},
  Proc. Symp. Pure Math. \textbf{3} (1961), 7--38.

\bibitem[AH2]{AH2}
\bysame, \emph{Riemann-{R}och theorems for differentiable manifolds},
  \href{http://dx.doi.org/10.1090/S0002-9904-1959-10344-X}{Bull. Amer. Math.
  Soc. \textbf{65} (1959)}, 276--281.

\bibitem[AHS]{AHS}
M.~F. Atiyah, N.~J. Hitchin, and I.~M. Singer, \emph{Self-duality in
  four-dimensional {R}iemannian geometry},
  \href{http://dx.doi.org/10.1098/rspa.1978.0143}{Proc. Roy. Soc. London Ser. A
  \textbf{362} (1978)}, no.~1711, 425--461.

\bibitem[AM]{AM}
F.~Arici and B.~Mesland, \emph{Toeplitz extensions in noncommutative topology
  and mathematical physics}, \href{http://arxiv.org/abs/arXiv:1911.05823}{{\tt
  arXiv:1911.05823}}.

\bibitem[APS1]{APS1}
M.~F. Atiyah, V.~K. Patodi, and I.~M. Singer, \emph{Spectral asymmetry and
  {R}iemannian geometry. {I}}, Math. Proc. Cambridge Philos. Soc. \textbf{77}
  (1975), 43--69.

\bibitem[APS2]{APS2}
\bysame, \emph{Spectral asymmetry and {R}iemannian geometry. {II}}, Math. Proc.
  Cambridge Philos. Soc. \textbf{78} (1975), no.~3, 405--432.

\bibitem[APS3]{APS3}
\bysame, \emph{Spectral asymmetry and {R}iemannian geometry. {III}}, Math.
  Proc. Cambridge Philos. Soc. \textbf{79} (1976), no.~1, 71--99.

\bibitem[AS1]{AS1}
M.~F. Atiyah and I.~M. Singer, \emph{The index of elliptic operators on compact
  manifolds}, Bull. Amer. Math. Soc. \textbf{69} (1963), 422--433.

\bibitem[AS2]{AS2}
\bysame, \emph{The index of elliptic operators. {I}}, Ann. of Math. (2)
  \textbf{87} (1968), 484--530.

\bibitem[AS3]{AS3}
\bysame, \emph{The index of elliptic operators. {III}},
  \href{http://dx.doi.org/10.2307/1970717}{Ann. of Math. (2) \textbf{87}
  (1968)}, 546--604.

\bibitem[AS4]{AS4}
\bysame, \emph{The index of elliptic operators. {IV}}, Ann. of Math. (2)
  \textbf{93} (1971), 119--138.

\bibitem[AS5]{AS5}
\bysame, \emph{The index of elliptic operators. {V}}, Ann. of Math. (2)
  \textbf{93} (1971), 139--149.

\bibitem[AS6]{AS6}
\bysame, \emph{Index theory for skew-adjoint {F}redholm operators}, Inst.
  Hautes {\'E}tudes Sci. Publ. Math. (1969), no.~37, 5--26.

\bibitem[AS7]{AS7}
\bysame, \emph{Dirac operators coupled to vector potentials}, Proceedings of
  the National Academy of Sciences \textbf{81} (1984), no.~8, 2597--2600.

\bibitem[ASeg]{ASeg}
M.~F. Atiyah and G.~B. Segal, \emph{The index of elliptic operators. {II}},
  Ann. of Math. (2) \textbf{87} (1968), 531--545.

\bibitem[B1]{B1}
R.~Bott, \emph{The periodicity theorem for the classical groups and some of its
  applications}, \href{http://dx.doi.org/10.1016/0001-8708(70)90030-7}{Advances
  in Math. \textbf{4} (1970)}, 353--411 (1970).

\bibitem[B2]{B2}
\bysame, \emph{The stable homotopy of the classical groups},
  \href{http://dx.doi.org/10.2307/1970106}{Ann. of Math. (2) \textbf{70}
  (1959)}, 313--337.

\bibitem[BdM]{BdM}
L.~Boutet~de Monvel, \emph{On the index of {T}oeplitz operators of several
  complex variables}, \href{http://dx.doi.org/10.1007/BF01410080}{Invent. Math.
  \textbf{50} (1978/79)}, no.~3, 249--272.

\bibitem[BH2]{BH2}
A.~Borel and F.~Hirzebruch, \emph{Characteristic classes and homogeneous
  spaces. {II}}, \href{http://dx.doi.org/10.2307/2372747}{Amer. J. Math.
  \textbf{81} (1959)}, 315--382.

\bibitem[BH3]{BH3}
\bysame, \emph{Characteristic classes and homogeneous spaces. {III}},
  \href{http://dx.doi.org/10.2307/2372969}{Amer. J. Math. \textbf{82} (1960)},
  491--504.

\bibitem[Bi1]{Bi1}
J.-M. Bismut, \emph{The {A}tiyah-{S}inger theorems: a probabilistic approach.
  {I}. {T}he index theorem},
  \href{http://dx.doi.org/10.1016/0022-1236(84)90101-0}{J. Funct. Anal.
  \textbf{57} (1984)}, no.~1, 56--99.

\bibitem[Bi2]{Bi2}
\bysame, \emph{The {A}tiyah-{S}inger index theorem for families of {D}irac
  operators: two heat equation proofs},
  \href{http://dx.doi.org/10.1007/BF01388755}{Invent. Math. \textbf{83}
  (1986)}, no.~1, 91--151.

\bibitem[BiF1]{BiF1}
J.-M. Bismut and D.~S. Freed, \emph{The analysis of elliptic families. {I}.
  {M}etrics and connections on determinant bundles}, Comm. Math. Phys.
  \textbf{106} (1986), no.~1, 159--176.

\bibitem[BiF2]{BiF2}
\bysame, \emph{The analysis of elliptic families. {II}. {D}irac operators, eta
  invariants, and the holonomy theorem}, Comm. Math. Phys. \textbf{107} (1986),
  no.~1, 103--163.

\bibitem[BS]{BS}
A.~Borel and J.-P. Serre, \emph{Le th\'eor\`eme de Riemann-Roch}, Bull. Soc.
  Math. France \textbf{86} (1958), 97--136.

\bibitem[Bu]{Bu}
U.~Bunke, \emph{Transgression of the index gerbe},
  \href{http://dx.doi.org/10.1007/s00229-002-0314-8}{Manuscripta Math.
  \textbf{109} (2002)}, no.~3, 263--287,
  \href{http://arxiv.org/abs/arXiv:math/0109052}{{\tt arXiv:math/0109052}}.

\bibitem[BV]{BV}
N.~Berline and M.~Vergne, \emph{A computation of the equivariant index of the
  {D}irac operator}, Bull. Soc. Math. France \textbf{113} (1985), no.~3,
  305--345.

\bibitem[BW]{BW}
R.~Brauer and H.~Weyl, \emph{Spinors in {$n$} {D}imensions},
  \href{http://dx.doi.org/10.2307/2371218}{Amer. J. Math. \textbf{57} (1935)},
  no.~2, 425--449.

\bibitem[Ca]{Ca}
E.~Cartan, \emph{Les groupes projectifs qui ne laissent invariante aucune
  multiplicit{\'e} plane}, Bulletin de la Soci{\'e}t{\'e} Math{\'e}matique de
  France \textbf{41} (1913), 53--96.

\bibitem[Ch]{Ch}
C.~C. Chevalley, \emph{The algebraic theory of spinors}, Columbia University
  Press, New York, 1954.

\bibitem[Che]{Che}
J.~Cheeger, \emph{{$\eta$}-invariants, the adiabatic approximation and conical
  singularities. {I}. {T}he adiabatic approximation}, J. Differential Geom.
  \textbf{26} (1987), no.~1, 175--221.

\bibitem[ChS]{ChS}
J.~Cheeger and J.~Simons,
  \href{http://dx.doi.org/10.1007/BFb0075216}{\emph{Differential characters and
  geometric invariants}}, Geometry and topology ({C}ollege {P}ark, {M}d.,
  1983/84), Lecture Notes in Math., vol. 1167, Springer, Berlin, 1985,
  pp.~50--80.

\bibitem[Cl]{Cl}
W.~K. Clifford, \emph{On the classification of geometric algebras},
  Mathematical Papers (1882), 397--401.

\bibitem[CS]{CS}
S.~S. Chern and J.~Simons, \emph{Characteristic forms and geometric
  invariants}, Ann. of Math. (2) \textbf{99} (1974), 48--69.

\bibitem[CZ]{CZ}
A.-P. Calder\'{o}n and A.~Zygmund, \emph{Singular integral operators and
  differential equations}, \href{http://dx.doi.org/10.2307/2372441}{Amer. J.
  Math. \textbf{79} (1957)}, 901--921.

\bibitem[D]{D}
P.~A.~M. Dirac, \emph{The quantum theory of the electron}, Proceedings of the
  Royal Society of London. Series A, Containing Papers of a Mathematical and
  Physical Character \textbf{117} (1928), no.~778, 610--624.

\bibitem[deR]{deR}
G.~De~Rham, \emph{Sur l'analysis situs des vari\'{e}t\'{e}s {\`a} {$n$}
  dimensions}, NUMDAM, [place of publication not identified], 1931.

\bibitem[DH]{DH}
J.~J. Duistermaat and G.~J. Heckman, \emph{On the variation in the cohomology
  of the symplectic form of the reduced phase space},
  \href{http://dx.doi.org/10.1007/BF01399506}{Invent. Math. \textbf{69}
  (1982)}, no.~2, 259--268.

\bibitem[Do]{Do}
S.~K. Donaldson, \emph{Atiyah's work on holomorphic vector bundles and gauge
  theories}. this volume.

\bibitem[E]{E}
C.~L. Epstein,
  \href{http://dx.doi.org/10.1142/9789812701398_0002}{\emph{Lectures on indices
  and relative indices on contact and {CR}-manifolds}}, Woods {H}ole
  mathematics, Ser. Knots Everything, vol.~34, World Sci. Publ., Hackensack,
  NJ, 2004, pp.~27--93.

\bibitem[F1]{F1}
Daniel~S. Freed, \emph{The Atiyah-Singer Index Theorem}, April 2021.
  \url{https://www.youtube.com/watch?v=AJHKp9kYm90}. CMSA Math-Science
  Literature Lecture Series.

\bibitem[F2]{F2}
\bysame, \emph{Classical {C}hern-{S}imons theory. {II}}, Houston J. Math.
  \textbf{28} (2002), no.~2, 293--310. Special issue for S. S. Chern.

\bibitem[F3]{F3}
\bysame, \emph{Determinants, torsion, and strings}, Comm. Math. Phys.
  \textbf{107} (1986), no.~3, 483--513.

\bibitem[FH]{FH}
Daniel~S. Freed and Michael~J. Hopkins, \emph{Reflection positivity and
  invertible topological phases}, Geometry \& Topology \textbf{25} (2021),
  no.~3, 1165--1130, \href{http://arxiv.org/abs/arXiv:1604.06527}{{\tt
  arXiv:1604.06527}}.

\bibitem[FHT]{FHT}
Daniel~S. Freed, Michael~J. Hopkins, and Constantin Teleman,
  \href{http://dx.doi.org/10.1093/acprof:oso/9780199534920.003.0019}{\emph{Consistent
  orientation of moduli spaces}}, The many facets of geometry, Oxford Univ.
  Press, Oxford, 2010, pp.~395--419.
  \href{http://arxiv.org/abs/arXiv:0711.1909}{{\tt arXiv:0711.1909}}.

\bibitem[FL]{FL}
Daniel~S. Freed and John Lott, \emph{An index theorem in differential
  {$K$}-theory}, \href{http://dx.doi.org/10.2140/gt.2010.14.903}{Geom. Topol.
  \textbf{14} (2010)}, no.~2, 903--966,
  \href{http://arxiv.org/abs/arXiv:0907.3508}{{\tt arXiv:0907.3508}}.

\bibitem[FS]{FS}
L.D. Faddeev and S.~L. Shatashvili, \emph{{Algebraic and Hamiltonian Methods in
  the Theory of Nonabelian Anomalies}},
  \href{http://dx.doi.org/10.1007/BF01018976}{Theor. Math. Phys. \textbf{60}
  (1985)}, 770--778.

\bibitem[FW]{FW}
D.~Friedan and P.~Windey, \emph{Supersymmetric derivation of the
  {Atiyah-Singer} index and the chiral anomaly},
  \href{http://dx.doi.org/10.1016/0550-3213(84)90506-6}{Nucl. Phys.
  \textbf{B235} (1984)},
395.
%%CITATION = NUPHA,B235,395;%%.

\bibitem[G]{G}
I.~M. Gelfand, \emph{On elliptic equations},
  \href{http://dx.doi.org/10.1070/RM1960v015n03ABEH004094}{Russian Math.
  Surveys \textbf{15} (1960)}, no.~3, 113--123.

\bibitem[Ge1]{Ge1}
E.~Getzler, \emph{Pseudodifferential operators on supermanifolds and the
  {A}tiyah-{S}inger index theorem}, Comm. Math. Phys. \textbf{92} (1983),
  no.~2, 163--178.

\bibitem[Ge2]{Ge2}
\bysame, \emph{A short proof of the local {A}tiyah-{S}inger index theorem},
  \href{http://dx.doi.org/10.1016/0040-9383(86)90008-X}{Topology \textbf{25}
  (1986)}, no.~1, 111--117.

\bibitem[Gi1]{Gi1}
P.~B. Gilkey, \emph{Curvature and the eigenvalues of the {L}aplacian for
  elliptic complexes},
  \href{http://dx.doi.org/10.1016/0001-8708(73)90119-9}{Advances in Math.
  \textbf{10} (1973)}, 344--382.

\bibitem[Gi2]{Gi2}
\bysame, \emph{The eta invariant for even-dimensional {${\rm PIN}_{{\rm c}}$}
  manifolds}, \href{http://dx.doi.org/10.1016/0001-8708(85)90119-7}{Adv. in
  Math. \textbf{58} (1985)}, no.~3, 243--284.

\bibitem[GJ]{GJ}
J.~Glimm and A.~Jaffe, \emph{Quantum physics: A Functional Integral Point of
  View}, second ed., Springer-Verlag, New York, 1987.

\bibitem[GK]{GK}
I.~C. Gohberg and M.~G. Krein, \emph{The basic results on the defect numbers,
  the root numbers and the indices of linear operators}, Magyar Tud. Akad. Mat.
  Fiz. Oszt. K\"{o}zl. \textbf{23} (1977), no.~3-4, 387--460. Translated from
  the Russian by K\'{a}roly Buz\'{a}si.

\bibitem[H1]{H1}
F.~Hirzebruch, \emph{Topological methods in algebraic geometry}, Classics in
  Mathematics, Springer-Verlag, Berlin, 1995. Translated from the German and
  Appendix One by R. L. E. Schwarzenberger, With a preface to the third English
  edition by the author and Schwarzenberger, Appendix Two by A. Borel, Reprint
  of the 1978 edition.

\bibitem[H2]{H2}
\bysame, \emph{A {R}iemann-{R}och theorem for differentiable manifolds},
  S\'{e}minaire {B}ourbaki, {V}ol. 5, Soc. Math. France, Paris, 1995, pp.~Exp.
  No. 177, 129--149.

\bibitem[H3]{H3}
\bysame, \emph{Some problems on differentiable and complex manifolds},
  \href{http://dx.doi.org/10.2307/1969629}{Ann. of Math. (2) \textbf{60}
  (1954)}, 213--236.

\bibitem[H4]{H4}
\bysame, \emph{Hilbert modular surfaces}, Enseign. Math. (2) \textbf{19}
  (1973), 183--281.

\bibitem[Hi]{Hi}
N.~J. Hitchin, \emph{Harmonic spinors}, Advances in Math. \textbf{14} (1974),
  1--55.

\bibitem[Ho1]{Ho1}
L.~H\"{o}rmander, \emph{Pseudo-differential operators},
  \href{http://dx.doi.org/10.1002/cpa.3160180307}{Comm. Pure Appl. Math.
  \textbf{18} (1965)}, 501--517.

\bibitem[Ho2]{Ho2}
\bysame, \emph{Pseudo-differential operators and hypoelliptic equations},
  Singular integrals ({P}roc. {S}ympos. {P}ure {M}ath., {V}ol. {X}, {C}hicago,
  {I}ll., 1966), 1967, pp.~138--183.

\bibitem[Ho3]{Ho3}
\bysame, \emph{The analysis of linear partial differential operators. {III}},
  Grundlehren der Mathematischen Wissenschaften [Fundamental Principles of
  Mathematical Sciences], vol. 274, Springer-Verlag, Berlin, 1985.
  Pseudodifferential operators.

\bibitem[Hod]{Hod}
W.~V.~D. Hodge, \emph{Harmonic {F}unctionals in a {R}iemannian {S}pace},
  \href{http://dx.doi.org/10.1112/plms/s2-38.1.72}{Proc. London Math. Soc. (2)
  \textbf{38} (1935)}, 72--95.

\bibitem[HS]{HS}
M.~J. Hopkins and I.~M. Singer, \emph{Quadratic functions in geometry,
  topology, and {M}-theory}, J. Diff. Geom. \textbf{70} (2005), 329--452,
\href{http://arxiv.org/abs/math/0211216}{{\tt arXiv:math/0211216}}.
%%CITATION = MATH/0211216;%%.

\bibitem[HZ]{HZ}
F.~Hirzebruch and D.~Zagier, \emph{The {A}tiyah-{S}inger theorem and elementary
  number theory}, Publish or Perish, Inc., Boston, Mass., 1974. Mathematics
  Lecture Series, No. 3.

\bibitem[J]{J}
K.~J\"{a}nich, \emph{Vektorraumb\"{u}ndel und der {R}aum der
  {F}redholm-{O}peratoren}, \href{http://dx.doi.org/10.1007/BF01360851}{Math.
  Ann. \textbf{161} (1965)}, 129--142.

\bibitem[JL]{JL}
J.~Jorgenson and S.~Lang,
  \href{http://dx.doi.org/10.1007/BFb0077194}{\emph{Basic analysis of
  regularized series and products}}, Lecture Notes in Mathematics, vol. 1564,
  Springer-Verlag, Berlin, 1993.

\bibitem[K]{K}
M.~Kac, \emph{Can one hear the shape of a drum?},
  \href{http://dx.doi.org/10.2307/2313748}{Amer. Math. Monthly \textbf{73}
  (1966)}, no.~4, part II, 1--23.

\bibitem[Kar]{Kar}
Max Karoubi, \emph{Alg{\`e}bres de {C}lifford et {$K$}-th\'{e}orie}, Ann. Sci.
  \'{E}cole Norm. Sup. (4) \textbf{1} (1968), 161--270.

\bibitem[KN]{KN}
J.~J. Kohn and L.~Nirenberg, \emph{An algebra of pseudo-differential
  operators}, \href{http://dx.doi.org/10.1002/cpa.3160180121}{Comm. Pure Appl.
  Math. \textbf{18} (1965)}, 269--305.

\bibitem[L]{L}
P.~L\'{e}vy, \emph{Wiener's random function, and other {L}aplacian random
  functions}, Proceedings of the {S}econd {B}erkeley {S}ymposium on
  {M}athematical {S}tatistics and {P}robability, 1950, University of California
  Press, Berkeley and Los Angeles, 1951, pp.~171--187.

\bibitem[LM]{LM}
H.~B. Lawson, Jr. and M.-L. Michelsohn, \emph{Spin geometry}, Princeton
  Mathematical Series, vol.~38, Princeton University Press, Princeton, NJ,
  1989.

\bibitem[Lo]{Lo}
Ya.~B. Lopatinski\u{\i}, \emph{On a method of reducing boundary problems for a
  system of differential equations of elliptic type to regular integral
  equations}, Ukrain. Mat. \v{Z}. \textbf{5} (1953), 123--151.

\bibitem[Lt1]{Lt1}
J.~Lott, \emph{Higher-degree analogs of the determinant line bundle},
  \href{http://dx.doi.org/10.1007/s00220-002-0686-3}{Comm. Math. Phys.
  \textbf{230} (2002)}, no.~1, 41--69,
  \href{http://arxiv.org/abs/arXiv:math/0106177}{{\tt arXiv:math/0106177}}.

\bibitem[Lt2]{Lt2}
\bysame, \emph{A proof of the axial anomaly}, Comm. Math. Phys. \textbf{97}
  (1985), no.~3, 371--379.

\bibitem[Mi1]{Mi1}
J.~W. Milnor, \emph{On the cobordism ring {$\Omega ^{\ast} $} and a complex
  analogue. {I}}, \href{http://dx.doi.org/10.2307/2372970}{Amer. J. Math.
  \textbf{82} (1960)}, 505--521.

\bibitem[Mi2]{Mi2}
\bysame, \emph{Lectures on the {$h$}-cobordism theorem}, Notes by L. Siebenmann
  and J. Sondow, Princeton University Press, Princeton, N.J., 1965.

\bibitem[Mik]{Mik}
S.~G. Mikhlin, \emph{Singular integral equations}, Amer. Math. Soc. Translation
  \textbf{1950} (1950), no.~24, 116. (Uspehi Matem. Nauk (N.S.), {\bf 3}
  (1948), 29--112).

\bibitem[MP]{MP}
S.~Minakshisundaram and \r{A} Pleijel, \emph{Some properties of the
  eigenfunctions of the {L}aplace-operator on {R}iemannian manifolds},
  \href{http://dx.doi.org/10.4153/cjm-1949-021-5}{Canad. J. Math. \textbf{1}
  (1949)}, 242--256.

\bibitem[MS]{MS}
H.~P. McKean, Jr. and I.~M. Singer, \emph{Curvature and the eigenvalues of the
  {L}aplacian}, J. Differential Geometry \textbf{1} (1967), no.~1, 43--69.

\bibitem[Mu]{Mu}
W.~M\"{u}ller, \emph{Signature defects of cusps of {H}ilbert modular varieties
  and values of {$L$}-series at {$s=1$}}, J. Differential Geom. \textbf{20}
  (1984), no.~1, 55--119.

\bibitem[N]{N}
F.~Noether, \emph{\"{U}ber eine {K}lasse singul\"{a}rer {I}ntegralgleichungen},
  \href{http://dx.doi.org/10.1007/BF01457974}{Math. Ann. \textbf{82} (1920)},
  no.~1-2, 42--63.

\bibitem[NAg]{NAg}
P.~Nelson and L.~Alvarez-Gaum{{\'e}}, \emph{Hamiltonian interpretation of
  anomalies}, Comm. Math. Phys. \textbf{99} (1985), no.~1, 103--114.

\bibitem[P1]{P1}
V.~K. Patodi, \emph{Curvature and the fundamental solution of the heat
  operator}, J. Indian Math. Soc. \textbf{34} (1970), no.~3-4, 269--285 (1971).

\bibitem[P2]{P2}
\bysame, \emph{An analytic proof of {R}iemann-{R}och-{H}irzebruch theorem for
  {K}aehler manifolds}, J. Differential Geometry \textbf{5} (1971), 251--283.

\bibitem[Pa]{Pa}
R.~S. Palais, \emph{Seminar on the {A}tiyah-{S}inger index theorem}, With
  contributions by M. F. Atiyah, A. Borel, E. E. Floyd, R. T. Seeley, W. Shih
  and R. Solovay. Annals of Mathematics Studies, No. 57, Princeton University
  Press, Princeton, N.J., 1965.

\bibitem[PS]{PS}
R.~S. Palais and R.~T. Seeley, \emph{Construction of $Int_k$}, Seminar on the
  {A}tiyah-{S}inger index theorem, With contributions by M. F. Atiyah, A.
  Borel, E. E. Floyd, R. T. Seeley, W. Shih and R. Solovay. Annals of
  Mathematics Studies, No. 57, Princeton University Press, Princeton, N.J.,
  1965, pp.~x+366. Chapter~XVI.

\bibitem[Q]{Q}
D.~Quillen, \emph{Determinants of {C}auchy-{R}iemann operators on {R}iemann
  surfaces}, Funktsional. Anal. i Prilozhen. \textbf{19} (1985), no.~1, 37--41,
  96. English translation: Functional Anal. Appl. 19 (1985), no. 1, 31--34.

\bibitem[Ri]{Ri}
B.~Riemann, \emph{Theorie der {A}bel'schen {F}unctionen},
  \href{http://dx.doi.org/10.1515/crll.1857.54.115}{J. Reine Angew. Math.
  \textbf{54} (1857)}, 115--155.

\bibitem[Ro]{Ro}
G.~Roch, \emph{Ueber die {A}nzahl der willk\"{u}rlichen {C}onstanten in
  algebraischen {F}unctionen},
  \href{http://dx.doi.org/10.1515/crll.1865.64.372}{J. Reine Angew. Math.
  \textbf{64} (1865)}, 372--376.

\bibitem[Roh]{Roh}
V.~A. Rohlin, \emph{New results in the theory of four-dimensional manifolds},
  Doklady Akad. Nauk SSSR (N.S.) \textbf{84} (1952), 221--224.

\bibitem[RS]{RS}
D.~B. Ray and I.~M. Singer, \emph{{$R$}-torsion and the {L}aplacian on
  {R}iemannian manifolds},
  \href{http://dx.doi.org/10.1016/0001-8708(71)90045-4}{Advances in Math.
  \textbf{7} (1971)}, 145--210.

\bibitem[S1]{S1}
I.~M. Singer, 2010. \url{https://www.youtube.com/watch?v=5FoaMcCJnmQ}.
  Interview of Isadore Singer for the MIT+150 Infinite History Project.

\bibitem[S2]{S2}
\bysame, \emph{The {$\eta$}-invariant and the index}, Mathematical aspects of
  string theory ({S}an {D}iego, {C}alif., 1986), Adv. Ser. Math. Phys., vol.~1,
  World Sci. Publishing, Singapore, 1987, pp.~239--258.

\bibitem[S3]{S3}
\bysame, \emph{Families of {D}irac operators with applications to physics}, no.
  Num\'{e}ro Hors S\'{e}rie, 1985, pp.~323--340. The mathematical heritage of
  \'{E}lie Cartan (Lyon, 1984).

\bibitem[Se1]{Se1}
R.~T. Seeley, \emph{Elliptic singular integral equations}, Singular {I}ntegrals
  ({P}roc. {S}ympos. {P}ure {M}ath., {C}hicago, {I}ll., 1966), Amer. Math.
  Soc., Providence, R.I., 1967, pp.~308--315.

\bibitem[Se2]{Se2}
\bysame, \emph{Singular integrals on compact manifolds},
  \href{http://dx.doi.org/10.2307/2372921}{Amer. J. Math. \textbf{81} (1959)},
  658--690.

\bibitem[Se3]{Se3}
\bysame, \emph{Integro-differential operators on vector bundles},
  \href{http://dx.doi.org/10.2307/1994203}{Trans. Amer. Math. Soc. \textbf{117}
  (1965)}, 167--204.

\bibitem[Se4]{Se4}
\bysame, \emph{Complex powers of an elliptic operator}, Singular {I}ntegrals
  ({P}roc. {S}ympos. {P}ure {M}ath., {C}hicago, {I}ll., 1966), Amer. Math.
  Soc., Providence, R.I., 1967, pp.~288--307.

\bibitem[Seg1]{Seg1}
G.~B. Segal, \emph{Equivariant {$K$}-theory}, Inst. Hautes {\'E}tudes Sci.
  Publ. Math. (1968), no.~34, 129--151.

\bibitem[Seg2]{Seg2}
\bysame, \emph{Faddeev's anomaly in Gauss's law}. preprint.

\bibitem[Seg3]{Seg3}
\bysame, \emph{The definition of conformal field theory}, Topology, geometry
  and quantum field theory, London Math. Soc. Lecture Note Ser., vol. 308,
  Cambridge Univ. Press, Cambridge, 2004, pp.~421--577.

\bibitem[Seg4]{Seg4}
\bysame, \emph{Michael Atiyah's work in algebraic topology}. this volume.

\bibitem[Sh]{Sh}
H.~Shimizu, \emph{On discontinuous groups operating on the product of the upper
  half planes}, \href{http://dx.doi.org/10.2307/1970201}{Ann. of Math. (2)
  \textbf{77} (1963)}, 33--71.

\bibitem[T1]{T1}
R.~Thom, \emph{Quelques propri{\'e}t{\'e}s globales des vari{\'e}t{\'e}s
  diff{\'e}rentiables}, Comment. Math. Helv. \textbf{28} (1954), 17--86.

\bibitem[T2]{T2}
\bysame, \emph{Espaces fibr\'{e}s en sph{\`e}res et carr\'{e}s de {S}teenrod},
  Ann. Sci. Ecole Norm. Sup. (3) \textbf{69} (1952), 109--182.

\bibitem[W]{W}
H.~Weyl, \emph{The classical groups}, Princeton Landmarks in Mathematics,
  Princeton University Press, Princeton, NJ, 1997. Their invariants and
  representations, Fifteenth printing, Princeton Paperbacks.

\bibitem[Wi1]{Wi1}
E.~Witten, \emph{Constraints on supersymmetry breaking},
  \href{http://dx.doi.org/10.1016/0550-3213(82)90071-2}{Nuclear Phys. B
  \textbf{202} (1982)}, no.~2, 253--316.

\bibitem[Wi2]{Wi2}
\bysame, \emph{Index of {D}irac operators}, Quantum Fields and Strings: a
  course for mathematicians, {V}ol. 1, 2 ({P}rinceton, {NJ}, 1996/1997), Amer.
  Math. Soc., Providence, RI, 1999, pp.~475--512.

\bibitem[Wi3]{Wi3}
\bysame, \emph{Global gravitational anomalies},
  \href{http://dx.doi.org/10.1007/BF01212448}{Commun. Math. Phys. \textbf{100}
  (1985)},
197.
%%CITATION = CMPHA,100,197;%%.

\bibitem[Wi4]{Wi4}
\bysame, \emph{Quantum field theory and the {J}ones polynomial}, Comm. Math.
  Phys. \textbf{121} (1989), no.~3, 351--399.

\bibitem[Wo]{Wo}
R.~Wood, \emph{Banach algebras and {B}ott periodicity},
  \href{http://dx.doi.org/10.1016/0040-9383(66)90035-8}{Topology \textbf{4}
  (1965/66)}, 371--389.

\end{thebibliography}


\clearpage
\section*{附录 A：审校批注汇编（Tutor Notes）}
\noindent 以下按节汇录正文中的全部审校批注（正文中以 \texttt{\% [Tutor Note: \dots]} 形式给出，
编译后不显示）。批注分四类：\textbf{易错}（符号、指标、归一化、定向/spin 结构等高风险处）、
\textbf{含糊}（作者以 obvious/standard 略去的证明实质）、\textbf{笔误}（原文疑似错误，照录并指出）、
\textbf{术语}（易混译名与概念边界）。

\subsection*{§1 引言}
\begin{enumerate}[leftmargin=2.2em,itemsep=2pt,topsep=2pt]
  \item 易错——“解构成一维复向量空间”依赖于重数被精确指定：若要求 $z_i$ 恰为 $n_i$ 重零点（当 $n_i>0$）或恰为 $|n_i|$ 阶极点（当 $n_i<0$），则全部解只相差常数因子 $c$；但黎曼–罗赫定理的“至少”提法取 $D=\sum n_iz_i$、$\sect D=H^0(\CP^1,\sO(D))$，此时当 $\deg D=\sum n_i\ge0$ 有 $\dim\sect D=\deg D+1=1+\sum n_i$，一般并不等于 $1$（对照式 \eqref{eq:1}）。原文亦未声明 $\sum n_i=0$：若 $\sum n_i\neq0$，$f$ 还在 $\infty$ 处有零点或极点，解空间仍是一维（相差一个常数因子）。
  \item 易错——“唯一约束”一句漏掉了次数条件：按阿贝尔（Abel）定理，$E=\CC/(\ZZ+\ZZ\tau )$ 上的除子 $D=\sum n_iz_i$ 是主除子当且仅当 $\deg D=\sum n_i=0$ 且 $\sum n_iz_i\in \ZZ+\ZZ\tau $（后者即除子的阿贝尔–雅可比像在雅可比簇 $J(E)$ 中为零）。原文只写出后一条件；例如取 $n_1=1$、$z_1=0$（$0$ 表 $E$ 中的原点），此时 $\sum n_iz_i=0\in \ZZ+\ZZ\tau $ 但 $\deg D=1\neq0$，并不存在以 $0$ 为唯一零点、别无极点的亚纯函数。对亏格为 $g$ 的一般曲面，“更多约束”正是阿贝尔定理的 $g$ 个周期条件。
  \item 含糊——“约束构成维数为亏格的向量空间”的精确含义是：对每个全纯微分 $\omega\in H^0(X,K_X)$ 都有周期条件 $\sum_i n_i\int_{z_0}^{z_i}\omega\equiv0$（模周期），共 $g=\dim H^0(X,K_X)=\dim\sect K$ 个独立复条件（由式 \eqref{eq:1} 取 $D=\sO$ 得 $\dim\sect K=g$，见 \S\ref{subsec:2.1}）。故“约束空间”指全纯微分空间的对偶 $H^0(X,K_X)^*$，它是 $g$ 维复向量空间；条件本身的取值落在 $g$ 维复环面（雅可比簇，Jacobian）$J(X)=H^0(X,K_X)^*/H_1(X;\ZZ)$ 中，固定度数的除子类构成 $J(X)$ 上的一个主齐性空间（torsor），并非线性空间。可查 [H1]。
  \item 含糊——“由拓扑来‘计数’”的实质是 Fredholm 指数：$\dbar$ 作用在线丛 $\sO(D)$ 的光滑截面上是椭圆算子，其指数 $\dim\ker\dbar-\dim\coker\dbar=\dim\sect D-\dim\sect{K-D}$ 恰为黎曼–罗赫公式 \eqref{eq:1} 的左端，只依赖 $\deg D$ 与 $g$ 这两个拓扑量；而单个空间 $\dim\sect D$ 的维数依赖于复结构。指数定理正是把“两个维数之差可以拓扑地计算”这一现象推广到一般椭圆算子；详见 \S\ref{subsec:2.1} 与 \S\ref{subsec:3.2}。
  \item 含糊——“于 1962–3 年提出并证明”是压缩的叙述：据本文 \S\ref{sec:2} 末尾的回忆，阿蒂亚与辛格的合作始于 1962 年 1 月（辛格访牛津），当年 3 月二人已得到狄拉克算子与指标公式，约九个月后完成第一证明；定理于 1963 年在 [AS1] 中宣布，第一证明（配边路线，仿照希策布鲁赫的号差与黎曼–罗赫证明）未由二人撰写，而发表于帕莱（Palais）主编的会议录 [Pa]；第二证明（$K$ 理论路线，仿照格罗滕迪克）于 1968 年发表于 [AS2]。此外还有以伪微分算子与热方程为工具的解析路线（[Se2,Se3]，本文 \S\ref{sec:6}，帕托迪（Patodi）、吉尔基（Gilkey））。按发表年份引用时，宜区分“1963 年宣布”与“1968 年发表证明”。
  \item 术语——Fredholm 指数定义为 $\dim\ker-\dim\coker$，但“有限维”并不是定义的一部分：闭流形上的椭圆算子经椭圆正则性（Sobolev 空间上的 Fredholm 性质）其核与余核才有限维，且与光滑截面上的核、余核一致；标准出处见本文 \S\ref{subsec:3.2} 以及 [Se2,Se3]、[Ho1]。余核须在适当的 Sobolev（或光滑）范畴中取，才能由正则性保证其有限维性。
  \item 术语——cobordism 与 bordism 均译“配边”：本文 \S\ref{sec:2} 的脚注指出托姆（Thom）、希策布鲁赫等人用“cobordism”代替“bordism”，阿蒂亚 [A10] 澄清了二者的关系，故文中两种拼法并存，不必据拼法区分理论。又，希策布鲁赫的“号差定理”即 signature theorem（号差 = signature，切勿与“符号”symbol 混同），它与黎曼–罗赫定理都是指数定理的特例。
  \item 含糊——“受格罗滕迪克的黎曼–罗赫定理启发”的实质：格罗滕迪克把黎曼–罗赫定理推广为相对版本（对复流形间的 proper 映射，见 \S\ref{subsec:2.2} 与 [H1]），其 $K$ 理论形式用 $f_!$ 与 Todd 类、陈特征的关系表达；阿蒂亚–辛格的第二证明把这一相对形式搬到微分拓扑中，在 $K^0(T^*X,T^*X\setminus0)$ 上使用 Thom 同构与差丛（difference bundle）构造，见 \S\ref{subsec:4.4} 与 [AS2]。
  \item 术语——“局部指数定理”（local index theorem）指热核 $e^{-t\Delta }$ 当 $t\to0$ 时的对角渐近展开给出指数密度，其积分等于指标；它与“整体拓扑不变量”的形式相对，但“局部”仅指公式的逐点形式，不表示指标本身可以局部化。本文 \S\ref{sec:6} 指出该定理的特殊情形由帕托迪、一般情形由吉尔基在 1970 年代初证明。
  \item 术语——[F1] 是 Freed 本人 2021 年 4 月的同名讲座视频（文献表中给出链接），并非论文；正文“a lecture based on this paper may be viewed at [F1]”即指此条目，引用时不要把它当作印刷文献。
\end{enumerate}

\subsection*{§2 来自代数几何与拓扑的先驱与动机}
\begin{enumerate}[leftmargin=2.2em,itemsep=2pt,topsep=2pt]
  \item 易错——源文件的小节标号与出现顺序不一致：按出现顺序依次是 subsec:2.1（黎曼–罗赫）、subsec:2.3（希策布鲁赫）、subsec:2.2（格罗滕迪克）、subsec:2.4（整性定理），而正文第 428、431 行的 \S\ref{subsec:2.3}、\S\ref{subsec:2.2} 正是按这套标号引用的。译文保留原标号、不按顺序重编，否则交叉引用会指向错误小节；排版时编号由出现顺序自动给出，仍是 2.1–2.4。
  \item 术语——"stable homotopy groups of Lie groups" 指经典群的稳定同伦群 $\pi_{n+k}(U(n))$、$\pi_{n+k}(O(n))$（$n\to\infty$ 的稳定范围），即 Bott 周期性所计算的对象，不是"每个李群自身的同伦群"。第 753–755 行明确说这里算的是酉群，两处指同一件事；Bott 的 {\it periodicity theorem} 正是由此得名。
  \item 术语——divisor 是整系数形式和 $D=\sum_x n_xx$（$n_x=\ord_x(D)\in\ZZ$，只有有限多个非零），并非集合意义上的"点集"；原句"每点附一个整数"是简化说法。 $\deg D=\sum_x\ord_x(D)$。另外 $\sect D$ 就是 $\sL(D)=H^0(X,\sO(D))$：源文件的宏定义为 \textbackslash{}sect{\#1}=\textbackslash{}sL(\#1)（第 265 行），故第 488 行写成的 $\sL(K)$ 与 $\sect K$ 完全同义。
  \item 易错——上句"至多有 $\ord_x(D)$ 阶极点"必须按整系数读：$\ord_x(D)>0$ 时是极点阶数的上界；$\ord_x(D)<0$ 时应读作"$x$ 处至少有 $-\ord_x(D)$ 阶零点"；而 $x\notin D$ 时 $\ord_x(D)=0$，即要求全纯。漏掉负系数情形就理解不了后面 $K-D$ 这类除子为何可能没有截面。
  \item 易错——定理左端其实是 $\sO(D)$ 的层上同调 Euler 示性数：由 Serre 对偶 $H^1(X,\sO(D))^\vee\cong H^0(X,K-D)$，即 $\dim\sect{K-D}=h^1(D)$，故左端 $=\chi(X,\sO(D))$。符号要点：对偶写的是 $K-D$；若误写成 $K+D$，$h^0(K+D)$ 与 $h^1(D)$ 不再对偶，(1) 即不成立。
  \item 易错——$\deg(K)=2g-2$，故 $\deg(D)>2g-2\iff\deg(K-D)<0$；负次数除子没有非零全纯截面（$\deg E<0\Rightarrow H^0(E)=0$），这才是 $\sect{K-D}=0$ 的理由。原文用 "it can be shown" 带过了这一步。另外 $g=\frac12\rank H^1(X;\ZZ)$ 用的是自由秩（$\dim_\RR H^1=2g$），不要把 $g$ 与 $\dim\sect K$ 混为一谈：后者等于 $g$ 是下面用 $D=\sO$ 代入 (1) 推出的。
  \item 易错——这里的 $\Euler(X)$ 是实二维闭曲面的欧拉示性数 $2-2g$，故 $1-\Euler(X)/2=g$。整性链条是：$\dim\sect K=g\in\ZZ$（层论/Serre 对偶）$\Rightarrow$ $\rank H^1=2g$ 为偶数 $\Rightarrow$ $\Euler/2\in\ZZ$。三步实为同一组等式，核对时注意不要把这个逻辑顺序读反。
  \item 含糊——曲面情形的推广（即 \eqref{eq:4} 的 Noether 公式）是关于 Euler 示性数 $\chi(\sO_X)$ 的等式，而不是逐项 $h^0$ 的公式；19 世纪末这三人的结果要到 Hirzebruch 的 \eqref{eq:7} 才被统一为"示性类积分"的形式。
  \item 术语——这条脚注讲的是术语源流：老文献（Thom、Hirzebruch 等）把 "bordism" 说成 "cobordism"。不同作者用法不一：bordism 多指流形在配边关系下的等价类与配边群 $\Omega_*^{SO}$，cobordism 有时指具体的配边 $W$（$\partial W=M\sqcup N$），有时也指同一等价关系。 Atiyah\textasciitilde{}\cite{A10} 澄清了这一用法。本节其余各处按中文惯例一律译作"配边"。
  \item 含糊——"可以用椭圆算子理论与 Dolbeault 定理证明"省略的实质是 Hodge 理论：$H^q(X,V)\cong H^{0,q}_{\bar\partial}(X,V)\cong\mathcal H^{0,q}(X,V)$ （调和 $(0,q)$-形式空间），而椭圆算子 $\dbar+\dbar^*$ 的核有限维（详见 \S\S\ref{subsec:3.1}--\ref{subsec:3.2}）。原文一句话带过了整个 Hodge–Dolbeault 同构，需要时可查 [H1] 与本文 \S\ref{subsec:3.1}。
  \item 术语——$\chi(X,V)$ 是层上同调的 Euler 示性数（交错和），不是流形的 Euler 数 $\Euler(X)$：例如 $n=1$、$V=\sO$ 时 $\chi(X,\sO)=1-g$，而 $\Euler(X)=2-2g$。\eqref{eq:4} 中的 $\chi(X)$ 是 $\chi(X,\sO_X)$ 的简写，尽管字母相同，含义与 $\Euler$ 不同。
  \item 易错——\eqref{eq:4} 是 Noether 公式 $\chi(\sO_X)=\frac1{12}(c_1^2+c_2)[X]$，其中 $c_1^2$ 表示 $c_1(X)\smile c_1(X)$，$c_2[X]=\Euler(X)$。由 $p_1=c_1^2-2c_2$、 $\Sign=p_1/3$ 可把它改写成 $\chi=(\Sign+\Euler)/4$。用 $K3$ 曲面核对：$\Sign=-16$、 $\Euler=24$，得 $\chi=(-16+24)/4=2$，与 $(0+24)/12=2$ 一致。若把分母 12 记成 24，就会与 $\Ahat$ 的 $-p_1/24$ 混淆。
  \item 易错——Todd 类的逐因子展开为 $\frac{y}{1-e^{-y}}=1+\frac y2+\frac{y^2}{12}+\cdots$，故 $\Todd(X)=1+\frac{c_1}{2}+\frac{c_1^2+c_2}{12}+\cdots$。本节出现三个互不相同的归一化，务必分清：Todd 的 4 维项 $\frac1{12}(c_1^2+c_2)$（见 \eqref{eq:4}）、 $L$-类的 4 维项 $\frac{p_1}{3}$（见 \eqref{eq:9}）、$\Ahat$-类的 4 维项 $-\frac{p_1}{24}$（见 \eqref{eq:22}）。$\frac12$ 因子（$c_1/2$）来自 $\frac{y}{1-e^{-y}}$ 与 $\Ahat$ 的换算，不是笔误。
  \item 含糊——"分裂原理"的实质：存在旗流形 $\pi:F(X)\to X$ 使 $\pi^*TX$ 分裂为线丛之和，且 $\pi^*:H^*(X)\hookrightarrow H^*(F(X))$ 为单射；于是先在分裂情形定义示性类/示性数，再证明结果与所取分裂无关。$\ch$ 的加性（$\ch(V\oplus W)=\ch V+\ch W$）与 $\Todd$ 对直和的乘性正是在此用到。
  \item 易错——中间杯积配对要在 $H^{2k}(X;\RR)$（实系数）上谈对称性与非退化性；因 $2k$ 是偶数，$\alpha_1\smile\alpha_2=\alpha_2\smile\alpha_1$， $\Sign$ 才有定义。若 $\dim X\equiv2\pmod4$，中间配对是交错的，"号差"恒为 0，故号差定理只对 $4k$ 维陈述；原文的 $4k$ 不可省。
  \item 术语——原文特意注明较老文献把 signature 叫作 "index"，这正是 Hirzebruch 号差定理旧称 "index theorem" 的由来；它与第 4 节起讨论的 Fredholm 指数 $\ind D$ 是两个概念，二者要等到指标定理才统一（号差等于号差算子的指标）。
  \item 易错——\eqref{eq:9} 的乘积上限是 $2k$：实 $4k$ 维流形的复化切丛 $TX\otimes\CC$ 分裂为 $4k$ 个复线丛，其 Chern 根成对出现 $y_i,-y_i$ （$i=1,\dots,2k$），$p_1=\sum y_i^2$。展开得 $L(X)=1+\frac{p_1}{3} +\frac{7p_2-p_1^2}{45}+\cdots$，故 $\dim X=4$ 时 $\Sign=p_1/3$。若上限误取 $k$ 或漏掉根的另一半，度数与符号都会错。
  \item 含糊——"只需在生成元上验证"省略了两件事：(i) Thom 定理给出 $\Omega_*^{SO}\otimes\QQ=\QQ[\CP^2,\CP^4,\dots]$，故偶射影空间是有理生成元（\cite{T1}）；(ii) 号差在定向配边下不变是 Thom 的定理 [T2，\S IV]，而号差对乘积是乘性的（$\Sign(M\times N)=\Sign M\cdot\Sign N$，由 Künneth 公式与中间次数配对得到）。两条合起来，才能从生成元推广到所有闭定向流形。
  \item 含糊——把"凝聚代数层"换成"全纯向量丛"而 $K(X)$ 不变，依赖 $X$ 光滑：此时每个凝聚层有有限长度的局部自由（向量丛）消解，故两者的 Grothendieck 群一致（Serre；[BS] 即 Borel–Serre 的黎曼–罗赫论文）。在一般概形上这一替换不成立。
  \item 术语——此处的 $K(X)$ 是凝聚层生成的自由阿贝尔群模去由短正合列生成的子群（Grothendieck 群），不是拓扑 $K$ 理论；下一段的 $f_!$ 是层论直接像的交错和。 \eqref{eq:11} 中的 $R^qf_*$ 是"高次直接像层"，不是同调群。
  \item 含糊——$f_!$ 良定义除需要短正合列给出的长正合列（保证 $\sum(-1)^qR^qf_*$ 只依赖 $K$ 类）之外，还需 $R^qf_*=0$ 对 $q\gg0$ 成立，这要求 $f$ 真（下一段因此明确假设 $f$ proper）；非逆紧态射的高次直接像未必凝聚，$f_!$ 未必落在 $K(Y)$ 中。
  \item 易错——\eqref{eq:13} 中 $L=\sO_Y(X)$，故 $L\inv=\sO_Y(-X)$， \eqref{eq:14} 给出 $f_!\sO_X=\sO_Y-L\inv$；若把 $L\inv$ 写成 $L$，符号就反了。 \eqref{eq:16} 的 $f_*(1)=y$ 是 Gysin 前推：$f_*:H^0(X)\to H^2(Y)$ 把 $1$ 送到法丛的 Thom 类 $c_1(L)=y\in H^2(Y)$（Euler 类 $f^*c_1(L)\in H^2(X)$ 则是它的拉回）。把 \eqref{eq:15} 与 \eqref{eq:16} 对比可见 $f_*\circ\ch\ne\ch\circ f_!$，相差的因子恰是 $\Todd(L)=\frac{y}{1-e^{-y}}$，即 \eqref{eq:128} 末行。
  \item 笔误——\eqref{eq:127}、\eqref{eq:128} 的标签命名是源文件修订时的旧编号残留（label 名与排版编号无关，不影响引用，故译文照录原标号）。另注意 \eqref{eq:128} 第一步是把 \eqref{eq:12} 除以 $Y$ 上的 $\Todd(Y)$：$\Todd(Y)$ 的常数项为 1、故可逆， $\frac{\Todd(X)}{f^*\Todd(Y)}$ 中分母是拉回到 $X$ 的类。
  \item 易错——\eqref{eq:17} 是 Koszul 复形（比较 \eqref{eq:13} 的除子情形），映射为与 $e$ 的缩并 $\iota(e)$，$k=\rank E$，复形在 $e\ne0$ 处正合，故其 $K$ 类以零截面 $X$ 为支撑。\eqref{eq:18} 的 $\wedge^\bullet(E^*)$ 是交错和 $\sum_{k=0}^{\rank E}(-1)^k\wedge^kE^*$；原文 $\sum\limits_{}$ 的求和下标留空，是省略记法而非笔误。漏掉 $(-1)^k$ 得到的是和 $\sum_k[\wedge^kE^*]$（其陈特征为 $\prod_i(1+e^{-y_i})$），不再是 Thom 类。
  \item 术语——genus 在本节有两种用法：(i) 曲线/曲面的亏格 $g$（拓扑不变量）；(ii) 示性类意义下的"亏格"，即上同调类在基本类上取值的数。第 710–711 行的 "Todd 亏格 $\Todd(X)[X]$" 属 (ii)：$\Todd(X)$ 是类（Todd 类），$\Todd(X)[X]$ 才是数（Todd 亏格）；同理有 $\Ahat$-类与 $\Ahat$ 亏格、$L$-类与 $L$ 亏格。注意 $\Ahat$ 亏格与 Todd 亏格是两个不同的不变量，仅在近复流形上由 \eqref{eq:23} 相联系。
  \item 易错——归一化核对：\eqref{eq:22} 的 $\Ahat(X)=\prod_{i=1}^{2k}\frac{y_i/2}{\sinh y_i/2}$ 展开为 $1-\frac{p_1}{24}+\cdots$。$K3$ 曲面：$p_1=-48$，故 $\Ahat(X)[X]=-p_1/24=2$（即 Dirac 算子的指数），而 $L(X)[X]=\Sign=p_1/3=-16$。对 spin 四维流形 $p_1=3\Sign$，于是 $\Ahat[X]=-\Sign/8$；Rokhlin 定理 $\Sign\equiv0\pmod{16}$ 等价于 $p_1\equiv0\pmod{48}$，正是它保证 $\Ahat[X]\in\ZZ$（且为偶数）。文献中若见到 $-p_1/48$ 之类系数，那是把 $p_1/48$ 的整除性误当成 $\Ahat$ 的归一化，或 $A$ 与 $\Ahat$ 旧记号混用所致；引用前须按 \eqref{eq:22} 逐项换算。
  \item 含糊——"$w_2$ 是 $c_1$ 的模 2 约化"只对近复流形成立（此时 $w_2(X)\equiv c_1(X)\pmod2$）；一般实流形没有 $c_1$，故脚注把问题放宽为"存在 $c\in H^2(X;\ZZ)$ 使 $c\equiv w_2(X)\pmod2$"。这个 $c$ 正是后文 \eqref{eq:19} 中"定向"条件的一般形式，也是 \eqref{eq:23} 中 $e^{c/2}$ 的来源。
  \item 易错——\eqref{eq:23} 的 $c_1/2$ 因子不是任意的：逐因子恒等式 $\frac{y}{1-e^{-y}}=e^{y/2}\frac{y/2}{\sinh(y/2)}$（因 $1-e^{-y}=2e^{-y/2}\sinh\frac y2$）说明 $e^{c_1/2}$ 恰好补上 Todd 的 $y_i$ 与 $\Ahat$ 的 $y_i/2$ 之间的差别；写成 $e^{c_1}$ 会差一个 $e^{c_1/2}$。校核到 4 次项： $e^{c_1/2}\Ahat=1+\frac{c_1}{2}+\frac{c_1^2+c_2}{12}$，与 \eqref{eq:5} 一致（用 $p_1=c_1^2-2c_2$）。
  \item 易错——$K(S^n)$ 的约化部分同构于 $\pi_{n-1}(U)$ 而非 $\pi_n(U)$： $\tilde K(S^n)\cong[S^n,BU\times\ZZ]\cong\pi_n(BU)\cong\pi_{n-1}(U)$。Bott 给出 $\pi_{n-1}(U)=\ZZ$（$n$ 偶）、$0$（$n$ 奇），周期性由此而来。原文"归约为酉群的稳定同伦群"不可读成"球面的同伦群由酉群给出"。
  \item 术语——"extraordinary cohomology" 即广义（上）同调：满足 Eilenberg–MacLane 公理中除维数公理外全部公理的上同调理论（$K$ 理论、配边理论等）；中文文献亦译"异常上同调"。维数公理失效正是它有别于奇异上同调之处。
  \item 术语——这里的"定向"不是流形的定向（只需 $w_1=0$），而是 $K$ 理论 Thom 类存在所需的 $w_2$-定向：$c\equiv w_2(X)-f^*w_2(Y)\pmod2$。近复流形取 $c=c_1(X)$ 即满足（$w_2$ 是 $c_1$ 的模 2 约化）；spin 流形取 $c=0$。$c$ 的作用是在 \eqref{eq:23} 中用 $e^{c/2}$ 代替 $e^{c_1/2}$（见第 734–736 行脚注）。
  \item 含糊——\eqref{eq:20} 就是 $K$ 理论 Thom 类的构造：$f_!(1)\in K(Y)$ 是法丛的 Thom 类，"零延拓"一步用到管状邻域定理，与 \eqref{eq:17} 的 Koszul 复形是同一件事（把 $E$ 上的复形限制/延拓）。张量 $\pi^*V$ 即得 $f_!(V)$。下文"证明是一个练习：比较 $K$ 理论与上同调的 Thom 同构"实质是要验证 $\ch(U_K)=\Todd(E)\smile U_H$（$U_K,U_H$ 分别为两种 Thom 类），细节见 [ABS，\S12]。
  \item 笔误——第 812 行 "for $z=0$ in\textasciitilde{}\eqref{eq:21}" 疑应为 $z=1$ （$K(X)$ 的单位元）：$z=0$ 时 \eqref{eq:21} 两端都是 0，推不出任何整性；只有 $z=1$ 时左端才给出 $\ch(f_!(1))$，再由 $Y=\pt$、$c=0$ 及 \eqref{eq:23} 得 $\ch(f_!(1))=f_*(\Todd(X))=\Ahat(X)[X]$。译文照录原文 $z=0$，此处指出。
  \item 易错——顺着上一条：$K(\pt)\cong\ZZ$，而 $\ch:K(\pt)\to H^*(\pt;\QQ)$ 是恒等映射，故 $f_!(1)$ 这个整数与 $\Ahat(X)[X]$ 的相等是"数"层面的等式；但 $f_!(1)$ 作为 $K$ 类存在，比"示性数恰为整数"强，这正是这一论证被称为"给出了 $\Ahat[X]$ 的拓扑解释"的原因。
  \item 含糊——"号差是调和形式空间的维数之差"省略了 Hodge 理论： $\Sign(X)$ 等于杯积形式限制在 $H^{2k}$ 的调和形式空间上所得对称型的正、负特征子空间维数之差（Hodge 分解 + 中间次数的 Hodge 星算子）。Hirzebruch 要的是 $\Ahat$ 对应的解析刻画，即某椭圆算子的核与余核维数之差，答案就是 Dirac 算子的指数 $\ind D$。
  \item 术语——"A-roof genus" 是 $\Ahat$ 亏格的口语读法（$\hat A$ 形似屋顶，读作 "A-roof" 或 "A-hat"），并非另一个不变量；它与 \eqref{eq:24} 的问句是同一个问题，只是这里要求的是分析/指标解释，而非用 Atiyah–Hirzebruch $K$ 理论给出的拓扑解释。
\end{enumerate}

\subsection*{§3 分析中的先驱}
\begin{enumerate}[leftmargin=2.2em,itemsep=2pt,topsep=2pt]
  \item 术语——对 $q\ge1$，(3.2) 的分子 $\ker[d:\Omega^q\to\Omega^{q+1}]$ 与分母 $\image[d:\Omega^{q-1}\to\Omega^q]$ 都是无限维向量空间（$q=0$ 时分子为常数、分母为 $0$），取商之后才有限维；有限性可由 Hodge 理论（(3.4) 的椭圆性）或 de Rham 定理证明，无法从 (3.2) 直接看出。
  \item 含糊——“存在唯一的‘最佳’形式”所略去的是 Hodge 分解 $\Omega^q(X)=\sH^q(X)\oplus\image d\oplus\image d^*$：存在性由在闭流形上极小化 $L^2$ 范数得到（椭圆算子 $\Delta$ 的 Fredholm 性及紧性），唯一性由 $\langle\Delta\omega,\omega\rangle=\|d\omega\|^2+\|d^*\omega\|^2$ 得到。原始文献为 [Hod]；椭圆理论的证明见 [Ho3]。
  \item 易错——Hodge 星号 $*$ 的归一化决定 (3.3) 的写法：若取 $\alpha\wedge *\beta=\langle\alpha,\beta\rangle\,\mathrm{vol}$，则伴随为 $d^*=(-1)^{np+n+1}*d*$（$p$ 为形式次数），故 $d(*\omega)=0$ 与 $d^*\omega=0$ 相差一个可逆的 Hodge 星号算子（可能差符号），两种写法等价；但若 $*$ 的约定整体差一个常数或符号，则 $d^*$ 与 $\Delta$ 的整体符号须重新核对（(3.4) 取 $\Delta=dd^*+d^*d$ 使 $\Delta$ 非负；几何分析文献有时把非负拉普拉斯算子记作 $-\Delta$，与 (3.4) 相差一个整体符号）。
  \item 易错——“在闭流形上 (3.4) 等价于 (3.3) 中的两个方程”依赖分部积分：由 $\int_X\langle\Delta\omega,\omega\rangle\,\mathrm{vol}=\|d\omega\|^2+\|d^*\omega\|^2$ 知 $\Delta\omega=0$ 当且仅当 $d\omega=d^*\omega=0$，而 $d^*\omega=0$ 与 $d(*\omega)=0$ 相差一个可逆的 Hodge 星号算子（可能差符号）。闭流形（无边）保证分部积分无边界项；带边流形上须补加边界条件。
  \item 术语——$\sH^q(X)$ 是调和 $q$-形式构成的空间（由 $\Delta\omega=0$ 定义），与上同调群 $H^q(X;\RR)$（由闭形式模恰当形式定义）是不同对象；“更强”的同构 $\sH^q(X)\xrightarrow{\;\cong\;}H^q(X;\RR)$ 才是 Hodge 定理的完整内容，$\dim\sH^q=b_q$ 只是其推论。
  \item 易错——“这两个论断都不能推广”指对一般的椭圆算子：$\dim\ker P$ 不再是某个贝蒂数，$\ker P$ 也不再与某个上同调群同构；保留下来的只有交错和 $\ind P=\dim\ker P-\dim\coker P$（核与余核一般也不等维，见 (3.6)），这就是 §3.2 的 Fredholm 指数。
  \item 易错——(3.5) 的 $\Euler(X)$ 是常系数实上同调的交错和，对应平凡系数 $\RR$ 的 de Rham 复形；而 (2.3) 的全纯欧拉示性数即使取平凡全纯线丛也给出 $\Todd(X)[X]$，一般不等于 $\Euler(X)$。作者“(Compare (2.3))”只是提醒同一“交错和”格式，不可读成数值相等。
  \item 易错——把 (3.6) 的余核与伴随的核等同，需要 $d+d^*$ 为 Fredholm 算子（闭值域、有限维核与余核），此时才有 $\coker T\cong\ker T^*$；对一般算子这仅在闭值域时成立，否则 $\ker T^*=(\image T)^\perp$ 不再同构于 $\coker T$。此处由椭圆理论保证。
  \item 易错——脚注中的 $\tau$ 是对合：由 $*^2=(-1)^{p(n-p)}$（$n=2\ell$）可得 $\tau^2\omega=(-1)^{p(n-p)}\,i^{p(p-1)+q(q-1)+2\ell}\omega=\omega$（$q=n-p$，且 $i$ 的幂为 $(-1)^p$），故 $\Omega^\pm$ 是 $\tau$ 的 $\pm1$ 特征空间。若把 $\tau$ 换成 $-\tau$（或把 $i$ 的幂次整体乘 $-1$），$\Omega^\pm$ 互换，(3.7) 的指数变号；此 $\zt$-分级不同于 (3.6) 的偶/奇次分级。核对 [AS3, §6] 的约定。
  \item 含糊——(3.7) 的指数等于号差，其表述要求 $\dim X$ 为 $4$ 的倍数（中间维数 $2k$ 为偶数，(2.9) 的配对才有非零号差）；若 $\dim X\equiv2\pmod4$，(2.9) 的配对在中间维数是交替的、号差为 $0$，(3.7) 的指数也应为 $0$。此外定向与 $\Omega^\pm$ 的选取须与 (2.9) 一致，否则差一个负号。
  \item 术语——复流形上 $\Omega^{p,q}$ 的第一个指标是全纯次数、第二个是反全纯次数，$\dbar$ 只增大第二个指标；故 (3.8) 是 $(0,q)$ 行，Dolbeault 上同调记作 $H^{0,q}_{\dbar}(X,V)$，与层上同调 $H^q(X,V)$ 的 $q$ 对应。若与 de Rham 的 $H^q_{\text{DR}}$ 混用次数，会差一位。
  \item 易错——此处 $n$ 改指复维数（$X$ 是 $n$ 维复流形），与本节开头 $n=\dim_{\RR}X$ 的用法不同；(3.8) 的末项 $\Omega^{0,n}$ 与 (3.9) 的求和上限 $n$ 都应读作复维数，请结合 $\dim X=1$ 的退化情形核对。
  \item 含糊——“Hodge 理论表明……”分两步：(i) 对任意紧 Hermitian 复流形，$\dbar$-Laplacian $\Delta_{\dbar}=\dbar\dbar^*+\dbar^*\dbar$ 的 Hodge 理论（$\dbar$ 复形是椭圆复形）已给出 $H^{0,q}(X,V)\cong\sH^{0,q}_{\dbar}(X,V)$；(ii) 凯勒条件用于 Kähler 恒等式 $\Delta_d=2\Delta_{\dbar}$，从而 $\dbar$-调和的 $(0,q)$-形式恰为 de Rham 意义下的复调和形式，并有 $H^r(X;\CC)=\bigoplus_{p+q=r}H^{p,q}(X)$。见 [H1]。
  \item 笔误——(3.9) 左端原文作 $\sum_{q=0}^n\dim\sH^{0,q}(X)$，但按上文“交错和”的表述与 (2.3) 的定义，疑应为 $\sum_{q=0}^n(-1)^q\dim\sH^{0,q}(X,V)$：既缺 $(-1)^q$，又缺系数丛 $V$，否则左端不可能等于 $\chi(X,V)=\Todd(X)\ch(V)[X]$。译文照录原文，未改 (3.9)。
  \item 术语——原文所谓 $P$ 的 symbol $\sigma(P)$ 实为主象征（principal symbol），只记录最高阶项；(3.10) 的全象征还含各低阶项。下文“symbols 的连续路径”也指主象征的路径。
  \item 易错——“象征可逆”是逐点条件：对每个 $x$ 与每个非零 $\theta\in T^*_xX$，$\sigma(P)(\theta,\dots,\theta)$ 可逆；它在 $\theta=0$ 处不可逆，故准确说法是 $\sigma(P)$ 在 $T^*X\setminus0$ 上给出丛同构。这要求 $\rank E^0=\rank E^1$；椭圆性是开条件，对主象征的 $C^0$ 扰动稳定。
  \item 含糊——“由椭圆理论可知核与余核有限维（用到紧致性）”所略去的是：椭圆正则性（$Pu\in C^\infty$ 推出 $u\in C^\infty$，见 [Se2] 与 [Ho3]）加上 Rellich–Kondrachov 定理（$X$ 紧时 $H^s\hookrightarrow H^{s-m}$ 为紧嵌入），由此 $P:H^s\to H^{s-m}$ 为 Fredholm 算子。$X$ 非紧时紧嵌入失效，核或余核可以是无限维的。
  \item 含糊——“适当 Sobolev 空间”指对每个实数 $s$，$P$ 延拓为 $P_s:H^s(X,E^0)\to H^{s-m}(X,E^1)$ 且为 Fredholm 算子；其核在 $s$ 充分大时与 $C^\infty$ 核一致，指数与 $s$ 无关。细节见 [AS2, §1] 与 [Ho3]。
  \item 术语——Fredholm 算子的定义关键是“闭值域”（不是“闭算子”）：$T$ 为 Fredholm 算子当且仅当它在 Calkin 代数（有界算子模紧算子构成的商代数）中可逆（Atkinson 定理），此时 $\coker T\cong\ker T^*$。只有有限维核与余核而不闭值域者不算 Fredholm 算子。
  \item 含糊——(3.12) 的符号脚注实质是行列式的约定问题：若在次数 $0,1$ 上写两项复形 $0\to V^0\to V^1\to0$，则 $\sum(-1)^q\dim V^q=\dim V^0-\dim V^1=\dim\ker-\dim\coker$ 恰为 (3.12)，故数值上并无歧义；作者的顾虑（“符号本应相反”）来自 $T\in V^1\otimes(V^0)^*$ 的对偶性以及 $\det T$ 与 $\det T^*$ 在行列式线意义下互为倒数，这在后文行列式线丛的讨论中会再次出现。
  \item 笔误——脚注末句原文作“the usual sign for the index (as in (3)) reduce to (30)”，其中 reduce 疑应为 reduces（主语是 the usual sign）；按印刷号即 (2.3) 与 (3.12)，译文按单数理解。
  \item 含糊——Fredholm 算子空间“有无界阶的非平凡同伦群”指 Atiyah–Jänich 定理：$\Fred(H^0,H^1)$ 同伦等价于 $\ZZ\times BU$（复 K 理论的分类空间），故由 Bott 周期性（周期 $2$）各阶同伦群在无穷多阶非平凡，且 $\pi_0\cong\ZZ$ 正由 (3.13) 的指数给出。见 [A4]；这正是 §4–§5 用 K 理论证明指数定理的出发点。
  \item 含糊——“紧算子平移保持 Fredholm 子空间且指数不变”即 Atkinson 定理的推论：$T$ Fredholm、$K$ 紧则 $T+K$ Fredholm 且 $\ind(T+K)=\ind T$；这也是下段“低阶项相对最高阶项为紧”能推出“指数只依赖象征”的原因。
  \item 含糊——“低阶项相对最高阶项为紧”即 Rellich 定理：$m$ 阶算子与 $(m-1)$ 阶算子之差 $H^s\to H^{s-m}$ 为紧算子（经由紧嵌入 $H^{s-m+1}\hookrightarrow H^{s-m}$）；等价地，$P$ 与只保留主象征的算子相差一个紧算子，故指数相同。参见 [Se2] 与 [Ho3]。
  \item 术语——“index”（指数）在此指解析的 Fredholm 指数 (3.12)；§4 将证明它等于拓扑指数（由主象征的同伦类经 K 理论推前给出）。注意 §2 提到旧文献把 signature 也称作 index，那与 Fredholm 指数不是同一概念，勿混。
  \item 术语——Fritz Noether（弗里茨·诺特）是 Emmy Noether（埃米·诺特）之弟，两人中文均译“诺特”，此处须保留名“弗里茨”以免混淆；[E,AM] 都把 Fredholm 指数的首创归于其 1920 年 Math. Ann. 论文 [N]。
  \item 术语——旧文献的“Noether property”（诺特性质）即今之 Fredholm 性（闭值域、有限维核与余核），而旧称“Fredholm”专指指数为 $0$ 者；“Fredholm alternative”（Fredholm 二择一）指 $Tu=f$ 要么对每个 $f$ 唯一可解、要么齐次方程有非平凡解，恰是指数为 $0$ 的 Fredholm 算子的性质。
  \item 笔误——原文此句“in the older literature the following terminology that is sometimes used”语法不完整，疑脱漏 is（应为“the following terminology is sometimes used”）；译文按其本意译出。
  \item 易错——$\sH$ 若理解为圆盘上的 Bergman 空间（$L^2$ 全纯函数），它并不能按傅里叶级数嵌入 $L^2(\cir)$；下一句“By Fourier series $\sH$ sits as a subspace in $\tsH$”表明这里实际是 Hardy 空间 $H^2$（边界值实现，Fourier 系数只含非负频率）。$i$ 为包含映射、$\pi$ 为正交投影（Szegő 投影）。
  \item 易错——定理的符号约定为 $\ind T_f=-\mathrm{wind}(f)$（绕数逆时针为正）；用 $f_n(z)=z^n$ 核对：$n>0$ 时 $T_{z^n}$ 是单侧移位的 $n$ 次幂，核为 $0$、余核为 $n$ 维，指数为 $-n=-\mathrm{wind}(z^n)$。若把 $\sH$ 取在圆盘外部，或把 $T_f$ 定义为 $i\circ M_f\circ\pi$，符号会改变。
\end{enumerate}

\subsection*{§4 指数定理及其证明}
\begin{enumerate}[leftmargin=2.2em,itemsep=2pt,topsep=2pt]
  \item 笔误——源文件中 subsec:4.5 与 subsec:4.4 互换：文档顺序第 4 小节（应用）带 subsec:4.5，第 5 小节（K 理论）带 subsec:4.4。因此本节开头“第二个证明（\S\ref{subsec:4.4}）”编译后印作 4.5，而“\S\ref{subsec:4.5} 列出应用”印作 4.4。译文按源码文档顺序保留原 label，不重排、不改名；印刷号与交叉引用本身均正确，仅为 label 命名与所在小节不符。
  \item 术语——"announced in\textasciitilde{}\cite{AS1} in\textasciitilde{}1963" 指 Bull. Amer. Math. Soc. 69 (1963) 的公告；第一个证明由 Palais 编成 Annals of Math. Studies 57 (1965) 的 Seminar 文集 [Pa]（撰稿人含 Atiyah、Borel、Floyd、Seeley、Shih、Solovay），并非 Atiyah–Singer 本人执笔；第二个证明见 Ann. of Math. 87 (1968) 系列第一篇 [AS2]。核对年代时注意“宣布年”与“正式发表年”的差别。
  \item 术语——Dirac 原文 [D] 在 Minkowski 时空（号差 $(1,3)$）上工作，其方程是一阶的，其平方为波算子；本节把时间方向解析延拓 $t\mapsto i x^n$（Wick 旋转），得到欧氏空间上的椭圆算子。原文说“the analogue”，要点正在于这一转动把双曲方程变为椭圆方程，从而进入 Fredholm 指数理论的适用范围。
  \item 易错——(4.1)–(4.4) 的符号约定是配套的：取几何学家常用的 $\Delta=-\sum_i\partial_i^2$ 时，Clifford 关系必须是 (4.4) 的 $\gamma^i\gamma^j+\gamma^j\gamma^i=-2\delta^{ij}$（即 $\Cliff_{-n}$）；若改用 $\Delta=+\sum_i\partial_i^2$，右端要换成 $+2\delta^{ij}$（即脚注中的 $\Cnp$）。两处符号必须同时翻转，否则 $D^2=\Delta$ 不成立。
  \item 术语——$\Cnm=\Cliff_{-n}$、$\Cnp=\Cliff_{+n}$ 是原文的记号（见脚注与 §5.4，源码 label 为 subsec:5.6）；实 Clifford 代数 $\Cliff_{\pm n}$ 的类型以 $n\bmod 8$ 为周期（[ABS] 的周期表），这决定了旋量表示 $\SS^0,\SS^1$ 在 $n\bmod 8$ 不同时的实/复/四元数性质。不要把 $\Cnm$ 的下标读成“负维数”。
  \item 含糊——“the spin group is a double cover of $\SO_n$, and a subgroup of the units in $\Cnm$”一句跳过两点：(i) 对 $n\ge3$，$\Spin_n\to\SO_n$ 是万有覆盖（$\pi_1(\SO_n)=\ZZ/2$），$n=2$ 时 $\Spin_2\cong U(1)$；(ii) 更准确地说 $\Spin_n$ 位于 $\Cnm$ 偶部分 $\Cliff^0_{-n}$ 的可逆元群中，这正是旋量丛可由主 $\Spin_n$-丛相伴构造的原因。详见 [ABS] 与 [Ch]。
  \item 笔误——(4.5)（源码 \eqref{eq:40}）原文作 \textbackslash{}frac{\textbackslash{}sinh y\_i/2}{1/2}；按 TeX 的结合规则，$\sinh y_i/2$ 读作 $(\sinh y_i)/2$，再除以分母 $1/2$ 得 $\sinh y_i$，并不是旋量特征标之差。作者本意应为 $\sinh(y_i/2)/(1/2)=2\sinh(y_i/2)$，因为 $\SS^0-\SS^1$ 的特征标等于 $\prod_i(e^{y_i/2}-e^{-y_i/2})=\prod_i 2\sinh(y_i/2)$。译文照录原文，未改公式。
  \item 笔误——紧接 (4.5) 的正文说 $y_1,\dots ,y_n$ 是 $\Spin_n$ 极大环面特征标的一组基，但 $\Spin_n$（$n$ 偶）的极大环面秩为 $n/2$，(4.5) 的乘积也只到 $n/2$；疑应为 $y_1,\dots ,y_{n/2}$。可对照 (2.24) 的 $\Ahat$ 亏格：那里 $X^{4k}$ 有 $2k=n/2$ 个 Chern 根。
  \item 易错——(4.5) 并不是 $1/\Ahat$：由 (2.24)，$\Ahat=\prod_i (y_i/2)/\sinh(y_i/2)$，其倒数的分母应为 $y_i/2$ 而非 $1/2$。准确地说 (4.5) 等于 $\prod_i 2\sinh(y_i/2)$，即旋量特征标之差；不要把“upside down variant”理解为严格取倒数。
  \item 含糊——“the spin condition was known to be related to the integrality of the $\Ahat$-genus, as explained in\textasciitilde{}\S\ref{subsec:2.4}”：所指是 Borel–Hirzebruch [BH3]（早期 [BH2] 差一个 2 的幂）利用 Milnor [Mi1] 的配边结果证明：若 $w_2(X)=0$，则 $\Ahat(X)[X]\in\ZZ$；更一般地，若 $w_2$ 是某个整上同调类的模 2 约化，则相应的 Todd 数为整数。见 §2.4 的脚注与 [H3]。
  \item 术语——“spin structure is a double cover $\Spin(X)\to\SO(X)$ together with a compatible principal $\Spin_n$-bundle structure”是 spin 结构的几何定义：它等价于把主 $\SO_n$-丛提升为主 $\Spin_n$-丛；存在性障碍是 $w_2(X)=0$，而两个 spin 结构之差由 $H^1(X;\ZZ/2)$ 分类。不要把主丛 $\Spin(X)$ 与旋量丛 $\SS_X$ 混为一谈：前者是主 $\Spin_n$-丛，后者是相伴的复向量丛。
  \item 易错——(4.6) 的 $D_X=c\circ\nabla$ 只是手征 Dirac 算子 $\SS^0_X\to\SS^1_X$，它不（形式）自伴；（形式）自伴的是作用在 $\SS^0_X\oplus\SS^1_X$ 上、以 $D_X$ 与 $D_X^*$ 为两个非对角块的分块算子。其自伴性来自 $c$ 关于度量是斜伴随、$\nabla$ 保度量，以及 $X$ 紧闭（可用 Stokes 定理分部积分验证）。指数按 (3.12) 取 $\dim\ker D_X-\dim\coker D_X$，与 (4.7) 的左端对应。
  \item 含糊——“The Levi-Civita connection on $X$ induces a covariant derivative $\nabla$ on each spinor bundle”省略的实质是 spin 联络：把 Levi-Civita 联络的 $\SO_n$-联络形式经 $\Spin_n\to\SO_n$ 的微分（同构 $\mathfrak{spin}_n\cong\mathfrak{so}_n$）提升到旋量丛，得到协变导数 $\nabla$；Dirac 算子的局部公式是 $D=\sum_i c(e_i)\nabla_{e_i}$（$e_i$ 为局部正交标架）。标准处理见 [ABS] 与 [Ch]。
  \item 术语——“only determined uniquely up to tensoring by a line”：$\Spin_n$ 的两个不可约复旋量表示只在相差一个一维表示的意义下唯一。张量一个线表示不改变手征分解的结构，却改变旋量丛与 Dirac 算子（相当于用该线丛扭化）；§5.4（源码 label 为 subsec:5.6）的 Clifford 线性 Dirac 算子在 $\Cliff_n$-模层面消除了这一非典范性。
  \item 含糊——“solutions $\psi$ to $D\psi=0$”中的 $D$ 是 $\SS^0_X\oplus\SS^1_X$ 上的自伴 Dirac 算子，而非 (4.6) 的手征算子 $D_X$；因此“调和旋量”按手征性自然分成 $\sH^+$ 与 $\sH^-$，它们分别对应 $D_X$ 的核与余核（经由 $D_X^*$ 的核）。这是 (4.7) 左端取 $\dim\sH^+-\dim\sH^-$ 的原因。
  \item 易错——(4.7) 的正负号依赖于 $\SS^0,\SS^1$ 的编号约定：交换两者会使左端整体变号。与 (4.7) 配套的约定是 $\SS^0$ 取“正”手征、$\Ahat$ 取 (2.24) 的 $\prod(y_i/2)/\sinh(y_i/2)$；若采用 $\Delta=+\sum\partial_i^2$ 或相反的 $\gamma$ 符号约定，需重新核对符号。参见 [ABS] 对 $\SS^\pm$ 的约定。
  \item 笔误——原文说“the DPhil thesis [Hi] of Atiyah's student Nigel Hitchin”，但文献表 [Hi] 是 N. J. Hitchin, Harmonic spinors, Adv. Math. 14 (1974) 1–55 的论文，并非学位论文本身（该文源自其 1972 年牛津 DPhil 论文）。引用时可注明“论文正式发表版”。
  \item 术语——analytic index（解析指数）指 (3.12) 的 Fredholm 指数 $\dim\ker P-\dim\coker P$；topological index（拓扑指数）在这里先以上同调形式给出，即 $\Todd(TX\otimes\CC)\smallsmile\phi^{-1}\ch\sigma(P)$ 的示性数；§4.5 的 $K$ 理论形式用 $K$ 理论 Thom 同构定义 t-ind。阅读时注意作者的措辞层次：$\ind P$、上同调拓扑指数、t-ind 在定理中相等。
  \item 易错——(4.10) 中必须是复化切丛的 Todd 类 $\Todd(TX\otimes\CC)$，它与 $\Ahat$ 的关系是 $\Todd(TX\otimes\CC)=\Ahat(X)^2$（复化后 Chern 根成对为 $\pm y_i$，故 $\Todd=\prod_i[(y_i/2)/\sinh(y_i/2)]^2$）。注意不要与复流形的 Todd 类混淆：后者按 (2.25) 为 $\Todd(X)=e^{c_1(X)/2}\Ahat(X)$，二者相差一个因子。另注意 $\phi$ 是 Thom 同构、$\phi\inv$ 才是回到 $X$ 上同调的映射。
  \item 术语——注 4.11 说“$K$ 理论公式（\S\ref{subsec:4.4}）更自然”：此处的 $K$ 理论公式指 §4.5（源码 label 为 subsec:4.4）末尾的 K 理论定理（定理 4.18，a-ind = t-ind）；由于 label 互换，\ref{subsec:4.4} 印刷为 4.5，但指向的正是 $K$ 理论小节。
  \item 含糊——“the cyclic module $K^0(X)\sigma_0$ ... is a subgroup of finite index”省略了论证的关键：解析指数与拓扑指数都是 $\Ell(X)\to\ZZ$ 的群同态；若子群 $K^0(X)\sigma_0$ 在 $\Ell(X)$ 中指数有限，则只需在该子群上验证等式（对任意 $\sigma$，有 $m\neq0$ 使 $m\sigma\in K^0(X)\sigma_0$，同态性给出 $m(\ind\sigma-\text{top}\ \sigma)=0$）。而 $K^0(X)\sigma_0$ 中的元素正是扭化号差算子的象征，故问题化归到 $D_W$。
  \item 含糊——“Bott periodicity, as realized in $K$-theory”：指复 $K$ 理论的 Bott 周期性 $K^0(X)\cong K^2(X)$（[B1]、[B2]），它给出 $K$ 理论的 Thom 同构 $K^0(X)\cong K^0(T^*X,T^*X\setminus 0)$（在约化 $K$ 群中即 suspension 同构），从而在 spin 流形上 $\Ell(X)$ 由 Dirac 象征循环生成。实情形的周期 8 见 [ABS] 的 Clifford 模表。
  \item 含糊——“This is straightforward for the cohomological formula”省略了配边不变性的来源：右端是示性数，其配边不变性来自 $\Todd(TX\otimes\CC)$、$\Ahat$ 等示性类在配边下的自然性（在配边 $Y$ 上由 Stokes 定理化为边界上的积分）。左端解析指数的配边不变性则靠下述“在 $Y$ 上补一个算子”的论证；细节见 [AS3] 与 [AB3]。
  \item 易错——“local elliptic boundary condition”指 Atiyah–Bott 意义下的局部（Lopatinski–Shapiro）边界条件，而非 APS 的整体谱投影条件；局部条件使 $D_U$ 在 $Y$ 上形式自伴且边界条件自伴，故 $\ind_YD_U=0$；而 (i) 式 $\ind_YD_U=\ind_XD_W$ 来自内部算子与其边界算子的指数关系（粘合/分离论证）。参见 [AB3]、[AS3]。
  \item 术语——(4.14) 中的 $p(x,\xi)$ 是伪微分算子的全象征（total symbol），而 (3.11) 与 (4.15) 中的 $\sigma(P)$ 是主象征（principal/top order symbol）；原文在“symbol”一词上两义并用，阅读时须按上下文区分。全象征依赖于坐标与截断，主象征才是流形上整体定义的。
  \item 术语——“Schwartz kernel smooth away from the diagonal”：伪微分算子的核 $K(x,y)$ 在对角线外光滑，在对角线附近有 $|x-y|^{-n-m}$ 型奇性；微分算子的核是 $\delta(x-y)$ 及其导数，支集在对角线上，故“off the diagonal 为零”。这一性质保证伪微分算子把 $C^\infty$ 映到 $C^\infty$，并在 Sobolev 空间上有界。
  \item 易错——(4.15) 用极限定义主象征，要求 $p$ 满足 $p(x,\lambda\xi)/\lambda^m$ 当 $\lambda\to\infty$ 有极限；对一般的全象征这并非自动成立（例如对数项会破坏齐次性）。因此“限制到使极限存在的伪微分算子”不是技术性小条件，而是主象征有定义的前提；主象征仅保留最高阶的齐次部分，低阶项与有界项都被丢弃。
  \item 含糊——“the surjectivity and continuity of the principal symbol map ... are crucial”的实质是两条定理：(i) 每个椭圆象征类都可由某个伪微分算子实现（满射性，需要 $X$ 紧致）；(ii) 算子同伦与象征同伦相容（主象征映射的连续性/稳定性），从而解析指数只依赖象征的同伦类。这正是 §4.2 末尾所提两条性质的来源，见 [AS2] §1 与 [AB1]。
  \item 术语——(4.7) 的 $\Ahat(X)[X]$ 是 $\Ahat$ 亏格的示性数：$\Ahat(X)$ 按 (2.24) 取 $\prod_i(y_i/2)/\sinh(y_i/2)$，其中 $p_1$ 的系数是 $-1/24$（即 $\Ahat=1-p_1/24+\cdots$）。若采用把 $\Ahat$ 换成与 $L$ 亏格混用的约定（$L=p_1/3$，$\Ahat=1-L/8+\cdots$），公式会差一个因子；核对时以 (2.24) 的 $\prod_i(y_i/2)/\sinh(y_i/2)$ 为准，参见 [H1]。
  \item 易错——“the more flexible techniques of Atiyah-Singer prove it for arbitrary closed complex manifolds, a powerful generalization from projective algebraic manifolds”：希策布鲁赫–黎曼–罗赫定理 2.7（thm:1）原本对射影代数流形成立，而 AS 的椭圆算子方法把它推广到任意闭复流形（无需射影/代数）。这是“解析指数=拓扑指数”的直接应用：$\dbar+\dbar^*$ 的指数就是 $\chi(X,V)$。
  \item 含糊——“given appropriate inequalities between $\dim X$ and $\rank E^0=\rank E^1$”原文未写出具体不等式；其意为：当 $E^0,E^1$ 为平凡丛时，象征是球面丛到 $\GL_k(\CC)$（$k=\rank E^0$）的映射，在 $\dim X$ 与 $k$ 满足适当不等式时该象征类平凡，指数因而为零。最清楚的特例是单个函数（$k=1$）。精确陈述请回到 [AS1] 核对维数与秩的条件。
  \item 术语——pushforward $f_!$ 在格罗滕迪克的黎曼–罗赫定理（(2.13)）中是层的前推 $f_!:K(X)\to K(Y)$；在拓扑范畴中 Atiyah–Hirzebruch 用 $K$ 理论 Thom 同构定义拓扑推前，使其在 $X$ 为复流形、$f$ 为复映射时与代数 $f_!$ 一致（见 (2.27) 附近的讨论）。
  \item 含糊——(4.17) 的箭头只标了 $\ind$，而紧接着的正文定义了两个同态 a-ind 与 t-ind；可把 (4.17) 理解为二者共同的“定义域与值域”式，两个同态在此后分别命名。此处不是笔误，只是表述从简。下文定理 4.18 写 a-ind = t-ind，可证此读法。
  \item 术语——$K$ 理论的 Thom 同构给出 $\Ell(X)=K^0(T^*X,T^*X\setminus 0)\cong K^0(X)$（“符号映射”把相对 $K$ 群与底流形的 $K$ 群等同起来）；于是 t-ind 可写成 $\Ell(X)\xrightarrow{\cong}K^0(X)\to K^0(\pt)=\ZZ$ 的复合，其中第二步是映到一点的推前。这正是 §4.2 上同调公式的 $K$ 理论对应物，其存在性依赖 Bott 周期性（即 $K$ 群中的 suspension/Thom 同构）。
  \item 含糊——定理 4.18（a-ind = t-ind）标为“Atiyah-Singer\textasciitilde{}1967”，而本节上文说第二个证明于 1968 年发表于 [AS2]（Ann. of Math. 87）；1967 或指宣布/投稿/预印年份，不宜判为笔误。核对文献年份时以 [AS2] 为准。
  \item 含糊——“Atiyah-Singer uniquely characterize 'index homomorphisms' ... by a short list of axioms”：所谓短公理表包括正规化（在一点或平凡象征上的取值）、切除（excision）与乘性（multiplicativity）等；证明的核心是用 $K$ 理论的 Thom 同构与 Bott 周期性把任意象征类化归到这些基本情形。细节见 [AS2] §1 与 §2。
\end{enumerate}

\subsection*{§5 主题的各种变奏}
\begin{enumerate}[leftmargin=2.2em,itemsep=2pt,topsep=2pt]
  \item 笔误——式 \eqref{eq:51} 左端求和指标为 $x\in \Fix(f)$，被加项却写作 $\epsilon _X$（下标为大写 $X$），按式 \eqref{eq:50} 应为 $\epsilon _x$。此式照录原文，引用时请按 $\epsilon _x$ 理解。
  \item 术语——"带符号计数（count with signs）"指按 $\sign\det(1-df_x)$ 加权；这里的 $f$ 必须是横截的（脚注假设 $1$ 不是 $df_x$ 的本征值，即 $\det(1-df_x)\neq0$），否则须借助等变局部化或热核方法处理非横截不动点。
  \item 易错——式 \eqref{eq:52} 的分子是\emph{分级}迹 $\Tr(f^0)-\Tr(f^1)$，分母取绝对值 $|\det(1-df_x)|$；符号信息全在分子的分级迹里。核验：对 de Rham 复形 $E^q=\Omega^q$，有 $\sum_q(-1)^q\tr\Lambda^q(df_x^*)=\det(1-df_x^*)=\det(1-df_x)$，故 $\nu_x=\sign\det(1-df_x)=\epsilon_x$，与式 \eqref{eq:50}–\eqref{eq:51} 一致。若误把分母写成不带绝对值的 $\det(1-df_x)$，则比值为 $1$ 而非 $\epsilon_x$，即逐点乘上 $\epsilon_x$（它可在同一连通分支内变号），并非仅在分支上整体差一个符号。
  \item 含糊——作者以"不评论证明"带过，其实质是 Atiyah–Bott 在 [AB1]（Ann. of Math. 86 (1967)）中用伪微分算子把 $P$ 局部化到不动点附近，并对分级迹作迹公式；所需技术是伪微分算子代数上的迹与象征演算。完整的证明见 [AB1]；椭圆\emph{复形}（而非单个算子）情形还需对每个 $f^i$ 逐项跟踪分级。
  \item 术语——"等变指数"指指数作为 $R(G)$ 中的虚拟表示；$g\in G$ 给出虚拟特征标值 $\ind_g(P)\in\CC$，即"特征标值指标"，且 $g=e$ 时 $\ind_e(P)=\ind P\in\ZZ$。注意它与"指标丛"（families 情形，取值于 $K^{\bullet }(S)$）是两回事：前者底空间是一点、对称群为 $G$，后者底空间是族参数空间 $S$。
  \item 含糊——"等变 $K$ 理论中的局部化定理"即 Atiyah–Segal 的局部化定理（[Seg1]，亦见 [ASeg]）：等变 $K$ 类由它在不动点集上的限制（或按法向权展开的贡献）决定。作者未写出公式；具体地，$\chi(g)=\ind_g(P)$ 可写成 $g$ 的不动点集 $\Fix(g)$ 上的求和，各项含 $\det(1-dg)$ 型的因子，这正是式 \eqref{eq:53} 的等变来源。
  \item 易错——把定理\textasciitilde{}\ref{thm:9} 用于外尔特征标公式时必须取 $T$ 的\emph{一般}元素（generic element）：一般性保证不动点孤立且 $1$ 不是 $dt_x$ 的本征值，即定理\textasciitilde{}\ref{thm:9} 的横截条件成立。若 $t$ 落在环面的真子群中，不动点集非离散，须改用等变局部化公式。
  \item 术语——"$h$-配边（$h$-cobordism）"指嵌入 $W$ 满足两侧包含映射都是同伦等价；"透镜空间（lens space）"是 $S^3$ 对有限循环群自由作用的商。此处结论的强度在于把拓扑的 $h$-配边关系升级为黎曼几何的等距，所用工具是号差算子的 Lefschetz 公式，见 [AB2]。
  \item 含糊——从 Lefschetz 公式到外尔特征标公式的推导被一句"leads to"带过。实质：旗流形 $G/T$ 上全纯线丛 $\mathscr{L}$ 的上同调实现 $G$ 的不可约表示（Borel–Weil），而 $t\in T$ 的 Lefschetz 数 $\sum_q(-1)^q\tr(t^*|H^q(G/T;\mathscr{L}))$ 正是该表示的特征标在 $t$ 处的值；不动点即 Weyl 群元素对应的点。详见 [AB2] 与 [AS3]。
  \item 易错——"$H^q(X;\sO_X)=0$ 对 $q>0$" 的全纯映射必有不动点：这是用 Lefschetz 数等于 $1$（只有 $H^0$ 贡献）而非用拓扑不动点性质；请核对 $\sO_X$ 是结构层、$H^0(X;\sO_X)=\CC$（$X$ 连通紧致）。
  \item 术语——"局部边界条件"指条件只依赖未知函数在边界点处的有限阶 jet：Dirichlet 条件只涉及 $0$ 阶 jet（$u(x)$），Neumann 条件涉及法向导数（$1$ 阶 jet）。与之相对的是"整体边界条件"（如 §7.2 的 APS 条件），后者涉及边界上算子的谱投影，因而不是局部的。
  \item 含糊——"Lopatinski 判别法"即今称的 Lopatinski–Shapiro 条件（亦称覆盖条件）：要求边界象征在法向的零空间维数恒定，且边界象征的各因子在法向"覆盖"靶向量丛。作者只给出 [Lo]（1953）与 [Ho1]（其 §20.1 疑有误，见下条）的指引，未写条件本身；核对时请注意不同文献对"椭圆边界条件"的定义在光滑性要求上略有差异。
  \item 笔误——原文 \cite[\S20.1]{Ho1} 疑误：[Ho1] 是 Hörmander 1965 年 CPAM 短文（pp. 501–517），不可能有 §20.1；疑应为 Hörmander 专著 [Ho3]（*The analysis of linear partial differential operators III*），请核对文献与节号。
  \item 易错——圆盘上 $\dbar$ 的核无限维，根源是 Cauchy–Riemann 方程作为边值问题是\emph{不适定}的（Lopatinski 条件失败），而不是 $\dbar$ 不椭圆。$\dbar$ 的主象征在余切丛的非零部分是可逆的（复线性），所以它是椭圆算子；毛病完全出在边界。
  \item 含糊——障碍的精确表述：作者说"$\sigma$ 到 $K^0(T^*X|_{\partial X})$ 的限制"，但该限制所在的群、以及"提升到适当的相对 $K$ 理论群"是哪个群，原文未写出。按 [AB3]（Bombay 1964 报告）与 [A2]（Seminar 附录 I）的叙述，障碍 $=\sigma$ 的边界限制类 $\in K^0(T^*X|_{\partial X})$；该障碍消失时可提升到相对群 $K^0(T^*X,T^*X|_{\partial X})$，其消失当且仅当存在局部椭圆边界条件。审校时请核对 [AB3] 的定义。
  \item 含糊——"Atiyah–Bott–Singer" 并无独立的三人论文：上文引的是 [AB3]（Atiyah–Bott，Bombay 报告）与 [A2]（Atiyah，Seminar 附录 I），可理解为二者的合并称呼；经典的 K 理论证明见 [AS2]。请勿据三人合名另找文献。
  \item 术语——此处 "Bott 周期性定理" 指复 $K$ 理论的 $2$ 周期性。Atiyah–Bott 用带边流形的相对 $K$ 群给出的新证明把周期性化归为边界椭圆算子的指数定理，见 [AB4]（Acta Math. 112 (1964)）。
  \item 术语——本小节只处理\emph{局部}椭圆边界条件：此时拓扑指数可无修正项地计算解析指数。若改用\emph{整体}边界条件（如 §7.2 的 APS 条件），指标公式须加上边界修正项（$\eta$ 不变量，其几何原型即号差亏损 signature defect），见 §7.3 与 [APS1]–[APS3]；两种边界问题不可混为一谈。
  \item 易错——反自伴 Fredholm 算子的整数指数必为零，因为 $\ker P^*=\ker P$ 且 $\coker P\cong\ker P^*$，正负指标相消；因此这里只能取模 $2$ 指标。这不是笔误。另注意：$D$（周期）与 $D'$（反周期）的核分别是 $1$ 维与 $0$ 维，故其模 $2$ 指标分别为 $1$ 与 $0$，恰好区分 $S^1$ 的两个 spin 结构。
  \item 术语——$\cir$ 上恰有两个 spin 结构：周期（Lie 群）结构与反周期（扭化、bounding）结构；相应的旋量丛截面可等同于周期或反周期函数。见 \S\ref{subsec:4.1} 与 [LM] I.3。学者常说的"两个 spin 结构"与 $\zt$ 指标的对应，正是 $KO^{-1}(\pt)\cong\zt$ 的几何体现。
  \item 术语——大写的 "Real $K$ 理论"（[A3]，*K-theory and reality*, 1966）指 Atiyah 的 KR 理论：带对合 $\tau$ 的空间上的 $K$ 理论，其元素是"实"丛在对合下的复共轭配对；不要与实 $K$ 理论 $KO$ 混为一谈。本节的实指标定理最终用 $KO$ 表述（见式 \eqref{eq:81}），KR 只是导出它的框架；二者的关系见 [A3]。
  \item 易错——"四维实 Dirac 算子的自然整数不变量是整数指数的一半"：复指标为 $\ind_\CC D=\Ahat=-\sigma/8$；实 Dirac 算子的 $KO^{-4}\cong\ZZ$ 指标与此相差因子（源自 $\Cliff_{-4}\cong\HH(2)$ 的模维数），故其自然整数不变量为 $\Ahat/2$。于是"该不变量为整数"等价于 $\Ahat$ 为偶数，即 $\sigma/8$ 为偶数，即 $16\mid\sigma$，这正是 Rokhlin 定理。核对 [AS5] 与 [LM] II.4 的归一化，别把因子 $2$ 弄丢。
  \item 含糊——作者自陈"对模 $2$ 指标不知道其他证明"，这凸显了 $KO$ 理论方法在 $\zt$ 指标上的不可替代性；热方程/$\eta$ 不变量方法（§6、§7）只直接给出整指标（或 $\eta$ 的实值修正），模 $2$ 指标需额外的拓扑输入。引用时请勿把这句话误读为"模 $2$ 指标没有证明"。
  \item 笔误——本小节在文档中排第 4 个（按编号应为 5.4），却带 label subsec:5.6；随后"族指数定理"排第 5 个（应为 5.5）而带 label subsec:5.4，"覆盖与 von Neumann 代数"排第 6 个（应为 5.6）而带 label subsec:5.5：三个 label 名与位置错位。LaTeX 按位置编号，故交叉引用仍指向正确小节、印出的编号也正确，问题仅在 label 名与位置不符；为不破坏全篇引用，此处照录源码 label 未改。
  \item 易错——式 \eqref{eq:79} 的 $\pm$ 约定决定了 $\Cliff_{\pm n}$ 的符号：$\Cliff_{+n}$ 中 $\gamma_i^2=+1$，$\Cliff_{-n}$ 中 $\gamma_i^2=-1$。Bott 周期表（$n\bmod 8$）与 $KO^{-n}$ 的取值都依赖这一约定；若把两个符号对调，$\Cliff_{\pm n}$ 的矩阵代数类型会移位，周期表随之整体平移。核对 [ABS]（*Clifford modules*, Topology 3 (1964)）与 [LM] I.4 的约定。
  \item 术语——"混合构造（mixing construction）"即纤维丛的相伴丛构造 $P\times_G V$：用主 $\Spin_n$ 丛 $\Spin(X)$ 把表示 $\Cnp$ 扭化为向量丛。它不是新的几何操作，只是把表示论数据升级为丛的标准手段。
  \item 易错——"$\Cnm$ 是 $\Cnp$ 的反代数"这一步用的是分级反代数同构 $\Cliff_{-n}\cong(\Cliff_{+n})^{\text{graded op}}$（Koszul 符号反乘法 $a\cdot b=(-1)^{|a||b|}ba$；复化后也可用 $\gamma ^i\mapsto i\gamma ^i$ 实现，参见 [LM] I.1），因此右 $\Cnp$ 模自动成为左 $\Cnm$ 模。这正是"Clifford 线性"的含义：算子与左 $\Cnm$ 作用分级交换，而单个 $\Spin_n$ 表示被整套 Clifford 生成元的作用取代。
  \item 术语——第 2 条"旋量表示只确定到与线丛张量积"（原文只在 §4.1 提过）：因为 $\Spin_n$ 的中心（$n$ 偶时为 $\zt$）与复化后矩阵代数的自同构允许把 $\SS^{\pm}$ 与一维表示作张量；故不能单靠表示论定义典范旋量丛，须额外选 spin 结构。$\Cliff$ 模构造的优点是：$\Cnp$ 作为 $\Spin_n$ 模虽是可约的，但带有典范的 Clifford 生成元作用，因此"典范"到无需选择。
  \item 术语——"Clifford 线性"意为算子与左 $\Cnm$ 模结构相容（分级交换），故核与余核都是 $\Cnm$ 模；因此指标落在 $\Cliff$ 模的 Grothendieck 群中，这由 [ABS] 等同于 $KO^{-n}(\pt)$。这一"以模代替表示"的技巧是 Atiyah–Singer III 之后的标准语言，见 [LM] II.7。
  \item 术语——"解析前推"与"拓扑前推"：$f_!$ 指沿 $f$ 的 umkehr（Gysin）映射；实 $K$ 理论中 $f\:X\to\pt$ 的定向由一个 spin 结构给出（$M\Spin$ 定向），见式 \eqref{eq:19} 的复情形类比。不要把这里的 $f_!$ 与前推 $f_*$（同调/上同调）混淆。
  \item 易错——式 \eqref{eq:81} 的周期表请按 $n\bmod 8$ 核对：$KO^{-0}=\ZZ$，$KO^{-1}=\zt$，$KO^{-2}=\zt$，$KO^{-3}=0$，$KO^{-4}=\ZZ$，$KO^{-5}=0$，$KO^{-6}=0$，$KO^{-7}=0$。表中 $n\equiv 1,2$ 给 $\zt$ 而 $n\equiv 0,4$ 给 $\ZZ$，这对应 $\pi _{n-1}O_\infty$ 的取值（$\pi_0O=\zt$，$\pi _1O=\zt$，$\pi _3O=\ZZ$，$\pi _7O=\ZZ$）。注意表只对 $n\ge1$ 陈述，$n=0$ 需另行补上 $KO^0(\pt)\cong\ZZ$。
  \item 易错——"$n=1$ 的模 $2$ 指标是特例"：$KO^{-1}(\pt)\cong\zt$，由 $\dim\ker D\bmod2$ 给出；对 $S^1$ 的两个 spin 结构分别给 $1$ 与 $0$（见 \S\ref{subsec:5.3}），正与表中 $n\equiv1$ 的 $\zt$ 相符。
  \item 术语——此处的"谱（spectrum）"是同伦论意义下的谱（$\Omega$-谱），与算子的谱毫无关系。同一节里两种"谱"并存是易混点：$\Fn(\sH_n)$ 构成的谱代表 $KO$ 或 $K$ 的 $\Omega$-谱，而"算子的谱"指本征值集。
  \item 易错——$\Fn(\sH_n)$ 的定义要求算子\emph{奇}且反自伴；去掉"奇"或去掉"反自伴"都会得到不同的分类空间（例如自伴情形对应 $KR$ 的其他分量）。这里 $\Cn$ 作用与"odd"条件共同保证 $\Fred$ 恰是 $KO^{-n}$（或 $K^{-n}$）的分类空间；详见 [AS6] 定理 A。
  \item 术语——"逆紧（proper）纤维丛"指纤维紧致、$f$ 逆紧；族指数定理要求纤维紧致以保证逐纤维的 Fredholm 性，而底空间 $S$ 可以非紧（此时 $K^{\bullet }(S)$ 用紧支撑或 $K$ 理论的一般定义）。不要把 proper 与"正常"或"真"混淆。
  \item 术语——"族指数（family index）"与"指标丛（index bundle）"：族指标是 $K^{\bullet }(S)$ 中的一个类；当 $\ker$、$\coker$ 的维数沿 $S$ 跳跃时它一般不是真实向量丛，而是以 $\ker-\coker$ 为纤维的\emph{虚拟}丛（K 理论类）。只有在维数局部常定的开集上才能谈论具体的指标丛；这正是必须用 $K$ 理论而非普通丛来陈述族指数定理的原因。
  \item 易错——"Fredholm 算子构成 $K$ 理论的分类空间"对复 $K$ 理论给 $\Fred\simeq \ZZ\times BU$（$K^0$），带分级的奇算子给 $K^{-1}$；因此族指标取值于 $K^{\bullet }(S)$ 而非单个 $K^0(S)$。实情形对应 $KO$ 与 8 周期，见 [AS6] 与 [J]。
  \item 易错——纤维取\emph{奇}维是关键的：复 Bott 周期性给出 $K^{-n}(S)\cong K^{-1}(S)$（$n$ 奇），而 $K^{-1}(S)=[S,U(\infty)]$ 经 $\det:U(\infty)\to U(1)$ 映到 $H^1(S;\ZZ)\cong[S,S^1]$，这才有"行列式映射"。若 $n$ 为偶数，最低阶分量将是 $K^0$ 的秩（映到 $H^0$），而不是行列式。
  \item 含糊——"相对黎曼 spin 结构"的定义放在脚注里，但族 Dirac 算子的自伴/反自伴性与 Fredholm 性需要"相对黎曼结构"（内积 + 水平分布）才能逐纤维定义；作者未给出细节。可查 [AS4]（*The index of elliptic operators IV*）与 [APS3] §3（自伴表述）。
  \item 易错——"$H^1(S;\ZZ)$ 无挠"才使"由在 $\cir$ 上的周期决定类"成立：一般地 $H^1(S;\ZZ)\cong\Hom(H_1(S),\ZZ)$ 确实无挠，故周期（对 $S^1$ 的配对）足以确定类；但对 $H^2$ 一般有挠部分，不能用 de Rham/周期方法决定。基变换（base change）即沿 $\varphi $ 拉回纤维丛与向量丛。
  \item 术语——"最低阶片"指 $K$ 理论类在 Atiyah–Hirzebruch 谱序列中的最低非零分量；行列式映射 $K^{-1}(S)\to H^1(S;\ZZ)$ 由 $\det:U(\infty)\to U(1)$ 诱导，把 $K$ 类送到 $U(1)$-类（即 $S\to S^1$ 的绕数/相位；等价于 $S\times S^1$ 上行列式线丛的第一陈类），勿与 $K^0\to H^2$ 的第一陈类混同。注意它\emph{不是}单射：核包含挠信息与高阶信息。
  \item 术语——$\RZ$ 在此是 Eilenberg–MacLane 空间 $K(\ZZ,1)\simeq S^1$，故 $H^1(S;\ZZ)\cong[S,\RZ]$；把同伦类"提升为具体映射"指 $\eta$ 不变量给出规范代表元（而不仅是同伦类），这正是"几何不变量"的含义。
  \item 术语——"模空间，或栈"：模空间 $\sM=F\inv (c)/\sG$ 一般是奇异的（当 $\sG$ 作用有稳定化子时），"栈"指保留稳定化子信息的精细版本；此时"虚拟切丛"由族的指数给出（形变复形的上同调），在基点处等于 $\ker dF_b-\coker dF_b$。
  \item 易错——原文 "its linearizations $dF_b$ at $b\in F\inv (c)$ fit together to a family ... parametrized by $\sB$"：严格地说族应由 $F\inv (c)/\sG$ 参数化，且只有当 $F$ 在 $\sG$-轨道横截方向上 Fredholm 时指标才有定义；作者用 $\sB$ 参数化是为简化。审校时可保留原文措辞，但须知这里隐含了"模 $\sG$ 的 Fredholm"假设。
  \item 含糊——"\eqref{eq:12} 的拓扑类比"指族情形的 Riemann–Roch 型公式 $\ch(f_!V)=f_*(\ch V\cdot \Todd(T(X/S)))$ 之类；作者未写出公式。热方程方法能算出 $\ch$ 像是因为热核的局部展开给出示性形式的积分，而要恢复 $K$ 理论中的挠信息（如指标丛本身）则超出热核方法的能力，见 [Bi2]。
  \item 术语——"II 型 von Neumann 代数"（type II von Neumann algebra）的关键特征是存在实值的维数函数（Murray–von Neumann 维数），可以取连续值；这正是把指数推广到实数的代数基础。与 $\ssA$ 上的正规迹 $\tau$ 配合，投影的 $\Gamma$-维数 $\dim_\Gamma=\tau(p)$ 就是 $L^2$ 意义下的"维数"。
  \item 易错——von Neumann 迹的归一化是本节最容易出错处：$\Gamma$ 无限时 $\ssA$ 为 II$_\infty$ 型，$\tau(1)=\infty$，不存在 $\tau(1)=1$ 的迹；应取 $\tau$ 使 $\dim_\Gamma(L^2(\Gamma))=1$（即把基本域维数归一化为 $1$），此时 $\dim_\Gamma(\ker\tD)=\tau(p)$ 是"每个基本域的维数"。若不归一化，定理会差一个体积因子，$\ind_\Gamma\tD=\ind D$ 不再成立。核对 [A5] 的约定。
  \item 术语——"$L^2$ 指数"与"$\Gamma$ 维数"：对投影 $p$（到 $\ker\tD$ 的正交投影），$\dim_\Gamma(\ker\tD)=\tau(p)\in[0,\infty)$ 可为非整数；这就是 $L^2$-index。它与普通指数 $\ind D\in\ZZ$ 的区别只在归一化与适用范围（非紧覆盖上），定理断言两者数值相等。
  \item 术语——"无分叉 Galois 覆盖"：拓扑覆盖自动无分叉，"Galois"即正规覆盖，其 deck 变换群 $\Gamma=\pi_1(X)/\pi_1(\tX)$ 自由作用，"unramified"是为区别于代数几何中带分歧的覆盖。定理所需的是 $\Gamma$ 自由作用且 $X=\tX/\Gamma$ 紧致。
  \item 易错——定理\textasciitilde{}\ref{thm:12} 断言 $\ind_\Gamma\tD=\ind D$：左端按基本域计数、右端是 $X$ 上的整数指数。等号成立依赖 $\Gamma$ 自由作用与 $X$ 紧致；deck 群总是自由作用，挠元同样可以自由作用（如 $\ZZ_2$ 对径作用于 $S^n$），故"$\Gamma$ 有挠元"并不等于"作用不自由"。若作用确有不动点，须改用轨道叶状空间上的指标理论（Connes 的叶状结构理论），见 [A5] 及后续文献。
\end{enumerate}

\subsection*{§6 热方程证明}
\begin{enumerate}[leftmargin=2.2em,itemsep=2pt,topsep=2pt]
  \item 术语——本节的“局部”指 $t\to0$ 时在每一点 $x$ 处成立的恒等式（局部指标密度），与“整体拓扑”（示性类在基本类 $[X]$ 上的配对）相对；下一节讨论的“几何不变量”指 $\eta$ 不变量、号差亏损等次级不变量，勿与本节的热系数混淆。
  \item 术语——\emph{热算子}是算子 $H_t=e^{-t\Delta}$（\eqref{eq:56}），\emph{热核}是它的积分核 $h_t(x,y)$（\eqref{eq:59}）。二者不可互换：只有热核才有逐点的 $t\to0$ 渐近展开，而热算子本身在 $t\to0$ 时只是趋于恒等算子。
  \item 易错——\eqref{eq:57} 的 $\Tr\Delta^{-s}$ 在无穷维时是形式写法：先取谱截断、再解析延拓（$\Delta$ 须有离散谱且特征值按 Weyl 律增长，闭流形上的椭圆算子满足）。零特征值的约定尤须注意：若把 $\ker\Delta$ 上的 $0^{-s}$ 排除（只对 $\lambda>0$ 求和），$\zeta$ 就少一项 $\dim\ker\Delta$；作者在正文（$s=0$ 处的取值 $=\dim\sH^0-\dim\sH^1$）用的是计入零特征值、取 $0^0=1$ 的约定。两种归一化相差 $\dim\ker\Delta$，引用时须核对。
  \item 易错——\eqref{eq:58} 的被积式是 $\frac{dt}{t}\,t^s=t^{s-1}dt$（不是 $t^{-s}$），收敛域也不是“仅 $\Re s>0$”：小 $t$ 端由 \eqref{eq:60} 的 $t^{-n/2}$ 控制，要求 $\Re s>n/2$（否则 $t\to0$ 处发散）；大 $t$ 端若 $\ker\Delta\neq0$，则 $\Tr e^{-t\Delta}\to\dim\ker\Delta$ 而不趋于 $0$，须先减去零特征值投影 $P_0$（改用 $\Tr(e^{-t\Delta}-P_0)$）才收敛。前置的 $1/\Gamma(s)$ 消去 $s=0,-1,\dots$ 处的极点；$\zD$ 的解析延拓即由此得到。
  \item 术语——$\zeta$ 正规化（zeta regularization）指用 $\zD(s)$ 的亚纯延拓给发散量赋值；例如 $\det\Delta:=e^{-\zeta_\Delta'(0)}$（$\zeta$ 正规化行列式）。这是取有限部分的\emph{定义}，不是普通行列式；本节只用到 $\zD$ 的亚纯延拓，尚未用到行列式。
  \item 笔误——原文作“the heat operator\textasciitilde{}$H_t$ in\textasciitilde{}\eqref{eq:56} is exists by basic elliptic theory”，多一“is”；译文按 exists 理解。
  \item 含糊——“basic elliptic theory”的实质是：闭流形上非负自伴椭圆算子的谱定理（本质自伴、离散谱）与半群理论给出 $e^{-t\Delta}$ 作为强连续压缩半群，且其核是热方程初值问题的基本解。完备叙述见 \cite{Se4} 与 Berline–Getzler–Vergne《Heat Kernels and Dirac Operators》第 2 章（本书无该文献键）。
  \item 术语——\emph{磨光算子}（smoothing operator）指核 $h_t$ 光滑、从而把分布映为光滑函数；它必是迹类紧算子，但反之不然，“紧”不蕴含“磨光”。
  \item 易错——\eqref{eq:59} 约定 $h_t(x,y)=(e^{-t\Delta}\delta_y)(x)$，配合热方程 $u_t=-\Delta u$（传播子 $e^{-t\Delta}$，$\Delta$ 非负）；若改用 $e^{+t\Delta}$，则 $t$ 换成 $-t$，渐近展开的方向随之反转。又 $h_t(x,y)$ 关于 $x,y$ 对称，但须以同一黎曼测度 $|dx|$、$|dy|$ 为前提。
  \item 易错——两端极限的拓扑不同：$t\to0$ 时 $e^{-t\Delta}\to I$ 仅在\emph{强}算子拓扑下成立（当 $\spec\Delta$ 无界时 $\|I-e^{-t\Delta}\|=1$，范数收敛失败），而 $t\to\infty$ 时在\emph{一致}（算子范数）拓扑下收敛到 $\ker\Delta$ 上的投影，速率由第一个正特征值控制。原文特意区分这两种拓扑，不可混用。
  \item 易错——原文说 $x\neq y$ 时 $h_t(x,y)$ “指数式地收敛到零”；精确速率由 Varadhan 公式给出（本书无该文献键）：$\log h_t(x,y)\sim -d(x,y)^2/4t$，$d$ 为黎曼距离。故“指数式”应读作 $e^{-d^2/4t}$ 这一尺度，而非任意的 $e^{-c/t}$。
  \item 笔误——原文作“Minakshisundarum”，而 \cite{MP} 条目作“Minakshisundaram”（S. Minakshisundaram），少一字母“a”；又原文称该文为 1948 年的论文，而 \cite{MP} 为 Canad. J. Math. \textbf{1} (1949), 242–256，年份不符。
  \item 笔误——\eqref{eq:60} 的求和项写作 $A_k(x)t^i$，指数 $t^i$ 中的 $i$ 未定义（且与 \eqref{eq:72} 的下标 $i$ 冲突），按标准热核展开应为 $t^k$：$h_t(x,x)\sim t^{-n/2}\sum_k A_k(x)t^k$。照录原文，特此指出。
  \item 易错——\eqref{eq:60} 的归一化：$t^{-n/2}$ 与 $A_0=(4\pi)^{-n/2}$（见 §\ref{subsec:6.2} 首段）配套，等价于 $h_t(x,x)\sim(4\pi t)^{-n/2}\sum_k a_k(x)t^k$、$a_0=1$；若采用 $(2\pi)^{-n/2}$ 等别的约定，全部 $A_k$ 会差一个常数因子。核对 \cite{MP}、\cite{MS} 的归一化。
  \item 含糊——“渐近展开与亚纯延拓等价”正是梅林变换的实质：把 $\int_0^\infty t^{s-1}\bigl(\Tr e^{-t\Delta}-\text{有限项}\bigr)dt$ 逐项积分，$t\in(0,1)$ 上的 $t^{s-1}t^{m}$ 项给出 $\frac{1}{s+m}$ 型极点，$t\in(1,\infty)$ 上给出整函数；于是 \eqref{eq:60} 等价于 $\zD$ 在全 $s$ 平面的亚纯延拓。极点的位置与留数（即 Seeley–DeWitt 系数）的严格证明见 \cite{MP,Se4}。
  \item 易错——\eqref{eq:61} 用到 $\ker P^*P=\ker P$ 与 $\ker PP^*\cong\coker P$；这需要 $P$ 为 Fredholm 算子（椭圆性保证闭值域与有限维核、余核）。若不假定 $P^*$ 为形式伴随，右端第二项应直接写成 $\dim\coker P$。
  \item 易错——\eqref{eq:62} 只在 $\lambda>0$ 时成立；$\lambda=0$ 处 $P$ 把 $\ker P^*P=\ker P$ 映入 $0$，指数的全部信息恰在 $\lambda=0$ 这一项。这正是 McKean–Singer 把“核的维数差”改写为“谱上的和之差”的出发点。
  \item 易错——\eqref{eq:63} 要求 $\chi(0)=1$ 以保住 $\lambda=0$ 的项，并要求 $\chi$ 衰减足够快以使两个和收敛；取 $\chi(\lambda)=\lambda^{-s}$ 时需 $\Re s\gg0$，才能先对谱截断求和再取极限（即交换求和与取迹）。
  \item 术语——$\zD$ 在 \eqref{eq:57} 中已含迹（$\zD=\Tr\Delta^{-s}$），而 \eqref{eq:64} 又对 $\zeta_{P^*P}$ 取迹，可见此处的 $\zeta_{P^*P}(s)$ 指算子 $(P^*P)^{-s}$（复幂），记号前后并不完全一致；勿据此认为原文重复取迹。
  \item 含糊——“explicit integral formulas\textasciitilde{}\cite{Se4} in terms of the symbol”指 Seeley 的复幂构造：$(P^*P)^{-s}$ 的核由象征 $\sigma(P^*P)^{-s}$ 的振荡积分给出，故 $\zeta(s)$ 是象征的积分。但正因如此，这些公式含象征的高阶导数，与 \eqref{eq:44}（Chern–Weil 形式只用曲率的低阶导数）形成作者所说的“mismatch”。
  \item 易错——\eqref{eq:65} 两端与 $t\in(0,\infty)$ 无关：$t>0$ 时 $\Tr e^{-tP^*P}$ 与 $\Tr e^{-tPP^*}$ 均可微（热核光滑、迹类），求导后由迹的循环性得 $\Tr\bigl(P^*P e^{-tP^*P}\bigr)=\Tr\bigl(PP^*e^{-tPP^*}\bigr)$，故导数为零。不能以直接取 $t\to0$ 或 $t\to\infty$ 的极限代替这一论证：两端的收敛性都须单独证明。
  \item 笔误——原文作“evaluation as $t\to\infty$ reproduces\textasciitilde{}\eqref{eq:65}”，但 $t\to\infty$ 时 $e^{-tP^*P}$ 与 $e^{-tPP^*}$ 分别收敛到 $\ker P^*P$、$\ker PP^*$ 上的投影，得到的是 \eqref{eq:61}（即 $\ind P=\dim\ker P^*P-\dim\ker PP^*$），而不是 \eqref{eq:65} 本身。疑应为 \eqref{eq:61}。
  \item 易错——\eqref{eq:66} 中的 $A^i_{n/2}(x)$ 要求 $n$ 为偶数（热系数 $A_k$ 只对非负整数 $k$ 有定义）；对奇数维 $n$ 此式需另行解释，原文未加声明。
  \item 术语——此处“热系数”即 Seeley–DeWitt 系数；\eqref{eq:66} 的被积式 $\tr\bigl[A^0_{n/2}-A^1_{n/2}\bigr]$ 是“局部指标密度”（local index density），积分后给出指数。注意区分：局部指数定理断言该密度存在并等于某个示性形式；“指标密度”只是被积式本身。
  \item 易错——\eqref{eq:67} 的交错和 $\sum_q(-1)^q\tr h_t\q$ 就是超迹 $\Trs=\Tr(\gamma\,\cdot)$：在 $\Omega^{\bullet}$ 上取 $\gamma=(-1)^{\deg}$（形式次数分级），即 de Rham 算子的手征性算子；若采用 $\Omega^{\pm}$（偶/奇形式）分级的 $\gamma=\pm1$ 约定，分项写法不同而交错和相同。若与 §2 的手征约定符号相反，则 $\Trs$ 整体变号，务必核对。
  \item 术语——“Gauss–Bonnet–Chern integrand”即欧拉形式 $\Euler(TX)$ 的被积式：$\chi(X)=\int_X \Euler(TX)$；原文以连字符写作“Gauss-Bonnet-Chern”，规范写法为 Gauss–Bonnet–Chern。对奇数维 $n$，McKean–Singer 证明该极限为零。
  \item 含糊——Patodi 的“virtuoso direct computation”指逐项算出 Seeley–DeWitt 系数 $A_k\q$ 并直接验证 $k<n/2$ 时交错和为零，本质上是曲率张量多项式的组合恒等式硬算（极为繁重）；见 \cite{P1} 以及 \cite{ABP} 的相应部分。
  \item 含糊——“由标准的拓扑论证，这蕴含 Atiyah–Singer 指数定理”的实质是配边论证（即本文 §\ref{subsec:4.2} 的“第一证明”）：扭化号差算子情形的指数定理经配边不变性与乘法性，可化归为一般椭圆算子的情形。细节见 \cite{ABP}；配边（第一）证明的原始文本见 \cite{Pa}（据本文 §\ref{sec:4} 开头的说明，Atiyah–Singer 的第一个证明未由他们自己写出，而是收入 Palais 编的 Seminar 卷；1968 年的 \cite{AS2} 是第二个、K 理论证明）。
  \item 术语——原文把 functor 与 sheaf 并列（“the functors (sheaves)”）：$\Met$ 与 $\Omega^q$ 都是 $\Mn\op\to\Set$ 的函子，且是层；范畴 $\Mn$ 的态射取\emph{局部微分同胚}，正是为了使度量的拉回只在局部有定义，从而得到层而非仅为预层。
  \item 易错——\eqref{eq:71} 的自变量是 $\lambda^2g$（而非 $\lambda g$）；常值缩放 $g\mapsto\lambda^2g$ 不改变 Levi-Civita 联络（Christoffel 符号不变），故曲率 2-形式 $\Omega$ 与由它构造的 Pontryagin/Chern–Weil 形式都不变，只有标量曲率按 $\lambda^{-2}$ 缩放。于是 Pontryagin 形式的任何多项式权均为 $0$；再结合定理\textasciitilde{}\ref{thm:13}（非负权、正则、自然的微分形式必为 Pontryagin 形式的多项式），即得“非零形式的权为零”，同时可见 Hirzebruch 的 $L$ 类 \eqref{eq:9} 与 $\Ahat$ 类 \eqref{eq:22} 的权为 $0$。
  \item 笔误——\eqref{eq:72} 中被求和的系数写作 $\omega ^{i,j}_{I,\alpha }$（带逗号），紧接着的正文却写作 $\omega ^{ij}_{I,\alpha }$（无逗号），二者指同一对象；照录原文，仅指出写法不一致。
  \item 含糊——“Weyl's theorem\textasciitilde{}\cite{W} on invariants of the orthogonal group”即正交群不变量的\emph{第一基本定理}：$O(n)$ 不变多项式环由内积 $g_{ij}$ 与体积形式 $\epsilon$ 生成。Gilkey 用它把 \eqref{eq:72} 中出现的全部 $g$ 的导数项化归为这些不变量的多项式，再借“权”的齐次性消去除 Pontryagin 形式之外的项；完整论证见 \cite{ABP}。
  \item 术语——“局部指数定理”含两层内容：(i) 极限 \eqref{eq:74} 存在；(ii) 它等于某个局部示性形式（此处为 $L$ 类，对 Dirac 算子则为 $\Ahat$ 形式）。被取极限的量 $\tr h^0_t-\tr h^1_t$ 是\emph{局部指标密度}；“局部指数定理”断言密度与示性形式恒等，而“整体指数定理”是其积分结果。
  \item 含糊——Mehler 公式（\cite[p.\textasciitilde{}19]{GJ}）给出谐振子热核的显式高斯核 $K_t(x,y)$；Getzler 的要点在于：把位似后的热核当 $t\to0$ 时直接与这一显式核比较，从而把渐近展开化为算子代数的极限。细节见 \cite{Ge2}。
  \item 易错——Getzler 的位似须同时缩放三类对象：$t\mapsto\varepsilon^2 t$、空间坐标按 $\varepsilon$ 缩放、Clifford 变量（奇变量）按 $\varepsilon$ 缩放（各版本指数略有出入，须核对 \cite{Ge2} 的约定）。若只缩放时空而不缩放 Clifford 变量，重标度后的热核不收敛，也得不到 $\Ahat$ 形式。
  \item 术语——“Duistermaat–Heckman exactness of stationary phase”指平稳相近似在等变（环面作用）情形下\emph{精确}：积分恰等于驻点局部贡献之和，无高阶修正项，即 DH 局部化公式。Atiyah 在 \cite{A6} 中把这一精确性形式地移置到自由环路空间上。
  \item 术语——此处的“特征函数”（characteristic function）是概率论意义（$\mathbb{E}[e^{izS}]$），与示性类（characteristic class）无关；中文两者都作“特征”，极易混淆。
  \item 易错——Lévy 公式 $\pi z/\sinh\pi z$ 与 $\Ahat$ 的对应靠 $2\pi$ 归一化：取 $z_i=y_i/2\pi$（$y_i$ 为曲率特征值），则 $\prod_i\frac{\pi z_i}{\sinh \pi z_i}=\prod_i\frac{y_i/2}{\sinh(y_i/2)}$，正是 \eqref{eq:22} 的 $\Ahat$ 类。注意此处 $\sinh$ 取实变量，$z$ 与曲率的对应含因子 $2\pi$，漏掉该因子会得到错误的归一化。
  \item 含糊——“a standard formula in the theory of Lie groups”指：对 $\exp(a+tX)$ 在 $t=0$ 展开可得 $d\exp_a=\frac{1-e^{-\ad a}}{\ad a}$（Maurer–Cartan 形式展开）；此处 $J(x)=(1-e^{-x})/x$，其级数从 $1$ 开始，故可逆。推导与约定见 \cite{BV}。
  \item 笔误——公式标签顺序与出现顺序不符：\eqref{eq:76} 在正文中先于 \eqref{eq:75} 出现，按排版编号二者将互换。正文引用无误，仅标签编号顺序错位，照录原文。
  \item 易错——$\Ahat$ 与 $J$ 的关系正是此谜的答案：$J(x)=\frac{1-e^{-x}}{x}$，故 $\Ahat(x)=\bigl(J(x)J(-x)\bigr)^{-1/2}=\frac{x/2}{\sinh(x/2)}$（与 \eqref{eq:22} 核对）。于是 Berline–Vergne 公式中由 Todd 级数之逆 $J$ 产生的因子，经曲率代入后恰好给出 $\Ahat$ 亏格；注意其中取了几何平均 $\bigl(J(x)J(-x)\bigr)^{-1/2}$，符号约定须与定向一致。
  \item 易错——\eqref{eq:75} 有三处约定须核对 \cite{BV}：(i) 曲率项写作 $\Omega_p/2$，这个 $1/2$ 正是 $\Ahat$ 级数中 $x/2$ 的来源；(ii) $\exp(-a)$ 的符号取决于指标准架与 Levi-Civita 联络的约定，取 $\exp(+a)$ 会使 $J$ 的自变量变号；(iii) 与 $a$ 的缩并用 Killing 型，结果是 $\mathfrak{o}(n)$ 值的 2 形式，还须与 $\Omega$ 配对。
  \item 含糊——“研究热核沿纤维的行为”指在 $\O(X)$ 上作等变（沿纤维的 $\O(n)$ 作用）热核的渐近展开并对纤维方向作局部化；正是在 $t\to0$ 的极限中，纤维方向的 Jacobian 取 \eqref{eq:75} 的形式，从而产生 $\Ahat$ 因子。细节见 \cite{BV}。
\end{enumerate}

\subsection*{§7 Dirac 算子的几何不变量}
\begin{enumerate}[leftmargin=2.2em,itemsep=2pt,topsep=2pt]
  \item 术语——标题中的「几何不变量」指依赖度量（对 Dirac 算子还依赖 spin 结构与联络）的不变量，与 §2–§6 中由示性数和 $K$ 理论给出的拓扑不变量相对；本节按三个层次展开：号差亏损（普通实数或 $\RZ$ 值）、$\eta$ 不变量、以及行列式线丛（带度量与联络的几何对象）。
  \item 含糊——脚注说几何不变量的归宿是微分 $K$ 理论「a more recent development」：其实际内容是把初级（示性数）与次级（陈–西蒙斯、$\eta$）不变量统一为同一个上同调理论中的类的 Cheeger–Simons 微分特征理论（本文 [ChS]、[F2]、[HS]），其算子代数版本见 [FL]；作者在此只提及而不展开。
  \item 含糊——"the index theorem for flat bundles" 指 APS 系列论文中的平坦丛指数定理（见 [APS2]、[APS3] 的相关讨论），它把带边指标定理推广到平坦 $\Gamma$-丛的情形；作者未注出处，阅读时不要把它与 \eqref{eq:92} 的 Dirac 算子指标定理混为一谈。
  \item 易错——\eqref{eq:84} 的归一化：全 Gauss–Bonnet 定理写作 $\int_X K\,d\mu_X=2\pi\chi(X)$，故被积式取 $\frac{K}{2\pi}$ 才恰给出 Euler 数；把 $\Euler$ 理解为 Euler 形式的积分时，这个 $2\pi$ 不可漏掉。
  \item 易错——\eqref{eq:85} 中 $\kappa$ 的符号取决于边界的外法向约定：对平面圆盘（$K=0$、$\kappa=1/r$，取外法向）有 $\Euler=1$；若改取内法向，边界项整体变号，须与 $X$ 的定向一并固定后再核对。
  \item 术语——"a neighborhood ... isometric to the cylinder $[0,\epsilon )\times \bX$ with its product metric" 即度量在边界附近为乘积度量（领状邻域）。这是本节所有带边公式的标准假设：只有此时边界算子 $A_Y$ 才与 $t$ 无关，\eqref{eq:90} 的分解才成立。
  \item 易错——\eqref{eq:86} 中的 $\omega$ 是 $p_1/3$ 的 Chern–Weil 4-形式，而 4 维时 $L$ 亏格的最低次项正是 $L_1=p_1/3$，故 \eqref{eq:86} 与 \eqref{eq:94} 在 4 维一致。切勿把它与 $\Ahat$ 的整性归一化混用：$\hat A_1=-p_1/24$（\eqref{eq:24} 与 \eqref{eq:99}），两者相差因子 $-8$。
  \item 易错——这里是本小节的关键对比：乘积度量使 \eqref{eq:85} 的边界项消失，却不能让 \eqref{eq:86} 成立。原因是 4 维的「边界项」并非某个 3-形式的局部积分（下文说 $\alpha$ 不是 $\int_Y\eta$ 的形式），而是本质上的谱量（非局部）；切勿照搬 Gauss–Bonnet 的直观。
  \item 含糊——"as follows easily from \eqref{eq:86}" 的实质：若 $Y=\bX=\bX'$，把 $X$ 与 $-X'$ 沿 $Y$ 粘合得闭 4-流形 $Z$，对 $Z$ 用 \eqref{eq:86} 即得 $\Sign(X)-\Sign(X')=\int_X\omega-\int_{X'}\omega$，故 $\alpha$ 只依赖 $Y$；此处用到号差在粘合下的可加性与积分的可加性。
  \item 易错——号差亏损 $\alpha$ 在定向反向下变号：$\Sign(\bar X)=-\Sign(X)$ 且 $\int_{\bar X}\omega=-\int_X\omega$，故 $\alpha(\bar Y)=-\alpha(Y)$。$\eta$ 不变量同样在流形反向下变号（其谱关于原点翻转，见 \eqref{eq:93}），故 $-\eta=\alpha$ 与这一奇偶性相容；对照 \eqref{eq:94} 时符号不要弄反。
  \item 含糊——"not of the form $\int_Y\eta$" 的论据是「号差亏损在有限覆盖下不具乘性」：若 $\alpha(Y)=\int_Y\eta$ 且 $\eta$ 由度量局部给出，则 $n$ 重覆盖 $Y'\to Y$ 将给出 $\alpha(Y')=n\,\alpha(Y)$；而 [APS1]（亦见 [A7] §10.3）说明号差亏损没有这种乘性。严格核对需计算 $\eta$ 在覆盖下的行为。
  \item 术语——Hilbert 模曲面 $X^0(K)=(\HH\times\HH)/\PSL_2(\sO)$ 是 Hilbert 模形式所对应的模曲面；它有 orbifold 点（由 $\sO$ 中的有限阶单位引起）和与理想类一一对应的端。截断端、挖去 orbifold 点邻域后得到的 $X(K)$ 是带边光滑流形，其边界正是号差亏损的计算对象。
  \item 含糊——透镜空间边界上的号差亏损「was known from the equivariant signature theorem §5.1」：实质是对轨形点的循环群稳定化子用 [A7] §10.3 的等变号差定理计算 $\Sign$（所得边界为透镜空间 $L(p,q)$，公式含 Dedekind 和）；作者未写出公式，查 [A7] §10.3 与 §5.1 的等变号差定理。
  \item 含糊——「desingularization of the cusp singularity in the cone on the boundary」指 Hirzebruch 用尖点奇点的解消（例外除子链的相交数由 Hirzebruch–Jung 连分数给出）把号差亏损写成连分数表达式（[H4] p.\textasciitilde{}222）；"continued fraction" 就是该解消的连分数。
  \item 含糊——「up to a numerical factor」中的因子需按 [H4] p.\textasciitilde{}231 的 $L$ 函数归一化核对（是否含 $\pi$、类数以及 $\sqrt d$ 的幂）；这一步是 $\eta$ 不变量与数论 $L$ 函数对应的核心，不要把它当作纯记号问题。
  \item 术语——号差亏损 $\alpha$ 与 $\eta$ 不变量的关系是 $\alpha(Y)=-\eta(Y)$，其中 $\eta(Y):=\eta(0)$ 是解析延拓值（本节末 \eqref{eq:94} 之后）。作者先用 $\alpha$ 的几何/数论性质（[H4]、[Sh]、[ADS]、[Mu]）引出问题，再用 $\eta$ 不变量给出其解析表达，这条线索不要读反。
  \item 含糊——局部椭圆边界条件的「拓扑障碍」见 [AB3]（Atiyah–Bott, The index problem for manifolds with boundary）与本文 §5.2：符号类 $\sigma\in K^0(T^*X,T^*X\setminus0)$ 在边界 $K^0(T^*X|_{\bX})$ 上的限制一般非零；只有当它消失（即 $\sigma$ 提升为相对 K 群中的类）时才存在局部边界条件。对号差算子这一限制非零，因此必须改用整体条件。
  \item 术语——$\dbar$ 是 Cauchy–Riemann 算子（$\partial/\partial\bar z$ 的一半），其核是全纯函数而非调和函数，故为无穷维；这正是需要边界条件把核截成有限维的原因。与调和形式的 Hodge 理论对照阅读更清楚。
  \item 易错——$a\in \RR\setminus \ZZ^{\ge0}$ 且条件取 $m>a$：若 $a$ 取非负整数，则 $m=a$ 的 Fourier 系数的归属需另行规定（谱在 $a$ 处有本征值，谱投影 $\lambda<a$ 与 $\lambda\le a$ 给出不同的边界条件）；这正是作者要求 $a$ 避开 $\ZZ^{\ge0}$ 的原因，核对 §5.2 与 [APS1] §1 的相应约定。
  \item 含糊——"it extends to a Fredholm operator on suitable Sobolev completions" 的实质：在 $H^s(\Omega)$ 上取 $\sH_a$ 的闭包（即边界 Fourier 系数在 $L^2$ 中的闭包），并验证该闭扩张的核、余核有限维且有闭值域；这正是 [APS1] §1 对「椭圆边界条件」所要求的两条性质。
  \item 易错——\eqref{eq:90} 中 $\sigma$ 是 $dt$ 的 Clifford 乘法，若 $dt$ 取单位 1-形式并采用黎曼约定则 $\sigma^2=-1$，从而 $\sigma\inv=-\sigma$；$A_Y=\sigma\inv D_Y$ 的自伴性、以及 \eqref{eq:91} 中谱 $\lambda$ 的符号都依赖这一约定，核对 §4 的 Clifford 约定。
  \item 易错——取 $a=0$ 时边界条件是「负谱空间」（$\lambda<0$）的闭包；若改取正谱空间，相当于取伴随算子或反转定向，指标公式 \eqref{eq:92} 中 $(\et{\bX}+h)/2$ 的符号要相应改变。用 \eqref{eq:92} 前先固定这个约定，并与 [APS1] 定理 4.14 对照。
  \item 含糊——"one must take into account $\ker D_Y$ separately" 的实质：零特征空间的归属（取谱投影 $P_{<0}$ 还是 $P_{\le0}$）会改变指标，其效应正是 \eqref{eq:92} 分子中的 $h=\dim\ker D_{\bX}$；所以 $h$ 项不是修正项，而是边界条件定义的一部分。
  \item 含糊——"The problem splits into two pieces: a cylinder near $\bX$ and $X\setminus\bX$" 与 "Gluing ... via a partition of unity" 的实质是 [APS1] §4 的两步：柱面上把 $D_X$ 化为以谱参数为变量的常微分算子族并显式计算其指标；内部用热核方法（如 §6.2 的局部指数定理）计算，最后用单位分解与热核的边界层渐近展开把两部分拼接。真正困难的是边界附近热核的渐近展开。
  \item 易错——\eqref{eq:92} 的归一化：分母的 $2$ 来自把 $\eta$ 的解析延拓与谱投影（正/负谱空间）拼成 Fredholm 指标，$(h+\eta)/2$ 中的 $h=\dim\ker D_{\bX}$ 是零模的贡献；若把 $h$ 取成别的量（如 $\dim\ker D_X^-$）或把边界条件取成另一侧投影，系数会变成 $(\eta-h)/2$ 或整体反号。核对 [APS1] 定理 4.14 的约定。
  \item 易错——\eqref{eq:93} 只对非零谱求和：零特征值若出现在求和中，$\sign 0$ 无定义且 $|\lambda|^{-s}$ 项发散，这正是 \eqref{eq:92} 中另设 $h$ 的原因。另注意 $\eta(0)$ 是解析延拓值，不是级数在 $s=0$ 处的和（级数只在 $\Re(s)$ 充分大时收敛）。
  \item 含糊——解析延拓的存在性、有限性与脚注中「在 $\Re(s)>-1/2$ 上全纯」的断言依赖 Seeley 的复幂理论（见 [Se4]，以及 [Se1]、[Se2]）与 [APS1] §2 的定理；一般的亚纯延拓可以有极点，$s=0$ 处正则并非自动成立，而是 APS 定理的实质内容之一。作者以一句「there is a meromorphic continuation」带过。
  \item 易错——$\eta$ 在流形反向时变号：反向使 $A_Y$ 变为 $-A_Y$（或复共轭），故 $\eta(-A_Y)=-\eta(A_Y)$，于是 $\eta(\bar Y)=-\eta(Y)$。这与 \eqref{eq:94} 中 $\Sign(X)=\int_X\omega-\et{\bX}$ 的奇偶性一致：$\Sign$ 与 $\int\omega$ 都变号时，$-\eta$ 保持同一变换规律。
  \item 易错——\eqref{eq:94} 中边界算子的公式 $(-1)^{k+q+1}(*d-d*)$ 须按 [APS2] 的约定理解：$*$ 为 Hodge 星、$d$ 为外微分，$k$ 来自 $X$ 的维数 $4k$（故 $\dim\bX=4k-1$），$q$ 为形式的次数。不同文献把它写成 $\pm(*d-d*)$ 或等价的 $\pm(d*+*d)$ 等形式，符号与 $k,q$ 的平移约定不同；用 \eqref{eq:94} 前请核对 [APS2] §4 与 [APS1] 的 $\eta$ 定义。
  \item 含糊——\eqref{eq:94} 与 \eqref{eq:92} 相比不见了 $h$ 项（\eqref{eq:92} 为 $(\et{\bX}+h)/2$）。这一「简化」依赖号差算子的特殊约定（零特征空间与谱投影的处理方式由号差算子的奇偶性决定），作者未加说明；严格核对须看 [APS1] 定理 4.14 与 [APS2] §4 中 $h$ 的定义，以及 $-\eta$ 即号差亏损 \eqref{eq:87} 这一事实。
  \item 术语——「谱不变量」指只依赖 $\bX$ 上号差算子谱（特征值）的量，而不像示性数那样是局部曲率多项式的积分。号差亏损的定义 \eqref{eq:87} 用的是 4-流形的号差与 $p_1$ 积分，二者都不是谱量，但差值却是谱量——这正是 APS 定理的惊人之处。
  \item 易错——\eqref{eq:95} 定义的是 $\RZ$ 值映射 $\xia{Y/S}$（模 $1$ 的等价类），不是实值函数；只有选定实值提升后才能谈 \eqref{eq:96} 的微分。$(\et{Y/S}+h)/2$ 的模 $1$ 不确定性恰好是族指标定理所给的整性，故 \eqref{eq:95} 把 $S\to\RZ$ 的\emph{同伦类}精化为具体的几何映射。
  \item 易错——\eqref{eq:96} 中 $d\xia{Y/S}$ 是对实值提升取的微分，而 $\xia{Y/S}$ 只定义到模 $1$；两者相容是因为 $\int_{Y/S}\omega$ 的积分在换提升时不变（整数部分被模 $1$ 吸收）。与 \eqref{eq:111} 比较时注意这里的 $d$ 是 $S$ 上的外微分，而非流形 $X$ 上的。
  \item 含糊——"Atiyah-Patodi-Singer [§7 APS3], in collaboration with Lusztig" 给出的另一种解析计算：谱流 = 绕数。其严格证明需要族的谱在 $\cir$ 上构成 $\cir\times\RR$ 的闭子集（由自伴性与逆紧性保证），从而相交数有定义；见 [APS3] §7（该文与 Lusztig 合作的结果）以及 [APS1] 中 $\eta$ 的解析性质。
  \item 术语——谱流与指标：谱流是 $\cir$ 参数族 Dirac 算子的整数不变量，可理解为 $\cir\times Y$ 上的族指标（或映射环面上的指标），也是 \eqref{eq:96} 的「曲率」在 $1$ 维底空间上的积分；不要把它与单个算子的 Fredholm 指标 \eqref{eq:34} 混为一谈。
  \item 易错——本例谱为 $\{s+n:n\in \ZZ\}$，$s$ 从 $0$ 走到 $1$ 时恰有一个分支穿过 $0$（$n=0$ 分支由 $0$ 移到 $1$，同时 $n=-1$ 分支由 $-1$ 移到 $0$），故谱流为 $1$（「从负到正」取正号）。不要因为 $s=0$ 与 $s=1$ 两端各有一个零特征值就误判为 $0$ 或 $2$。
  \item 术语——初级不变量（示性数、指标）是整值拓扑量；次级不变量是其「积分」的依赖联络的修正，取 $\RZ=\RR/\ZZ$ 值。总测地曲率 $\int\kappa/2\pi$ 就是最简单的次级不变量：它与和乐角只差 $2\pi$ 的整数倍，故模 $\ZZ$ 后等于 Levi-Civita 联络的和乐，这解释了为什么 CS 型次级不变量天然是 $\RZ$ 值的。
  \item 易错——Chern–Weil 的曲率约定：$\Omega = d\Theta + \frac{1}{2}[\Theta,\Theta]$ 与 $\Omega = d\Theta + [\Theta,\Theta]$ 两种写法取决于 $\Theta$ 的取值是否含因子；$p(\Omega)$ 的符号与系数随之改变。对照本文 §3 与 [CS] 时务必先固定 $\Theta = g^{-1}dg$（左不变）还是 $\Theta = dg\,g^{-1}$（右不变）的约定。
  \item 含糊——"Refinements exist only if the periods of $c_p$ are integers"：这是必要条件，其充分性等价于 $c_p$ 落在 $H^{2k}(BG;\ZZ)\to H^{2k}(BG;\RR)$ 的像中（可查 [ChS] 的 Cheeger–Simons 理论）。作者未说明 "periods" 指在闭 $2k$ 维流形上的配对值，也未给出构造：$d\,\mathrm{CS}(\Theta) = p(\Omega)$，其中 $\mathrm{CS}$ 由沿 $\Theta_t=(1-t)\Theta_0+t\Theta_1$ 的积分给出。
  \item 笔误——原文作「1972 年」陈–西蒙斯引入该不变量，但 [CS] 于 1974 年才刊出（1971 年提交）；此年份应指提交/预印本，引用时以 [CS] 的发表年 1974 为准。
  \item 笔误——\eqref{eq:97} 写作 $p_1(W)[W]$，但本段刚把 4-流形记作 $X$，$W$ 未定义；\eqref{eq:99} 又写成 $p_1(TW)[W]$。疑原文应以 $X$（或 $TX$）替代 $W$；照录原文，请核对 [H4] 或 §2 中示性数记号。
  \item 含糊——$\Gamma (Y)$ 的共形不变性与「共形浸入 $Y^3\to\EE^4$ 的障碍」出自 [CS] §6；这里 $\Gamma$ 是 Levi-Civita 联络的 CS 不变量（3 维、$\mathrm{SO}_3$ 丛），与 \eqref{eq:98} 的 $\xia Y$ 相差一个归一化因子（见例\textasciitilde{}\ref{thm:22} 中 $-24\xia Y$ 与 $\Gamma(Y)$ 的比较）。
  \item 术语——微分上同调（differential cohomology）的类同时携带上同调类（初级）、联络/曲率（次级）与挠信息；[ChS] 是 Cheeger–Simons 微分特征，[F2] 是 Freed 的经典 Chern–Simons 理论，[HS] 是 Hopkins–Singer 关于二次函数与微分上同调的表述。与 \eqref{eq:98} 的 $\RZ$ 值 $\eta$ 不变量对照，可见它应由微分 $K$ 理论（而非微分上同调）承载，见注\textasciitilde{}\ref{thm:23}(2)。
  \item 易错——\eqref{eq:98} 中 $\pmod1$ 的归一化：$\eta$ 与 $h$ 都是实数，除以 $2$ 后模 $1$；若改用 $\eta/2$（不除 $h$）或把 $\eta$ 的定义 \eqref{eq:93} 换成不取 $\sign\lambda$ 的版本，结果会整体相差 $h/2$ 或符号。核对 [APS1]/[APS2] 的 $\eta$ 归一化。
  \item 易错——"depends smoothly" 的断言依赖族指标定理（[APS3]）：$\eta$ 不变量本身只是 $\RR/\ZZ$ 值，其光滑性指模 $1$ 后的光滑映射；脚注中号差算子情形的 $\RR$ 值光滑性来自调和形式的上同调解释，二者不同，不要混淆。
  \item 易错——\eqref{eq:99} 的系数 $-1/24$ 来自 $\hat A_1=-p_1/24$（与 \eqref{eq:24} 的整性定理配套）；不要把它与 \eqref{eq:86} 的 $p_1/3$ 相混。$\hat A$ 亏格的整性（在 spin 流形上）是这里「$-24\xia Y$ 与 $\Gamma(Y)$ 可比较」的根据。
  \item 含糊——"a similar comparison in more general circumstances leads to a spin bordism invariant"：作者未指出该配边不变量的确切群与计算；其含义是次级不变量的比较在（一般维数或带 $G$-丛的）配边群上给出 $\RZ$ 值不变量（见 [APS2] 及 [Gi2]；注意 $\Omega^{Spin}_3=0$，3 维情形本身不产生非平凡同态）。清分母这一步正是 $\RZ$ 值与有理系数示性数不能直接相乘的技术后果。
  \item 术语——Adams $e$ 不变量是稳定同伦论中标架配边 $\Omega^{fr}_*$ 的不变量，[APS2] §4 说明它可由 $\eta$ 不变量构造；而 "framed bordism" 与本文 §2 的配边记号一致，勿与 spin 配边混淆。
  \item 易错——\eqref{eq:101} 只是形式记号（原文用引号标出 $=$）：无穷乘积 $\prod\lambda_m$ 一般发散，且若 $\lambda_1=0$ 则平凡地为零；故必须（i）剔除零特征值，(ii) 用 $\zeta$ 正规化（\eqref{eq:104}）赋值。$X=\cir$ 时 $\prod m^2$ 发散这一点正说明需要正规化。
  \item 笔误——\eqref{eq:102}、\eqref{eq:103} 中 $\lambda^{-s}$ 与 $\log\lambda$ 均未标出下标 $m$，而求和指标为 $m$；疑应为 $\lambda_m^{-s}$、$\log\lambda_m$。照录原文，请核对 [RS] 或 §6.1 的 $\zeta$ 函数定义。
  \item 含糊——$\zeta_\Delta(s)$ 的解析延拓存在、且 $s=0$ 为正则点，这依赖 Seeley 的复幂/热核理论（[Se4]）与 [RS] §1；作者以「use the regularity」一句带过。另注意 $\det\Delta$ 无尺度不变性：替换 $\Delta\mapsto c\Delta$ 会引入 $e^{+\zeta_\Delta(0)\log c}$ 型因子。
  \item 易错——行列式线 \eqref{eq:106} 的定义要求 $\dim V^0=\dim V^1$：此时 $\Det T$ 是 $\Det V^0\to\Det V^1$ 的映射（即一条线上的元素），而不是数；若 $V^0=V^1$，$\det T$ 才是通常的数值行列式。对一般的 Fredholm 算子，$V^0,V^1$ 无穷维，\eqref{eq:123} 不再有意义，必须改用 $\ker$ 与 $\coker$ 的有限维性。
  \item 含糊——Quillen [Q] 在 $\Fred$ 上构造行列式线丛 $\pi:\Det\to\Fred$：对 $S$ 上的 Fredholm 算子族 $T:S\to\Fred$，拉回 $T^*\Det\to S$ 的纤维为 $\det(\ker T_s)\otimes\det(\coker T_s)^*$（在指标局部常值的分层上粘合），且 $\det$ 给出连续截面；该族的同构类正是 \eqref{eq:108} 中的 $\Det\in H^2(S;\ZZ)$（拓扑部分）。这里「next topological invariant」指该 $H^2$ 类，而非数值指标 \eqref{eq:34}。
  \item 易错——\eqref{eq:110} 的相位约定：$e^{-2\pi i\xia{X_\varphi(\epsilon)}(V)}$ 中 $\xi$ 取 \eqref{eq:98} 的 $\RZ$ 值，$2\pi i$ 因子使其与和乐的 $U(1)$ 值匹配；若把 $\xi$ 换成 $-\xi$ 或 $\eta$ 而非 $(\eta+h)/2$，会得到 $e^{+2\pi i\cdots}$ 或差一个符号。核对 [BiF2] 的符号约定。
  \item 含糊——"the \emph{adiabatic limit} in \eqref{eq:110} was introduced in [Wi3]"：绝热极限 $\epsilon\to0$ 使纤维方向的度量相对底空间变得巨大（$g_\cir/\epsilon^2$），谱的间隙行为决定极限存在；其严格分析见 [BiF2]（Bismut–Freed 的和乐定理），其他几何解释见 [Che]、[S2]。作者未说明极限为何收敛。
  \item 易错——\eqref{eq:111} 的 $2\pi i$ 与 $\omega$ 的 Chern–Weil 归一化必须匹配：$\omega$ 代表 $\Ahat\ch$ 时，曲率形式按 $\curv = 2\pi i\,\omega$ 的约定给出第一 Chern 类；对照 \eqref{eq:96} 的 $d\xi=\int_{Y/S}\omega$（无 $2\pi i$）可见两处因子的来源不同。
  \item 术语——Quillen 度量与 [BiF1] 的协变导数用的是不同的 $\zeta$ 函数：前者用 $-\zeta'(0)$（\eqref{eq:104}），后者用沿族参数方向的 $\zeta$ 函数（即对族 Dirac 算子作超迹/热核展开），二者合起来才给出 Hermite 线丛的结构；不要把它们混为同一个 $\zeta$ 正规化。
\end{enumerate}

\subsection*{§8 反常与指数理论}
\begin{enumerate}[leftmargin=2.2em,itemsep=2pt,topsep=2pt]
  \item 术语——本节“anomaly（反常）”至少有四义，读时须分清：(i) 手征/微扰反常，由 $\frac{1}{2\pi}\int F$ 型积分表达的对称性破缺；(ii) 整体反常，即把射影对象提升为线性对象时的拓扑阻碍，由和乐测得（[Wi3]）；(iii) 引力反常，见 [AgW]；(iv) 反常流入，用高一维的经典作用量抵消低维边界上的反常。作者在本节把“反常”统一提升为量子理论射影性的度量，见 \eqref{eq:114} 与 §\ref{subsec:8.4}。
  \item 含糊——作者用三元组 $(\sP,p,H)$ 而不用希尔伯特空间来定义量子系统，正是为了把相位自由度从结构里剔除：$\sH$ 只是坐标选择，$\sP$ 才是物理态空间；\eqref{eq:112} 与 \eqref{eq:113} 两处的讨论都只依赖射线的相对相位。脚注提到的 $C^*$ 代数框架（代数量子力学）被作者以“较简单的语境已经足够”一句带过，实质是本节不需要算子代数的一般理论。
  \item 含糊——“Define $\sP$ as the limit of $\PH$ over all $(\sH,\theta )\in \cP$”中的 limit 用词不严：$\cP$ 是连通的群胚，且其自同构在 $\PH$（以及 $\End\sH$）上作用平凡，故此图表上的 limit 与 colimit 同构，$\sP$ 按哪一种读都得到同一空间；作者沿用 limit 一词只是行文方便，不必按范畴论的严格意义理解。
  \item 易错——\eqref{eq:112} 的跃迁概率只依赖两条射线的相对相位：把 $\psi _i$ 换成 $e^{i\alpha _i}\psi _i$ 不改变 $|\langle \psi _1,\psi _2\rangle|^2$，却改变振幅 \eqref{eq:116} 的相位，这正是射影性的数学内容；此外 $p$ 的定义要求 $\psi _i$ 取单位范数，否则须除以范数，归一化因子极易漏掉。
  \item 含糊——“A basic theorem of Wigner asserts”省略的是：$\PH$ 上保持跃迁概率 $p$ 的任一（双射）变换必由 $\sH$ 的酉算子或反酉算子诱导，故 \eqref{eq:113} 中 $q$ 满射、\eqref{eq:114} 的扩张存在。本文参考文献未收 Wigner 原文；标准证明见 Wigner 1931 年专著 \emph{Group Theory}（中译本《群论与量子力学》），或 Weinberg \emph{The Quantum Theory of Fields} 卷一 §2.2。
  \item 术语——“反常”在此定义为 \eqref{eq:114} 这个 $\zt$ 分次的群扩张，而不只是一个中心扩张：反酉（时间反演型）对称性的存在使扩张带有 $\zt$ 分次 $\epsilon $。若 $\epsilon $ 非平凡，$G_\sH$ 不再是 $G$ 上的中心扩张，而是“$\zt$ 分次中心扩张”，此时同构类不由 $H^2(G;\TT)$ 单独决定，还需记录 $\epsilon $。
  \item 含糊——“the symmetry persists in the limit”的实质是：经典极限下对称群 $G$ 线性地作用，而量子系统中它只射影地作用，$G_\sH$ 就是这一射影作用的线性化；作者没有给出“经典极限存在且对称性保持”的判别法，这在一般情形下是困难的（需半经典分析）。
  \item 易错——第 2 条里“把 $S$ 积掉”要求先选定 $S$ 上的测度，$L^2$ 截面空间才良定义；测度会改变所得量子系统，但反常（提升的阻碍）与测度无关，它只依赖射影丛结构。另注意“反常是阻碍”与“反常是特征”两种说法的区别：不积分 $S$ 时它恰是族指标定理中的那个类，这正是 §\ref{subsec:8.4} 中“反常本身是场论”的雏形。
  \item 笔误——注 \ref{thm:27}(3) 中原文作 “in quantum field \textasciitilde{}(\S\ref{subsec:8.4})”，疑脱 theory 一词，应为 “in quantum field theory”（量子场论）；译文按补足后的意思翻译，原文照录。
  \item 含糊——“The isomorphism class of the central extension lives in $H^2(G;\TT)$”是中心扩张的分类定理，作者未给证明：$G$ 离散时用抽象群上同调，$G$ 为李群时须用连续（光滑）上同调，否则会出现非连续的“怪”扩张——这正是“The type of cohomology depends on the type of group”一句的含义。作者以“假定 $G$ 有限”回避了这一点；§\ref{subsec:8.4} 把 $G$ 换成配边范畴时，群上同调就被替换为配边谱的上同调。
  \item 易错——\eqref{eq:115} 的 $\hbar$ 只出现在演化算子 $\evo{k+1}k$ 中（它们通常含因子 $e^{-iH(t_{k+1}-t_k)/\hbar}$）；概率对 $\hbar$ 的依赖是间接的。别把 $\hbar$ 误写进 \eqref{eq:116} 的振幅相位，也别把 $\hbar$ 与 $2\pi$（例如瞬子相位 $e^{2\pi i\theta}$ 中的 $2\pi$）混为一谈。
  \item 术语——注 \ref{thm:28} 中“诸线 $\sL$ 构成数据参数空间上的线丛”把本节的射影性讨论与 §\ref{subsec:8.3} 接上：在后者中这些线丛具体化为狄拉克算子族的行列式线丛 \eqref{eq:122}。注意线丛作为拓扑对象只记 $c_1$，作为带度量与联络的线丛才记录反常的几何信息（[F3]、[BiF1]）。
  \item 含糊——“soft”指本节只用到射影空间、群扩张等不依赖具体分析的拓扑/范畴事实，类比于对一般 Fredholm 算子的指标理论（只需算子的拓扑性质即可定义指数）。作者预告“这不是类比”：一旦换成威克转动后的旋量场，反常就由族指标定理精确计算，见 \eqref{eq:118} 与 \eqref{eq:119}。
  \item 含糊——[AS7] 是《美国科学院院刊》上的四页短文，只陈述结果；“we first summarize their work”所省略的细节（族指标、超渡、行列式线丛）散见阿蒂亚–辛格的系列长文与 [S3]、[F3]；本节的 \eqref{eq:117}–\eqref{eq:119} 就是这些细节的提要。另注意 [Lt2] 是 Lott 的“轴反常的一个证明”，不要与 [Lt1] 的指标 gerbe 混淆。
  \item 易错——为何要除以 $\sG_0$ 而非整个 $\sG$：若 $g\in\sG$ 固定联络 $A$ 且 $g(p_0)=1$，则由 $A$ 的平行性 $g$ 必恒等于 $1$，故 $\sG_0$ 在 $\sA$ 上自由作用；对整个 $\sG$，使 $A$ 不动的常值（中心）变换使作用不自由，$\sA/\sG$ 也不再是好的丛。由 $g\mapsto g(p_0)$ 得 $\sG/\sG_0\cong G$，于是 \eqref{eq:117} 中 $\pi $ 的纤维恰是 $G$。$\sG_0$ 依赖 $p_0$ 的选取，但 $\pi $ 的同构类与 $\ind$ 不依赖。
  \item 术语——\eqref{eq:118} 中 $K^{-2n}\cong K^0$ 用的是复 $K$ 理论的 Bott 周期性（周期 $2$）：$-2n$ 为偶数故同构于 $K^0$。注意指数里出现 $2n$，与本节开头所设 $\dim X=n$ 不一致，详见下一条批注。
  \item 易错——归一化：本文取陈特征 $\ch_1=\frac{i}{2\pi}\Tr F$，故第一陈类的实表示由 $\frac{i}{2\pi}\Tr F$ 给出、周期为整数；若改用物理学家的实值场强 $F$ 书写，则对应 $\frac{1}{2\pi}\int F\in \ZZ$。$2\pi$ 因子写错时反常只差整数倍，却会翻转“整体反常是否为零”的判定；此处的 $2\pi$ 与 $\hbar$ 无关，切勿混同。
  \item 笔误——本段三处 $2n$ 与开头“$\dim X=n$”不相容：\eqref{eq:118} 的指数是 $K^{-2n}$，\eqref{eq:119} 的 $\omega $ 是 $(2n+2)$ 形式、超渡后是 $(2n+1)$ 形式。若改按 $\dim X=2n$ 读则全部自洽：$(2n+2)$ 形式沿 $2n$ 维纤维积分恰得基上的 $2$ 形式 $=\ch_1$，族指数落在 $K^{-2n}\cong K^0$。疑原文的 $n$ 与 $2n$ 混用（请核对 [AS7] 中流形维数所采用的记号），照录原文。
  \item 含糊——“They also relate their result to determinants”一句跳过的实质是：$c_1$ 等于行列式线丛的同构类（拓扑层面），而要在几何层面得到该线丛的度量与联络，须用 [BiF1] 的 Bismut–Freed 联络与 [F3] 的分析；本文 §\ref{subsec:7.5} 已讨论这一构造。仅凭 \eqref{eq:119} 只能得到 $c_1$，得不到“平坦截面”的阻碍（见下段）。
  \item 易错——\eqref{eq:121} 的指数是 $-(\psi ,D_s\psi )$，既无 $i$ 也无 $\hbar$：这是已威克转动到黎曼号差的欧几里得作用量；回到洛伦兹号差时指数为 $e^{iS/\hbar}$，而“相位取 $2\pi i$ 的整数倍”这类量子化条件在欧几里得表述中表现为 $\det D$ 的相位/符号。不要把 $e^{iS/\hbar}$ 的 $i$ 与 $2\pi i$ 中的 $2\pi$ 的归一化混在一起，否则行列式的相位会差一个与 $\hbar$ 无关的整数倍。
  \item 术语——\eqref{eq:121} 中的 $(\psi ,D_s\psi )$ 不是希尔伯特空间上的埃尔米特内积，而是旋量丛的实（Majorana 型、电荷共轭型）双线性型 $\bar\psi D_s\psi $；其反对称性等价于 $D_s$ 关于相应埃尔米特内积自伴、而该双线性型由电荷共轭 $C$ 给出（于是 $CD_s$ 反对称）。换用不同的实结构（或改变旋量丛定向）会使该型乘上符号，从而 $\pfaff$ 差一个符号——这是行列式线丛相位模糊的根源。
  \item 含糊——“the pfaffian reduces to a determinant”省略的是：当旋量丛按手征性分裂为 $\sS=\sS^+\oplus \sS^-$、狄拉克算子呈非对角形式 $D=\begin{pmatrix}0&D^-\\D^+&0\end{pmatrix}$ 时，$\pfaff D=\det D^+$（至多相差一个与实结构、定向有关的符号）；“阿蒂亚–辛格框架”之所以能作此简化，是因为那里的算子正是这种手征形式。若不指定实结构，$\pfaff$ 本身只在相差一个相位的意义下有定义。
  \item 术语——行列式线丛的纤维是 $\det\Ker D_s\otimes(\det\Coker D_s)^{-1}$；本文的 \eqref{eq:121} 直接给出正规化行列式作为截面，而在一般的族上定义 $\det D_{X/S}$ 须先作 $\zeta$ 正规化（[F3]）并赋予 Quillen 度量与 [BiF1] 的联络。作者把 \eqref{eq:121} 当作定义，是为了避开这些分析细节。
  \item 易错——第 1 条中“就是反常”应理解为：反常的完整数据是线丛 + 度量 + 联络，而不只是它的拓扑类 $c_1$。只记 $c_1$ 只能看到“能否取非零截面”，看不到“能否取平坦截面”；后者才是（下一段 [F3] 意义下的）反常。
  \item 含糊——“Such integrals are problematic analytically”在注 \ref{thm:27}(2) 处已被点过一次：对度量/联络的积分没有勒贝格测度，路径积分在分析上未定义；但“反常”与“平凡化”可以严格定义为线丛及其截面，故作者能在数学上谈论它。这正是 §\ref{subsec:8.4} 把反常提升为场论的出发点。
  \item 易错——局部与整体的区分在此是：局部反常用 \eqref{eq:119} 的 $2$ 形式（自然联络的曲率，闭而不精确）度量，整体反常用其和乐（沿参数空间中非平凡闭路的相位）度量。若讨论反常流入（anomaly inflow），方向是：$(n+1)$ 维体上的陈–西蒙斯型作用量的变分精确抵消 $n$ 维边界上的规范变分；把体与边的定向弄反，会使抵消变成加倍。本文不展开流入机制，可参 [Wi3]、[AgW]。
  \item 术语——gerbe（中文或作“格贝”，常保留原文）是线丛的下一级对象，可理解为“局部由线丛粘合而成的层”，其“曲率”是 $3$ 形式；把射影丛提升为向量丛的阻碍比第一陈类“高一级”，故由狄拉克格贝而非线丛记录。[Lt1]、[Bu] 讨论指标 gerbe 及其超渡；[FS] 是 Faddeev–Shatashvili 的论文（切勿误读为本文作者 Freed 与他人合著），[Seg2]、[NAg] 讨论其在物理中的实现。
  \item 含糊——引文出自 [A12]（阿蒂亚自述）。“the combinatorial approach is so far the only one that has been made to work (for the Jones polynomials)”指 Reshetikhin–Turaev/量子群给出的组合构造：在阿蒂亚写这段话的年代，陈–西蒙斯路径积分尚无严格定义，而组合构造可以严格化。对照 [Wi4] 中威滕用路径积分给出不变量，二者的关系即“公理 vs 构造”。
  \item 易错——号差为何是配边不变量：$\Sign(X)=\langle L(X),[X]\rangle$，$L$ 亏格由 Pontryagin 类给出，配边关系下该配对不变；更本质地，\eqref{eq:124} 是整配边群 $\Omega _{4k}(\SO)$ 上的 $\ZZ$ 值同态，而任何到 $\ZZ$ 的同态在挠类上自动为零，故号差只看到自由部分（$\Omega _{4k}(\SO)\otimes\QQ$ 由复射影空间的乘积生成）。定理 \ref{thm:3} 的证明正依赖于此。
  \item 易错——这里的 $d$ 是空间维数，TQFT 的时空维数是 $d+1$；§\ref{subsec:8.3} 的 $n$ 是时空维数，故 $n=d+1$。后文 §\ref{subsec:8.4} 说反常是 $(n+1)$ 维场论，因为它比被反常化的 $n$ 维理论高一维。这三处维数极易混，读时务必对表：$d$（对象维数）、$d+1$（态射/时空维数）、$n=d+1$、$n+1$（反常理论维数）。
  \item 含糊——“categorifies”的确切含义是：$\Bd$ 的 $\pi_0$（对象模“存在态射”）给出 $\Omega _d$，而 $\Bd$ 还保留了对象之间的态射与复合，即保留了配边的全部信息。严格说 $\Omega _d$ 视切结构而定，见第 2 条。
  \item 术语——tangential structure（切结构）包括定向、$\Spin$ 结构、标架等，它决定配边群的类型；正文未指定切结构处写的 $\Omega _d$ 严格应为非定向配边群 $\Omega _d(\O)$，而 \eqref{eq:124} 中有定向，写作 $\Omega _{4k}(\SO)$。后文 $M\!\Spin$ 对应的正是 $\Spin$ 切结构下的配边谱。
  \item 含糊——定义 \ref{thm:31} 只写了函子性，未写出全部相容性：还须要求 $F(\emptyset)=\mathbf{1}$（空流形对应基域）、$F(Y_0\amalg Y_1)\cong F(Y_0)\otimes F(Y_1)$（不交并对张量积）、以及粘合配边对应态射的复合；这些正是“symmetric monoidal functor”一词的内容。现代严格化版本要求 $\Bd$ 为 $(\infty,1)$-范畴、$F$ 为幺半 $(\infty,1)$-函子（配边假设），作者以“with suitable modifications and extensions it is believed to apply widely”一句带过。
  \item 术语——Atiyah–Segal 公理系统中，Segal 的原始公理针对二维共形场论（[Seg3]，保留共形/度量结构），Atiyah 把它搬到拓扑情形（[A11]，只需配边结构，并要求 $F(Y)$ 有限维）。二者的差别在于是否保留度量/共形数据；拓扑情形相当于共形权为零的极端，故配分函数只剩拓扑信息。
  \item 易错——法标架的选取会改变同痕类 $\partial X\approx\cir\times \cir$，从而改变 $F(X)$ 的取值：这正是“标架纽结不变量”与“无标架纽结不变量”的区别。琼斯多项式要成为无标架不变量，须再乘/除以一个依赖标架的因子（陈–西蒙斯理论中表现为一个相位）；作者说“of this type”是就标架依赖的意义而言。
  \item 含糊——“the trivializations of the anomaly must be coherent in the background data”是本节最关键也最含糊的一句：局部性要求反常的平凡化在配边的粘合下相容（即平凡化构成背景数据范畴上的一个函子性对象），这正是阿蒂亚–西格尔公理的内容；作者未证明“局部性 $\Leftrightarrow$ 函子性”，该等价是现代公理化场论（[FHT]、[FH]）的工作。
  \item 术语——“可逆场论”的确切含义是：$F$ 经由线与可逆线性映射构成的 Picard 群胚 $\Line_k\subset \Vk$ 分解，于是所有态空间都是一维、所有态射都可逆，$F$ 只记录相位而无动力学自由度；物理上这对应可逆相/SPT 相。反常正是这种场论，但一般不是拓扑的：它依赖背景度量与联络（见 \eqref{eq:120} 的数据），故只能定义在带相应结构的流形上——这与 §\ref{subsec:8.1} 中“振幅落在未被平凡化的线里”是同一件事的两种说法。
  \item 易错——脚注强调：$n$ 维理论的反常只定义在维数 $\le n$ 的流形上，故它不是完整的 $(n+1)$ 维理论。把它延拓为完整的 $(n+1)$ 维可逆理论需要额外数据（例如把 $n$ 维理论视为 $\sB$ 的表示）。作者在下一条注中称之为“once-categorified invertible $n$-dimensional field theory”，切勿与“$(n+1)$ 维理论”混为一谈。
  \item 含糊——注 \ref{thm:34} 的“cohomological interpretation”指 [FHT] 的如下对应：可逆场论 = 谱的映射（或稳定同伦意义下的一个类），于是把“一维表示”提升为上同调类、把次数 $+1$，正好对应“取 $(n+1)$ 维配边谱的一个类”。作者说这是“roadmap”而非证明，实质是 [FHT] 的定理；这正是下一段用稳定同伦论重述的原因。
  \item 易错——反常相位常写成 $e^{2\pi i\eta(D)}$，准确说是 $e^{2\pi i\xi}$，其中 $\xi(Y)=\frac{1}{2}\bigl(\eta(D)+\dim\ker D\bigr)$：对闭 $\Spin$ 流形 $Y$，$\xi(Y)$ 只在模 $1$ 意义下良定义（换度量或算子会差整数），故指数化后相位有意义；仅当 $Y$ 是某个 $\Spin$ 流形 $W$ 的边界、且维数使 $\int_W\Ahat\in\ZZ$ 时，APS 指标定理才给出 $\xi(Y)\equiv\int_W\Ahat\pmod 1$ 为整数（相位为 $1$）。若忘了 $\frac{1}{2}\dim\ker$ 或谱流修正而直接用 $e^{2\pi i\eta}$，相位不闭环——这正是整体反常的和乐。
  \item 含糊——“conjectural as a mathematical assertion until more foundations are laid”指 [FH, Conjecture 9.70] 依赖于配边谱与路径积分的对应尚未完全严格化（需把 $\eta$ 不变量、指标密度、$M\!\Spin\to KO$ 的 Atiyah–Bott–Shapiro 映射（[ABS]）组装成一个可逆场论）。对旋量场，$M\!\Spin\to KO$ 把 $n$ 维 $\Spin$ 配边类送到 $KO^{-n}$ 中的类，其与 $\eta$ 的关系是 APS 定理的“谱侧”表述；作者说“Implicit in it are ... index theorems”，正是本节的落脚点。
\end{enumerate}



\clearpage
\section*{附录 B：核心术语对照表}
\noindent 下表给出本译本采用的英--中术语对应（按原文出现的大致顺序排列）。

\begin{longtable}{@{}p{0.42\linewidth}p{0.52\linewidth}@{}}
\toprule\textbf{English} & \textbf{中文译名}\<本地路径>
index & 指数（Fredholm 指数：$\dim\ker-\dim\coker$）\\
elliptic operator & 椭圆算子\\
pseudodifferential operator & 伪微分算子\\
principal symbol & 主象征\\
cokernel & 余核\\
cohomology / homology & 上同调 / 同调\\
de Rham cohomology & de Rham 上同调\\
Dolbeault cohomology & Dolbeault 上同调\\
Hodge theory & Hodge 理论\\
harmonic form & 调和形式\\
Laplacian & 拉普拉斯算子\\
self-adjoint & 自伴的\\
compact operator & 紧算子\\
Fredholm operator & Fredholm 算子\\
vector bundle / line bundle & 向量丛 / 线丛\\
section (of a bundle) & 截面\\
connection / curvature & 联络 / 曲率\\
holonomy & 和乐（和乐群）\\
characteristic class & 示性类\\
Chern class / Chern character & 陈类 / 陈特征\\
Pontryagin class & Pontryagin 类\\
Euler class & 欧拉类\\
Todd class / Todd genus & Todd 类 / Todd 亏格\\
$\Ahat$-genus & $\Ahat$ 亏格\\
signature & 号差\\
Riemann--Roch theorem & 黎曼--罗赫定理\\
sheaf / coherent sheaf & 层 / 凝聚层\\
divisor & 除子\\
genus (of a surface) & 亏格\\
spin structure & spin 结构\\
spinor / spinor bundle & 旋量 / 旋量丛\\
Clifford algebra / module & Clifford 代数 / 模\\
Dirac operator & Dirac 算子\\
chirality / chiral & 手征性 / 手征的\\
heat kernel & 热核\\
zeta regularization & 泽塔（$\zeta$）正规化\\
spectral flow & 谱流\\
eta invariant & $\eta$ 不变量\\
determinant line bundle & 行列式线丛\\
path integral & 路径积分\\
partition function & 配分函数\\
fermion / boson & 费米子 / 玻色子\\
gauge theory / gauge field & 规范场论 / 规范场\\
anomaly & 反常\\
chiral anomaly & 手征反常\\
global anomaly & 整体反常\\
gravitational anomaly & 引力反常\\
anomaly inflow & 反常流入\\
moduli space & 模空间\\
$K$-theory / Bott periodicity & $K$ 理论 / Bott 周期性\\
classifying space & 分类空间\\
equivariant & 等变的\\
fixed point & 不动点\\
von Neumann algebra & 冯·诺依曼代数\\
bordism / cobordism & 配边\\
topological field theory (TQFT) & 拓扑场论\\
invertible field theory & 可逆场论\\
adiabatic limit & 绝热极限\\
APS boundary condition & APS 边界条件\\
signature defect & 号差亏损\\
secondary invariant & 次级不变量\\
Chern--Simons form & 陈--西蒙斯形式\\
projective representation & 射影表示\\
central extension & 中心扩张\\
\bottomrule\end{longtable}


\end{document}
